% !TeX program = xelatex
% 译本：《分形集的几何》(The Geometry of Fractal Sets, K. J. Falconer,
% Cambridge Tracts in Mathematics 85, Cambridge University Press, 1985).
% 编译：xelatex -> biber -> xelatex -> xelatex
\documentclass[
  BookTitleCN={分形集的几何},
  BookTitleEN={The Geometry of Fractal Sets},
  OriginalAuthor={K. J. Falconer},
  OriginalEdition={Cambridge Tracts in Mathematics 85},
  OriginalPublisher={Cambridge University Press},
  OriginalYear={1985},
  Translator={Fingolfin},
  ModelUsed={GLM-5.3-Flash},
  TranslationDate={11 Sep 2026},
  BibStyle={authoryear}
]{mathtranslation}

\raggedbottom

\addbibresource{references.bib}

% 参考文献与引文中作者名之间使用中文连接词
\DeclareDelimFormat{multinamedelim}{~与~}
\DeclareDelimFormat{finalnamedelim}{~与~}

% ================================================================
%  与原书一致的编号机制
% ================================================================
% 原书中定义、定理、引理、推论共用一个按章计数的计数器
% (如 Theorem 1.10、Lemma 1.3、Corollary 2.4)，此处重新声明.
\makeatletter
\let\definition\relax\let\enddefinition\relax
\let\lemma\relax\let\endlemma\relax
\let\theorem\relax\let\endtheorem\relax
\let\corollary\relax\let\endcorollary\relax
\let\proposition\relax\let\endproposition\relax
\let\remark\relax\let\endremark\relax
\let\example\relax\let\endexample\relax
\let\problem\relax\let\endproblem\relax
% 原计数器一并清除, 否则 \newtheorem 因 \c@definition 已定义而报错
\csundef{c@definition}\csundef{c@lemma}\csundef{c@theorem}\csundef{c@corollary}%
\csundef{c@proposition}\csundef{c@remark}\csundef{c@example}\csundef{c@problem}
% amsmath 的计数器检查点列表 \cl@@ckpt 中残留被清除的计数器, 一并滤除
\def\MT@elt#1{\ifcsname c@#1\endcsname\noexpand\@elt{#1}\fi}
{\let\@elt\MT@elt\xdef\cl@@ckpt{\cl@@ckpt}}
\theoremstyle{mtfang}
\newtheorem{definition}{定义}[chapter]
\newtheorem{lemma}[definition]{引理}
\newtheorem{theorem}[definition]{定理}
\newtheorem{corollary}[definition]{推论}
\newtheorem{proposition}[definition]{命题}
\newtheorem{remark}[definition]{注}
\theoremstyle{mtsong}
\newtheorem{example}[definition]{例}
\newtheorem{problem}[definition]{问题}

% 清除旧 [section] 声明留在 \cl@section 重置链中的残留
\@removefromreset{definition}{section}\@removefromreset{lemma}{section}%
\@removefromreset{theorem}{section}\@removefromreset{corollary}{section}%
\@removefromreset{proposition}{section}\@removefromreset{remark}{section}%
\@removefromreset{example}{section}\@removefromreset{problem}{section}

% 原书所有定理类环境共用一个计数器; cleveref 依计数器判定引用类型.
% 通过 \crefalias 在各环境开始处设置别名, 使 \cref 能区分定理/引理/推论等.
\AtBeginEnvironment{definition}{\crefalias{definition}{definition}}
\AtBeginEnvironment{lemma}{\crefalias{definition}{lemma}}
\AtBeginEnvironment{theorem}{\crefalias{definition}{theorem}}
\AtBeginEnvironment{corollary}{\crefalias{definition}{corollary}}
\AtBeginEnvironment{proposition}{\crefalias{definition}{proposition}}
\AtBeginEnvironment{remark}{\crefalias{definition}{remark}}
\AtBeginEnvironment{example}{\crefalias{definition}{example}}
\AtBeginEnvironment{problem}{\crefalias{definition}{problem}}
\crefname{exercise}{习题}{习题}

% 公式与插图按章编号 (原书: (1.12)、Fig. 8.5).
\numberwithin{equation}{chapter}
\counterwithin{figure}{chapter}

% 习题按章编号 (原书: Exercise 1.6), 习题环境在章末统一给出.
\counterwithout{exercise}{section}
\counterwithin{exercise}{chapter}
\renewcommand{\theexercise}{\thechapter.\arabic{exercise}}
\makeatother

% 章首的 \cleardoublepage 插入的空白页不排页眉页脚
\renewcommand{\cleardoublepage}{%
  \clearpage
  \if@twoside
    \ifodd\c@page\else
      \thispagestyle{empty}%
      \hbox{}\newpage
    \fi
  \fi
}
\makeatother

% cleveref 多重引用使用中文连接词 (命令由 cleveref 在 \begin{document} 时定义)
\AtBeginDocument{%
  \renewcommand{\crefpairconjunction}{ 与 }%
  \renewcommand{\crefmiddleconjunction}{、}%
  \renewcommand{\creflastconjunction}{ 与 }%
}

\begin{document}

\maketitle

\frontmatter
% !TeX root = ../main.tex

\phantomsection
\chapter*{译者说明}
\addcontentsline{toc}{chapter}{译者说明}

本书为 K.~J.~Falconer 所著 \emph{The Geometry of Fractal Sets}（《分形集的几何》，Cambridge Tracts in Mathematics 85，Cambridge University Press，1985 年初版，1986 年附勘误平装重印本）的中文翻译与重排版.

原书是关于整数维数与分数维数集合之几何测度论的一部严谨而自足的专著，以 Hausdorff 测度与 Hausdorff 维数为主线，系统阐述了 $s$-集的密度、切线性质、网测度方法、投影定理、Kakeya 问题以及分形集的各类例子.本次翻译以原书 1985 年初版为底本，参照 1986 年平装重印本的勘误.

翻译遵循以下原则：

\begin{enumerate}
    \item 以原书为正文内容、数学公式、章节结构与编号的最高依据；
    \item 数学正确性优先于语言润色，不增删或改写正文内容；
    \item 数学术语在全文中保持一致，新术语在首次出现处以楷体标出并附英文原词；
    \item 定理、引理、推论沿用原书的按章统一编号，公式与插图亦按章编号，习题按原书格式附于各章末尾；
    \item 原书的参考文献列表完整保留，正文中的文献引用一律通过文献数据库实现交叉引用；原书参考文献列表中方括号所注的引用章节，以注记形式保留在各条目之后.
\end{enumerate}

原书正文中的个别明显排印错误（如符号表中的 ``Grossmann manifold'' 应为 Grassmann 流形等）已在相应位置以译者注说明.书中发现的其余可疑之处，均核对原书影印件后处理.

\vspace{0.5cm}
\hfill Fingolfin

\hfill \TranslationDate

% !TeX root = ../main.tex

\chapter*{前言}
\addcontentsline{toc}{chapter}{前言}

这本专著对具有分数维数与整数维数 Hausdorff 维数的集合的数学理论给出了严谨而自足的论述.它主要关注几何理论而非应用.书中的许多内容此前只能见于原始数学论文，而这些论文往往高度专门、令人困惑，并且使用陈旧的记号.我希望通过撰写本书使这些材料更易于获取，同时也为使用分形集合的读者提供一份有用而精确的论述.

本书虽然主要为纯数学家而写，但我希望它对若干类型或多或少有素养且有较高要求的读者也有用处.在最基本的层次上，本书可供在其他数学或科学学科中遇到分形的读者作为参考书使用：将主要定理与推论同基本定义结合起来阅读，即可得到那些已被严格建立的性质的精确表述.

要想获得这一学科的全貌，或者在初次阅读时，可以只跟随基本的解说与定理的陈述而略去详细的证明.非专业的数学工作者或许还可以略去~1.1~节的细节，该节从技术的观点建立了一般测度的性质.

要充分领会书中的细节，需要具备初等数学分析与一般拓扑的实用知识.毫无疑问，构成这一理论核心的某些证明既困难又冗长，但为了完整洞察其中数学的美妙与精巧，掌握它们是十分值得的.

本书着重于该学科的基本工具，例如 Vitali 覆盖定理、网测度与位势论方法.

我们所需要的测度与 Hausdorff 测度的性质在第 1 章的前两节中建立.全书着重于把测度就其自身用于估计集合的大小，而不是将其作为定义积分的一个步骤.积分只在后面各章中作为一个方便的工具使用；大体上，关于积分的直观认识就完全够用了.

在所采用的普遍性水平上不可避免地做了一些折中.我们在 $n$ 维 Euclid 空间中工作，尽管许多思想同样适用于更一般的度量空间.在某些情形，高维类似物的证明要复杂得多，此时定理只在二维情形给出证明，并对其推广给出参考文献.类似地，有时也只给出一维或二维的证明，如果它们包含一般情形的本质思想，同时又能使记号得到简化.我们还把注意力限制在与数值维数 $s$ 相应的 Hausdorff 测度上，而不是任意函数对应的 Hausdorff 测度.

书中不少证明相对原始形式做了简化.进一步的缩短无疑也是可能的，但作者对完美性的追求不得不与在有限时间内完成本书的要求相权衡！

虽然本专著本质上是自足的，书中还是简要描述了这一工作的各种变形与推广，并给出了完整的参考文献.更多的变形与推广可以在各章末尾的习题中找到.

在此谨向所有以任何方式帮助过本专著写作的人们致谢.我特别感谢 Roy Davies 教授对书稿的细致批评，并允许我使用其未发表的资料；感谢 Hallard Croft 博士提出的详尽建议以及在校对中给予的帮助.我也非常感谢 B.~B.~Mandelbrot 教授、J.~M.~Marstrand 教授、P.~Mattila 教授和 C.~A.~Rogers 教授有益的意见与讨论.

我还要感谢 Maureen Woodward 夫人与 Rhoda Rees 夫人为书稿打字，并感谢 Cambridge University Press 的 David Tranah 与 Sheila Shepherd 陪伴本书走过了出版的各个阶段.最后，我必须感谢我的妻子 Isobel 拨冗阅读了本书的早期草稿，以及她始终如一的鼓励与支持.

% !TeX root = ../main.tex

\chapter*{导言}
\addcontentsline{toc}{chapter}{导言}

整数维数与分数维数集合的几何测度论自本世纪初以来一直由纯数学家们发展.近来，分形集合在科学中的重要性急剧上升.Mandelbrot（\citeyear{mandelbrot1975,mandelbrot1977,mandelbrot1982}） 开创了用分形集合为从分子到天文尺度的各种科学现象建模的先河，例如：粒子的布朗运动、流体中的湍流、植物的生长、地理上的海岸线与地面、宇宙中星系的分布，甚至股票市场价格的波动.分数维数集合也出现在纯数学的诸多分支中，例如数论与非线性微分方程.更多的例子在 Mandelbrot 的科学性、哲学性与图示性的论述中有所描述.于是，起源于纯数学的概念在科学中找到了大量应用，而这些应用又为数学家提供了进一步问题的丰富源泉.本专著主要关心这类集合的几何理论，而非其应用.

\newterm{fractal}{分形}{fractal}（fractal）一词由 Mandelbrot（\citeyear{mandelbrot1975}）从意为``破碎的''拉丁词 \emph{fractus} 导出，他给出了分形的一个``尝试性定义'':Hausdorff 维数严格大于拓扑维数的集合；但他同时指出这一定义并不令人满意，因为它排除了某些高度不规则的集合，而这些集合按其精神显然应当视为分形.此前数学家们曾以多种方式称呼这类集合——``分数维数集合''、``Hausdorff 测度集合''、``具有精细结构的集合''或``不规则集合''.对这类集合的任何严格研究也必须包含对那些拓扑维数与 Hausdorff 维数相等的集合的考察，哪怕只是为了将它们从讨论中排除出去.因此，我在本书中收录这类正则集合（光滑曲线与曲面等）时并不感到抱歉.

为了推广点、曲线、曲面分别具有维数 $0$、$1$、$2$ 的观念，人们提出了许多估计``细的''或``高度不规则''集合的``大小''或``维数''的方法.由 Hausdorff 测度定义的 Hausdorff 维数，从数学家的观点看具有压倒性的优点：Hausdorff 测度是一个测度（即在两两不交集合的可数族上具有可加性）.

遗憾的是，即便是相对简单的集合，其 Hausdorff 测度与维数也可能很难计算；特别是这些量的下界往往不易得到.这已成为物理应用中的一个相当大的缺陷，并导致 Hausdorff 维数定义的若干变体被采用，其中某些变体的采用甚至是无意的.

\textcite{hurewiczwallman1941}、\textcite{kahane1976}、\textcite{mandelbrot1982}（第 39 节）与 Tricot（\citeyear{tricot1981,tricot1982}） 综述并比较了这些 Hausdorff 维数的替代定义.它们包括熵（见 \textcite{hawkes1974}）、相似维数（见 \textcite{mandelbrot1982}），以及 \textcite{johnsonrogers1982} 的局部维数与局部测度.基于任何这样的定义都不难写出一本类似性质的书，但 Hausdorff 测度与维数无疑是研究得最广泛、使用得最广泛的.

通过定义外测度把区间长度的概念推广到更复杂的实数集合，这一想法出现得出人意料地晚.\textcite{borel1895} 利用测度估计集合的大小，从而构造了某些病态函数.这些思想被 \textcite{lebesgue1904} 采纳，成为其积分理论的基础概念.\textcite{caratheodory1914} 引入了更一般的 ``Carathéodory 外测度''.特别地，他在 $n$ 维 Euclid 空间中定义了 ``$1$ 维''或``线性''测度，并指出 $s$ 维测度对其他整数 $s$ 也可以类似地定义.\textcite{hausdorff1919} 指出 Carathéodory 的定义对非整数 $s$ 同样有价值.他通过证明 Cantor 著名的``中三分之一''集在 $s=\log 2/\log 3=0.6309\ldots$ 时具有正而有限的 $s$ 维测度来说明这一点.分数维数集合的概念由此诞生，而 Hausdorff 之名也被用于称呼相应的维数与测度.

从那时起，关于 Hausdorff 测度及 Hausdorff 可测集合几何的大量结果被发现.\textcite{rogers1970} 的著作对内在的测度论给出了出色的论述，而 \textcite{federer1969} 的学术巨著则给出了测度几何的一个非常一般的处理，并延伸到曲面面积与同调积分理论的问题.

关于 Hausdorff 测度及其几何的许多工作归功于 Besicovitch，读者在本书中将反复遇到他的名字.事实上，多年来几乎所有关于 Hausdorff 测度的已发表工作都带有他的名字，其中大量工作包含了极为精巧的论证.近年来他的学生们又做出了许多重要的进一步贡献.\textcite{burkill1971} 与 \textcite{taylor1975} 的讣告文章可以让读者对 Besicovitch 在这一学科中的影响之巨大有所体会.

显然，Besicovitch 曾打算写一本题为 \emph{The Geometry of Sets of Points}（点集的几何）的几何测度论著作，它在许多方面很可能与本卷相似.Besicovitch 于 1970 年去世后，Roy Davies 教授在 Helen Alderson 博士（1972 年去世）的协助下，整理出了可能构成 Besicovitch``第 1 章''的一个版本.这一章注定没有任何后续，但它对本卷的前面部分产生了相当大的影响.

在第 1 章中我们定义 Hausdorff 测度并研究其基本性质，说明在某些直截了当的情形下如何计算集合的 Hausdorff 维数与测度.

我们特别感兴趣的是维数为 $s$ 的 $s$-集，即 $s$ 维 Hausdorff 测度非零但有限的集合.对仅受如此弱条件限制的一类集合而言，其几何必然只能是一般点邻域性质的研究.因此接下来的三章讨论局部性质：集合在一点处的密度，以及集合绕其每一点的方向分布，即切线的存在性问题.分数维数集合与整数维数集合分别处理.分数维数集合必然是分形，但整数维数集合中正则的``曲线状''或``曲面状''集合与不规则的``分形''集合之间存在显著对照.

第 5 章引入强有力的网测度技术.这使我们能够证明任何具有无穷 $s$ 维 Hausdorff 测度的集合都包含一个 $s$-集，从而按要求把 $s$-集理论推广到更一般的集合.网测度还被用于研究集合的 Descartes 积的 Hausdorff 测度.

下一章处理集合到低维子空间的投影.在这一部分工作中，位势论方法作为直接几何方法的替代被引入.

第 7 章讨论 Kakeya 问题，即求在其中可以转动单位长线段的最小测度集合.文中讨论了若干变体，并通过对偶性与前一章的投影定理以及调和分析联系起来.

最后一章收录了若干杂例，用以说明书中前面遇到的一些思想.

参考文献列于书末并按年份引用.更多的文献目录见 Federer（\citeyear{federer1947,federer1969}）与 \textcite{rogers1970} 与 \textcite{mandelbrot1982}.

% !TeX root = ../main.tex

\chapter*{符号表}
\addcontentsline{toc}{chapter}{符号表}

鉴于本书论题的范围，特别是最后一章的范围，记号的使用不可能完全一致.一般而言，符号在首次引入时定义；本节仅作为一个粗略的指南.

我们完全在 $n$ 维 Euclid 空间 $\mathbb{R}^n$ 中工作.$\mathbb{R}^n$ 的点有时按向量的意义理解，用小写字母 $x,y,z$ 等表示.偶尔我们也用 $(x,y)$ 表示 Descartes 坐标.大写字母 $E,F,\Gamma$ 等用于表示 $\mathbb{R}^n$ 的子集，花体大写字母 $\mathcal{C},\mathcal{V},\mathcal{I}$ 用于表示集合族.我们采用如下约定：集合包含符号 $\subset$ 允许取等的可能性.集合 $E$ 的直径记作 $|E|$；但在意义明确时，模长符号也按通常方式表示向量的长度，例如 $|x-y|$ 表示点 $x$ 与 $y$ 之间的距离.常数 $b,c,c_1,\varepsilon,\delta$ 与指标 $i,j,k$ 也用（可带下标的）小写字母表示.

下表可用作对较常用记号的提醒.

\bigskip
\noindent
\begin{tabular}{@{}p{0.22\textwidth}@{\quad}p{0.70\textwidth}@{}}
\multicolumn{2}{@{}l}{\headfont 集合}\\[2pt]
$\mathbb{R}^n$ & $n$ 维 Euclid 空间.\\
$B_r(x)$ & 以 $x$ 为心、$r$ 为半径的闭圆盘或闭球.\\
$S_r(x,\theta,\phi)$ & 角为 $\phi$、半径为 $r$ 的扇形.\\
$C_r(x,I)$ & 双扇形.\\
$R(x,y)$ & 以 $x$ 与 $y$ 为心的圆对所围的公共区域.\\
$G_{n,k}$ & $\mathbb{R}^n$ 的 $k$ 维子空间的 Grassmann 流形.\footnotemark\\
$L(a,b),L(E)$ & 直线集合.\\
$\bar{E},\operatorname{int}E$ & $E$ 的拓扑闭包与内部.\\
$[E]_\delta$ & $E$ 的 $\delta$-平行体.\\[8pt]
\multicolumn{2}{@{}l}{\headfont 映射}\\[2pt]
$\operatorname{proj}_\theta,\operatorname{proj}_\Pi$ & 到方向 $\theta$ 的直线上的正交投影，相应地，到平面 $\Pi$ 上的正交投影.\\
$\hat{f},\hat{\mu}$ & 函数 $f$ 与测度 $\mu$ 的 Fourier 变换.\\
$f\circ g$ & 映射的复合：先 $g$ 后 $f$.\\
\end{tabular}

\medskip
\noindent
\begin{tabular}{@{}p{0.22\textwidth}@{\quad}p{0.70\textwidth}@{}}
\multicolumn{2}{@{}l}{\headfont 测度等}\\[2pt]
$\mathcal{H}^s$ & $s$ 维 Hausdorff 测度或外测度.\\
$\mathcal{L}^n$ & $n$ 维勒贝格测度.\\
$\mathcal{M}^s$ & $s$ 维可比较网测度.\\
$\mathcal{H}^s_\delta,\ \mathcal{M}^s_\delta$ & 构造 $\mathcal{H}^s$ 与 $\mathcal{M}^s$ 时所用的 $\delta$-外测度.\\
$\mathscr{L}(\Gamma)$ & 曲线 $\Gamma$ 的长度.\\
$\dim E$ & $E$ 的 Hausdorff 维数.\\
$\phi_t,\ C_t,\ I_t$ & $t$-位势、容量、能量.\\[8pt]
\multicolumn{2}{@{}l}{\headfont 密度}\\[2pt]
$D^s(E,x)$ & $E$ 在 $x$ 处的密度.\\
$\underline{D}^s(E,x),\ \bar{D}^s(E,x)$ & 下密度、上密度.\\
$\bar{D}^s_c(E,x)$ & 上凸密度.\\
$\underline{D}^s(E,x,\theta,\phi)$ & 下角密度等.\\
\end{tabular}
\footnotetext{原书此处印作 ``Grossmann manifold'', 应为 Grassmann 流形——译者注.}


\makecontents

\mainmatter
% !TeX root = ../main.tex
% 第 1 章 测度与维数 (对应原书 pp. 1--19)

\chapter{测度与维数}

\section{基本测度论}

本节简明扼要地叙述我们所需要的基本测度论.更完整的论述可参见 \textcite{kingmantaylor1966} 或 \textcite{rogers1970}.

设 $X$ 为任意集合.(我们随后将取 $X$ 为 $n$ 维 Euclid 空间 $\mathbb{R}^n$.)称 $X$ 的子集的一个非空族 $\mathcal{S}$ 为 \newterm{sigma-field}{$\sigma$-域}{sigma-field}，如果 $\mathcal{S}$ 在取补与可数并下封闭（即若 $E\in\mathcal{S}$，则 $X\setminus E\in\mathcal{S}$；且若 $E_1,E_2,\ldots\in\mathcal{S}$，则 $\bigcup_{j=1}^{\infty}E_j\in\mathcal{S}$）.一点初等集合论表明，$\sigma$-域在可数交与集合差下也封闭，并且 $X$ 与空集 $\emptyset$ 属于 $\mathcal{S}$.

集合序列 $\{E_j\}$ 的\newterm{limit-lower}{下极限}{limit, lower}与\newterm{limit-upper}{上极限}{limit, upper}定义为
\[
  \varliminf_{j\to\infty}E_j=\bigcup_{k=1}^{\infty}\bigcap_{j=k}^{\infty}E_j
\]
与
\[
  \varlimsup_{j\to\infty}E_j=\bigcap_{k=1}^{\infty}\bigcup_{j=k}^{\infty}E_j.
\]
于是 $\varliminf E_j$ 由那些除有限多个 $E_j$ 之外都属于每个 $E_j$ 的点组成，而 $\varlimsup E_j$ 由那些属于无穷多个 $E_j$ 的点组成.由这些定义的形式显然可见，若每个 $E_j$ 都属于 $\sigma$-域 $\mathcal{S}$，则 $\varliminf E_j,\ \varlimsup E_j\in\mathcal{S}$.若 $\varliminf E_j=\varlimsup E_j$，则用 $\lim E_j$ 表示这个公共值；当 $\{E_j\}$ 为递增或递减集合序列时总会出现这种情形.

设 $\mathcal{C}$ 为 $X$ 的子集的任意族.称由 $\mathcal{C}$ \newterm{sigma-field-generated}{生成的 $\sigma$-域}{sigma-field generated by}，记作 $\mathcal{S}(\mathcal{C})$，为所有包含 $\mathcal{C}$ 的 $\sigma$-域之交.直接验证可知 $\mathcal{S}(\mathcal{C})$ 本身是一个 $\sigma$-域，可以视为包含 $\mathcal{C}$ 的``最小''的 $\sigma$-域.

\newterm{measure}{测度}{measure}是定义在 $X$ 的子集的某个 $\sigma$-域 $\mathcal{S}$ 上、取值于 $[0,\infty]$ 的函数 $\mu$，使得
\begin{equation}\label{eq:1.1}
  \mu(\emptyset)=0
\end{equation}
并且
\begin{equation}\label{eq:1.2}
  \mu\Bigl(\bigcup_{j=1}^{\infty}E_j\Bigr)=\sum_{j=1}^{\infty}\mu(E_j)
\end{equation}
对 $\mathcal{S}$ 中每个由两两不交的集合构成的可数序列 $\{E_j\}$ 成立.

由 \eqref{eq:1.2} 可知 $\mu$ 是递增的集合函数，即若 $E\subset E'$ 且 $E,E'\in\mathcal{S}$，则
\[
  \mu(E)\leq\mu(E').
\]

\begin{theorem}[\newterm{continuity-of-measures}{测度的连续性}{continuity of measures}]\label{thm:1.1}
设 $\mu$ 是 $X$ 的子集的 $\sigma$-域 $\mathcal{S}$ 上的测度.

(a) 若 $E_1\subset E_2\subset\cdots$ 是 $\mathcal{S}$ 中的递增集合序列，则
\[
  \mu\Bigl(\lim_{j\to\infty}E_j\Bigr)=\lim_{j\to\infty}\mu(E_j).
\]

(b) 若 $F_1\supset F_2\supset\cdots$ 是 $\mathcal{S}$ 中的递减集合序列且 $\mu(F_1)<\infty$，则
\[
  \mu\Bigl(\lim_{j\to\infty}F_j\Bigr)=\lim_{j\to\infty}\mu(F_j).
\]

(c) 对 $\mathcal{S}$ 中的任意集合序列 $\{F_j\}$，
\[
  \mu\Bigl(\varliminf_{j\to\infty}F_j\Bigr)\leq\varliminf_{j\to\infty}\mu(F_j).
\]
\end{theorem}

\begin{proof}
(a) 我们可以把 $\bigcup_{j=1}^{\infty}E_j$ 表为不交并 $E_1\cup\bigcup_{j=2}^{\infty}(E_j\setminus E_{j-1})$.于是由 \eqref{eq:1.2}，
\[
\begin{aligned}
  \mu\Bigl(\lim_{j\to\infty}E_j\Bigr)
    &=\mu\Bigl(\bigcup_{j=1}^{\infty}E_j\Bigr)
     =\mu(E_1)+\sum_{2}^{\infty}\mu(E_j\setminus E_{j-1})\\
    &=\lim_{k\to\infty}\Bigl[\mu(E_1)+\sum_{2}^{k}\mu(E_j\setminus E_{j-1})\Bigr]
     =\lim_{k\to\infty}\mu\Bigl(E_1\cup\bigcup_{2}^{k}(E_j\setminus E_{j-1})\Bigr)
     =\lim_{k\to\infty}\mu(E_k).
\end{aligned}
\]

(b) 若令 $E_j=F_1\setminus F_j$，则 $\{E_j\}$ 如 (a) 中所述.由于 $\bigcap_{j}F_j=F_1\setminus\bigcup_{j}E_j$，
\begin{align*}
  \mu\Bigl(\lim_{j\to\infty}F_j\Bigr)
  &=\mu\Bigl(\bigcap_{j=1}^{\infty}F_j\Bigr)
  =\mu(F_1)-\mu\Bigl(\bigcup_{j}E_j\Bigr)
  =\mu(F_1)-\lim_{j\to\infty}\mu(E_j)\\
  &=\lim_{j\to\infty}\bigl(\mu(F_1)-\mu(E_j)\bigr)
  =\lim_{j\to\infty}\mu(F_j).
\end{align*}

(c) 现在令 $E_k=\bigcap_{j=k}^{\infty}F_j$.则 $\{E_k\}$ 是 $\mathcal{S}$ 中的递增集合序列，故由 (a)，
\[
  \mu\Bigl(\lim_{j\to\infty}F_j\Bigr)
  =\mu\Bigl(\bigcup_{k=1}^{\infty}E_k\Bigr)
  =\lim_{k\to\infty}\mu(E_k)
  \leq\varliminf_{j\to\infty}\mu(F_j). \qedhere
\]
\end{proof}

接下来我们引入 \newterm{outer measure}{外测度}{outer measure}，它本质上是把性质 \eqref{eq:1.2} 减弱为次可加性的``测度''.确切地说，集合 $X$ 上的外测度 $\nu$ 是定义在 $X$ 的所有子集上、取值于 $[0,\infty]$ 的函数，使得
\begin{equation}\label{eq:1.3}
  \nu(\emptyset)=0,
\end{equation}
\begin{equation}\label{eq:1.4}
  \nu(A)\leq\nu(A')\qquad\text{若 } A\subset A',
\end{equation}
并且
\begin{equation}\label{eq:1.5}
  \nu\Bigl(\bigcup_{1}^{\infty}A_j\Bigr)\leq\sum_{1}^{\infty}\nu(A_j)
  \qquad\text{对 $X$ 的任意子集 } \{A_j\}.
\end{equation}

外测度之所以有用，是因为总存在一个子集的 $\sigma$-域，外测度在其上表现得如测度一般；对定义合理的外测度，这个 $\sigma$-域可以相当大.

称 $X$ 的子集 $E$ 为 \newterm{measurable set}{$\nu$-可测的}{measurable set}（即关于外测度 \newterm{outer-measurable}{可测}{outer measurable}），如果它对 $X$ 的每个子集都可加分解，即若
\begin{equation}\label{eq:1.6}
  \nu(A)=\nu(A\cap E)+\nu(A\setminus E)
\end{equation}
对所有``检验集''$A\subset X$ 成立.注意，为证集合 $E$ 是 $\nu$-可测的，只需验证
\begin{equation}\label{eq:1.7}
  \nu(A)\geq\nu(A\cap E)+\nu(A\setminus E),
\end{equation}
因为相反的不等式已含于 \eqref{eq:1.5} 中.显然可验证：若 $\nu(E)=0$，则 $E$ 是 $\nu$-可测的.

\begin{theorem}\label{thm:1.2}
设 $\nu$ 是一个外测度.$\nu$-可测集的族 $\mathcal{M}$ 构成一个 $\sigma$-域，且 $\nu$ 在 $\mathcal{M}$ 上的限制是一个测度.
\end{theorem}

\begin{proof}
显然 $\emptyset\in\mathcal{M}$，故 $\mathcal{M}$ 非空.又由 \eqref{eq:1.6} 的对称性，$A\in\mathcal{M}$ 当且仅当 $X\setminus A\in\mathcal{M}$，故 $\mathcal{M}$ 在取补下封闭.

为证 $\mathcal{M}$ 在可数并下封闭，设 $E_1,E_2,\ldots\in\mathcal{M}$，并设 $A$ 为任意集合.对 $E_1,E_2,\ldots$ 依次应用 \eqref{eq:1.6}，取相应的检验集，得
\[
\begin{aligned}
  \nu(A)
    &=\nu(A\cap E_1)+\nu(A\setminus E_1)\\
    &=\nu(A\cap E_1)+\nu\bigl((A\setminus E_1)\cap E_2\bigr)+\nu(A\setminus E_1\setminus E_2)\\
    &=\cdots\\
    &=\sum_{j=1}^{k}\nu\Bigl(\Bigl(A\setminus\bigcup_{i=1}^{j-1}E_i\Bigr)\cap E_j\Bigr)
      +\nu\Bigl(A\setminus\bigcup_{j=1}^{k}E_j\Bigr).
\end{aligned}
\]
于是
\[
  \nu(A)\geq\sum_{j=1}^{k}\nu\Bigl(\Bigl(A\setminus\bigcup_{i=1}^{j-1}E_i\Bigr)\cap E_j\Bigr)
  +\nu\Bigl(A\setminus\bigcup_{j=1}^{\infty}E_j\Bigr)
\]
对一切 $k$ 成立，从而
\begin{equation}\label{eq:1.8}
  \nu(A)\geq\sum_{j=1}^{\infty}\nu\Bigl(\Bigl(A\setminus\bigcup_{i=1}^{j-1}E_i\Bigr)\cap E_j\Bigr)
  +\nu\Bigl(A\setminus\bigcup_{j=1}^{\infty}E_j\Bigr).
\end{equation}
另一方面，
\[
  A\cap\bigcup_{j=1}^{\infty}E_j=\bigcup_{j=1}^{\infty}\Bigl(\Bigl(A\setminus\bigcup_{i=1}^{j-1}E_i\Bigr)\cap E_j\Bigr),
\]
故利用 \eqref{eq:1.5}，
\[
  \nu(A)\leq\nu\Bigl(A\cap\bigcup_{j=1}^{\infty}E_j\Bigr)+\nu\Bigl(A\setminus\bigcup_{j=1}^{\infty}E_j\Bigr)
  \leq\sum_{j=1}^{\infty}\nu\Bigl(\Bigl(A\setminus\bigcup_{i=1}^{j-1}E_i\Bigr)\cap E_j\Bigr)
  +\nu\Bigl(A\setminus\bigcup_{j=1}^{\infty}E_j\Bigr)\leq\nu(A),
\]
其中末一步用到 \eqref{eq:1.8}.由此得 $\bigcup_{j=1}^{\infty}E_j\in\mathcal{M}$，故 $\mathcal{M}$ 是一个 $\sigma$-域.

现在设 $E_1,E_2,\ldots$ 为 $\mathcal{M}$ 中两两不交的集合.在 \eqref{eq:1.8} 中取 $A=\bigcup_{j=1}^{\infty}E_j$，得
\[
  \nu\Bigl(\bigcup_{j=1}^{\infty}E_j\Bigr)\geq\sum_{j=1}^{\infty}\nu(E_j),
\]
再与 \eqref{eq:1.5} 结合，便知 $\nu$ 在 $\mathcal{M}$ 上是一个测度.
\end{proof}

称外测度 $\nu$ 是 \emph{正则的}\newterm{regular-outer-measure}{正则的}{regular outer measure}，若对每个集合 $A$ 都存在包含 $A$ 的 $\nu$-可测集 $E$，使得 $\nu(A)=\nu(E)$.

\begin{lemma}\label{lem:1.3}
若 $\nu$ 是正则外测度，$\{A_j\}$ 是任意递增集合序列，则
\[
  \lim_{j\to\infty}\nu(A_j)=\nu\Bigl(\lim_{j\to\infty}A_j\Bigr).
\]
\end{lemma}

\begin{proof}
对每个 $j$，选取 $\nu$-可测集 $E_j$ 使 $E_j\supset A_j$ 且 $\nu(E_j)=\nu(A_j)$.于是利用 \eqref{eq:1.4} 与 \cref{thm:1.1}(c)，
\[
  \nu\Bigl(\lim A_j\Bigr)=\nu\Bigl(\varliminf A_j\Bigr)\leq\nu\Bigl(\varliminf E_j\Bigr)\leq\varliminf\nu(E_j)=\lim\nu(A_j).
\]
相反的不等式由 \eqref{eq:1.4} 得出.
\end{proof}

现在设 $(X,d)$ 为度量空间.（就我们的目的而言，$X$ 通常为 $n$ 维 Euclid 空间 $\mathbb{R}^n$，$d$ 为通常的距离函数.）属于由 $X$ 的闭子集生成的 $\sigma$-域的集合称为该空间的 \newterm{Borel sets}{Borel 集}{Borel set}.Borel 集包括开集（作为闭集的补集）、\newterm{fsigma-set}{$F_{\sigma}$-集}{F$\sigma$-set}（即可数个闭集之并）、\newterm{gdelta-set}{$G_{\delta}$-集}{G$\delta$-set}（即可数个开集之交）等.

称 $X$ 上的外测度 $\nu$ 为 \newterm{metric outer measure}{度量外测度}{metric outer measure}，若只要 $E$ 与 $F$ \emph{正分离}（\newterm{positively separated}{正分离}{positively separated}），即只要
\[
  d(E,F)=\inf\{d(x,y):x\in E,\,y\in F\}>0,
\]
就有
\begin{equation}\label{eq:1.9}
  \nu(E\cup F)=\nu(E)+\nu(F).
\end{equation}
我们将证明，若 $\nu$ 是度量外测度，则 $\nu$-可测集的族包含 Borel 集.证明基于下述形式的 \newterm{caratheodory-lemma}{Carathéodory 引理}{Caratheodory's lemma}.

\begin{lemma}\label{lem:1.4}
设 $\nu$ 是 $(X,d)$ 上的度量外测度.设 $\{A_j\}_{1}^{\infty}$ 是 $X$ 的子集的递增序列，$A=\lim_{j\to\infty}A_j$，并设 $d(A_j,A\setminus A_{j+1})>0$ 对每个 $j$ 成立.则 $\nu(A)=\lim_{j\to\infty}\nu(A_j)$.
\end{lemma}

\begin{proof}
只需证明
\begin{equation}\label{eq:1.10}
  \nu(A)\leq\lim_{j\to\infty}\nu(A_j),
\end{equation}
因为相反的不等式由 \eqref{eq:1.4} 得出.令 $B_1=A_1$，$B_j=A_j\setminus A_{j-1}$（$j\geq 2$）.若 $j+2\leq i$，则 $B_j\subset A_j$ 且 $B_i\subset A\setminus A_{i-1}\subset A\setminus A_{j+1}$，故 $B_i$ 与 $B_j$ 正分离.于是，把 \eqref{eq:1.9} 应用 $(m-1)$ 次得
\[
  \nu\Bigl(\bigcup_{k=1}^{m}B_{2k-1}\Bigr)=\sum_{k=1}^{m}\nu(B_{2k-1}),\qquad
  \nu\Bigl(\bigcup_{k=1}^{m}B_{2k}\Bigr)=\sum_{k=1}^{m}\nu(B_{2k}).
\]
我们可以假定这两个级数都收敛——否则，由于 $\bigcup_{k=1}^{m}B_{2k-1}$ 与 $\bigcup_{k=1}^{m}B_{2k}$ 都含于 $A_{2m}$，将有 $\lim_{j\to\infty}\nu(A_j)=\infty$.于是
\[
  \nu(A)=\nu\Bigl(\bigcup_{j=1}^{\infty}A_j\Bigr)
  =\nu\Bigl(A_j\cup\bigcup_{k=j+1}^{\infty}B_k\Bigr)
  \leq\nu(A_j)+\sum_{k=j+1}^{\infty}\nu(B_k)
  \leq\lim_{i\to\infty}\nu(A_i)+\sum_{k=j+1}^{\infty}\nu(B_k).
\]
由于该级数当 $j\to\infty$ 时趋于 $0$，故 \eqref{eq:1.10} 得证.
\end{proof}

\begin{theorem}\label{thm:1.5}
若 $\nu$ 是 $(X,d)$ 上的度量外测度，则 $X$ 的所有 Borel 子集都是 $\nu$-可测的.
\end{theorem}

\begin{proof}
由于 $\nu$-可测集构成一个 $\sigma$-域，而 Borel 集是含 $X$ 的闭子集的最小 $\sigma$-域，故只需证当 $E$ 为闭集、$A$ 为任意集合时 \eqref{eq:1.7} 成立.

设 $A_j$ 为 $A\setminus E$ 中与 $E$ 的距离至少为 $1/j$ 的点的集合.则 $d(A\cap E,A_j)\geq 1/j$，故
\begin{equation}\label{eq:1.11}
  \nu(A\cap E)+\nu(A_j)=\nu\bigl((A\cap E)\cup A_j\bigr)\leq\nu(A)
\end{equation}
对每个 $j$ 成立，因为 $\nu$ 是度量外测度.集合序列 $\{A_j\}$ 递增，且由于 $E$ 闭，$A\setminus E=\bigcup_{j=1}^{\infty}A_j$.于是，只要 $d(A_j,A\setminus E\setminus A_{j+1})>0$ 对一切 $j$ 成立，\cref{lem:1.4} 便给出 $\nu(A\setminus E)\leq\lim_{j\to\infty}\nu(A_j)$，再由 \eqref{eq:1.11} 即得 \eqref{eq:1.7}.但若 $x\in A\setminus E\setminus A_{j+1}$，则存在 $z\in E$ 使 $d(x,z)<1/(j+1)$，于是当 $y\in A_j$ 时
\[
  d(x,y)\geq d(y,z)-d(x,z)>1/j-1/(j+1)>0.
\]
故 $d(A_j,A\setminus E\setminus A_{j+1})>0$，如所需.
\end{proof}

还有一类重要的集合，与 Borel 集不同，它们由闭集的并与交显式地定义.若 $(X,d)$ 是度量空间，则 \newterm{Souslin sets}{Souslin 集}{Souslin set} 是如下形式的集合：
\[
  E=\bigcup_{i_1i_2\ldots}\ \bigcap_{k=1}^{\infty}E_{i_1i_2\ldots i_k},
\]
其中对每个由正整数构成的有限序列 $\{i_1,i_2,\ldots,i_k\}$，$E_{i_1i_2\ldots i_k}$ 都是闭集.注意，尽管 $E$ 由可数多个闭集构造而成，但其中的并是对连续统多个整数无穷序列取的.（每个闭集在该表达式中出现许多次.）

可以证明每个 Borel 集都是 Souslin 集，而且若底度量空间完备，则 Souslin 集的任何连续像仍是 Souslin 集.此外，若 $\nu$ 是度量空间 $(X,d)$ 上的外测度，则只要闭集都 $\nu$-可测，Souslin 集也就都 $\nu$-可测.由 \cref{thm:1.5} 可得：若 $\nu$ 是 $(X,d)$ 上的度量外测度，则 Souslin 集是 $\nu$-可测的.我们只顺带提及 Souslin 集.其测度论方面的更详细论述见 \textcite{rogers1970}，行家还可参阅 \textcites{rogers1980}.

\section{Hausdorff 测度}


从本书余下部分起，我们在 $n$ 维 Euclid 空间 $\mathbb{R}^n$ 中工作，不过应当强调，这里所说的许多内容在一般度量空间的框架下同样成立.

若 $U$ 是 $\mathbb{R}^n$ 的非空子集，我们定义 $U$ 的\newterm{diameter}{直径}{diameter}为 $|U|=\sup\{|x-y|:x,y\in U\}$.若 $E\subset\bigcup_{i}U_i$ 且 $0<|U_i|\leq\delta$ 对每个 $i$ 成立，则称 $\{U_i\}$ 为 $E$ 的一个 \newterm{delta-cover}{$\delta$-覆盖}{$\delta$-cover}.

设 $E\subset\mathbb{R}^n$，$s$ 为非负数.对 $\delta>0$ 定义
\begin{equation}\label{eq:1.12}
  \mathcal{H}^s_{\delta}(E)=\inf\sum_{i=1}^{\infty}|U_i|^s,
\end{equation}
其中下确界取遍 $E$ 的所有（可数）$\delta$-覆盖 $\{U_i\}$.容易验证 $\mathcal{H}^s_{\delta}$ 是 $\mathbb{R}^n$ 上的外测度.

为得到 $E$ 的 Hausdorff $s$ 维\newterm{hausdorff-outer-measure}{外测度}{Hausdorff outer measure}，令 $\delta\to 0$，即
\begin{equation}\label{eq:1.13}
  \mathcal{H}^{s}(E)=\lim_{\delta\to 0}\mathcal{H}^{s}_{\delta}(E)=\sup_{\delta>0}\mathcal{H}^{s}_{\delta}(E).
\end{equation}
该极限存在（但可能为无穷），因为当 $\delta$ 减小时 $\mathcal{H}^s_{\delta}$ 增大.易见 $\mathcal{H}^s$ 是外测度，而且它还是度量外测度.事实上，若 $\delta$ 小于正分离的集合 $E$ 与 $F$ 之间的距离，则 $E\cup F$ 的 $\delta$-覆盖中的任何集合都不能同时与 $E$ 和 $F$ 相交，于是
\[
  \mathcal{H}^s_{\delta}(E\cup F)=\mathcal{H}^s_{\delta}(E)+\mathcal{H}^s_{\delta}(F),
\]
从而对 $\mathcal{H}^s$ 也有类似的等式.$\mathcal{H}^s$ 在 $\mathcal{H}^s$-可测集构成的 $\sigma$-域上的限制称为 \newterm{Hausdorff measure}{Hausdorff $s$ 维测度}{Hausdorff measure}；由 \cref{thm:1.5}，该 $\sigma$-域包含 Borel 集（事实上还包含 Souslin 集）.

注意，若把 \eqref{eq:1.12} 中的下确界改为取遍 $E$ 的由凸集构成的 $\delta$-覆盖而非任意集合的 $\delta$-覆盖，则得到 Hausdorff 测度的一个等价定义，因为任何集合都含于一个同直径的凸集之中.类似地，取开集或取闭集构成的 $\delta$-覆盖有时也很方便.在各种情形下，尽管对 $\delta>0$ 所得的 $\mathcal{H}^s_{\delta}$ 的值可能不同，但极限 $\mathcal{H}^s$ 的值是一样的，见 \textcite{davies1956}.（然而，若下确界取遍由球构成的 $\delta$-覆盖，则会得到一个不同的测度；\textcite{besicovitch1928a} 第 3 章把这类``\newterm{spherical-hausdorff-measure}{球面 Hausdorff 测度}{Hausdorff measure, spherical}''与 Hausdorff 测度作了比较.）

对任意 $E$，显然当 $s$ 从 $0$ 增到 $\infty$ 时 $\mathcal{H}^s(E)$ 不增.此外，若 $s<t$，则
\[
  \mathcal{H}^s_{\delta}(E)\geq\delta^{s-t}\,\mathcal{H}^t_{\delta}(E),
\]
由此可得：若 $\mathcal{H}^t(E)$ 为正，则 $\mathcal{H}^s(E)$ 为无穷.于是存在唯一的值 $\dim E$，称为 $E$ 的 \newterm{Hausdorff dimension}{Hausdorff 维数}{Hausdorff dimension}，使得
\begin{equation}\label{eq:1.14}
  \mathcal{H}^{s}(E)=\infty\ \ \text{若}\ 0\leq s<\dim E,
  \qquad
  \mathcal{H}^{s}(E)=0\ \ \text{若}\ \dim E<s<\infty.
\end{equation}

若 $C$ 是 $\mathbb{R}^n$ 中边长为 $1$ 的方体，则用显然的方式把 $C$ 分成 $k^n$ 个边长为 $1/k$ 的小方体即可看出：若 $\delta\geq k^{-1}n^{\frac{1}{2}}$，则
\[
  \mathcal{H}^n_{\delta}(C)\leq k^{n}\bigl(k^{-1}n^{\frac{1}{2}}\bigr)^{n}\leq n^{\frac{1}{2}n},
\]
从而 $\mathcal{H}^n(C)<\infty$.于是当 $s>n$ 时 $\mathcal{H}^s(C)=0$，并且由于 $\mathbb{R}^n$ 可表示为可数多个这类方体之并，$\mathcal{H}^s(\mathbb{R}^n)=0$.由此可知对任意 $E\subset\mathbb{R}^n$ 有 $0\leq\dim E\leq n$.同样显然，若 $E\subset E'$，则 $\dim E\leq\dim E'$.

若 $\mathcal{H}^s$-可测集 $E\subset\mathbb{R}^n$ 满足 $0<\mathcal{H}^s(E)<\infty$，则称 $E$ 为一个 \newterm{s-set}{$s$-集}{$s$-set}；$1$-集有时也称为\newterm{linearly-measurable-set}{线性可测集}{linearly measurable set}.显然，$s$-集的 Hausdorff 维数等于 $s$，但重要的是要认识到，$s$-集远比 Hausdorff 维数为 $s$ 的可测集特殊.事实上，\textcite{besicovitch1942} 证明了任何集合都可以表示为连续统多个同维数集合的不交并.本书的大部分内容致力于研究 $s$-集的几何性质.

把 \eqref{eq:1.12} 中的 $|U_i|^s$ 换成 $h(|U_i|)$（其中 $h$ 是某个右连续的递增正函数），便得到 Hausdorff 测度定义的推广.我们的许多结果对这些更一般的测度都有直接的类似物，不过有时以代数上的简洁性为代价.此时集合 $E$ 的 Hausdorff``维数''可以更确切地理解为对函数的一个分划：一类函数使 $E$ 的测度为零，另一类使其为无穷（见 \textcite{rogers1970}）.若把 $|U_i|^s$ 换成 $h(U_i)$，其中 $h$ 只是集合 $U_i$ 的函数，甚至也能得到某些进展（见 \textcites{davies1969,daviessamuels1974}）.

下面我们证明 $\mathcal{H}^s$ 是正则测度，以及一个有用的推论：$s$-集可以从下方用闭子集逼近.这一证明由 \textcite{besicovitch1938} 给出，他还在 \citeyear{besicovitch1954} 证明了 \cref{thm:1.6}(b) 中有限性条件的必要性.

\begin{theorem}\label{thm:1.6}
(a) 若 $E$ 是 $\mathbb{R}^n$ 的任意子集，则存在包含 $E$ 的 $G_{\delta}$-集 $G$ 使 $\mathcal{H}^s(G)=\mathcal{H}^s(E)$.特别地，$\mathcal{H}^s$ 是正则外测度.

(b) 任何 $\mathcal{H}^s$-可测且 $\mathcal{H}^s$-测度有限的集合都包含一个测度相等的 $F_{\sigma}$-集，从而包含一个与它之差的测度可任意小的闭集.
\end{theorem}

\begin{proof}
(a) 若 $\mathcal{H}^s(E)=\infty$，则 $\mathbb{R}^n$ 就是测度相等之开集，故设 $\mathcal{H}^s(E)<\infty$.对每个 $i=1,2,\ldots$，选取 $E$ 的开 $2/i$-覆盖 $\{U_{ij}\}_j$ 使
\[
  \sum_{j=1}^{\infty}|U_{ij}|^s<\mathcal{H}^s_{1/i}(E)+1/i.
\]
则 $E\subset G$，其中 $G=\bigcap_{i=1}^{\infty}\bigcup_{j=1}^{\infty}U_{ij}$ 是一个 $G_{\delta}$-集.由于 $\{U_{ij}\}_j$ 是 $G$ 的 $2/i$-覆盖，$\mathcal{H}^s_{2/i}(G)\leq\mathcal{H}^s_{1/i}(E)+1/i$，令 $i\to\infty$ 即得 $\mathcal{H}^s(E)=\mathcal{H}^s(G)$.由于 $G_{\delta}$-集是 $\mathcal{H}^s$-可测的，故 $\mathcal{H}^s$ 是正则外测度.

(b) 设 $E$ 为 $\mathcal{H}^s$-可测且 $\mathcal{H}^s(E)<\infty$.利用 (a)，可找到包含 $E$ 的开集 $O_1,O_2,\ldots$，使得
\[
  \mathcal{H}^s\Bigl(\bigcap_{i=1}^{\infty}O_i\setminus E\Bigr)
  =\mathcal{H}^s\Bigl(\bigcap_{i=1}^{\infty}O_i\Bigr)-\mathcal{H}^s(E)=0.
\]
$\mathbb{R}^n$ 的任何开子集都是 $F_{\sigma}$-集，故设对每个 $i$ 有 $O_i=\bigcup_{j=1}^{\infty}F_{ij}$，其中 $\{F_{ij}\}_j$ 是递增的闭集序列.于是由 $\mathcal{H}^s$ 的连续性，
\[
  \lim_{j\to\infty}\mathcal{H}^s(E\cap F_{ij})=\mathcal{H}^s(E\cap O_i)=\mathcal{H}^s(E).
\]
因此，给定 $\varepsilon>0$，可找到 $j_i$ 使
\[
  \mathcal{H}^s(E\setminus F_{ij_i})<2^{-i}\varepsilon\qquad(i=1,2,\ldots).
\]
若 $F$ 为闭集 $\bigcap_{i=1}^{\infty}F_{ij_i}$，则
\[
  \mathcal{H}^s(F)\geq\mathcal{H}^s(E\cap F)\geq\mathcal{H}^s(E)-\sum_{i=1}^{\infty}\mathcal{H}^s(E\setminus F_{ij_i})>\mathcal{H}^s(E)-\varepsilon.
\]
由于 $F\subset\bigcap_{i=1}^{\infty}O_i$，故 $\mathcal{H}^s(F\setminus E)\leq\mathcal{H}^s\bigl(\bigcap_{i=1}^{\infty}O_i\setminus E\bigr)=0$.由 (a)，$F\setminus E$ 含于某个满足 $\mathcal{H}^s(G)=0$ 的 $G_{\delta}$-集 $G$.于是 $F\setminus G$ 是含于 $E$ 的 $F_{\sigma}$-集，且
\[
  \mathcal{H}^s(F\setminus G)\geq\mathcal{H}^s(F)-\mathcal{H}^s(G)>\mathcal{H}^s(E)-\varepsilon.
\]
对 $\varepsilon=\frac{1}{2},\frac{1}{3},\frac{1}{4},\ldots$ 取这样的 $F_{\sigma}$-集的可数并，便得到含于 $E$ 且与 $E$ 测度相等的 $F_{\sigma}$-集.
\end{proof}

下一条引理断言：任何用足够小的集合构成的覆盖去估计一个集合的 Hausdorff 测度的尝试，所得结果都不会比实际的 Hausdorff 测度小太多.

\begin{lemma}\label{lem:1.7}
设 $E$ 为 $\mathcal{H}^s$-可测且 $\mathcal{H}^s(E)<\infty$，$\varepsilon>0$.则存在只依赖于 $E$ 与 $\varepsilon$ 的 $\rho>0$，使得对任何满足 $0<|U_i|\leq\rho$ 的 Borel 集族 $\{U_i\}_{i=1}^{\infty}$ 有
\[
  \mathcal{H}^{s}\Bigl(E\cap\bigcup_{i}U_i\Bigr)<\sum_{i}|U_i|^{s}+\varepsilon.
\]
\end{lemma}

\begin{proof}
由 $\mathcal{H}^s$ 作为 $\mathcal{H}^s_{\delta}$ 当 $\delta\to 0$ 时极限的定义，可选取 $\rho$ 使得对 $E$ 的任何 $\rho$-覆盖 $\{W_i\}$ 有
\begin{equation}\label{eq:1.15}
  \mathcal{H}^{s}(E)<\sum|W_i|^{s}+\tfrac{1}{2}\varepsilon.
\end{equation}
给定满足 $0<|U_i|\leq\rho$ 的 Borel 集 $\{U_i\}$，可找到 $E\setminus\bigcup_{i}U_i$ 的 $\rho$-覆盖 $\{V_i\}$ 使
\[
  \mathcal{H}^{s}\Bigl(E\setminus\bigcup_{i}U_i\Bigr)+\tfrac{1}{2}\varepsilon>\sum|V_i|^{s}.
\]
由于 $\{U_i\}\cup\{V_i\}$ 便是 $E$ 的一个 $\rho$-覆盖，故由 \eqref{eq:1.15}
\[
  \mathcal{H}^{s}(E)<\sum|U_i|^{s}+\sum|V_i|^{s}+\tfrac{1}{2}\varepsilon.
\]
从而
\[
  \mathcal{H}^{s}\Bigl(E\cap\bigcup_{i}U_i\Bigr)
  =\mathcal{H}^{s}(E)-\mathcal{H}^{s}\Bigl(E\setminus\bigcup_{i}U_i\Bigr)
  <\sum|U_i|^{s}+\sum|V_i|^{s}+\tfrac{1}{2}\varepsilon-\sum|V_i|^{s}+\tfrac{1}{2}\varepsilon
  =\sum|U_i|^{s}+\varepsilon. \qedhere
\]
\end{proof}

本节最后，我们证明一个关于``一致 Lipschitz''映射所联系的集合的测度的简单引理.

\begin{lemma}\label{lem:1.8}
设 $\psi:E\to F$ 为满射，且对常数 $c$ 满足
\[
  |\psi(x)-\psi(y)|\leq c\,|x-y|\qquad(x,y\in E).
\]
则 $\mathcal{H}^s(F)\leq c^s\mathcal{H}^s(E)$.
\end{lemma}

\begin{proof}
对每个 $i$ 有 $|\psi(U_i\cap E)|\leq c\,|U_i|$.于是若 $\{U_i\}$ 是 $E$ 的 $\delta$-覆盖，则 $\{\psi(U_i\cap E)\}$ 是 $F$ 的 $c\delta$-覆盖.又 $\sum_{i}|\psi(U_i\cap E)|^s\leq c^s\sum_{i}|U_i|^s$，故 $\mathcal{H}^s_{c\delta}(F)\leq c^s\mathcal{H}^s_{\delta}(E)$，令 $\delta\to 0$ 即得结论.
\end{proof}

\section{覆盖结果}

Vitali 覆盖定理是几何测度论中最有用的工具之一.给定覆盖某集合 $E$ 的一个``足够大''的集合族，Vitali 定理从中选出一个不交子族，使之覆盖 $E$ 的几乎全部.

我们在下一条引理中阐述 Vitali 结果之证明所体现的基本原理，不过是在一个简化的场景中.称一个集合族是\newterm{semidisjoint}{半不交的}{semidisjoint}，若族中没有任何成员包含于另一个不同的成员.

\begin{lemma}\label{lem:1.9}
设 $\mathcal{C}$ 是含于 $\mathbb{R}^n$ 的有界子集中的一族球.则可找到有限或可数的不交子族 $\{B_i\}$，使得
\begin{equation}\label{eq:1.16}
  \bigcup_{B\in\mathcal{C}}B\subset\bigcup_{i}B_i',
\end{equation}
其中 $B_i'$ 是与 $B_i$ 同心、半径为其 $5$ 倍的球.此外，可使族 $\{B_i'\}$ 为半不交的.
\end{lemma}

\begin{proof}
我们归纳地选取 $\{B_i\}$.令 $d_0=\sup\{|B|:B\in\mathcal{C}\}$，并从 $\mathcal{C}$ 中选取满足 $|B_1|\geq\frac{1}{2}d_0$ 的 $B_1$.设 $B_1,\ldots,B_m$ 已经选定，令 $d_m=\sup\{|B|:B\in\mathcal{C},\ B\ \text{与}\ \bigcup_{1}^{m}B_i\ \text{不交}\}$.若 $d_m=0$，过程终止.否则，从 $\mathcal{C}$ 中选取与 $\bigcup_{1}^{m}B_i$ 不交且 $|B_{m+1}|\geq\frac{1}{2}d_m$ 的 $B_{m+1}$.这些球当然不交；我们断言它们还有所需的覆盖性质.设 $B\in\mathcal{C}$，则要么 $B=B_i$ 对某个 $i$ 成立，要么 $B$ 与某个满足 $2|B_i|\geq|B|$ 的 $B_i$ 相交.若非如此，$B$ 本应先于第一个满足 $2|B_m|<|B|$ 的球 $B_m$ 被选出.（注意，把体积求和可知 $\sum|B_i|^n<\infty$，故若选出了无穷多个球，则当 $i\to\infty$ 时 $|B_i|\to 0$.）无论哪种情形都有 $B\subset B_i'$，即得 \eqref{eq:1.16}.为使 $\{B_i'\}$ 半不交，只需在 $B_i'\subset B_j'$（对某个 $j\neq i$）时把 $B_i$ 从子族中去掉即可，注意 $B_i'$ 只能含于有限多个 $B_j'$.
\end{proof}

称集合族 $\mathcal{V}$ 是 $E$ 的一个 \newterm{Vitali class}{Vitali 类}{Vitali class}，若对每个 $x\in E$ 与 $\delta>0$，存在 $U\in\mathcal{V}$ 使 $x\in U$ 且 $0<|U|\leq\delta$.


\begin{theorem}[\newterm{vitali-covering-theorem}{Vitali 覆盖定理}{Vitali covering theorem}]\label{thm:1.10}
(a) 设 $E$ 是 $\mathbb{R}^n$ 的 $\mathcal{H}^s$-可测子集，$\mathcal{V}$ 是 $E$ 的由闭集构成的 Vitali 类.则可从 $\mathcal{V}$ 中选出（有限或可数的）不交序列 $\{U_i\}$，使得或者 $\sum|U_i|^s=\infty$，或者 $\mathcal{H}^s\bigl(E\setminus\bigcup_{i}U_i\bigr)=0$.

(b) 若 $\mathcal{H}^s(E)<\infty$，则给定 $\varepsilon>0$，还可要求
\[
  \mathcal{H}^{s}(E)\leq\sum_{i}|U_i|^{s}+\varepsilon.
\]
\end{theorem}

\begin{proof}
固定 $\rho>0$；不妨设对所有 $U\in\mathcal{V}$ 有 $|U|\leq\rho$.归纳地选取 $\{U_i\}$.令 $U_1$ 为 $\mathcal{V}$ 的任一成员.设 $U_1,\ldots,U_m$ 已经选定，令 $d_m$ 为 $\mathcal{V}$ 中那些与 $U_1,\ldots,U_m$ 都不交的 $U$ 的 $|U|$ 的上确界.若 $d_m=0$，则 $E\subset\bigcup_{1}^{m}U_i$，于是 (a) 得证，过程终止.否则，令 $U_{m+1}$ 为 $\mathcal{V}$ 中与 $\bigcup_{1}^{m}U_i$ 不交且 $|U_{m+1}|\geq\frac{1}{2}d_m$ 的集合.

设过程无限进行下去且 $\sum|U_i|^s<\infty$.对每个 $i$，令 $B_i$ 为圆心在 $U_i$ 中、半径为 $3|U_i|$ 的球.我们断言对每个 $k>1$，
\begin{equation}\label{eq:1.17}
  E\setminus\bigcup_{1}^{k}U_i\subset\bigcup_{k+1}^{\infty}B_i.
\end{equation}
事实上，若 $x\in E\setminus\bigcup_{1}^{k}U_i$，则存在不与 $U_1,\ldots,U_k$ 相交且 $x\in U$ 的 $U\in\mathcal{V}$.由于 $|U_i|\to 0$，存在 $m$ 使 $|U|>2|U_m|$.由 $\{U_i\}$ 的选取方法，$U$ 必与某个满足 $k<i<m$ 且 $|U|\leq 2|U_i|$ 的 $U_i$ 相交.由初等几何 $U\subset B_i$，故得 \eqref{eq:1.17}.于是若 $\delta>0$，只要 $k$ 大到足以保证对 $i>k$ 有 $|B_i|\leq\delta$，便有
\[
  \mathcal{H}^{s}_{\delta}\Bigl(E\setminus\bigcup_{1}^{\infty}U_i\Bigr)
  \leq\mathcal{H}^{s}_{\delta}\Bigl(E\setminus\bigcup_{1}^{k}U_i\Bigr)
  \leq\sum_{k+1}^{\infty}|B_i|^{s}
  =6^{s}\sum_{k+1}^{\infty}|U_i|^{s}.
\]
从而对一切 $\delta>0$ 有 $\mathcal{H}^s_{\delta}\bigl(E\setminus\bigcup_{1}^{\infty}U_i\bigr)=0$，故 $\mathcal{H}^s\bigl(E\setminus\bigcup_{1}^{\infty}U_i\bigr)=0$，即 (a) 得证.

为证 (b)，不妨设证明开始时选取的 $\rho$ 就是 \cref{lem:1.7} 中对应于 $\varepsilon$ 与 $E$ 的那个数.若 $\sum|U_i|^s=\infty$，则 (b) 显然.否则，由 (a) 与 \cref{lem:1.7},
\[
  \mathcal{H}^{s}(E)
  =\mathcal{H}^{s}\Bigl(E\setminus\bigcup_{i}U_i\Bigr)+\mathcal{H}^{s}\Bigl(E\cap\bigcup_{i}U_i\Bigr)
  =0+\mathcal{H}^{s}\Bigl(E\cap\bigcup_{i}U_i\Bigr)
  <\sum|U_i|^{s}+\varepsilon. \qedhere
\]
\end{proof}

\newterm{covering-theorems}{覆盖定理}{covering theorems}本身也得到了广泛的研究，它在调和分析以及几何测度论中都具有特殊的重要性.de Guzmán（\citeyear{deguzman1975,deguzman1981}） 的两本书描述了针对非常一般的集合类与测度类的结果，书中还包含进一步的文献.覆盖原理的一种处理方法归功于 Besicovitch（\citeyear{besicovitch1945a,besicovitch1946,besicovitch1947}）；其中第一篇文章包含了对密度（如本书 2.2~节所描述的）的应用.

\section{勒贝格测度}

我们把 \newterm{lebesgue-measure}{勒贝格测度}{Lebesgue measure}作为 $\mathbb{R}^n$ 中\newterm{volume}{体积}{volume}的通常定义的推广而得到（``体积''在 $\mathbb{R}^1$ 中指长度，在 $\mathbb{R}^2$ 中指面积）.

设 $C$ 是 $\mathbb{R}^n$ 中如下形式的\newterm{coordinate-block}{坐标块}{coordinate block}：
\[
  C=[a_1,b_1)\times[a_2,b_2)\times\cdots\times[a_n,b_n),
\]
其中对每个 $i$ 有 $a_i<b_i$.按显然的方式定义 $C$ 的体积为
\[
  V(C)=(b_1-a_1)(b_2-a_2)\cdots(b_n-a_n).
\]
设 $E\subset\mathbb{R}^n$，令
\begin{equation}\label{eq:1.18}
  \mathcal{L}^{n}(E)=\inf\sum_{i}V(C_i),
\end{equation}
其中下确界取遍 $E$ 的所有由块构成的序列 $\{C_i\}$ 的覆盖.易见 $\mathcal{L}^n$ 是 $\mathbb{R}^n$ 上的外测度，称为\newterm{lebesgue-outer-measure}{勒贝格外测度}{Lebesgue outer measure}.进一步，若 $E$ 是任一块，则 $\mathcal{L}^n(E)$ 与 $E$ 的体积一致；这只需把 \eqref{eq:1.18} 中的和用有限和逼近，再用包含 $C_i$ 之面的平面剖分 $E$ 即得.由于任何块 $C_i$ 都可以剖分成小子块而使 \eqref{eq:1.18} 中的和不变，故对任意 $\delta>0$，把下确界取遍 $E$ 的 $\delta$-覆盖即可.因而 $\mathcal{L}^n$ 是 $\mathbb{R}^n$ 上的度量外测度.$\mathcal{L}^n$ 在 $\mathcal{L}^n$-可测集（即勒贝格可测集，由 \cref{thm:1.5} 它包含 Borel 集）上的限制称为勒贝格 $n$ 维测度或\newterm{n-dim-lebesgue-measure}{体积}{n-dimensional Lebesgue measure}.

显然，$\mathbb{R}^1$ 上 $\mathcal{L}^1$ 与 $\mathcal{H}^1$ 的定义一致.可以预期，当 $n>1$ 时 $\mathbb{R}^n$ 上的外测度 $\mathcal{L}^n$ 与 $\mathcal{H}^n$ 也有联系；事实上它们只差一个常数倍.为说明这一点，我们需要下述著名的几何结果——``等直径不等式''（\newterm{isodiametric inequality}{等直径不等式}{isodiametric inequality}），即给定直径时体积最大的集合是球.利用对称化或其他方法的证明可在任何关于凸性的书中找到，例如 \textcite{eggleston1958}；另见\cref{exe:1.6}.

\begin{theorem}\label{thm:1.11}
直径为 $d$ 的闭凸集的 $n$ 维体积至多为 $\pi^{\frac{1}{2}n}\bigl(\tfrac{1}{2}d\bigr)^{n}/\bigl(\tfrac{1}{2}n\bigr)!$，即直径为 $d$ 的球的体积.
\end{theorem}

\begin{theorem}\label{thm:1.12}
若 $E\subset\mathbb{R}^n$，则 $\mathcal{L}^n(E)=c_n\mathcal{H}^n(E)$，其中 $c_n=\pi^{\frac{1}{2}n}/2^{n}\bigl(\tfrac{1}{2}n\bigr)!$.特别地，$c_1=1$，$c_2=\pi/4$.
\end{theorem}

\begin{proof}
给定 $\varepsilon>0$，可用闭凸集族 $\{U_i\}$ 覆盖 $E$ 使 $\sum|U_i|^n<\mathcal{H}^n(E)+\varepsilon$.由 \cref{thm:1.11}，$\mathcal{L}^n(U_i)\leq c_n|U_i|^n$，故
\[
  \mathcal{L}^n(E)\leq\sum\mathcal{L}^n(U_i)<c_n\mathcal{H}^n(E)+c_n\varepsilon,
\]
从而 $\mathcal{L}^n(E)\leq c_n\mathcal{H}^n(E)$.

反之，设 $\{C_i\}$ 为覆盖 $E$ 的坐标块族，满足
\begin{equation}\label{eq:1.19}
  \sum_{i}V(C_i)<\mathcal{L}^{n}(E)+\varepsilon.
\end{equation}
把这些块稍稍放大而保持该不等式，可设它们都是开块.对每个 $i$，含于 $C_i$ 中半径至多为 $\delta$ 的闭球构成 $C_i$ 的一个 Vitali 类.由 Vitali 覆盖定理（\cref{thm:1.10}(a)），存在 $C_i$ 中直径至多为 $\delta$ 的不交球 $\{B_{ij}\}_j$，使 $\mathcal{H}^n\bigl(C_i\setminus\bigcup_{j=1}^{\infty}B_{ij}\bigr)=0$，从而 $\mathcal{H}^n_{\delta}\bigl(C_i\setminus\bigcup_{j=1}^{\infty}B_{ij}\bigr)=0$.由于 $\mathcal{L}^n$ 是 Borel 测度，
\[
  \sum_{j=1}^{\infty}\mathcal{L}^n(B_{ij})=\mathcal{L}^n\Bigl(\bigcup_{j=1}^{\infty}B_{ij}\Bigr)\leq\mathcal{L}^n(C_i).
\]
于是
\[
\begin{aligned}
  \mathcal{H}^{n}_{\delta}(E)
    &\leq\sum_{i=1}^{\infty}\mathcal{H}^{n}_{\delta}(C_i)
     \leq\sum_{i=1}^{\infty}\sum_{j=1}^{\infty}\mathcal{H}^{n}_{\delta}(B_{ij})
      +\sum_{i=1}^{\infty}\mathcal{H}^{n}_{\delta}\Bigl(C_i\setminus\bigcup_{j=1}^{\infty}B_{ij}\Bigr)\\
    &\leq\sum_{i=1}^{\infty}\sum_{j=1}^{\infty}|B_{ij}|^{n}
     =\sum_{i=1}^{\infty}\sum_{j=1}^{\infty}c_n^{-1}\mathcal{L}^{n}(B_{ij})
     \leq c_n^{-1}\sum_{i=1}^{\infty}\mathcal{L}^{n}(C_i)
     <c_n^{-1}\mathcal{L}^{n}(E)+c_n^{-1}\varepsilon,
\end{aligned}
\]
其中末一步用了 \eqref{eq:1.19}.于是对一切 $\varepsilon$ 与 $\delta$ 有 $c_n\mathcal{H}^n_{\delta}(E)\leq\mathcal{L}^n(E)+\varepsilon$，从而 $c_n\mathcal{H}^n_{\delta}(E)\leq\mathcal{L}^n(E)$.
\end{proof}

勒贝格测度论中的经典结果之一是勒贝格密度定理.我们后面的许多工作都源于为 Hausdorff 测度建立类似定理的尝试.读者不妨给出一个证明，作为使用 Vitali 覆盖定理的练习；此外，该定理也是\cref{thm:2.2} 的简单推论.

\begin{theorem}[\newterm{lebesgue-density-theorem}{勒贝格密度定理}{Lebesgue density theorem}]\label{thm:1.13}
设 $E$ 是 $\mathbb{R}^n$ 的 $\mathcal{L}^n$-可测子集.则 $E$ 在 $x$ 处的\newterm{lebesgue-density}{勒贝格密度}{Lebesgue density}
\begin{equation}\label{eq:1.20}
  \lim_{r\to 0}\frac{\mathcal{L}^{n}\bigl(E\cap B_r(x)\bigr)}{\mathcal{L}^{n}\bigl(B_r(x)\bigr)}
\end{equation}
存在，且当 $x\in E$ 时等于 $1$，当 $x\notin E$ 时等于 $0$，除一个 $\mathcal{L}^n$-测度为 $0$ 的 $x$ 集合外.（$B_r(x)$ 表示以 $x$ 为心、$r$ 为半径的闭球；与通常一样，$r$ 取正值趋于 $0$.）
\end{theorem}

\section{Hausdorff 维数与测度的计算}

确定一个集合的 Hausdorff 维数往往是困难的，而求出甚至估计其 Hausdorff 测度则更为困难.在已被考虑过的情形中，通常都是下估计难以获得.本章最后分析某些集合的维数与测度（\newterm{hausdorff-measure-calculation}{Hausdorff 测度的计算}{Hausdorff measure, calculation of}）；更多的例子将在全书各处出现.读者将会看到，对一切 $s$ 与 $n$ 的值，$\mathbb{R}^n$ 中存在极其丰富多样的 $s$-集，因此这里要描述的一般理论有着广泛的适用性.

实数中最熟悉的具有非整数 Hausdorff 维数的集合是 \newterm{cantor-set}{Cantor 集}{Cantor set}.令 $E_0=[0,1]$，
\[
  E_1=[0,\tfrac{1}{3}]\cup[\tfrac{2}{3},1],\qquad
  E_2=[0,\tfrac{1}{9}]\cup[\tfrac{2}{9},\tfrac{1}{3}]\cup[\tfrac{2}{3},\tfrac{7}{9}]\cup[\tfrac{8}{9},1],
\]
等等，其中 $E_{j+1}$ 由 $E_j$ 中每个区间去掉（开）中间三分之一而得；见图~\ref{fig:1.1}.于是 $E_j$ 由 $2^j$ 个长度为 $3^{-j}$ 的区间组成.Cantor 集是完备（闭且自稠密的）集合 $E=\bigcap_{j=0}^{\infty}E_j$.（在这一构造中出现的闭区间之族构成一个``网''（\newterm{net}{网}{net}），即任何两个这样的区间要么不交，要么一个包含于另一个.集合之网的思想在本书中频繁出现.）等价地，除去可数多个点不计，$E$ 就是 $[0,1]$ 中三进展开不含数字 $1$ 的数的集合.我们将显式计算 $E$ 的 Hausdorff 维数与测度；这种基本类型的计算可推广到相当复杂的集合.

\begin{figure}[htbp]
  \centering
  \begin{tikzpicture}[xscale=6.2]
    \foreach \j/\lst in {0/{0},1/{0,2},2/{0,2,6,8},3/{0,2,6,8,18,20,24,26},4/{0,2,6,8,18,20,24,26,54,56,60,62,72,74,78,80}} {
      \pgfmathsetmacro{\ys}{-1.55*\j}
      \begin{scope}[yshift=\ys cm]
        \pgfmathtruncatemacro{\N}{3^\j}
        \foreach \i in \lst {
          \pgfmathsetmacro{\xa}{\i/\N}
          \pgfmathsetmacro{\xb}{(\i+1)/\N}
          \ifnum\j>2
            \draw[dashed] (\xa,0)--(\xb,0);
            \draw[dashed] (\xa,0.05)--(\xb,0.05);
          \else
            \draw (\xa,0)--(\xb,0);
            \draw (\xa,0.05)--(\xb,0.05);
          \fi
        }
        \node[right] at (1.05,-0.02) {$E_{\j}$};
      \end{scope}
    }
    \node[right] at (1.05,-7.85) {$\vdots$};
  \end{tikzpicture}
  \caption{Cantor 集}
  \label{fig:1.1}
\end{figure}

\begin{theorem}\label{thm:1.14}
Cantor 集 $E$ 的 Hausdorff 维数为 $s=\log 2/\log 3=0.6309\ldots$，而且 $\mathcal{H}^s(E)=1$.
\end{theorem}

\begin{proof}
由于 $E$ 可用构成 $E_j$ 的 $2^j$ 个长度为 $3^{-j}$ 的区间覆盖，立即得到
$\mathcal{H}^s_{3^{-j}}(E)\leq 2^{j}3^{-sj}=2^{j}2^{-j}=1$.
令 $j\to\infty$ 得 $\mathcal{H}^s(E)\leq 1$.

为证相反的不等式，我们证明：若 $\mathcal{I}$ 是任何覆盖 $E$ 的区间族，则
\begin{equation}\label{eq:1.21}
  1\leq\sum_{I\in\mathcal{I}}|I|^{s}.
\end{equation}
把每个区间稍稍放大并利用 $E$ 的紧性，只需对 $\mathcal{I}$ 为有限个闭区间构成的情形证明 \eqref{eq:1.21}.再作一步简化，可取每个 $I\in\mathcal{I}$ 为包含 $E$ 的构造中出现的一对网区间 $J$ 与 $J'$ 的最小区间.（$J$ 与 $J'$ 不必是同一个 $E_j$ 中的区间.）设 $J$ 与 $J'$ 是这样的最大区间，则 $I$ 由 $J$、随后是 $E$ 的补集中的某个区间 $K$、再随后是 $J'$ 拼成.由 $E_j$ 的构造可见
\begin{equation}\label{eq:1.22}
  |J|,\ |J'|\leq|K|.
\end{equation}
于是
\[
  |I|^{s}=\bigl(|J|+|K|+|J'|\bigr)^{s}
  \geq\bigl(\tfrac{3}{2}(|J|+|J'|)\bigr)^{s}
  =2\bigl(\tfrac{1}{2}|J|^{s}+\tfrac{1}{2}|J'|^{s}\bigr)
  \geq|J|^{s}+|J'|^{s},
\]
这里用了函数 $t^s$ 的凹性及 $3^s=2$ 的事实.于是把 $I$ 换成两个子区间 $J$ 与 $J'$ 不会使 \eqref{eq:1.21} 中的和增大.如此进行下去，经过有限多步后，便得到由长度为 $3^{-j}$ 的相等区间构成的 $E$ 的覆盖.这样的覆盖必包含 $E_j$ 的所有区间，故既然 \eqref{eq:1.21} 对这一覆盖成立，它对原来的覆盖 $\mathcal{I}$ 也成立.
\end{proof}

Cantor 集构造中所用的因子 $\frac{1}{3}$ 并无特殊之处.若取 $E_0$ 为单位区间，并通过从 $E_j$ 中每个区间的中心去掉比例 $1-2k$ 而得到 $E_{j+1}$，则用与上面类似的论证（把 \eqref{eq:1.22} 换成 $|J|,\ |J'|\leq|K|k/(1-2k)$）可以证明 $\mathcal{H}^{s}\bigl(\bigcap_{1}^{\infty}E_j\bigr)=1$，其中 $s=\log 2/\log(1/k)$.

我们可以在更高维中以类似方式构造不规则子集.例如，取 $E_0$ 为 $\mathbb{R}^2$ 中的单位正方形，去掉除四个角上边长为 $k$ 的正方形以外的一切，得到 $E_1$；如此继续，在第 $j$ 步把 $E_{j-1}$ 中每个正方形换成四个边长为 $k^{j}$ 的角上正方形而得到 $E_j$（构造的前几步见图~\ref{fig:1.2}）.同样的计算给出 $\mathcal{H}^{s}\bigl(\bigcap_{1}^{\infty}E_j\bigr)$ 的正的上下界，其中 $s=\log 4/\log(1/k)$.要求出此类情形中测度的精确值则需要更精细的论证，我们不作进一步讨论.

\begin{figure}[htbp]
  \centering
  \begin{tikzpicture}[scale=2.6]
    \fill[gray!35] (0,0) rectangle (1,1);
    \draw (0,0) rectangle (1,1);
    \node[above] at (0.5,1) {$1$};
    \node[below] at (0.5,0) {$E_0$};
    \begin{scope}[xshift=1.7cm]
      \fill[gray!35] (0,0) rectangle (0.4,0.4);
      \fill[gray!35] (0.6,0) rectangle (1,0.4);
      \fill[gray!35] (0,0.6) rectangle (0.4,1);
      \fill[gray!35] (0.6,0.6) rectangle (1,1);
      \draw (0,0) grid (1,1);
      \draw (0,0) rectangle (1,1);
      \node[above] at (0.5,1) {$k$};
      \node[below] at (0.5,0) {$E_1$};
    \end{scope}
    \begin{scope}[xshift=3.4cm]
      \foreach \qx in {0,0.6} {\foreach \qy in {0,0.6} {
        \foreach \ax in {0,0.24} {\foreach \ay in {0,0.24} {
          \fill[gray!35] (\qx+\ax,\qy+\ay) rectangle ++(0.16,0.16);
        }}
      }}
      \draw (0,0) grid (1,1);
      \draw (0,0) rectangle (1,1);
      \node[above] at (0.5,1) {$k^2$};
      \node[below] at (0.5,0) {$E_2$};
    \end{scope}
  \end{tikzpicture}
  \caption{平面上一个不规则集合之构造的前几个阶段}
  \label{fig:1.2}
\end{figure}

作为替代，我们描述实直线上 Cantor 构造的一个推广.设 $s$ 是严格介于 $0$ 与 $1$ 之间的数；所构造的集合将具有维数 $s$.设 $E_0$ 表示单位区间；通过对 $E_j$ 的每个区间 $I$ 指定 $E_{j+1}\cap I$，归纳地定义集合 $E_0\supset E_1\supset E_2\cdots$，每个都是有限个闭区间之并.设 $I$ 是这样的区间，令 $m\geq 2$ 为整数，并令 $J_1,J_2,\ldots,J_m$ 为 $I$ 的等长且等距的闭子区间，其长度由
\begin{equation}\label{eq:1.23}
  |J_i|^{s}=\frac{1}{m}|I|^{s}
\end{equation}
给出，且 $J_1$ 的左端点与 $I$ 的左端点重合，$J_m$ 的右端点与 $I$ 的右端点重合.于是
\begin{equation}\label{eq:1.24}
  m|J_i|+(m-1)d=|I|\qquad(1\leq i\leq m),
\end{equation}
其中 $d$ 是相邻两个区间 $J_i$ 之间的间隔.定义 $E_{j+1}$ 使 $E_{j+1}\cap I=\bigcup_{1}^{m}J_i$.注意，$m$ 的值在 $E_j$ 的不同区间上可以不同，因此 $E_j$ 可以包含多种不同长度的区间.

集合 $E=\bigcap_{j=0}^{\infty}E_j$ 是单位区间的一个完备无处稠密子集.下述分析将出现在 Davies 即将发表的一篇文章中.

\begin{theorem}\label{thm:1.15}
若 $E$ 是如上描述的集合，则 $\mathcal{H}^s(E)=1$.
\end{theorem}

\begin{proof}
$E$ 的构造中所用的区间，即某个 $E_j$ 的分支区间，称为\newterm{net interval}{网区间}{net interval}.对 $F\subset E$，令
\begin{equation}\label{eq:1.25}
  \mu(F)=\inf\sum_{I\in\mathcal{I}}|I|^{s},
\end{equation}
其中下确界取遍 $F$ 的所有由网区间之族 $\mathcal{I}$ 构成的覆盖.则 $\mu$ 是 $E$ 的子集上的外测度（事实上还是 Borel 测度）.注意，若坚持要求 $\mathcal{I}$ 对任意 $\delta>0$ 都是 $F$ 的 $\delta$-覆盖，$\mu$ 的值并不改变，因为利用 \eqref{eq:1.23} 总可以把一个网区间 $I$ 换成若干更小的网区间而不改变 \eqref{eq:1.25} 中的和.

设 $\mathcal{I}$ 是 $E$ 的由网区间构成的覆盖.为求 $\sum_{I\in\mathcal{I}}|I|^s$ 的下界，可设族 $\mathcal{I}$ 有限（因为每个网区间相对于紧集 $E$ 是开的），还可设 $\mathcal{I}$ 中的区间两两不交（利用网的性质可去掉含于其他区间之中的那些区间）.设 $J$ 是 $\mathcal{I}$ 中最短的区间之一；设 $J$ 是某个 $E_j$ 的分支区间.则 $J\subset I$，$I$ 为 $E_{j-1}$ 中的某个区间.由于 $\mathcal{I}$ 是 $E$ 的不交覆盖，$E_j\cap I$ 的所有其他区间必都属于 $\mathcal{I}$.若把这些区间换成单个区间 $I$，则由 \eqref{eq:1.23}，$\sum|I|^s$ 的值不变.如此继续，用较大的区间替换网区间之组而不改变和的值，最终到达单个区间 $[0,1]$.于是 $\sum_{I\in\mathcal{I}}|I|^s=|[0,1]|^s=1$，特别地
\[
  \mu(E)=1.
\]
以完全相同的方式可知，若 $J$ 是任何网区间，则
\begin{equation}\label{eq:1.26}
  \mu(J\cap E)=|J|^{s}.
\end{equation}
接下来证明对任意区间 $J$ 有
\begin{equation}\label{eq:1.27}
  \mu(J\cap E)\leq|J|^{s}.
\end{equation}
必要时收缩 $J$，只需在 $J\subset[0,1]$、$J$ 的端点属于 $E$，而且（通过逼近）与含于 $J$ 中的网区间的端点重合的假设下证明.设 $I$ 是包含 $J$ 的最小网区间；设 $I$ 是某个 $E_j$ 的区间.设 $J$ 与 $E_{j+1}\cap I$ 的分支区间中的 $J_q,J_{q+1},\ldots,J_r$ 相交，其中 $1\leq q<r\leq m$.（由 $I$ 的极小性，这样的区间至少有两个.）我们断言
\begin{equation}\label{eq:1.28}
  |J_q\cap J|^{s}+|J_{q+1}|^{s}+\cdots+|J_{r-1}|^{s}+|J_r\cap J|^{s}\leq|J|^{s}.
\end{equation}
若 $J_q\cap J$ 不是整个 $J_q$，或 $J_r\cap J$ 不是整个 $J_r$，则把 $J$ 稍微放大时，\eqref{eq:1.28} 左端比右端增长得快.因此只需对 $J$ 为包含 $J_q$ 与 $J_r$ 的最小区间的情形证明 \eqref{eq:1.28}.此时 \eqref{eq:1.28} 成为
\begin{equation}\label{eq:1.29}
  k|J_i|^{s}\leq|J|^{s}=\bigl(k|J_i|+(k-1)d\bigr)^{s},
\end{equation}
其中 $k=r-q+1$.由 \eqref{eq:1.23} 与 \eqref{eq:1.24}，当 $k=m$ 时上式成立；当 $k=1$ 时显然成立且两种情形都取等号.求导两次可见 \eqref{eq:1.29} 右端作为 $k$ 的函数是凹函数\footnote{原书此处印作 ``convex function of $k$''（$k$ 的凸函数）.由于 $0<s<1$，\eqref{eq:1.29} 右端作为 $k$ 的函数实为凹函数，且论证恰好需要凹性（凹函数在区间上的最小值在端点取得），故此处按数学实际译为``凹函数''.——译者注.}，故 \eqref{eq:1.29} 对 $1\leq k\leq m$ 成立，从而 \eqref{eq:1.28} 得证.

最后，若 $J_q\cap J$ 或 $J_r\cap J$ 不是单个网区间，则可重复上述过程，把 $J_q\cap J$ 与 $J_r\cap J$ 换成更小的网区间，得到一个与 \eqref{eq:1.28} 类似但涉及 $E_{j+2}$ 而非 $E_{j+1}$ 的区间的表达式.如此继续，最终发现 $|J|^s$ 至少等于覆盖 $J\cap E$ 且含于 $J$ 内的两两不交的网区间长度之 $s$ 次幂的和.于是对任意区间 $J$，由 \eqref{eq:1.26} 即得 \eqref{eq:1.27}.

由于把 \eqref{eq:1.25} 中下确界改为取遍任意 $\delta>0$ 的 $\delta$-覆盖 $\mathcal{I}$ 时 \eqref{eq:1.25} 仍成立，故 $\mathcal{H}^s(E)\leq\mu(E)$.另一方面，由 \eqref{eq:1.27}，
\[
  \mu(E)\leq\sum\mu(J_i\cap E)\leq\sum|J_i|^{s}
\]
对 $E$ 的任何覆盖 $\{J_i\}$ 成立，故 $\mu(E)\leq\mathcal{H}^s(E)$.综上所述 $\mathcal{H}^s(E)=\mu(E)=1$.
\end{proof}

在更高维中的类似构造涉及方体或立方体的嵌套序列而非区间.同样的方法可求出 Hausdorff 维数并估计相应的 Hausdorff 测度.

\cref{thm:1.15} 的基本方法也可用于求其他相关类型集合的维数.例如，若在 $E$ 的构造中每个 $I$ 内的区间 $J_1,\ldots,J_m$ 只是``近乎相等''或``近乎等距''，该方法经调整后仍可求出 $E$ 的维数.类似地，若由 $E_j$ 得到 $E_{j+1}$ 时方程 \eqref{eq:1.23} 与 \eqref{eq:1.24} 只在``$j\to\infty$ 的极限''下成立，仍有可能求出 $E$ 的维数.

另一个求集合维数的有用技巧是把它稍加``畸变''而得到一个维数已知的集合，然后应用\cref{lem:1.8}.读者可参看\cref{thm:8.15}(a)，那里给出了示例.

\textcite{eggleston1952} 求出了由交过程形成的非常一般的集合的 Hausdorff 维数；其结果被 \textcite{peyriere1977} 推广.近来 \textcite{daviesfast1978} 描述了一种有趣而有力的方法.其他相关构造见 \textcites{randolph1941,erdos1946,ravetz1954,besicovitchtaylor1954,beardon1965,best1942,ciglervolkmann1963,wegmann1971b}，其中最后三篇是同一作者更早工作的延续.估计 Hausdorff 测度的另一种方法将在 8.3~节中描述.

\begin{exercises}
\item 证明：若 $\mu$ 是集合的 $\sigma$-域 $\mathcal{M}$ 上的测度，$E_j\in\mathcal{M}\ (1\leq j<\infty)$，则只要 $\mu\bigl(\bigcup_{1}^{\infty}E_j\bigr)<\infty$，就有
$\mu\bigl(\varlimsup_{j\to\infty}E_j\bigr)\geq\varliminf_{j\to\infty}\mu(E_j)$.
\item 设 $\nu$ 是度量空间 $(X,d)$ 上的外测度，且每个 Borel 集都 $\nu$-可测.证明 $\nu$ 是度量外测度.
\item 证明 $\mathbb{R}^n$ 上的外测度 $\mathcal{H}^s$ 平移不变，即 $\mathcal{H}^s(x+E)=\mathcal{H}^s(E)$，其中 $x+E=\{x+y:y\in E\}$.进而证明：$x+E$ 是 $\mathcal{H}^s$-可测的当且仅当 $E$ 是 $\mathcal{H}^s$-可测的.类似地，证明 $\mathcal{H}^s(cE)=c^s\mathcal{H}^s(E)$，其中 $cE=\{cy:y\in E\}$.
\item 对一般测度 $\mu$ 证明下述形式的 Vitali 覆盖定理：设 $E$ 是 $\mathbb{R}^n$ 的 $\mu$-可测子集且 $\mu(E)<\infty$.若 $\mathcal{V}$ 是 $E$ 的（由可测集构成的）Vitali 类，则存在不交集合 $U_1,U_2,\ldots\in\mathcal{V}$ 使 $\mu\bigl(E\setminus\bigcup_{i}U_i\bigr)=0$.
\item 用 Vitali 覆盖定理证明勒贝格密度定理.（对每个 $\alpha<1$ 与 $\rho>0$，考虑球类
$\mathcal{V}=\{B_r(x):x\in E,\ r\leq\rho\ \text{且}\ \mathcal{L}^n(B_r(x)\cap E)\leq\alpha\mathcal{L}^n(B_r(x))\}$.）
\item 证明直径为 $d$ 的平面凸集 $U$ 的面积至多为 $\frac{1}{4}\pi d^2$.（方法之一：取 $U$ 边界上一点为极坐标原点，使 $U$ 的面积为 $\frac{1}{2}\int r(\phi)^2\,\mathrm{d}\phi$，并注意对每个 $\phi$ 有 $r(\phi)^2+r\bigl(\phi+\frac{1}{2}\pi\bigr)^2\leq d^2$.）
\item 用勒贝格密度定理推出 \newterm{steinhaus-theorem}{Steinhaus 的结果}{Steinhaus' theorem}：若 $E$ 是测度为正的实数的勒贝格可测集，则差集 $\{y-x:x,y\in E\}$ 包含区间 $(-h,h)$.更一般地证明：若 $E$ 与 $E'$ 都是具有正勒贝格测度的可测集，则 $\{y-x:x\in E,\ y\in E'\}$ 包含一个区间.
\item 设 $\mu$ 是 $\mathbb{R}^n$ 上的 Borel 测度，$E$ 是 $\mu$-可测集且 $0<\mu(E)<\infty$.证明：(a) 若对 $x\in E$ 有 $\varlimsup_{r\to 0}r^{-s}\mu\bigl(B_r(x)\cap E\bigr)<c<\infty$，则 $\mathcal{H}^s(E)>0$；(b) 若对 $x\in E$ 有 $\varlimsup_{r\to 0}r^{-s}\mu\bigl(B_r(x)\cap E\bigr)>c>0$，则 $\mathcal{H}^s(E)<\infty$.（(a) 用 Hausdorff 测度的定义；(b) 用\cref{exe:1.4} 中形式的 Vitali 覆盖定理.）
\item 设 $E$ 是 $0$ 与 $1$ 之间其十进展开不含奇数数字的数的集合.试给出 $E$ 的 Hausdorff 维数与测度你能得到的最好上下估计.（事实上 $E$ 是一个 $s$-集，其中 $s=\log 5/\log 10$.这个例子旨在说明求 Hausdorff 测度时可能出现的一些困难，它比 Cantor 集要更棘手一些.处理此类问题的一种方法将在 8.3~节中描述.）
\end{exercises}

% !TeX root = ../main.tex
% 第 2 章 基本密度性质 (对应原书 pp. 20--27)

\chapter{基本密度性质}\label{chap:2}

\section{引言}

回顾一下，称 $\mathbb{R}^n$ 的子集 $E$ 为一个 $s$-集（$0\leq s\leq n$），如果 $E$ 是 $\mathcal{H}^s$-可测的且 $0<\mathcal{H}^s(E)<\infty$.一般我们排除 $s=0$ 的情形，因为对 $s=0$ 常常需要不同于其他 $s$ 值的专门处理.不过，由于 $0$-集恰好就是有限点集，其性质是直截了当的.

接下来三章讨论 $s$-集的局部性质，特别是密度与切线的存在性问题.由于测度性质在可数并下保持，其中一些结果可以用于 $\mathcal{H}^s$-测度为\newterm{sigma-finite-measure}{$\sigma$-有限的 $\mathcal{H}^s$-可测集}{sigma-finite measure}（即作为 $s$-集的可数并而形成的集合）.$\mathcal{H}^s$-测度非 $\sigma$-有限的 $s$ 维集合很难把握，除非求出其有限测度的子集；这将在\cref{chap:5} 中讨论.

通常把 $s$ 视为固定的；在没有歧义时，``测度''、``可测''、\newterm{almost-all}{``几乎所有的''}{almost all}（即``除一个测度为零的集合外''）等术语均指测度 $\mathcal{H}^s$.

首先定义在本书发展中起主要作用的基本集合密度.这些密度是勒贝格密度 \eqref{eq:1.20} 的自然类似物，尽管它们的行为可能截然不同.这些密度刻画了一个集合的局部测度与``期望''测度之比.

设 $B_r(x)$ 表示以 $x$ 为心、$r$ 为半径的\newterm{ball}{闭球}{ball}，于是 $|B_r(x)|=2r$.称
\[
  \bar{D}^{s}(E,x)=\varlimsup_{r\to 0}\frac{\mathcal{H}^{s}\bigl(E\cap B_r(x)\bigr)}{(2r)^{s}}
\]
与
\[
  \underline{D}^{s}(E,x)=\varliminf_{r\to 0}\frac{\mathcal{H}^{s}\bigl(E\cap B_r(x)\bigr)}{(2r)^{s}}
\]
分别为 $s$-集 $E$ 在点 $x\in\mathbb{R}^n$ 处的\newterm{upper density}{上密度}{density, upper}与\newterm{lower density}{下密度}{density, lower}（有时也称为上、下球面密度或圆形密度）.若 $\bar{D}^s(E,x)=\underline{D}^s(E,x)$，则称 $E$ 在 $x$ 处的\newterm{density}{密度}{density}存在，并用 $D^s(E,x)$ 表示这个公共值.

若点 $x\in E$ 满足 $\underline{D}^s(E,x)=\bar{D}^s(E,x)=1$，则称 $x$ 为 $E$ 的\newterm{regular point}{正则点}{regular point}；否则称 $x$ 为\newterm{irregular point}{非正则点}{irregular point}.若 $s$-集 $E$ 的 $\mathcal{H}^s$-几乎所有的点都是正则点，则称 $E$ 是\newterm{regular}{正则的}{regular set}；若几乎所有的点都是非正则点，则称 $E$ 是\newterm{irregular}{非正则的}{irregular set}.

刻画正则 $s$-集并给出 $s$-集密度的界，是本书的两个重要目标.勒贝格密度定理在这里没有类似物：$s$-集在几乎所有的点处密度为 $1$ 一般并不成立.事实上，该学科的主要结果之一是：除非 $s$ 为整数，$s$-集不可能是正则的.然而，若 $s$ 是整数，$s$-集可分解为一个正则部分与一个非正则部分.非常粗略地说，正则 $s$-集看起来像 $\mathbb{R}^n$ 中某个 $s$ 维流形的可测子集，而非正则集则可能表现得像 $n$ 个 Cantor 型集合的 Descartes 积——适当选取这些集合使所得集合具有所需的维数.

虽然我们主要关心球面密度，但引入 $s$-集 $E$ 在 $x$ 处的\newterm{upper convex density}{上凸密度}{density, upper convex}也很方便，其定义为
\begin{equation}\label{eq:2.1}
  \bar{D}^{s}_{\mathrm{c}}(E,x)=\lim_{r\to 0}\Bigl\{\sup\frac{\mathcal{H}^{s}(E\cap U)}{|U|^{s}}\Bigr\},
\end{equation}
其中上确界取遍所有满足 $x\in U$ 且 $0<|U|\leq r$ 的凸集 $U$.任何集合都含于一个等直径的凸集之中，故这等价于取遍所有满足 $x\in U$ 且 $0<|U|\leq r$ 的集合 $U$ 取上确界.

由于 $B_r(x)$ 是凸的，且当 $x\in U$ 时有 $U\subset B_r(x)$，其中 $r=|U|$，故有关系
\begin{equation}\label{eq:2.2}
  2^{-s}\,\bar{D}^{s}_{\mathrm{c}}(E,x)\leq\bar{D}^{s}(E,x)\leq\bar{D}^{s}_{\mathrm{c}}(E,x).
\end{equation}

在证明关于密度的基本结果之前，陈述一些经常应用的一般性原则或许是有价值的.

第一，我们常常需要知道由度量或测度性质定义的点集关于适当的测度是可测的.在本学科的实践中，验证可测性几乎总是形式上的手续.所遇到的集合通常可以用已知的可测集经由 $\lim_{\rho\to 0}$、$\varlimsup_{\rho>0}E_\rho$ 之类的组合表达出来；随后在检验集合 $E_\rho$ 的定义时发现，它就等于对正有理数 $\rho$ 取的可数交 $\bigcap_{\rho\in\mathbf{Q}^+}E_\rho$.\cref{lem:2.1} 给出了可测性的若干示范，但在其后当例行验证已足够时，我们常常不再明确提及而默认所涉集合可测.

第二，许多证明涉及证明由某个性质（如密度条件）定义的点集测度为零.典型地，需要证明集合 $E=\{x:\psi(x)>0\}$ 的 Hausdorff 测度为零，其中 $\psi$ 是某个实值函数.标准的方法是证明对每个 $\alpha>0$，集合 $\{x:\psi(x)>\alpha\}$ 测度为零，并指出可数并 $E=\bigcup_{j=1}^{\infty}\{x:\psi(x)>1/j\}$ 亦必如此.于是只需证明：对一切 $\alpha,\varepsilon>0$ 有 $\mathcal{H}^s\{x:\psi(x)>\alpha\}<\varepsilon$.

第三，提一下``一致化''的思想.若要求一个性质在集合 $E$ 的几乎所有的点处成立，只需证明：对一切 $\varepsilon>0$，存在子集 $F\subset E$ 使 $E\setminus F$ 的测度小于 $\varepsilon$，且该性质在整个 $F$ 上成立.这样做的好处是 $F$ 常可取得比 $E$ 正则得多.例如，利用\cref{thm:1.6} 一般可取 $F$ 为闭集，而且还可选取 $F$ 使各种密度（举例来说）在 $F$ 上一致收敛.在本书中，我们一般用字母 $F$ 表示这个``工作集''，即剥去 $E$ 最剧烈的不规则性后所得的集合.很多时候，如果一个结果能在合理的拓扑与一致性假设下对某个集合证明成立，那么把它推广到完全一般的情形就纯粹是技术性的事了.

\section{初等密度界}

本节开始研究 $s$-集的基本密度性质，证明一些对一切 $s$ 值都成立的结果.这类结果首先由 Besicovitch（\citeyear{besicovitch1928a,besicovitch1938}） 对线性可测集给出，后一篇文章包含了改进的证明.\textcite{besicovitch1945a,marstrand1954a,wallin1969} 指出，把这些工作推广到 $\mathbb{R}^n$ 中任意 $s$ 与 $n$ 的 $s$-集并无困难.

首先验证密度作为 $x$ 的函数的可测性.有时谈\newterm{measurable-function}{函数的可测性}{measurable function}比谈由函数定义的集合的可测性更简单.称函数 $f:\mathbb{R}^n\to\mathbb{R}$ 为可测的（相应地，\newterm{borel-measurable-function}{Borel 可测的}{Borel-measurable function}、\newterm{upper-semicontinuous-function}{上半连续的}{upper semicontinuous function}），如果 $\{x:f(x)<c\}$ 对每个 $c$ 都是可测集（相应地，Borel 集、开集）；把 $<$ 换成 $\leq$、$>$ 或 $\geq$ 也得到可测性的等价定义.由连续性的开集定义可知，$\mathbb{R}^n$ 上的连续函数都是 Borel 可测的.

\begin{lemma}\label{lem:2.1}
设 $E$ 是一个 $s$-集.

(a) 对每个 $r$，$\mathcal{H}^s\bigl(E\cap B_r(x)\bigr)$ 是 $x$ 的上半连续函数，从而是 Borel 可测函数.

(b) $\underline{D}^s(E,x)$ 与 $\bar{D}^s(E,x)$ 是 $x$ 的 Borel 可测函数.
\end{lemma}

\begin{proof}
(a) 给定 $r,\alpha>0$，写
\[
  F=\{x:\mathcal{H}^{s}\bigl(E\cap B_r(x)\bigr)<\alpha\}.
\]
设 $x\in F$.当 $\varepsilon\searrow 0$ 时，$B_{r+\varepsilon}(x)$ 递减到 $B_r(x)$，故由 $\mathcal{H}^s$ 的上连续性，
\[
  \mathcal{H}^{s}\bigl(E\cap B_{r+\varepsilon}(x)\bigr)\searrow\mathcal{H}^{s}\bigl(E\cap B_r(x)\bigr).
\]
于是可找到 $\varepsilon$ 使 $\mathcal{H}^s\bigl(E\cap B_{r+\varepsilon}(x)\bigr)<\alpha$，从而当 $|y-x|\leq\varepsilon$ 时 $B_r(y)\subset B_{r+\varepsilon}(x)$ 且 $\mathcal{H}^s\bigl(E\cap B_r(y)\bigr)<\alpha$.故 $F$ 是 $\mathbb{R}^n$ 的开子集.这对一切 $\alpha$ 成立，因此 $\mathcal{H}^s\bigl(E\cap B_r(x)\bigr)$ 是 $x$ 的上半连续函数.

(b) 利用 (a)，
\[
  \{x:\mathcal{H}^{s}\bigl(E\cap B_r(x)\bigr)<\alpha(2r)^{s}\}
\]
是开集，于是给定 $\rho>0$，
\[
  F_{\rho}=\{x:\mathcal{H}^{s}\bigl(E\cap B_r(x)\bigr)<\alpha(2r)^{s}\ \text{对某个}\ r<\rho\}
\]
作为这类集合之并也是开的.现在
\[
  \{x:\underline{D}^{s}(E,x)<\alpha\}=\bigcap_{\rho>0}F_{\rho};
\]
由于 $F_\rho$ 随 $\rho$ 减小而增大，这个交可以只对正有理数 $\rho$ 取.于是 $\{x:\underline{D}^s(E,x)<\alpha\}$ 是 $G_{\delta}$-集，从而对每个 $\alpha$ 都是 Borel 集，即 $\underline{D}^s(E,x)$ 是 $x$ 的 Borel 可测函数.类似的论证建立 $\bar{D}^s(E,x)$ 的可测性.
\end{proof}

这条引理使我们能够断言：对任何 $s$-集 $E$，形如 $\{x:D^s(E,x)>\alpha\}$ 的集合都是 $\mathcal{H}^s$-可测的.证明的另一个推论是：形如
\[
  \{x:\mathcal{H}^{s}\bigl(E\cap B_r(x)\bigr)<\alpha(2r)^{s}\ \text{对某个}\ r\leq\rho\}
\]
的集合是 $\mathcal{H}^s$-可测的（作为开集）；稍作变化便可用``对每个''替换``对某个''.

我们先对凸密度得到上密度的界，然后用 \eqref{eq:2.2} 推出圆形密度的相应结果.下述定理当 $E$ 为紧集时是显然的，但它成立得更广.

\begin{theorem}\label{thm:2.2}
若 $E$ 是 $\mathbb{R}^n$ 中的 $s$-集，则对 $\mathcal{H}^s$-几乎所有的 $x\notin E$ 有 $\bar{D}^s_{\mathrm{c}}(E,x)=0$.
\end{theorem}

\begin{proof}
固定 $\alpha>0$；我们证明可测集
\[
  F=\{x\notin E:\bar{D}^{s}_{\mathrm{c}}(E,x)>\alpha\}
\]
测度为零.由 $\mathcal{H}^s$ 的正则性，给定 $\delta>0$，可找到闭集 $E_1\subset E$ 使 $\mathcal{H}^s(E\setminus E_1)<\delta$.对 $\rho>0$ 令
\begin{equation}\label{eq:2.3}
  \mathcal{V}=\bigl\{U:\ U\ \text{为闭凸集},\ 0<|U|\leq\rho,\ U\cap E_1=\emptyset
  \ \text{且}\ \mathcal{H}^{s}(E\cap U)>\alpha|U|^{s}\bigr\}.
\end{equation}
利用 \eqref{eq:2.1} 与 $E_1$ 的闭性，$\mathcal{V}$ 是 $F$ 的由闭集构成的 Vitali 类，于是可用 Vitali 定理（\cref{thm:1.10}(a)）找到 $\mathcal{V}$ 中的不交序列 $\{U_i\}$，使或者 $\sum|U_i|^s=\infty$，或者 $\mathcal{H}^s\bigl(F\setminus\bigcup U_i\bigr)=0$.

但由 \eqref{eq:2.3}，
\[
  \sum|U_i|^{s}
  <\frac{1}{\alpha}\sum\mathcal{H}^{s}(E\cap U_i)
  =\frac{1}{\alpha}\mathcal{H}^{s}\Bigl(E\cap\bigcup_{i}U_i\Bigr)
  \leq\frac{1}{\alpha}\mathcal{H}^{s}(E\setminus E_1)
  <\frac{\delta}{\alpha}<\infty,
\]
因为 $\{U_i\}$ 两两不交且都与 $E_1$ 不交.于是 $\mathcal{H}^s\bigl(F\setminus\bigcup U_i\bigr)=0$，从而
\[
  \mathcal{H}^{s}_{\rho}(F)
  \leq\mathcal{H}^{s}_{\rho}\Bigl(F\setminus\bigcup U_i\Bigr)+\mathcal{H}^{s}_{\rho}\Bigl(F\cap\bigcup U_i\Bigr)
  \leq\mathcal{H}^{s}\Bigl(F\setminus\bigcup U_i\Bigr)+\sum|U_i|^{s}
  <0+\frac{\delta}{\alpha}.
\]
这对一切 $\delta>0$ 与一切 $\rho>0$ 成立，故 $\mathcal{H}^s(F)=0$，如所需.
\end{proof}

\begin{theorem}\label{thm:2.3}
若 $E$ 是 $\mathbb{R}^n$ 中的 $s$-集，则在 $\mathcal{H}^s$-几乎所有的 $x\in E$ 处 $\bar{D}^s_{\mathrm{c}}(E,x)=1$.
\end{theorem}

\begin{proof}
(a) 我们用 Hausdorff 测度的定义证明在 $E$ 中几乎处处有 $\bar{D}^s_{\mathrm{c}}(E,x)\geq 1$.取 $\alpha<1$ 与 $\rho>0$，令
\begin{equation}\label{eq:2.4}
  F=\bigl\{x\in E:\mathcal{H}^{s}(E\cap U)<\alpha|U|^{s}\ \text{对一切满足}\ x\in U,\ |U|\leq\rho\ \text{的凸集}\ U\bigr\}.
\end{equation}
则 $F$ 是 $E$ 的 Borel 子集.对任意 $\varepsilon>0$，可找到 $F$ 的由凸集构成的 $\rho$-覆盖 $\{U_i\}$ 使
\[
  \sum|U_i|^{s}<\mathcal{H}^{s}(F)+\varepsilon.
\]
于是，设每个 $U_i$ 都含有 $F$ 的某个点，并利用 \eqref{eq:2.4}，
\[
  \mathcal{H}^{s}(F)
  \leq\sum_{i}\mathcal{H}^{s}(F\cap U_i)
  \leq\sum_{i}\mathcal{H}^{s}(E\cap U_i)
  <\alpha\sum_{i}|U_i|^{s}
  <\alpha\bigl(\mathcal{H}^{s}(F)+\varepsilon\bigr).
\]
由于 $\alpha<1$ 且外面的不等式对一切 $\varepsilon>0$ 成立，故 $\mathcal{H}^s(F)=0$.对任意 $\rho>0$ 都可定义这样的 $F$，故由定义，在 $E$ 中几乎处处有 $\bar{D}^s_{\mathrm{c}}(E,x)\geq\alpha$.这对一切 $\alpha<1$ 成立，于是在 $E$ 中几乎处处有 $\bar{D}^s_{\mathrm{c}}(E,x)\geq 1$.

(b) 我们用 Vitali 方法证明几乎处处有 $\bar{D}^s_{\mathrm{c}}(E,x)\leq 1$.给定 $\alpha>1$，令 $F=\{x\in E:\bar{D}^s_{\mathrm{c}}(E,x)>\alpha\}$，则 $F$ 是 $E$ 的可测子集.令 $F_0=\{x\in F:\bar{D}^s_{\mathrm{c}}(E\setminus F,x)=0\}$.则由\cref{thm:2.2} 有 $\mathcal{H}^s(F\setminus F_0)=0$.利用凸密度的定义，当 $x\in F_0$ 时
\[
  \bar{D}^{s}_{\mathrm{c}}(F,x)\geq\bar{D}^{s}_{\mathrm{c}}(E,x)-\bar{D}^{s}_{\mathrm{c}}(E\setminus F,x)>\alpha.
\]
于是
\begin{equation}\label{eq:2.5}
  \mathcal{V}=\{U:U\ \text{为闭凸集且}\ \mathcal{H}^{s}(F\cap U)>\alpha|U|^{s}\}
\end{equation}
是 $F_0$ 的一个 Vitali 类，故由\cref{thm:1.10}(b)，给定 $\varepsilon>0$，可找到 $\mathcal{V}$ 中的不交序列 $\{U_i\}_i$ 使 $\mathcal{H}^s(F_0)\leq\sum|U_i|^s+\varepsilon$.由 \eqref{eq:2.5}，
\[
  \mathcal{H}^{s}(F)=\mathcal{H}^{s}(F_0)
  <\frac{1}{\alpha}\sum\mathcal{H}^{s}(F\cap U_i)+\varepsilon
  \leq\frac{1}{\alpha}\mathcal{H}^{s}(F)+\varepsilon.
\]
该不等式对任意 $\varepsilon>0$ 成立，故当 $\alpha>1$ 时 $\mathcal{H}^s(F)=0$，如所需.
\end{proof}

\textcite{freilich1966} 与 \textcite{daviessamuels1974} 对出乎意料一般的 Hausdorff 型测度得到了与\cref{thm:2.2,thm:2.3} 类似的结果.

这两个定理对圆形密度（在我们的发展中更为重要）的类似物，用 \eqref{eq:2.2} 立即得到.

\begin{corollary}\label{cor:2.4}
若 $E$ 是 $\mathbb{R}^n$ 中的 $s$-集，则在 $E$ 外 $\mathcal{H}^s$-几乎所有的 $x$ 处 $D^s(E,x)=0$.
\end{corollary}

\begin{corollary}\label{cor:2.5}
若 $E$ 是 $\mathbb{R}^n$ 中的 $s$-集，则在几乎所有的 $x\in E$ 处
\[
  2^{-s}\leq\bar{D}^{s}(E,x)\leq 1.
\]
\end{corollary}

\cref{cor:2.4} 有若干重要推论.其一是：$s$-集的可测子集的密度，在该子集的几乎所有点处与原集合的密度相等.这个结果常被用于准备性的场合：于是当考察某个具有特定密度性质的子集时，我们可以弃置集合的其余部分，使该密度性质几乎处处成立.

\begin{corollary}\label{cor:2.6}
设 $F$ 是 $s$-集 $E$ 的可测子集.则在几乎所有的 $x\in F$ 处
$\underline{D}^s(F,x)=\underline{D}^s(E,x)$ 且 $\bar{D}^s(F,x)=\bar{D}^s(E,x)$.
\end{corollary}

\begin{proof}
记 $H=E\setminus F$.由\cref{cor:2.4}，在 $F$ 中几乎所有的 $x$ 处 $D^s(H,x)=0$.对这些 $x$，
\[
  \underline{D}^{s}(E,x)=\underline{D}^{s}(F,x)+D^{s}(H,x)=\underline{D}^{s}(F,x)
  \quad\text{且}\quad
  \bar{D}^{s}(E,x)=\bar{D}^{s}(F,x)+D^{s}(H,x)=\bar{D}^{s}(F,x). \qedhere
\]
\end{proof}

\begin{corollary}\label{cor:2.7}
设 $E=\bigcup_{j}E_j$ 是 $s$-集的可数不交并且 $\mathcal{H}^{s}(E)<\infty$.则对任意 $k$，在几乎所有的 $x\in E_k$ 处
\[
  \underline{D}^{s}(E_k,x)=\underline{D}^{s}(E,x)
  \quad\text{且}\quad
  \bar{D}^{s}(E_k,x)=\bar{D}^{s}(E,x).
\]
\end{corollary}

\begin{proof}
对 $E$ 的子集 $E_k$ 应用\cref{cor:2.6}.
\end{proof}

\begin{corollary}\label{cor:2.8}
设 $E$ 是 $s$-集.若 $E$ 是正则的（相应地，非正则的），则 $E$ 的任何测度为正的可测子集也是正则的（相应地，非正则的）.
\end{corollary}

\begin{proof}
由\cref{cor:2.6} 与正则性的定义立得.
\end{proof}

\begin{corollary}\label{cor:2.9}
正则集（相应地，非正则集）与可测集之交是正则集（相应地，非正则集）.正则集与非正则集之交测度为零.
\end{corollary}

\begin{proof}
前两个断言由\cref{cor:2.8} 得出.最后的断言成立，因为这样的交必既是正则集又是非正则集.
\end{proof}

下一条推论使我们能够独立地处理 $s$-集的正则部分与非正则部分.

\begin{corollary}[\newterm{decomposition-theorem}{分解定理}{decomposition theorem}]\label{cor:2.10}
若 $E$ 是 $s$-集，则 $E$ 的正则点之集是一个正则集，而 $E$ 的非正则点之集是一个非正则集.
\end{corollary}

\begin{proof}
由\cref{lem:2.1}，正则点之集与非正则点之集都可测，故由\cref{cor:2.6} 即得本推论.
\end{proof}

对 $\mathbb{R}^n$（$n\geq 2$）中的 $s$-集，我们还能得到上角密度的界.\newterm{density-angular}{角密度}{density, angular}由 \textcite{besicovitch1929,besicovitch1938} 引入，用于研究 $s$-集的切线性质.设 $\theta$ 为单位向量，$\phi$ 为一个角；令 $S(x,\theta,\phi)$ 为顶点在 $x$、轴沿方向 $\theta$ 的\newterm{cone}{闭单向无穷锥}{cone}，由那些使线段 $[x,y]$ 与 $\theta$ 的夹角至多为 $\phi$ 的点 $y$ 组成.记 $S_r(x,\theta,\phi)=B_r(x)\cap S(x,\theta,\phi)$ 为相应的半径 $r$ 的\newterm{spherical-sector}{球面扇形}{spherical sector}.角密度的定义与球面密度类似，只是把 $B_r(x)$ 换成\newterm{sector}{扇形}{sector} $S_r(x,\theta,\phi)$.于是
\[
  \bar{D}^{s}(E,x,\theta,\phi)=\varlimsup_{r\to 0}\frac{\mathcal{H}^{s}\bigl(E\cap S_r(x,\theta,\phi)\bigr)}{(2r)^{s}}
\]
与
\[
  \underline{D}^{s}(E,x,\theta,\phi)=\varliminf_{r\to 0}\frac{\mathcal{H}^{s}\bigl(E\cap S_r(x,\theta,\phi)\bigr)}{(2r)^{s}}
\]
分别为 $E$ 在 $x$ 处的\newterm{upper angular density}{上角密度}{density, upper angular}与\newterm{lower angular density}{下角密度}{density, lower angular}.一次例行验证（见\cref{lem:2.1}）即可建立角密度及其相关集合的一切所需的可测性.

比较球面扇形的直径与半径，由定义可得
\[
  \bar{D}^{s}(E,x,\theta,\phi)\leq 2^{-s}|S_1(x,\theta,\phi)|^{s}\,\bar{D}^{s}_{\mathrm{c}}(E,x).
\]
于是由\cref{thm:2.3} 立即推出：给定 $\theta$ 与 $\phi$，对几乎所有的 $x\in E$，
\[
  \bar{D}^{s}(E,x,\theta,\phi)\leq 2^{-s}|S_1(x,\theta,\phi)|^{s}.
\]
计算扇形的直径，上式成为
\[
  \bar{D}^{s}(E,x,\theta,\phi)\leq
  \begin{cases}
    2^{-s} & \bigl(0<\phi\leq\tfrac{1}{6}\pi\bigr)\\
    (\sin\phi)^{s} & \bigl(\tfrac{1}{6}\pi\leq\phi\leq\tfrac{1}{2}\pi\bigr)\\
    1 & \bigl(\tfrac{1}{2}\pi\leq\phi\leq\pi\bigr).
  \end{cases}
\]
当 $\phi\geq\frac{1}{2}\pi$ 时，通过修改\cref{thm:2.3} 证明的 (a) 部分，可得到上角密度的正的下界，见\cref{exe:2.3}.于是
\[
  \bar{D}^{s}(E,x,\theta,\phi)\geq
  \begin{cases}
    0 & \bigl(0<\phi<\tfrac{1}{2}\pi\bigr)\\
    2^{-s} & \bigl(\tfrac{1}{2}\pi\leq\phi\leq\pi\bigr).
  \end{cases}
\]
下密度的估计要难得到一些；我们将在后面结合切线性质在某些情形下予以考虑.

\textcite{morgan1935,gillis1935,besicovitch1938,dickinson1939} 给出了密度与角密度取极端值的一些集合的构造.

\begin{exercises}
\item 设 $E$ 是 Cantor 集，证明对一切 $x$ 有 $\underline{D}^s(E,x)\leq 2^{-s}$，其中 $s=\log 2/\log 3$.进而证明 $E$ 是非正则的.
\item 定义 $\mathbb{R}^n$ 中 $s$-集 $E$ 在 $x$ 处的上立方密度为
$\varlimsup_{r\to 0}\mathcal{H}^s\bigl(E\cap S_r(x)\bigr)/|S_r(x)|^s$，
其中 $S_r(x)$ 表示以 $x$ 为心、边平行于坐标轴、边长为 $r$ 的立方体.\footnote{原书此处印作 $\varlimsup_{r\to\infty}$.因所证结论为``在几乎所有的 $x\in E$ 处''的局部密度估计，且当 $r\to\infty$ 时左端对任何有限测度之集恒趋于 $0$（例如取 $E$ 为 Cantor 集），$r\to\infty$ 应为 $r\to 0$ 之排印误，故按 $r\to 0$ 翻译.——译者注.}
证明在几乎所有的 $x\in E$ 处，$E$ 的上立方密度介于 $2^{-s}n^{-s/2}$ 与 $1$ 之间.
\item 设 $E$ 是 $\mathbb{R}^n$（$n\geq 2$）中的 $s$-集，$\theta$ 是单位向量，证明对几乎所有的 $x\in E$ 有 $2^{-s}\leq\bar{D}^s\bigl(E,x,\theta,\tfrac{1}{2}\pi\bigr)$.（仿照\cref{thm:2.3}(a) 的证明；需要利用 $\mathcal{H}^s$ 的正则性获得一个闭集来工作.）进而证明当 $\phi\geq\frac{1}{2}\pi$ 时 $2^{-s}\leq\bar{D}^s(E,x,\theta,\phi)$.
\item 设 $\psi:\mathbb{R}^n\to\mathbb{R}^n$ 是 Jacobian 不取零值的连续可微变换.若 $E$ 是 $s$-集，证明对一切 $x$ 有 $\bar{D}^s\bigl(\psi(E),\psi(x)\bigr)=\bar{D}^s(E,x)$；对下密度也有类似结果.
\item 用\cref{cor:2.4} 与\cref{thm:1.12} 证明勒贝格密度定理（\cref{thm:1.13}）.
\end{exercises}

% !TeX root = ../main.tex
% 第 3 章 整数维数集合的结构 (对应原书 pp. 28--53)

\chapter{整数维数集合的结构}\label{chap:3}

\section{引言}

本章讨论 $s$ 为整数时 $\mathbb{R}^n$ 中 $s$-集的密度与切线结构.由\cref{cor:2.10} 知 $s$-集分裂为一个正则部分与一个非正则部分，而且我们发现这两类集合表现出显著不同的性质.我们的目标之一是把正则集刻画为可求长曲线或曲面的可数并的子集，从而把测度论的观念与描述性的拓扑观念联系起来.

我们详细叙述 $\mathbb{R}^2$ 中线性可测集（即 $1$-集）的理论.这一工作几乎完全归功于 Besicovitch（\citeyear{besicovitch1928a,besicovitch1938}），后一篇文章包含了一些改进的证明以及进一步的结果.Besicovitch 的证明大多难以超越——除了相对次要的改进，以及在表述上或许可以改进之处.确实，Besicovitch 所用的某些几何方法精巧到令人惊讶：它们当初居然能被想出来.\textcite{deguzman1981} 也描述了本章的部分工作.

\section{曲线与连续统}

正则 $1$-集与可求长曲线紧密相关.事实上，除去一个测度为零的集合不计，正则 $1$-集是可求长曲线的可数族的子集.本节致力于研究曲线，主要从拓扑的观点出发，并与具有有限线性测度的连续统相联系.我们在 $\mathbb{R}^n$ 中工作，因为这一理论并不比平面曲线的情形更复杂.

\newterm{curve}{曲线}{curve}（或 \newterm{jordan-curve}{Jordan 曲线}{Jordan curve}）$\Gamma$ 是连续单射 $\psi:[a,b]\to\mathbb{R}^n$ 的像，其中 $[a,b]\subset\mathbb{R}$ 为闭区间.任何曲线都是\newterm{continuum}{连续统}{continua}，即紧连通集.这是因为紧连通集的连续像既紧又连通.特别地，任何曲线都是 Borel 集，从而是 $\mathcal{H}^s$-可测的.此外，紧集之间的连续双射有连续的逆，故曲线也可以定义为闭区间在同胚意义下的像.

曲线 $\Gamma$ 的\newterm{length}{长度}{length}定义为
\begin{equation}\label{eq:3.1}
  \mathscr{L}(\Gamma)=\sup\sum_{i=1}^{m}|\psi(t_i)-\psi(t_{i-1})|,
\end{equation}
其中上确界取遍 $[a,b]$ 的一切分割 $a=t_0<t_1<\cdots<t_m=b$.若 $\mathscr{L}(\Gamma)<\infty$（即 $\psi$ 具有\newterm{bounded-variation}{有界变差}{bounded variation}），则称 $\Gamma$ 是\newterm{rectifiable}{可求长的}{rectifiable curve}.

注意我们的定义排除了\newterm{curve-self-intersecting}{自交曲线}{curve, self-intersecting}；下面的引理把它们涵盖在内.

\begin{lemma}\label{lem:3.1}
设 $\psi:[a,b]\to\mathbb{R}^n$ 为连续映射，$\psi(a)\neq\psi(b)$.则 $\psi[a,b]$ 包含一条连接 $\psi(a)$ 与 $\psi(b)$ 的曲线.
\end{lemma}

\begin{proof}
对 $[a,b]$ 的每个重点 $x$，令 $I_x$ 是满足 $\psi(t_1)=\psi(t_2)=x$ 的最大（闭）区间 $[t_1,t_2]$.令 $\mathcal{I}$ 表示不被其他此类区间包含的那些区间之族.则 $\mathcal{I}$ 由可数多个两两不交的真闭区间组成.于是可以构造连续满射 $f:[a,b]\to[0,1]$，使 $f(a)=0$，$f(b)=1$，并且当 $t_1\leq t_2$ 时 $f(t_1)\leq f(t_2)$，其中等号成立当且仅当 $t_1$ 与 $t_2$ 位于 $\mathcal{I}$ 的同一区间中.定义 $\psi_0:[0,1]\to\mathbb{R}^n$ 如下：若对某个 $I_x\in\mathcal{I}$ 有 $f^{-1}(u)=I_x$，则令 $\psi_0(u)=x$；否则令 $\psi_0(u)=\psi(f^{-1}(u))$.容易验证 $\psi_0$ 是连续单射，且 $\psi_0(0)=\psi(a)$，$\psi_0(1)=\psi(b)$.
\end{proof}

\eqref{eq:3.1} 中的和不会因添加分割点而减小.显然，若把 $\psi[a,b]$ 分成两条曲线 $\psi[a,c]$ 与 $\psi[c,b]$（$a<c<b$），则新曲线长度之和等于原曲线的长度，即 $\mathscr{L}$ 在有公共端点的曲线上是可加的.

\newterm{curve-arc-length}{用弧长参数化可求长曲线}{curve, parameterized by arc-length}总是可能的，即把 $\Gamma$ 表示为函数 $\psi_0:[0,\mathscr{L}(\Gamma)]\to\mathbb{R}^n$ 的像，使 $\psi_0[0,t]$ 的长度为 $t$.这可以由令 $\psi_0(t)$ 为满足 $\mathscr{L}(\psi[a,u])=t$ 的唯一一点 $\psi(u)$ 而实现.若 $\psi$ 用弧长表示可求长曲线 $\Gamma$，则由 \eqref{eq:3.1} 可见
\begin{equation}\label{eq:3.2}
  |\psi(t_1)-\psi(t_2)|\leq|t_1-t_2|.
\end{equation}
特别地，这蕴含 $\psi$ 是绝对连续函数.

关于曲线定义的进一步讨论见 \textcite{burkillburkill1970}（第 246 页）或 \textcite{pelling1977}.

\begin{lemma}\label{lem:3.2}
若 $\Gamma$ 是曲线，则 $\mathcal{H}^1(\Gamma)=\mathscr{L}(\Gamma)$.
\end{lemma}

\begin{proof}
设 $\Gamma$ 是连接 $z$ 与 $w$ 的曲线.若 $\operatorname{proj}$ 表示从 $\mathbb{R}^n$ 到过 $z$ 与 $w$ 的直线上的正交投影，则当 $x,y\in\mathbb{R}^n$ 时 $|\operatorname{proj}x-\operatorname{proj}y|\leq|x-y|$.由\cref{lem:1.8}，
\[
  \mathcal{H}^1(\Gamma)\geq\mathcal{H}^1(\operatorname{proj}\Gamma)\geq\mathcal{H}^1([z,w])=\mathscr{L}^1([z,w])=|z-w|,
\]
因为 $\operatorname{proj}\Gamma\supset[z,w]$.

现在设 $\Gamma$ 由 $\psi:[a,b]\to\mathbb{R}^n$ 定义.由上面的注记，对任意 $t$ 与 $u$ 有 $\mathcal{H}^1(\psi[t,u])\geq|\psi(t)-\psi(u)|$.于是若 $a=t_0<t_1<\cdots<t_m=b$ 是 $[a,b]$ 的任一分割，则
\[
  \sum_{i}|\psi(t_i)-\psi(t_{i-1})|\leq\sum_{i}\mathcal{H}^1\bigl(\psi[t_{i-1},t_i]\bigr)=\mathcal{H}^1(\Gamma),
\]
因为 $\Gamma$ 的弧 $\psi[t_{i-1},t_i]$ 除端点外互不相交.故 $\mathscr{L}(\Gamma)\leq\mathcal{H}^1(\Gamma)$.

最后，设 $\mathscr{L}(\Gamma)<\infty$，并设 $\psi$ 用弧长参数化 $\Gamma$.由于 $\psi$ 是从 $[0,\mathscr{L}(\Gamma)]$ 到 $\Gamma$ 的满足 \eqref{eq:3.2} 的满射，\cref{lem:1.8} 蕴含
$\mathcal{H}^1(\Gamma)\leq\mathcal{H}^1([0,\mathscr{L}(\Gamma)])=\mathscr{L}(\Gamma)$.
\end{proof}

这条引理的一个推论是：若 $\Gamma$ 是可求长曲线，则当 $s<1$ 时 $\mathcal{H}^s(\Gamma)$ 为无穷，当 $s>1$ 时为零（见 \eqref{eq:1.14}）.

\begin{corollary}\label{cor:3.3}
设 $\psi$ 用弧长参数化可求长曲线 $\Gamma$.若 $E$ 是 $[0,\mathscr{L}(\Gamma)]$ 的任意勒贝格可测子集，则 $\psi(E)$ 是 $\mathbb{R}^n$ 的 $\mathcal{H}^1$-可测子集，且 $\mathcal{H}^1(\psi(E))=\mathscr{L}^1(E)$.
\end{corollary}

\begin{proof}
由\cref{lem:3.2}，可加集函数 $\mathcal{H}^1(\psi(\cdot))$ 与 $\mathscr{L}^1(\cdot)$ 在闭区间上一致，故由通常的测度扩张过程（例如见 \textcite{kingmantaylor1966}），它们在 $[0,\mathscr{L}(\Gamma)]$ 的勒贝格可测子集上一致.
\end{proof}

这条推论使得直线上的勒贝格测度的所有常用结果都可以移植到曲线上.例如，由勒贝格密度定理（\cref{thm:1.13}）可得：若 $E$ 是曲线 $\Gamma$ 的 $\mathcal{H}^1$-可测子集，则在几乎所有的 $x\in E$ 处，当 $\mathcal{H}^1(I)\to 0$ 时有 $\mathcal{H}^1(E\cap I)/\mathcal{H}^1(I)\to 1$，其中 $I$ 是 $\Gamma$ 中包含 $x$ 的子弧.

我们不再进一步使用这条推论.

我们经常要用到下一条引理：

\begin{lemma}\label{lem:3.4}
设 $E$ 是包含 $x$ 与 $y$ 的连续统.若 $|x-y|=\rho$，则 $\mathcal{H}^1\bigl(E\cap B_\rho(x)\bigr)\geq\rho$.特别地，$\mathcal{H}^1(E)\geq|E|$.
\end{lemma}

\begin{proof}
令 $f:\mathbb{R}^n\to[0,\infty)$ 定义为 $f(z)=|z-x|$.则 $f$ 是满足
\begin{equation}\label{eq:3.3}
  |f(z)-f(w)|\leq|z-w|\qquad(z,w\in\mathbb{R}^n)
\end{equation}
的连续映射.集合 $f\bigl(E\cap B_\rho(x)\bigr)$ 包含区间 $[0,\rho]$；否则，对某个 $r$（$0<r<\rho$），$E=(E\cap B_r(x))\cup(E\setminus B_r(x))$ 便成为 $E$ 的由两个不交闭集构成的分解.对映射 $f$ 应用\cref{lem:1.8}，
\[
  \mathcal{H}^{1}\bigl(E\cap B_{\rho}(x)\bigr)
  \geq\mathcal{H}^{1}\bigl(f(E\cap B_{\rho}(x))\bigr)
  \geq\mathcal{H}^{1}([0,\rho])=\mathscr{L}^{1}([0,\rho])=\rho. \qedhere
\]
\end{proof}

现在可以推出可求长曲线的基本密度性质.

\begin{lemma}\label{lem:3.5}
可求长曲线是正则 $1$-集.
\end{lemma}

\begin{proof}
若 $\Gamma$ 可求长，则 $\mathscr{L}(\Gamma)<\infty$，故由\cref{lem:3.2}，$\mathcal{H}^1(\Gamma)<\infty$.由于 $\Gamma$ 至少含两个不同的点，\cref{lem:3.4} 蕴含 $\mathcal{H}^1(\Gamma)>0$，故 $\Gamma$ 是一个 $1$-集.

设 $x$ 是 $\Gamma$ 上非端点的一点，并设 $x$ 把 $\Gamma$ 分成两条可求长子曲线 $\Gamma_-$ 与 $\Gamma_+$.当 $\rho$ 充分小时，由\cref{lem:3.4} 有 $\mathcal{H}^1\bigl(\Gamma_-\cap B_\rho(x)\bigr)\geq\rho$ 与 $\mathcal{H}^1\bigl(\Gamma_+\cap B_\rho(x)\bigr)\geq\rho$，从而 $\mathcal{H}^1\bigl(\Gamma\cap B_\rho(x)\bigr)\geq 2\rho$.于是在 $\Gamma$ 上除端点外的一切点处 $D^1(\Gamma,x)\geq 1$.结合\cref{cor:2.5}，这说明在几乎所有的 $x\in\Gamma$ 处 $D^1(\Gamma,x)$ 存在且等于 $1$.
\end{proof}

接下来讨论可求长曲线的切线的存在性.称 $\mathbb{R}^n$ 中的 $s$-集 $E$ 在 $x$ 处沿方向 $\pm\theta$ \newterm{tangent}{有切线}{tangent}，如果 $\bar{D}^s(E,x)>0$，而且对每个角 $\phi>0$，
\begin{equation}\label{eq:3.4}
  \lim_{r\to 0}r^{-s}\mathcal{H}^{s}\bigl(E\cap(B_r(x)\setminus S_r(x,\theta,\phi)\setminus S_r(x,-\theta,\phi))\bigr)=0
\end{equation}
（沿方向 $\pm\theta$ 过 $x$ 的直线当然就是切线）.显然，$s$-集在它的每一点处至多有一条切线.

我们给出 \textcite{besicovitch1944} 关于可求长曲线在几乎所有点处有切线的初等证明.另一种途径是诉诸有界变差函数可微性的标准结果，见 \textcite{kingmantaylor1966}（9.1~节）.

\begin{lemma}\label{lem:3.6}
设 $\Gamma$ 是端点为 $x$ 与 $y$ 的可求长曲线，$\phi$ 为正角.设 $E$ 是 $\Gamma$ 上这样的点的集合：$x$ 属于 $\Gamma$ 的任意小的一对子弧，这对子弧所对的弦彼此的夹角大于 $2\phi$.则 $\mathcal{H}^1(E)\leq\bigl(\mathscr{L}(\Gamma)-|x-y|\bigr)/(1-\cos\phi)$.
\end{lemma}

\begin{proof}
设 $L$ 表示过 $x$ 与 $y$ 的直线，并设 $\mathcal{V}$ 为 $\Gamma$ 的这样的闭子弧之族：子弧所对的弦与 $L$ 的夹角大于 $\phi$.由引理的条件，$\mathcal{V}$ 是 $E$ 的 Vitali 类.于是利用覆盖定理（\cref{thm:1.10}(b)），对任意 $\varepsilon>0$，可找到 $\Gamma$ 的属于 $\mathcal{V}$ 的不交子弧的有限族 $\Gamma_1,\ldots,\Gamma_m$ 使
\begin{equation}\label{eq:3.5}
  \mathcal{H}^{1}(E)\leq\sum_{i=1}^{m}|\Gamma_i|+\varepsilon\leq\sum_{i=1}^{m}\mathscr{L}(\Gamma_i)+\varepsilon.
\end{equation}
若 $\Gamma_0',\Gamma_1',\ldots,\Gamma_m'$ 是（可能退化为空的）补弧，则正交投影到 $L$ 上并利用投影不增长度的平凡事实，可见
\[
  \cos\phi\sum_{i}\mathscr{L}(\Gamma_i)+\sum_{i}\mathscr{L}(\Gamma_i')\geq|x-y|.
\]
由于 $\sum_{i}\mathscr{L}(\Gamma_i)+\sum_{i}\mathscr{L}(\Gamma_i')=\mathscr{L}(\Gamma)$，用 \eqref{eq:3.5} 得
\[
  \mathcal{H}^{1}(E)-\varepsilon\leq\sum_{i}\mathscr{L}(\Gamma_i)\leq\bigl(\mathscr{L}(\Gamma)-|x-y|\bigr)/(1-\cos\phi).
\]
$\varepsilon$ 任意，证明完成.
\end{proof}

\begin{corollary}\label{cor:3.7}
若 $\phi>0$，$E$ 是可求长曲线 $\Gamma$ 上这样的点的集合：$x$ 属于 $\Gamma$ 的任意小的一对子弧，这对子弧所对的弦彼此的夹角大于 $2\phi$，则 $\mathcal{H}^1(E)=0$.
\end{corollary}

\begin{proof}
给定 $\varepsilon>0$，由 \eqref{eq:3.1} 可在 $\Gamma$ 上（按次序，且 $x_0$ 与 $x_m$ 为 $\Gamma$ 的端点）找到点 $x_0,x_1,\ldots,x_m$ 使
\[
  \mathscr{L}(\Gamma)<\sum_{i=1}^{m}|x_i-x_{i-1}|+\varepsilon.
\]
记 $\Gamma_i$ 为 $\Gamma$ 介于 $x_{i-1}$ 与 $x_i$ 之间的部分.依次对每个 $\Gamma_i$ 应用\cref{lem:3.6} 得
\begin{align*}
  \mathcal{H}^{1}(E)
  &=\sum_{i=1}^{m}\mathcal{H}^{1}(E\cap\Gamma_i)
  \leq\sum_{i=1}^{m}\bigl(\mathscr{L}(\Gamma_i)-|x_i-x_{i-1}|\bigr)/(1-\cos\phi)\\
  &=\bigl(\mathscr{L}(\Gamma)-\sum_{i=1}^{m}|x_i-x_{i-1}|\bigr)/(1-\cos\phi)
  <\varepsilon/(1-\cos\phi).
\end{align*}
故 $\mathcal{H}^1(E)=0$.
\end{proof}

\begin{theorem}\label{thm:3.8}
可求长曲线 $\Gamma$ 在它的几乎所有点处有切线.
\end{theorem}

\begin{proof}
由于连续统 $\Gamma$ 至少有两个点，\cref{lem:3.4} 蕴含在所有 $x\in\Gamma$ 处 $\bar{D}^s(E,x)\geq\frac{1}{2}$.

设 $\psi:[a,b]\to\mathbb{R}^n$ 是 $\Gamma$ 的一个定义函数.由\cref{cor:3.7}，对几乎所有的 $x\in\Gamma$ 可找到单位向量 $\theta$，使得给定 $\phi>0$，当 $|u-t|<\varepsilon$ 时（这里 $\psi(t)=x$）有
\[
  \psi(u)\in S(x,\theta,\phi)\cup S(x,-\theta,\phi).
\]
此外，若 $\rho$ 为正数，则当 $|u-t|\geq\varepsilon$ 时 $\psi(u)\notin B_\rho(x)$（否则可选取满足 $|u_i-t|\geq\varepsilon$ 且 $\psi(u_i)\to x$ 的数列 $\{u_i\}$；利用序列紧性，存在 $u$ 使 $|u-t|\geq\varepsilon$ 且 $\psi(u)=x=\psi(t)$，这就给出了 $\Gamma$ 的一个重点）.于是
\[
  \Gamma\cap\bigl(B_{\rho}(x)\setminus S_{\rho}(x,\theta,\phi)\setminus S_{\rho}(x,-\theta,\phi)\bigr)=\emptyset,
\]
从而 \eqref{eq:3.4} 作为平凡推论成立.
\end{proof}

曲线的切线常用可微性来理解.\cref{thm:3.8} 证明的一个微小变形表明：由函数 $\psi$ 定义的可求长曲线在 $\psi(t)$ 处有切线，只要导数 $\psi'(t)$ 存在且不为零.

Besicovitch（\citeyear{besicovitch1934a,besicovitch1956,besicovitch1957,besicovitch1960}） 还考虑了如何把切线的定义调整后用于具有无穷线性测度的曲线与集合.

含于可求长曲线的可数并中的 $1$-集称为 \newterm{Y-set}{$Y$-集}{$Y$-set}.

\cref{lem:3.5} 与 \cref{thm:3.8} 有针对 $Y$-集的对应物.

\begin{corollary}\label{cor:3.9}
$Y$-集是正则 $1$-集.
\end{corollary}

\begin{proof}
把\cref{cor:2.8} 应用于\cref{lem:3.5}：含于可求长曲线中的 $1$-集几乎处处密度为 $1$.由\cref{cor:2.7}，$Y$-集亦然.
\end{proof}

\begin{corollary}\label{cor:3.10}
$Y$-集在它的几乎所有点处有切线.
\end{corollary}

\begin{proof}
由\cref{cor:3.9}，$Y$-集几乎处处具有正的下密度.把\cref{cor:2.6} 应用于\cref{thm:3.8}：含于可求长曲线中的 $1$-集几乎处处满足 \eqref{eq:3.4}，故由\cref{cor:2.7}，$Y$-集在几乎所有点处有切线.
\end{proof}

最终我们将证明正则 $1$-集是一个 $Y$-集加上一个测度为零的集合，从而由\cref{cor:3.10} 可得任何正则 $1$-集几乎处处有切线.

具有下一条引理所描述性质的集合有时称为 \newterm{lipschitz-set}{Lipschitz 集}{Lipschitz set}.

\begin{lemma}\label{lem:3.11}
设 $E$ 是 $\mathbb{R}^n$ 的有界子集，使得当 $x,y\in E$ 时线段 $[x,y]$ 与固定直线 $L$ 的夹角至多为 $\phi<\frac{1}{2}\pi$.则 $E$ 含于一条可求长曲线之中.
\end{lemma}

\begin{proof}
不妨设 $E$ 是闭集.设 $\Pi(t)$ 为与 $L$ 垂直、距某原点距离为 $t$ 的超平面，并设 $a$ 与 $b$ 是使 $\Pi(t)$ 与 $E$ 相交的 $t$ 的极端值.于是对每个 $t$，$\Pi(t)$ 至多含 $E$ 的一个点；若存在就用 $\psi(t)$ 表示这个点，若不存在且 $a<t<b$，则令 $\psi(t)$ 为 $\Pi(t)$ 上位于连接两侧离 $\Pi(t)$ 最近的 $E$ 中的点的线段上的点.当 $t_1\neq t_2$ 时，线段 $[\psi(t_1),\psi(t_2)]$ 与 $L$ 的夹角至多为 $\phi$，故 $|\psi(t_1)-\psi(t_2)|\leq|t_1-t_2|/\cos\phi$.由 \eqref{eq:3.1}，$\psi$ 定义了一条可求长曲线.
\end{proof}

接下来我们证明，具有有限 $\mathcal{H}^1$-测度的连续统除去一个零测集之外是可求长曲线的可数族，从而是正则的.遵循 \textcite{besicovitch1938}，我们以弧连通作为中间阶段.回忆：称集合 $E$ 是\newterm{arcwise-connected}{弧连通的}{arcwise connected}，如果任给 $z,w\in E$，存在位于 $E$ 内连接 $z$ 与 $w$ 的曲线.紧的弧连通集显然是连续统.我们需要下述逆命题.

\begin{lemma}\label{lem:3.12}
满足 $\mathcal{H}^1(E)<\infty$ 的连续统 $E$ 是弧连通的.
\end{lemma}

\begin{proof}
这本质上是\cref{lem:3.4} 的推论.取 $z,w\in E$.由于 $E$ 连通，$E$ 是\newterm{chain-connected}{链连通的}{chain connected}，即对每个 $\varepsilon>0$ 可在 $E$ 中找到点链 $z=x_0,x_1,\ldots,x_m=w$ 使 $|x_i-x_{i-1}|\leq\varepsilon$（$1\leq i\leq m$）.（这可以通过证明与 $z$ 链连通的点之集在 $E$ 中既开又闭而得.）必要时删去链的中间点，可设当 $|i-j|\geq 2$ 时 $|x_i-x_j|>\varepsilon$.于是 $\mathbb{R}^n$ 中没有点落在多于两个球 $B_{\frac{1}{2}\varepsilon}(x_i)$ 之中，故设 $m\geq 2$ 并用\cref{lem:3.4} 得
\[
  2\mathcal{H}^{1}(E)\geq\sum_{i=0}^{m}\mathcal{H}^{1}\bigl(E\cap B_{\frac{1}{2}\varepsilon}(x_i)\bigr)\geq\frac{1}{2}m\varepsilon.
\]
于是，若 $\Gamma_\varepsilon$ 表示连接 $x_0,x_1,\ldots,x_m$ 所得的折线曲线（必然不自交），则
\begin{equation}\label{eq:3.6}
  \mathscr{L}(\Gamma_{\varepsilon})=\sum_{i=1}^{m}|x_i-x_{i-1}|\leq m\varepsilon\leq 4\mathcal{H}^{1}(E).
\end{equation}
对 $0\leq t\leq 1$，令 $\psi_\varepsilon(t)$ 为 $\Gamma_\varepsilon$ 上沿其长度比例 $t$ 处的点，使 $\Gamma_\varepsilon$ 介于 $z=\psi_\varepsilon(0)$ 与 $\psi_\varepsilon(t)$ 之间的部分长度为 $t\mathscr{L}(\Gamma_\varepsilon)$.于是当 $0\leq t_1<t_2\leq 1$ 时，由 \eqref{eq:3.6}，
\begin{equation}\label{eq:3.7}
  |\psi_{\varepsilon}(t_1)-\psi_{\varepsilon}(t_2)|\leq\mathscr{L}\bigl(\psi_{\varepsilon}[t_1,t_2]\bigr)\leq|t_2-t_1|\,4\mathcal{H}^{1}(E).
\end{equation}
曲线族 $\{\Gamma_\varepsilon\}_{\varepsilon\leq 1}$ 都位于 $\mathbb{R}^n$ 的一个有界子集之中，故由 $[0,1]$ 到 $\mathbb{R}^n$ 的函数族 $\{\psi_\varepsilon\}_{\varepsilon\leq 1}$ 一致有界，且由 \eqref{eq:3.7} 等度连续.由 Arzelà--Ascoli 定理（例如见 \textcite{dunfordschwartz1958}，第四部分 6.7 节），这样的函数族是序列紧的，即存在序列 $\varepsilon(j)\to 0$ 与连续函数 $\psi:[0,1]\to\mathbb{R}^n$ 使 $\psi_{\varepsilon(j)}$ 在 $[0,1]$ 上一致收敛于 $\psi$.令 $\psi[0,1]=\Gamma$.显然 $\psi(0)=z$，$\psi(1)=w$.又若 $x\in\Gamma$，则给定 $\delta>0$，可找到 $\varepsilon(j)<\frac{1}{2}\delta$ 使 $\Gamma_{\varepsilon(j)}$ 与球 $B_{\frac{1}{2}\delta}(x)$ 相交.于是 $B_\delta(x)$ 中有 $E$ 的某个点（即 $\Gamma_{\varepsilon(j)}$ 的一个顶点）.由于 $E$ 是闭的，我们得出 $x\in E$，故 $\Gamma\subset E$.若 $\psi$ 是单射，则 $\Gamma$ 是曲线.否则，由\cref{lem:3.1} 可把 $\Gamma$ 缩减为连接 $z$ 与 $w$ 的一条曲线.
\end{proof}

\begin{lemma}\label{lem:3.13}
任何满足 $\mathcal{H}^1(E)<\infty$ 的紧弧连通集 $E$ 由可求长曲线的可数并再加上一个 $\mathcal{H}^1$-测度为零的集合组成.
\end{lemma}

\begin{proof}
归纳地定义曲线序列 $\{\Gamma_j\}$.令 $\Gamma_1$ 为 $E$ 中连接 $E$ 的两个最远点的曲线.设曲线 $\Gamma_1,\Gamma_2,\ldots,\Gamma_k$ 已经定义.令 $x$ 为 $E$ 中距 $\bigcup_{1}^{k}\Gamma_j$ 最远（距离记为 $d_k$）的点（由于所涉集合是紧的，该最大值可以取得）.若 $d_k=0$，过程终止，结论得证.否则，令 $\Gamma_{k+1}$ 为 $E$ 中连接 $x$ 与 $\bigcup_{1}^{k}\Gamma_j$ 的曲线，且除一个端点外 $\Gamma_{k+1}$ 与 $\bigcup_{1}^{k}\Gamma_j$ 不交.由\cref{lem:3.4},
\begin{equation}\label{eq:3.8}
  \mathcal{H}^{1}(\Gamma_{k+1})\geq d_k.
\end{equation}
于是
\begin{equation}\label{eq:3.9}
  \sum_{j=1}^{\infty}d_j\leq\sum_{1}^{\infty}\mathcal{H}^{1}(\Gamma_j)\leq\mathcal{H}^{1}(E)<\infty,
\end{equation}
故 $d_j\to 0$.若 $x\in E$ 且 $x$ 到 $\bigcup_{1}^{\infty}\Gamma_j$ 的距离 $d$ 为正，则对某个 $j$ 有 $d_j<d$，而 $x$ 本应被选为 $\Gamma_{j+1}$ 的端点，矛盾.故 $E$ 是 $\bigcup_{1}^{\infty}\Gamma_j$ 的闭包.由\cref{lem:3.2}, $\mathscr{L}(\Gamma_j)=\mathcal{H}^1(\Gamma_j)\leq\mathcal{H}^1(E)<\infty$，故每个 $\Gamma_j$ 都是可求长曲线.最后，对每个 $k$，$\bigcup_{1}^{k}\Gamma_j$ 是闭集，故
\[
  \mathcal{V}=\Bigl\{B_r(x):x\in E\setminus\bigcup_{1}^{k}\Gamma_j
  \ \text{且}\ B_r(x)\cap\bigcup_{1}^{k}\Gamma_j=\emptyset\Bigr\}
\]
是 $E\setminus\bigcup_{1}^{k}\Gamma_j$ 的由球构成的 Vitali 类.设 $x$ 为某个 $B\in\mathcal{V}$ 的圆心.由于 $x$ 位于 $\bigcup_{k+1}^{\infty}\Gamma_j$ 的闭包中，有该并集的点任意接近 $x$.这些点必须经由 $\bigcup_{k+1}^{\infty}\Gamma_j$ 中的一列弧连接到 $\bigcup_{1}^{k}\Gamma_j$ 中必在 $B$ 之外的一点，故由\cref{lem:3.4},
\[
  \frac{1}{2}|B|\leq\mathcal{H}^{1}\Bigl(B\cap\bigcup_{k+1}^{\infty}\Gamma_j\Bigr)
\]
对任何 $B\in\mathcal{V}$ 成立.于是给定 $\varepsilon>0$，可用 Vitali 定理（\cref{thm:1.10}(b)）从 $\mathcal{V}$ 中选出不交球列 $\{B_i\}$ 使
\[
  \mathcal{H}^{1}\Bigl(E\setminus\bigcup_{1}^{k}\Gamma_j\Bigr)
  \leq\sum_{i}|B_i|+\varepsilon
  \leq 2\sum_{i}\mathcal{H}^{1}\Bigl(B_i\cap\bigcup_{k+1}^{\infty}\Gamma_j\Bigr)+\varepsilon
  \leq 2\mathcal{H}^{1}\Bigl(\bigcup_{k+1}^{\infty}\Gamma_j\Bigr)+\varepsilon.
\]
从而
\[
  \mathcal{H}^{1}\Bigl(E\setminus\bigcup_{1}^{k}\Gamma_j\Bigr)\leq 2\mathcal{H}^{1}\Bigl(\bigcup_{k+1}^{\infty}\Gamma_j\Bigr).
\]
由 \eqref{eq:3.9}，当 $k\to\infty$ 时右端趋于零，故 $\mathcal{H}^1\bigl(E\setminus\bigcup_{1}^{\infty}\Gamma_j\bigr)=0$，如所需.
\end{proof}

上面的引理看起来已由 \textcite{wazewski1927} 独立于 Besicovitch 证明.

\begin{theorem}\label{thm:3.14}
设 $E$ 是满足 $\mathcal{H}^1(E)<\infty$ 的连续统.则 $E$ 由可求长曲线的可数并再加上一个 $\mathcal{H}^1$-测度为零的集合组成.
\end{theorem}

\begin{proof}
合并\cref{lem:3.12,lem:3.13}.
\end{proof}

\begin{corollary}\label{cor:3.15}
设 $E$ 是满足 $\mathcal{H}^1(E)<\infty$ 的连续统.则 $E$ 是正则的，且在几乎所有点处有切线.
\end{corollary}

\begin{proof}
由\cref{thm:3.14}, $E$ 是一个 $Y$-集加上一个零测集，故结论由\cref{cor:3.9,cor:3.10} 立即得到.（零测集不影响密度与角密度的值.）
\end{proof}

人们还建立了若干其他关于把正则 $1$-集嵌入曲线的结果.例如，\cref{thm:3.14} 中的例外点集可以视为一条可能自交的可求长曲线的长度为 $0$ 的子集.类似地，任何有限测度的连续统 $E$ 可嵌入一条长度至多为 $2\mathcal{H}^1(E)$ 的（可自交的）可求长曲线（见\cref{exe:3.5}）.

本节最后证明一个最初归功于 Gołąb 的关于连续统线性测度\newterm{golab-semicontinuity}{半连续性的定理}{semicontinuity theorem, Gołąb's}.这一结果在几何测度论中相当重要，在几何问题中它被用来断言具有指定性质的测度最大的连续统或曲线的存在性.我们采用 \textcite{blaschke1916} 著名的选择定理；\textcite{rogers1970} 在更广泛的框架下讨论了该定理.

若 $E\subset\mathbb{R}^n$，则 $E$ 的 $\delta$-\newterm{delta-parallel-body}{平行体}{delta-parallel body}是距 $E$ 的距离不超过 $\delta$ 的点构成的闭集，即
\[
  [E]_{\delta}=\Bigl\{x\in\mathbb{R}^n:\inf_{y\in E}|x-y|\leq\delta\Bigr\}.
\]
\newterm{hausdorff-metric}{Hausdorff 度量}{Hausdorff metric} $\delta$ 定义在 $\mathbb{R}^n$ 的所有非空紧子集之族上：
\[
  \delta(E,F)=\inf\bigl\{\delta:E\subset[F]_{\delta}\ \text{且}\ F\subset[E]_{\delta}\bigr\}.
\]
验证 $\delta$ 是一个度量是简单的练习.

\begin{theorem}[\newterm{blaschke-selection-theorem}{Blaschke 选择定理}{Blaschke selection theorem}]\label{thm:3.16}
设 $\mathcal{C}$ 是位于 $\mathbb{R}^n$ 的有界部分 $B$ 中的非空紧集的一个无穷族.则存在 $\mathcal{C}$ 中互不相同的集合的序列 $\{E_j\}$，在 Hausdorff 度量下收敛于一个非空紧集 $E$.
\end{theorem}

\begin{proof}
首先从 $\mathcal{C}$ 中产生一个 Cauchy 序列.设 $\{E_{1,i}\}_i$ 是 $\mathcal{C}$ 中互不相同的集合的任意序列.对每个 $k>1$，如下定义 $\{E_{k-1,i}\}_i$ 的无穷子列 $\{E_{k,i}\}_i$：令 $\mathcal{B}_k$ 为覆盖 $B$ 的直径至多为 $1/k$ 的闭球的有限族.$E_{k-1,i}$ 中的每一个都与这些球的某个特定组合相交，故由抽屉原理（不同的组合只有有限多种），必存在 $\{E_{k-1,i}\}_i$ 的无穷子族 $\{E_{k,i}\}_i$，它们都恰好与 $\mathcal{B}_k$ 中同样的球相交.若 $F$ 是 $\mathcal{B}_k$ 中这一特定组合的球之并，则对一切 $i$ 有 $E_{k,i}\subset F\subset[E_{k,i}]_{1/k}$，于是 $\delta(E_{k,i},F)\leq 1/k$，从而对一切 $i,j$ 有 $\delta(E_{k,i},E_{k,j})\leq 2/k$.令 $E_i=E_{i,i}$，立即看到
\begin{equation}\label{eq:3.10}
  \delta(E_i,E_j)\leq 2/\min\{i,j\},
\end{equation}
故 $\{E_i\}_i$ 是 Cauchy 序列.

现在
\[
  E=\bigcap_{j=1}^{\infty}\overline{\bigcup_{i=j}^{\infty}E_i}
\]
是非空紧集，因为它是非空紧集的递减序列之交.（横线表示取闭包.）由 \eqref{eq:3.10}, $\overline{\bigcup_{i=j}^{\infty}E_i}\subset[E_j]_{2/j}$，故对一切 $j$ 有 $E\subset[E_j]_{2/j}$.另一方面，若 $x\in E_j$，则由 \eqref{eq:3.10}，当 $i\geq j$ 时 $x\in[E_i]_{2/j}$，于是当 $k\geq j$ 时 $x\in[\overline{\bigcup_{i=k}^{\infty}E_i}]_{2/j}$.选取 $y_k\in\overline{\bigcup_{i=k}^{\infty}E_i}$ 使 $|x-y_k|\leq 2/j$.由序列紧性，$\{y_k\}$ 的子列收敛于某个 $y\in\mathbb{R}^n$，且 $|x-y|\leq 2/j$.但 $y\in\bigcap_{k=1}^{\infty}\overline{\bigcup_{i=k}^{\infty}E_i}=E$，故 $x\in[E]_{2/j}$.于是 $E_j\subset[E]_{2/j}$，从而 $\delta(E,E_j)\leq 2/j$.故 $\{E_j\}$ 在 Hausdorff 度量下收敛于 $E$.
\end{proof}

注意，上述证明的后半部分本质上就是 $\delta$ 是完备度量的一个证明.

每对点都由唯一的可求长道路连接的连续统称为\newterm{tree}{树}{tree}.下一条引理使我们能够对树的序列证明半连续性定理，并通过逼近把结果推广到一般连续统.

\begin{lemma}\label{lem:3.17}
设 $F\subset\mathbb{R}^n$ 是满足 $\mathcal{H}^1(F)<\infty$ 的树.则给定 $\delta>0$，可写 $F=\bigcup_{i=1}^{k}F_i$，其中各 $F_i$ 是连续统，满足

(a) $\sum_{i=1}^{k}|F_i|\leq\sum_{i=1}^{k}\mathcal{H}^1(F_i)=\mathcal{H}^1(F)$;

(b) 对一切 $i$ 有 $|F_i|\leq\delta$;

(c) $k\leq 3\delta^{-1}\mathcal{H}^1(F)+1$.
\end{lemma}

\begin{proof}
若 $|F|\leq\delta$，可取 $k=1$，$F_1=F$，故设 $|F|>\delta$.设 $x$ 为 $F$ 的任一点，令 $m$ 为 $F$ 中的点到 $x$ 的道路距离的上确界；则由\cref{lem:3.4} 有 $m\leq\mathcal{H}^1(F)<\infty$，而 $m>\frac{1}{2}\delta$，否则将有 $|F|\leq\delta$.设 $y$ 为 $F$ 中与 $x$ 的道路距离大于 $m-\frac{1}{6}\delta$ 的一点，并设 $z$ 为连接 $y$ 与 $x$ 的唯一道路上与 $x$ 的道路距离为 $m-\frac{1}{2}\delta$ 的点.点 $z$ 确定了 $F$ 到两个子树 $F_1$ 与 $F'$ 的一个分割，$z$ 是它们仅有的公共点，其中 $F_1$ 由 $F$ 中到 $x$ 的连接经过 $z$ 的那些点组成，见图~\ref{fig:3.1}. $F_1$ 的每一点与 $z$ 的道路距离在 $\frac{1}{2}\delta$ 之内，故 $|F_1|\leq\delta$.由\cref{lem:3.4}, $|F_1|\leq\mathcal{H}^1(F_1)$，且 $\mathcal{H}^1(F_1)$ 大于 $y$ 到 $z$ 的道路距离，后者至少为 $\bigl(m-\frac{1}{6}\delta\bigr)-\bigl(m-\frac{1}{2}\delta\bigr)=\frac{1}{3}\delta$.

若 $|F'|>\delta$，则对树 $F'$ 重复这一过程以分离出子树 $F_2$，如此继续，直至剩下直径至多为 $\delta$ 的树 $F_k$.结论的 (a) 与 (b) 是直接的，而 (c) 由 (a) 得出，因为当 $1\leq i\leq k-1$ 时 $\frac{1}{3}\delta\leq\mathcal{H}^1(F_i)$.
\end{proof}

\begin{figure}[htbp]
  \centering
  \includegraphics[width=0.5\textwidth]{fig3-1.png}
  \caption{引理 3.17 证明中树 $F$ 的分解}
  \label{fig:3.1}
\end{figure}

\begin{theorem}\label{thm:3.18}
设 $\{E_j\}$ 是 $\mathbb{R}^n$ 中在 Hausdorff 度量下收敛于紧集 $E$ 的连续统序列.则 $E$ 是连续统，且
\[
  \mathcal{H}^{1}(E)\leq\varliminf_{j\to\infty}\mathcal{H}^{1}(E_j).
\]
\end{theorem}

\begin{proof}
若 $E$ 不连通，则有分解 $E=G_1\cup G_2$，其中 $G_1$ 与 $G_2$ 是非空不交闭集，且它们之间的距离严格为正.于是当 $j$ 充分大时，$E_j\subset[G_1]_{\frac{1}{4}\delta}\cup[G_2]_{\frac{1}{4}\delta}$，其中 $[G_1]_{\frac{1}{4}\delta}$ 与 $[G_2]_{\frac{1}{4}\delta}$ 是含有 $E_j$ 的点的不交平行体，这与 $E_j$ 的连通性矛盾.故 $E$ 是连续统.

为证不等式，不妨设对一切 $j$ 有 $\mathcal{H}^1(E_j)\leq c<\infty$.对每个 $j$ 选取 $E_j$ 的有限子集 $S_j$，使 $S_j$ 也在 Hausdorff 度量下收敛于 $E$.（令 $\delta_j\to 0$，取有限集 $S_j$ 使 $E_j$ 的每一点都落在 $S_j$ 中某点的 $\delta_j$ 邻域内.）由\cref{lem:3.12}, 各 $E_j$ 弧连通，故可找到树 $F_j$ 使 $S_j\subset F_j\subset E_j$.（逐次添加``分支''构造 $F_j$，把 $S_j$ 的每一点用 $E_j$ 中的弧连接到已构造的 $F_j$ 部分.）于是在 Hausdorff 度量下 $F_j\to E$.

固定 $\delta>0$.利用\cref{lem:3.17}, 可分解每棵树 $F_j$ 使 $F_j=\bigcup_{i=1}^{k}F_{ji}$，其中对一切 $i,j$ 有 $|F_{ji}|\leq\delta$，且对一切 $j$ 有 $\sum_{i=1}^{k}|F_{ji}|\leq\mathcal{H}^1(F_j)$.（当然可以对每个 $j$ 用同一个满足 $k\leq 3\delta^{-1}c+1$ 的 $k$.）对 $\{F_{ji}\}_j$ 依次对每个 $i$ 应用 Blaschke 选择定理（\cref{thm:3.16}），通过取子列并重新编号，可设 $\{F_{ji}\}_j$ 在 Hausdorff 度量下收敛于非空紧集 $H_i$（$1\leq i\leq k$）.当然对一切 $i$ 有 $|H_i|\leq\delta$，且 $E\subset\bigcup_{1}^{k}H_i$.于是
\[
  \mathcal{H}^{1}_{\delta}(E)\leq\sum_{i}|H_i|
  =\lim_{j\to\infty}\sum_{i}|F_{ji}|
  \leq\varliminf_{j\to\infty}\mathcal{H}^{1}(F_j)
  \leq\varliminf_{j\to\infty}\mathcal{H}^{1}(E_j),
\]
令 $\delta\to 0$ 即得结论.
\end{proof}

下面的选择结果在几何问题中常常有用.


\begin{corollary}\label{cor:3.19}
设 $\mathcal{C}$ 是 $\mathcal{H}^1$-测度至多为 $c$ 且位于 $\mathbb{R}^n$ 的有界部分中的连续统的无穷族.则存在 $\mathcal{C}$ 中互不相同的集合的序列，在 Hausdorff 度量下收敛于满足 $\mathcal{H}^1(E)\leq c$ 的连续统 $E$.
\end{corollary}

\begin{proof}
用\cref{thm:3.16} 从 $\mathcal{C}$ 中得到收敛于非空紧集 $E$ 的序列 $\{E_j\}$.\cref{thm:3.18} 保证 $E$ 是满足 $\mathcal{H}^1(E)\leq c$ 的连续统.
\end{proof}

\section{密度与正则 $1$-集的刻画}

现在把注意力限于平面的子集.本节导出不规则 $1$-集的下密度的一个本质的上界，从而刻画正则与不规则 $1$-集.

若一个 $1$-集与每条可求长曲线的交都具有 $\mathcal{H}^1$-测度零，则称它为 \newterm{Z-set}{$Z$-集}{$Z$-set}.

由\cref{cor:3.9} 可得不规则集是 $Z$-集.我们将证明逆命题也成立.

\begin{lemma}\label{lem:3.20}
$Z$-集与具有有限 $\mathcal{H}^1$-测度的连续统的交测度为零.
\end{lemma}

\begin{proof}
由\cref{thm:3.14} 与 $Z$-集的定义立得.
\end{proof}

本节的主要部分致力于证明 $Z$-集的上密度几乎处处严格小于 $1$，从而是非正则的.为此我们使用 \textcite{besicovitch1938}（第 15 节）的``圆对''思想.（\newterm{circle-pair}{圆对}{circle-pair}是由半径相等的两个圆盘构成的图形，每个圆盘的圆心位于另一个圆盘的圆周上.）证明复杂但相当精妙.我们遵循 \textcite{morserandolph1944}（第 9 节）中给出的 Besicovitch 证明的推广形式.首先需要一个关于从连续统中除去圆盘内部的拓扑结果.回忆：若圆盘族中没有任何成员包含于另一个成员，则称该族是半不交的.

\begin{lemma}\label{lem:3.21}
设 $E$ 是 $\mathbb{R}^2$ 中的连续统.设 $\{B_i\}_{1}^{\infty}$ 是含于 $E$ 中的闭圆盘的可数半不交族，且对任意 $d>0$，满足 $|B_i|\geq d$ 的 $i$ 只有有限多个.则若 $\Gamma_i$ 是 $B_i$ 的圆周，
\[
  F=\Bigl(E\setminus\bigcup_{i}B_i\Bigr)\cup\bigcup_{i}\Gamma_i
\]
是连续统.
\end{lemma}

\begin{proof}
注意 $F=(E\setminus\bigcup_{i}\operatorname{int}B_i)\cup\bigcup_{i}\Gamma_i$，故为证 $F$ 闭，只需证 $\overline{\bigcup_{i}\Gamma_i}\subset F$.设 $x\in\overline{\bigcup_{i}\Gamma_i}$ 而 $x\notin E\setminus\bigcup_{i}\operatorname{int}B_i$，则 $x\in\operatorname{int}B_k$，对某个 $k$.设 $d>0$ 为 $x$ 到 $B_k$ 圆周的距离.若 $\Gamma_j$ 与 $B_{\frac{1}{2}d}(x)$ 相交，则不仅 $B_j$ 与 $B_{\frac{1}{2}d}(x)$ 相交，而且它与 $\mathbb{R}^2\setminus B_k$ 也相交，因为 $B_j\not\subset B_k$.于是 $|\Gamma_j|\geq\frac{1}{2}d$，故只有有限多个圆 $\Gamma_j$ 能与 $B_{\frac{1}{2}d}(x)$ 相交.既然 $x\in\overline{\bigcup_i\Gamma_i}$，我们得出对某个 $k$，$x\in\overline{\bigcup_{i=1}^{k}\Gamma_i}=\bigcup_{i=1}^{k}\Gamma_i\subset F$，如所需.

为证 $F$ 连通，设 $F=F_1\cup F_2$，其中 $F_1,F_2$ 是不交闭集，并由此导出矛盾.令
\[
  E_1=F_1\cup\bigcup_{\{i:\Gamma_i\subset F_1\}}B_i,\qquad
  E_2=F_2\cup\bigcup_{\{i:\Gamma_i\subset F_2\}}B_i.
\]
每个 $\Gamma_i$ 都含于 $E_1$ 或 $E_2$ 之一，故 $E_1\cup E_2=E$.由于圆盘 $\{B_i\}$ 半不交且对每个 $i$ 有 $F\cap B_i\subset\bigcup_{j}\Gamma_j$，容易验证 $E_1\cap E_2=\emptyset$.集合 $E_1$ 是闭的：若 $x\in\overline{\bigcup_{\{i:\Gamma_i\subset F_1\}}B_i}$，则 $x$ 要么是对某个 $k$ 为 $\bigcup_{1}^{k}B_i$ 中点列的极限，要么是圆周落在闭集 $F_1$ 中且半径趋于 $0$ 的一列圆盘的极限点.无论哪种情形都有 $x\in F_1$.类似地 $E_2$ 闭，故 $E=E_1\cup E_2$ 是 $E$ 的由不交闭集构成的分解.由 $E$ 的连通性，$E_1$ 或 $E_2$ 之一为空，从而 $F_1$ 或 $F_2$ 之一为空.
\end{proof}

$Z$-集下密度上界之证明的实质包含在下一条引理中.记 $R(x,y)$ 为以 $x$ 与 $y$ 为心的圆对的\newterm{common-region}{公共区域}{common region}，即 $R(x,y)=\operatorname{int}\bigl(B_{|x-y|}(x)\cap B_{|x-y|}(y)\bigr)$.粗略地说，我们将证明：若以 $1$-集 $E$ 中的点为心的圆对的公共区域包含了 $E$ 的很大一部分，则可以用在很大程度上位于 $E$ 内的曲线把 $E$ 的各个分支连接起来.

\begin{lemma}\label{lem:3.22}
设 $E$ 是 $\mathbb{R}^2$ 中的 $1$-集，$\alpha>0$.设 $E_0$ 是 $E$ 的紧子集，$\mathcal{H}^1(E_0)>0$，使得当 $x_1,x_2\in E_0$ 时 $\mathcal{H}^1\bigl(E\cap R(x_1,x_2)\bigr)\geq\alpha|x_1-x_2|$.则存在连续统 $H$ 使 $0<\mathcal{H}^1(H\cap E)\leq\mathcal{H}^1(H)<\infty$.
\end{lemma}

\begin{proof}
利用通常的一致化过程，由\cref{cor:2.5} 可找到 $\rho_1>0$ 与紧集 $F\subset E_0$，$\mathcal{H}^1(F)>0$，使
\begin{equation}\label{eq:3.11}
  \mathcal{H}^{1}\bigl(B_r(x)\cap E\bigr)\leq 2\cdot 2r
  \quad\text{若}\ x\in F\ \text{且}\ 0<r\leq\rho_1.
\end{equation}
进一步，由\cref{cor:2.4,cor:2.5}，存在点 $y\in F$ 与数 $\rho$（$0<\rho\leq\frac{1}{10}\rho_1$）使
\begin{equation}\label{eq:3.12}
  \mathcal{H}^{1}\bigl((E\setminus F)\cap B_r(y)\bigr)<2r\cdot 10^{-3}\alpha\qquad(0<r\leq 3\rho)
\end{equation}
且
\begin{equation}\label{eq:3.13}
  \mathcal{H}^{1}\bigl(F\cap B_\rho(y)\bigr)>\frac{1}{4}\cdot 2\rho=\frac{1}{2}\rho.
\end{equation}
必要时缩小 $\rho$，还可设圆盘 $B_\rho(y)$ 的圆周 $\Gamma$ 含有 $F$ 的某个点.

令 $\mathcal{C}$ 为闭圆盘族
\begin{equation}\label{eq:3.14}
  \mathcal{C}=\bigl\{B_r(x):x\in F\cap B_\rho(y),\ 0<r<2\rho
  \ \text{且}\ \mathcal{H}^{1}\bigl((E\setminus F)\cap B_r(x)\bigr)\geq\alpha r\bigr\}.
\end{equation}
由\cref{lem:1.9}，可在 $\mathcal{C}$ 中找到有限或可数的不交闭圆盘族 $\{B_i\}$，使 $\bigcup_{B\in\mathcal{C}}B\subset\bigcup_i B_i'$，其中 $B_i'$ 是与 $B_i$ 同心、半径为其 $5$ 倍的闭圆盘.此外，可使圆盘 $\{B_i'\}$ 半不交.记 $\Gamma_i$ 为 $B_i'$ 的圆周.令
\[
  G=\bigl(F\cap B_\rho(y)\bigr)\cup\Gamma\cup\bigcup_{i}B_i',
\]
并定义
\begin{equation}\label{eq:3.15}
  H=\Bigl(G\setminus\bigcup_{i}B_i'\Bigr)\cup\bigcup_{i}\Gamma_i,
\end{equation}
见图~\ref{fig:3.2}.证明的其余部分分为六步，验证 $H$ 是具有所述性质的连续统.

(a) $G$ 是闭的：由于 $F\cap B_\rho(y)$ 与 $\Gamma$ 闭，只需证 $\overline{\bigcup_{i}B_i'}\subset G$.若 $z\in\overline{\bigcup_{i}B_i'}$，则要么 $z\in\bigcup_{i=1}^{k}B_i'\subset G$ 对某个有限的 $k$ 成立，要么 $z$ 是圆盘子列 $\{B_i'\}_i$ 中点的极限.把面积求和可知 $\sum_i|B_i|^2<\infty$，故当 $i\to\infty$ 时 $|B_i'|\to 0$；于是 $z$ 是这些圆盘的圆心的极限点，而这些圆心都位于闭集 $F\cap B_\rho(y)$ 之中.

(b) $G$ 连通：设 $G=G_1\cup G_2$，其中 $G_1,G_2$ 是非空不交闭集，并由此导出矛盾.$G$ 的任何连通子集必含于 $G_1$ 或 $G_2$ 之一，故设 $\Gamma\subset G_1$.每个 $B_i'$ 的圆心都在 $B_\rho(y)$ 内，故若 $B_i'$ 含有 $\Gamma$ 上或 $\Gamma$ 外的点，它必与 $\Gamma$ 相交，从而含于 $G_1$.于是 $G_2\subset\operatorname{int}B_\rho(y)$.

\begin{figure}[htbp]
  \centering
  \includegraphics[width=0.6\textwidth]{fig3-2.png}
  \caption{引理 3.22 证明中的集合 $G$ 与 $H$}
  \label{fig:3.2}
\end{figure}

设 $G_1'$ 为把 $\mathbb{R}^2\setminus B_\rho(y)$ 并入 $G_1$ 后所得的集合，则 $G_1'$ 闭、连通且与 $G_2$ 不交.由于 $G_1'$ 与 $G_2$ 都含 $F$ 的点，设 $x_1\in F\cap G_1'$ 与 $x_2\in F\cap G_2$ 为使距离 $r=|x_1-x_2|$ 最小的点.该下确界可以取得且为正.由于 $\Gamma\subset G_1'$ 含 $F$ 的点，$0<r<2\rho$.公共区域 $R(x_1,x_2)$ 与 $F$ 不交，否则 $r$ 可进一步缩小，故由引理的假设
\[
  0<\alpha r
  \leq\mathcal{H}^{1}\bigl(E\cap R(x_1,x_2)\bigr)
  =\mathcal{H}^{1}\bigl((E\setminus F)\cap R(x_1,x_2)\bigr)
  \leq\mathcal{H}^{1}\bigl((E\setminus F)\cap B_r(x_2)\bigr).
\]
由于 $x_2\in F\cap B_\rho(y)$ 且 $r<2\rho$，$B_r(x_2)\in\mathcal{C}$，故 $B_r(x_2)\subset\bigcup_i B_i'\subset G$.于是
\[
  B_r(x_2)=(B_r(x_2)\cap G_1')\cup(B_r(x_2)\cap G_2)
\]
是 $B_r(x_2)$ 的由不交闭集构成的分解，$x_1$ 属于第一个，$x_2$ 属于第二个.显然荒谬，因为 $B_r(x_2)$ 连通.

(c) $H$ 是连续统：由 (a) 与 (b)，$G$ 是连续统，故由\cref{lem:3.21}，$H$ 也是连续统.

(d) $\sum_i|B_i'|\leq\frac{1}{8}\rho$：利用 \eqref{eq:3.14}、$B_i$ 的不交性与 \eqref{eq:3.12}，
\[
  \sum_{i}|B_i'|
  =5\sum_{i}|B_i|
  \leq\frac{10}{\alpha}\sum_{i}\mathcal{H}^{1}\bigl((E\setminus F)\cap B_i\bigr)
  \leq\frac{10}{\alpha}\mathcal{H}^{1}\bigl((E\setminus F)\cap B_{3\rho}(y)\bigr)
  \leq\frac{10}{\alpha}\cdot 6\rho\,10^{-3}\alpha
  \leq\frac{1}{8}\rho.
\]

(e) $\mathcal{H}^1(H)<\infty$：由于 $H\subset E\cup\Gamma\cup\bigcup_i\Gamma_i$，利用 (d) 得
\[
  \mathcal{H}^{1}(H)\leq\mathcal{H}^{1}(E)+2\pi\rho+\pi\sum_{i}|B_i'|<\infty.
\]

(f) $\mathcal{H}^1(H\cap E)>0$：由 $H$ 的定义，
\begin{align*}
  \mathcal{H}^{1}(H\cap E)
  &\geq\mathcal{H}^{1}(G\cap E)-\mathcal{H}^{1}\Bigl(E\cap\bigcup_{i}B_i'\Bigr)
  \geq\mathcal{H}^{1}\bigl(F\cap B_\rho(y)\bigr)-\sum_{i}\mathcal{H}^{1}\bigl(E\cap B_i'\bigr)\\
  &\geq\mathcal{H}^{1}\bigl(F\cap B_\rho(y)\bigr)-2\sum_{i}|B_i'|
  \geq\frac{1}{2}\rho-\frac{1}{4}\rho=\frac{1}{4}\rho>0,
\end{align*}
其中用到 \eqref{eq:3.11}（注意 $|B_i'|<20\rho\leq 2\rho_1$），而最后的估计由 \eqref{eq:3.13} 与 (d) 提供.
\end{proof}

关于 $Z$-集下密度的定理现在是圆对几何的一个简单推论.我们利用如下事实：以下密度与法向凸密度都很高的点为心的圆对，必在其公共区域中包含一个正测度子集.

\begin{theorem}\label{thm:3.23}
设 $E$ 是 $\mathbb{R}^2$ 中的 $Z$-集.则在几乎所有的 $x\in E$ 处 $\underline{D}^1(E,x)\leq\frac{3}{4}$.
\end{theorem}

\begin{proof}
设对某个 $\alpha>0$，集合 $E_1=\{x:D^1(E,x)>\frac{3}{4}+\alpha\}$ 具有正测度.于是利用\cref{thm:2.3}，按通常方式可找到紧 $1$-集 $E_2\subset E_1$ 与 $\rho>0$ 使
\begin{equation}\label{eq:3.16}
  \mathcal{H}^{1}\bigl(E\cap B_r(x)\bigr)>\bigl(\tfrac{3}{4}+\alpha\bigr)2r
  \qquad(x\in E_2,\ 0<r\leq\rho)
\end{equation}
且
\begin{equation}\label{eq:3.17}
  \mathcal{H}^{1}(E\cap U)<(1+\alpha)|U|
  \qquad(x\in E_2\cap U,\ 0<|U|\leq 3\rho).
\end{equation}
设 $E_0$ 是 $E_2$ 的紧子集，满足 $0<\mathcal{H}^1(E_0)<\infty$ 且 $|E_0|\leq\rho$.若 $x_1,x_2\in E_0$，则 $r=|x_1-x_2|\leq\rho$，于是注意到任何圆与 $E$ 的交测度为零，并利用 \eqref{eq:3.16}、\eqref{eq:3.17}（取 $U=B_r(x_1)\cap B_r(x_2)$）得
\begin{align*}
  \mathcal{H}^{1}\bigl(E\cap R(x_1,x_2)\bigr)
  &=\mathcal{H}^{1}\bigl(E\cap B_r(x_1)\bigr)+\mathcal{H}^{1}\bigl(E\cap B_r(x_2)\bigr)
  -\mathcal{H}^{1}\bigl(E\cap(B_r(x_1)\cup B_r(x_2))\bigr)\\
  &\geq 2\bigl(\tfrac{3}{4}+\alpha\bigr)2r-(1+\alpha)3r=\alpha r=\alpha|x_1-x_2|.
\end{align*}
于是对 $E_0\subset E$ 应用\cref{lem:3.22}，推得存在连续统 $H$ 使 $0<\mathcal{H}^1(H\cap E)\leq\mathcal{H}^1(H)<\infty$.由于 $E$ 是 $Z$-集，由\cref{lem:3.20} 这不可能.故当 $\alpha>0$ 时 $\mathcal{H}^1(E_1)=0$，如所需.
\end{proof}

\begin{corollary}\label{cor:3.24}
设 $E$ 是 $\mathbb{R}^2$ 中的\newterm{irregular-1-set}{不规则 $1$-集}{irregular 1-set}.则在几乎所有的 $x\in E$ 处 $\underline{D}^1(E,x)\leq\frac{3}{4}$.
\end{corollary}

\begin{proof}
由\cref{cor:3.9}，$E$ 的任何正测度子集都不含于任何可求长曲线.故 $E$ 是 $Z$-集，结论由\cref{thm:3.23} 立即得到.
\end{proof}

我们已经证明：就下密度而言，正则与不规则 $1$-集的性质截然不同——正则集的下密度为 $1$，而不规则集的下密度几乎处处至多为 $\frac{3}{4}$.不可能构造出在正测度集上下密度严格介于 $\frac{3}{4}$ 与 $1$ 之间的 $1$-集.

$\underline{D}^1(E,x)$ 的精确的本质上界仍未确定.虽然 $\frac{3}{4}$ 是迄今得到的最佳值，但人们普遍相信正确答案是 $\frac{1}{2}$；确实可以构造出使 $D^1(E,x)=\frac{1}{2}$ 几乎处处成立的集合.有趣的是，\textcite{besicovitch1928a} 最初给出的界是 $1-10^{-2576}$，其后在 1938 年的文章中才把它降到 $\frac{3}{4}$.

\textcite{besicovitch1938} 还发展了他的圆对方法以证明不规则集的密度几乎处处不存在，另见 \textcites{besicovitchwalker1931,gillis1934b,morserandolph1944}.

现在可以刻画正则与不规则 $1$-集了.

\begin{theorem}\label{thm:3.25}
不规则 $1$-集是 $Z$-集.\newterm{regular-1-set}{正则 $1$-集}{regular 1-set}由一个 $Y$-集再加上一个 $\mathcal{H}^1$-测度为零的集合组成.
\end{theorem}

\begin{proof}
若 $E$ 不规则，则 $E$ 与每条可求长曲线 $\Gamma$ 的交测度为零，否则 $E\cap\Gamma$ 将是一个 $Y$-集，从而是正则的.

若 $E$ 正则，则 $D^1(E,x)=1$ 几乎处处成立，于是\cref{thm:3.23} 结合\cref{cor:2.6} 蕴含：$E$ 的任何正测度可测子集与某条可求长曲线的交具有正测度.我们用这一事实定义可求长曲线序列 $\{\Gamma_j\}$.选取 $\Gamma_1$ 使
\[
  \mathcal{H}^{1}(\Gamma_1\cap E)\geq\frac{1}{2}\sup\bigl\{\mathcal{H}^{1}(\Gamma\cap E):\Gamma\ \text{可求长}\bigr\}.
\]
设 $\Gamma_1,\ldots,\Gamma_k$ 已经选出且 $E_k=E\setminus\bigcup_{1}^{k}\Gamma_j$ 具有正测度，则令 $\Gamma_{k+1}$ 为满足
\begin{equation}\label{eq:3.18}
  \mathcal{H}^{1}(\Gamma_{k+1}\cap E_k)\geq\frac{1}{2}\sup\bigl\{\mathcal{H}^{1}(\Gamma\cap E_k):\Gamma\ \text{可求长}\bigr\}>0
\end{equation}
的可求长曲线.

只有当对某个 $k$，曲线 $\{\Gamma_j\}_{1}^{k}$ 覆盖了 $E$ 的几乎全部时过程才终止，此时结论显然.否则，
\[
  \infty>\mathcal{H}^{1}(E)\geq\sum_{k}\mathcal{H}^{1}(\Gamma_{k+1}\cap E_k),
\]
故 $\mathcal{H}^1(\Gamma_{k+1}\cap E_k)\to 0$.若 $\mathcal{H}^1\bigl(E\setminus\bigcup_{1}^{\infty}\Gamma_j\bigr)>0$，则可找到可求长曲线 $\Gamma$ 使 $\mathcal{H}^1\bigl(\Gamma\cap(E\setminus\bigcup_{1}^{\infty}\Gamma_j)\bigr)=d>0$.但对某个 $k$ 有 $\mathcal{H}^1(\Gamma_{k+1}\cap E_k)<\frac{1}{2}d$，而按 \eqref{eq:3.18}，$\Gamma$ 本应被优先选为 $\Gamma_{k+1}$.故 $\mathcal{H}^1\bigl(E\setminus\bigcup_{1}^{\infty}\Gamma_j\bigr)=0$，于是 $E$ 由 $Y$-集 $\bigcup_{j}(E\cap\Gamma_j)$ 与一个零测集组成.
\end{proof}

不规则 $1$-集作为 $Z$-集，与每条可求长曲线的交测度为零.然而 \textcite{besicovitch1928a}（第 43 节）表明，不规则集可以嵌入非可求长曲线的可数并之中.

本节最后指出不规则集的``破碎''性质.回忆：若集合的任何两点都不位于同一连通分支中，则称该集合是完全不连通的.于是，给定集合中的任意一对点，存在由两个不交闭子集构成的分解，每个子集含其中一点.

\begin{corollary}\label{cor:3.26}
不规则 $1$-集 $E$ 是\newterm{totally-disconnected}{完全不连通的}{totally disconnected}.
\end{corollary}

\begin{proof}
若 $x$ 与 $y$ 同属于 $E$ 的某个闭连通子集 $F$，则由\cref{lem:3.4} 有 $|x-y|\leq\mathcal{H}^1(F)<\infty$.但 $E$ 是 $Z$-集，故由\cref{lem:3.20} 有 $\mathcal{H}^1(F)=\mathcal{H}^1(E\cap F)=0$.于是实际上 $x=y$.
\end{proof}

\section{切线性质}

与密度性质一样，正则集与不规则集的切线性质截然不同.回忆切线的定义 \eqref{eq:3.4}，正则集的结果已经清楚：

\begin{theorem}\label{thm:3.27}
$\mathbb{R}^n$ 中的正则 $1$-集 $E$ 在几乎所有的点处有切线.
\end{theorem}

\begin{proof}
由\cref{thm:3.25}，$E$ 由一个 $Y$-集与一个零测集组成，而由\cref{cor:3.10}，任何 $Y$-集在几乎所有点处有切线.
\end{proof}

为研究平面中不规则集的切线性质，我们使用 \textcite{besicovitch1938}（第 10 节）的一个精巧论证.首先需要一个几何引理.

\begin{lemma}\label{lem:3.28}
设 $\theta$ 是 $\mathbb{R}^2$ 中与直线 $L$ 垂直的单位向量.设 $P$ 是边与方向 $\pm\theta$ 成角 $\phi$ 的平行四边形，$y$ 与 $z$ 是 $P$ 的一对相对顶点，如图~\ref{fig:3.3}.则 $|y-z|\leq d/\sin\phi$，其中 $d$ 是 $P$ 到 $L$ 上投影的长度.
\end{lemma}

\begin{figure}[htbp]
  \centering
  \includegraphics[width=0.4\textwidth]{fig3-3.png}
  \caption{引理 3.28 中的平行四边形}
  \label{fig:3.3}
\end{figure}

\begin{proof}
设 $w$ 为 $P$ 的第三个顶点.则 $|y-z|\leq|y-w|+|w-z|$，且 $d=(|y-w|+|w-z|)\sin\phi$，由此得结论.
\end{proof}

下一条定理本质上是 Besicovitch 关于不规则集的切线结果.证明的想法是把不规则点集约化为一条可求长曲线的子集，而后者必有零测度.

\begin{theorem}\label{thm:3.29}
设 $E$ 是 $\mathbb{R}^2$ 中的不规则 $1$-集.则给定 $\theta$ 与 $0<\phi<\pi/2$，
\begin{equation}\label{eq:3.19}
  \bar{D}^{1}(E,x,\theta,\phi)+\bar{D}^{1}(E,x,-\theta,\phi)\geq\frac{1}{6}\sin\phi
\end{equation}
对几乎所有的 $x\in E$ 成立.
\end{theorem}

\begin{proof}
取 $\rho,\delta_+,\delta_->0$，并令 $F_0=F_0(\delta_+,\delta_-,\rho)$ 为 $E$ 中使
\begin{equation}\label{eq:3.20}
  \left.
  \begin{aligned}
    \mathcal{H}^{1}\bigl(E\cap S_r(x,\theta,\phi)\bigr)&\leq 2r\delta_+\\
    \mathcal{H}^{1}\bigl(E\cap S_r(x,-\theta,\phi)\bigr)&\leq 2r\delta_-
  \end{aligned}
  \right\}
\end{equation}
对一切 $r\leq\rho$ 成立的 $x$ 之集.照例 $F_0$ 可测；我们将证明若 $F_0$ 有正测度，则 $\delta_++\delta_-$ 不能太小.

若 $\mathcal{H}^1(F_0)>0$，则由 $\mathcal{H}^1$ 的正则性可找到 $F_0$ 的正测度紧子集 $F_1$.进一步，由\cref{thm:2.3} 或直接由 $\mathcal{H}^1$ 的定义，给定 $\eta>0$，可找到闭凸集 $U$ 使 $0<|U|\leq\rho$ 且
\begin{equation}\label{eq:3.21}
  \mathcal{H}^{1}(F)>(1-\eta)|U|,
\end{equation}
其中 $F=F_1\cap U$.从此我们只在 $U$ 内工作.

由于 $F$ 闭，可选取 $y_1$ 与 $z_1$ 为 $F$ 中连接线段与 $\theta$ 的夹角不超过 $\phi$ 的最远点对之一，使得
\[
  r_1=|y_1-z_1|
  =\sup\bigl\{|y-z|:z\in F\cap S(y,\theta,\phi)\ \text{且}\ y\in F\cap S(z,-\theta,\phi)\bigr\},
\]
其中 $S(y,\theta,\phi)$ 表示无穷扇形.$r_1$ 的极大性保证
\[
  F\cap S(y_1,\theta,\phi)=F\cap S_{r_1}(y_1,\theta,\phi)
\]
且
\[
  F\cap S(z_1,-\theta,\phi)=F\cap S_{r_1}(z_1,-\theta,\phi).
\]
由于 $r_1\leq\rho$ 且 $F\subset E$，由 \eqref{eq:3.20} 得
\begin{equation}\label{eq:3.22}
  \left.
  \begin{aligned}
    \mathcal{H}^{1}\bigl(F\cap S(y_1,\theta,\phi)\bigr)&\leq 2r_1\delta_+\\
    \mathcal{H}^{1}\bigl(F\cap S(z_1,-\theta,\phi)\bigr)&\leq 2r_1\delta_-.
  \end{aligned}
  \right\}
\end{equation}
令 $P_1$ 为闭平行四边形
\[
  P_1=S(y_1,\theta,\phi)\cap S(z_1,-\theta,\phi)
\]
并令 $Q_1$ 为开区域
\[
  Q_1=\operatorname{int}\bigl\{S(y_1,\theta,\phi)\cup S(z_1,-\theta,\phi)\bigr\},
\]
见图~\ref{fig:3.4}.由 \eqref{eq:3.22},
\[
  \mathcal{H}^{1}(F\cap Q_1)\leq 2r_1(\delta_++\delta_-).
\]

\begin{figure}[htbp]
  \centering
  \includegraphics[width=0.42\textwidth]{fig3-4.png}
  \caption{\cref{thm:3.29} 证明中的扇形与区域}
  \label{fig:3.4}
\end{figure}

集合 $U\setminus Q_1$ 紧且至多有两个分支，因此只要它在其同一分支中含有连接线段与 $\theta$ 的夹角至多为 $\phi$ 的一对 $F$ 中的点，我们就可用同样的方法得到位于 $Q_1$ 之外的平行四边形 $P_2$ 与开区域 $Q_2$.于是对 $j=1,2,\ldots$，我们在 $U\setminus\bigcup_{i=1}^{j-1}Q_i$ 的同一分支中找到 $F$ 的点 $y_j$ 与 $z_j$，使 $r_j=|y_j-z_j|$ 尽可能大，且
\begin{equation}\label{eq:3.23}
  \mathcal{H}^{1}(F\cap Q_j)\leq 2r_j(\delta_++\delta_-),
\end{equation}
其中
\[
  Q_j=\operatorname{int}\bigl\{S(y_j,\theta,\phi)\cup S(z_j,-\theta,\phi)\bigr\},
  \qquad
  P_j=S(y_j,\theta,\phi)\cap S(z_j,-\theta,\phi).
\]
这一过程或者无限进行下去，或者直到没有合适的点对为止.

设 $L$ 为与 $\theta$ 垂直的直线.若 $j>i$，则 $P_j$ 与 $Q_i$ 不交，故平行四边形 $\{P_j\}$ 在 $L$ 上的投影是内部不交的线段.用\cref{lem:3.28} 容易验证
\begin{equation}\label{eq:3.24}
  \sum_{j}r_j=\sum_{j}|y_j-z_j|\leq\frac{3|U|}{\sin\phi};
\end{equation}
因子 `3' 顾及了至多两个 $P_j$ 在 $U$ 两侧重叠的可能性.

设 $y$ 与 $z$ 为 $F\setminus\bigcup_{i}Q_i$ 中两个不同的点.若 $y$ 与 $z$ 位于某个 $Q_j$ 的两侧，则 $[y,z]$ 与 $\theta$ 的夹角至少为 $\phi$.若不然，且 $[y,z]$ 与 $\theta$ 的夹角小于 $\phi$，则 $y$ 与 $z$ 在每个阶段都本是 $y_j$ 与 $z_j$ 的候选者，故选取过程不会终止，且对一切 $j$ 有 $|y-z|\leq|y_j-z_j|$.这是不可能的，因为 \eqref{eq:3.24} 蕴含 $|y_j-z_j|\to 0$.于是我们得出：连接 $F\setminus\bigcup_{i}Q_i$ 中任意两点对线段与 $L$ 的夹角至多为 $\frac{1}{2}\pi-\phi$.由\cref{lem:3.11}, $F\setminus\bigcup_{i}Q_i$ 含于一条可求长曲线之中.由于 $F\setminus\bigcup_{i}Q_i$ 是 $E$ 的子集，而由\cref{thm:3.25} $E$ 是 $Z$-集，故
\[
  \mathcal{H}^{1}\Bigl(F\setminus\bigcup_{i}Q_i\Bigr)=0.
\]
于是利用 \eqref{eq:3.23}、\eqref{eq:3.24} 与 \eqref{eq:3.21},
\[
  \mathcal{H}^{1}(F)
  \leq\sum_{j}\mathcal{H}^{1}(F\cap Q_j)
  \leq 2(\delta_++\delta_-)\sum_{j}r_j
  \leq\frac{6(\delta_++\delta_-)|U|}{\sin\phi}
  \leq\frac{6(\delta_++\delta_-)}{(1-\eta)\sin\phi}\,\mathcal{H}^{1}(F).
\]
对任意 $\eta>0$ 都存在使上式成立的满足 $\mathcal{H}^1(F)>0$ 的 $F$，故
\[
  (\delta_++\delta_-)\geq\frac{1}{6}\sin\phi.
\]
于是若 $\delta_++\delta_-<\frac{1}{6}\sin\phi$，则由 \eqref{eq:3.20} 定义的集合 $F_0(\delta_+,\delta_-,\rho)$ 必测度为零，证明完成.
\end{proof}

\begin{corollary}\label{cor:3.30}
设 $E$ 是 $\mathbb{R}^2$ 中的不规则 $1$-集.则对几乎所有的 $x\in E$,
\begin{equation}\label{eq:3.25}
  \bar{D}^{1}(E,x,\theta,\phi)+\bar{D}^{1}(E,x,-\theta,\phi)\geq\frac{1}{6}\sin\phi>\frac{1}{10}\phi
\end{equation}
对一切 $\theta$ 与一切 $0<\phi<\frac{1}{2}\pi$ 成立.
\end{corollary}

\begin{proof}
由\cref{thm:3.29}, \eqref{eq:3.19} 对 $\{(\theta,\phi)\}$ 的某个可数稠密子集对几乎所有的 $x$ 成立.于是通过逼近，对几乎所有的 $x$，\eqref{eq:3.19} 对一切 $\theta$ 与 $\phi$ 成立.
\end{proof}

\eqref{eq:3.19} 或 \eqref{eq:3.25} 右端的系数当然可以改进.事实上，Besicovitch 的证明（用了略优于\cref{lem:3.28} 与 \eqref{eq:3.24} 的估计）给出 $\frac{1}{3}\sin\phi$，而且当 $\phi\geq\frac{1}{3}\pi$ 时给出 $\frac{1}{2}\sin\phi$.Besicovitch 猜测最佳系数可能大至 $2^{-\frac{1}{3}}$（Besicovitch (1982a, 第 11 节) 中的一个例子表明它不能比这更大）.\footnote{原书参考文献表中并无 Besicovitch (1982a) 这一条目，疑为 \citeyear{besicovitch1928a} 之误印，请读者鉴察.——译者注.}然而，任何正值都足以导出主要的切线结果.

\begin{corollary}\label{cor:3.31}
在平面不规则 $1$-集的几乎所有点处不存在切线.
\end{corollary}

\begin{proof}
由\cref{cor:3.30} 与切线的定义 \eqref{eq:3.4} 立得.
\end{proof}

如我们所见，正则 $1$-集可以作为可求长曲线的子集得到.不规则集的例子也许不那么显然，因此我们现在给出一个不规则 $1$-集的简单构造.

\begin{theorem}\label{thm:3.32}
$\mathbb{R}^2$ 中存在不规则 $1$-集.
\end{theorem}

\begin{proof}
设 $f$ 是 $[0,1)$ 上由
\[
  f(.x_1x_2\ldots)=.y_1y_2\ldots
\]
定义的实函数，其中 $x_1x_2\ldots$ 与 $y_1y_2\ldots$ 以 $4$ 为底书写（约定 $x_i$ 不最终全等于 $3$），使得 $y_i=5-x_i\pmod{4}$.（见图~\ref{fig:3.5}.）令
\[
  E=\{(x,f(x)):0\leq x<1\}.
\]
$E$ 是 Borel 集，且由于 $E$ 到 $x$ 轴的投影长度为 $1$，得 $1\leq\mathcal{H}^1(E)$.另一方面，对每个 $k$，$E$ 可用 $4^k$ 个边长为 $4^{-k}$ 的正方形覆盖，故 $\mathcal{H}^1(E)\leq\sqrt{2}$.于是 $E$ 是一个 $1$-集.容易看出任何可求长曲线（它在几乎处处有切线）与 $E$ 的交长度为零.故 $E$ 是 $Z$-集，是不规则的.另外，不规则性也可由 $E$ 在与 $x$ 轴成 $\pm 45^{\circ}$ 的两个方向上投影成勒贝格测度为零的集合这一事实得到；见\cref{cor:6.14}.
\end{proof}

\begin{figure}[htbp]
  \centering
  \includegraphics[width=0.42\textwidth]{fig3-5.png}
  \caption{\cref{thm:3.32} 中函数 $y=f(x)$ 的图像}
  \label{fig:3.5}
\end{figure}

\section{高维集合}

把该理论推广到 $\mathbb{R}^n$（$n\geq 3$）的整数维数子集，直到 Besicovitch 的第一篇基础性论文 \citeyear{besicovitch1928a} 发表约 47 年之后才完成.这是因为必须克服相当大的技术困难；例如，``有限 $\mathcal{H}^s$-测度的连续统是可求长曲线之并''这一定理当 $s$ 为大于 $1$ 的整数时没有类似物.

称 $\mathbb{R}^n$ 中的 $s$-集为\newterm{countably rectifiable}{可数可求长的}{countably rectifiable}，如果它形如
\[
  \bigcup_{j=1}^{\infty}f_j(E_j)\cup G,
\]
其中 $\mathcal{H}^s(G)=0$，且每个 $f_j$ 都是从 $\mathbb{R}^s$ 的有界子集 $E_j$ 到 $\mathbb{R}^n$ 的 \newterm{lipschitz-mapping}{Lipschitz 函数}{Lipschitz mapping}.（Lipschitz 函数 $f$ 要求对某个常数 $c$ 有 $|f(t)-f(u)|\leq c|t-u|$.）这是 Besicovitch 的``$Y$-集加上一个零测集''的直接类似物.

$\mathbb{R}^n$ 中的 $s$-集 $E$ 称为在 $x$ 处有切线，如果 $\bar{D}^s(E,x)>0$ 且存在过 $x$ 的 $s$ 维平面 $\Pi$，使对一切 $\phi>0$,
\begin{equation}\label{eq:3.26}
  \lim_{r\to 0}r^{-s}\mathcal{H}^{s}\bigl(E\cap(B_r(x)\setminus S(x,\Pi,\phi))\bigr)=0,
\end{equation}
其中 $S(x,\Pi,\phi)$ 表示使 $[y,x]$ 与 $\Pi$ 的夹角至多为 $\phi$ 的点 $y\in\mathbb{R}^n$ 之集.若 $\underline{D}^s(E,x)>0$ 且把 \eqref{eq:3.26} 中的 `lim' 换成 `$\varliminf$'，则称 $E$ 在 $x$ 处有\newterm{weak tangent}{弱切线}{weak tangent}.

\begin{theorem}\label{thm:3.33}
设 $E$ 是 $\mathbb{R}^n$ 中的 $s$-集，$s$ 为整数.则下述命题等价：

(a) $E$ 是正则的.

(b) $E$ 是可数可求长的.

(c) $E$ 在几乎所有的点处有切线.
\end{theorem}

如我们所见，Besicovitch（\citeyear{besicovitch1928a,besicovitch1938}） 对 $\mathbb{R}^2$ 中的 $1$-集证明了这一点.\textcite{federer1947} 在一篇极其复杂的文章中证明了 (b) 与 (c) 的等价性，而由这一点及切线的定义可得 (b) 蕴含 (a).（Federer 的 restrictedness 概念等价于切线的存在性.）\textcite{moore1950} 完成了 $\mathbb{R}^n$ 中 $1$-集的证明，\textcite{marstrand1961} 完成了 $\mathbb{R}^3$ 中 $2$-集的证明，而 \textcite{mattila1975a} 最终完成了 $\mathbb{R}^n$ 中 $s$-集的证明.对更一般的测度 $\mu$ 也得到了此类性质的结果.这里必须把正则性单纯定义为密度 $\lim_{r\to 0}r^{-s}\mu(B_r(x))$ $\mu$-几乎处处存在，而不涉及极限的实际值.上文提到的 Federer 与 Moore 的文章考虑了这样的推广，\textcite{morserandolph1944} 考虑了 $\mathbb{R}^2$ 中 $s=1$ 的子集，\textcite{marstrand1964} 则考虑了一般情形.特别地，Marstrand 证明对一般测度，正则性蕴含弱切线（按显然的方式定义）几乎处处存在；自然的猜想是 `弱'一词可以省去.

这些证明中的许多都与\cref{chap:6} 中讨论的投影性质的高维类似物紧密相关.

当 $s$ 是整数且 $1<s<n$ 时，把曲线理论推广到 $\mathbb{R}^n$ 中的 $s$ 维\newterm{surface}{曲面}{surface}存在相当的困难.这一理论远比曲线理论复杂；事实上，对``曲面''与``面积''等术语所赋予的精确含义需要格外小心.

$\mathbb{R}^3$ 中的曲面已被研究多年，很长时间里公认的曲面面积定义是 \newterm{lebesgue-frechet-area}{Lebesgue--Fréchet 定义}{Lebesgue-Fréchet area}，即作为内接多面体面积的上确界，与曲线长度的定义相类比.正是 \textcite{besicovitch1945b} 证明了这一定义是全然不适当的，它甚至不具备测度的基本可加性质.他提出最令人满意的面积定义是 Hausdorff 测度 $\mathcal{H}^2$，并着手由此出发处理曲面面积的各类问题.

由此产生的一个问题是如何把 Gołąb 半连续性定理（\cref{thm:3.18}）推广到曲面.显然没有直接的类似物，因为任何曲面都可以在 Hausdorff 度量下用一列可求长曲线逼近，而后者理应有零面积.\textcite{besicovitch1949} 证明了：若 $\{E_j\}$ 是收敛于参数曲面 $E$ 的参数曲面序列，且极限曲面 $E$ 满足合理的光滑性条件，则 $\mathcal{H}^2(E)\leq\varliminf_{j\to\infty}\mathcal{H}^2(E_j)$.把 Gołąb 定理向高维集合的一个更近的推广归功于 \textcite{vitushkin1966}，另见 \textcite{ivanov1975}.我们如下定义 $\mathbb{R}^n$ 的紧子集 $E$ 的\newterm{variations}{变差}{variations}.令 $V^0(E)$ 为 $E$ 的连通分支的个数，于是当 $V^0(E)=1$ 时 $E$ 是连续统.对 $1\leq s\leq n$，令 $V^s(E)=c\int V^0(E\cap\Pi)\,\mathrm{d}\Pi$，其中 $\Pi$ 取遍一切 $(n-s)$ 维平面，积分关于所有这类平面的空间上的自然不变测度（出现于积分几何中的那个）进行.规范化常数 $c$ 只依赖于 $n$ 与 $s$，其选取使得当 $E$ 是 $s$ 维方体时 $V^s(E)$ 就是其体积.于是对充分光滑的曲面或多面体曲面，$V^s$ 就是 $s$ 维体积.这些变差起初用于凸集的研究，但也可以卓有成效地用于一般紧集.\newterm{vitushkin-semicontinuity}{Vitushkin 半连续性定理}{semicontinuity theorem, Vitushkin's}断言：若 $\{E_j\}$ 是 $\mathbb{R}^n$ 中在 Hausdorff 度量下收敛于 $E$ 的紧集序列，则只要数 $\{V^i(E_j)\}$ 对 $0\leq i<s$ 一致有界，就有 $V^s(E)\leq\varliminf_{j\to\infty}V^s(E_j)$.当 $s=1$ 时它化为 Gołąb 定理.

一个特别有意义的问题是求\newterm{minimal-surface}{``极小曲面''}{minimal surface}，即张成 $\mathbb{R}^3$ 中一条闭曲线的面积最小的曲面.这需要对``曲面''作仔细的定义，随后问题便适合用 Hausdorff 测度来处理.\newterm{plateau-problem}{``Plateau 问题''}{Plateau problem}——面积最小的曲面是否总存在——最终由 \textcite{reifenberg1960} 以肯定的方式解决.（该文也列出了更早的文献.）\textcite{besicovitch1950} 对其中一些专题作了综述.

若 $E$ 是 $\mathbb{R}^n$ 中满足 $\mathcal{L}^n(E)>0$ 的有界勒贝格可测子集，$E$ 的\newterm{density-boundary}{密度边界}{density boundary} $\partial E$ 由 $\mathbb{R}^n$ 中使 $E$ 的勒贝格密度或者不存在、或者取 $0$ 与 $1$ 以外的值的那些点组成.勒贝格密度定理（\cref{thm:1.13}）蕴含 $\mathcal{L}^n(\partial E)=0$.然而在相反的方向上，当 $E$ 的边界高度不规则时，可以证明 $\mathcal{H}^{n-1}(\partial E)>0$ 这一非平凡结果.

人们还设计了多种其他曲面测度.除已提到者外，常用的还包括积分几何测度与 \newterm{de-gorgi-perimeter}{de Giorgi 周长}{de Gorgi perimeter}；见 \textcite{federer1969}，那里详细定义并考察了这类测度.

\begin{exercises}
\item 设 $\psi:[a,b]\to\mathbb{R}^n$ 是（不必连续的）有界变差函数，即
$\operatorname{var}[a,b]=\sup\sum_{i=1}^{m}|\psi(t_i)-\psi(t_{i-1})|<\infty$，
其中上确界取遍 $[a,b]$ 的一切分割 $a=t_0<t_1<\cdots<t_m=b$.证明 $\psi[a,b]$ 含于一条（可能自交的）可求长曲线.（提示：用 $\operatorname{var}[a,t]$ 定义一个从 $[0,\operatorname{var}[a,b]]$ 到 $\mathbb{R}^n$ 的连续函数.）进而证明 $\mathcal{H}^1(\psi[a,b])<\infty$.
\item 证明任何弧连通集都是连通的.给出 $\mathbb{R}^2$ 中一个不是弧连通的连续统的例子.（比较\cref{lem:3.12}.）
\item 验证 Hausdorff 度量是一个度量.
\item 给出 $\mathbb{R}^2$ 中在 Hausdorff 度量下收敛于 $E$ 的连续统序列 $\{E_j\}$ 的例子，使 $\mathcal{H}^1(E_j)\to\infty$ 而 $\mathcal{H}^1(E)<\infty$.
\item 设 $E$ 是满足 $\mathcal{H}^1(E)<\infty$ 的连续统.证明存在树 $F\subset E$ 使 $\mathcal{H}^1(E\setminus F)$ 任意小.进而证明 $E$ 含于一条长度至多为 $2\mathcal{H}^1(E)$ 的（可能自交的）曲线之中.
\item 设 $E$ 是不规则 $1$-集.给定 $\varepsilon,c>0$，证明存在 $\delta_0>0$，使得对任何满足 $\mathcal{H}^1(F)<c$ 的连续统 $F$，当 $0<\delta<\delta_0$ 时有 $\mathcal{H}^1\bigl(E\cap[F]_{\delta}\bigr)<\varepsilon$，其中 $[F]_{\delta}$ 是 $F$ 的 $\delta$-平行体.
\item 证明存在 Hausdorff 维数为 $2$ 的紧集 $F\subset\mathbb{R}^2$，使得含于 $F$ 中的每个 $1$-集都是不规则的.
\item 在 $\mathbb{R}^2$ 中构造一个完全不连通的正则 $1$-集.（比较\cref{cor:3.26}.）
\end{exercises}

% !TeX root = ../main.tex
% 第 4 章 非整数维数集合的结构 (对应原书 pp. 54--63)

\chapter{非整数维数集合的结构}\label{chap:4}

\section{引言}

本章考察 $s$ 非整数时 $\mathbb{R}^n$ 中 $s$-集的局部性质.基本结果是：任何这样的集合都是非正则的，即它的下圆形密度或下球面密度在几乎所有的点处严格小于 $1$.事实上，更强的结果——它的密度在几乎所有的点处不存在——也已被建立.与之前一样，我们还考察适当定义的切线的存在性，并证明存在这种切线的点之集必有零测度.

对平面的子集而言，这一工作完全归功于 Marstrand（\citeyear{marstrand1954a,marstrand1955}），前一篇文章给出了非常完整的论述.与整数维数的集合一样，高维类似物呈现出巨大的困难；其自然的推广最终由 \textcite{marstrand1964} 证明.

\section{\texorpdfstring{$0<s<1$ 的 $s$-集}{0<s<1 的 s-集}}

首先考虑 $\mathbb{R}^n$ 中 $s$ 严格小于 $1$ 的 \newterm{s-set-lt1}{$s$-集}{$s$-set (s < 1)}.此时基本性质（包括密度几乎处处不存在）相对容易得到.

关于这类集合的下述拓扑观察有时是有用的.

\begin{lemma}\label{lem:4.1}
$\mathbb{R}^n$ 中满足 $0<s<1$ 的 $s$-集 $E$ 是完全不连通的.
\end{lemma}

\begin{proof}
设 $x$ 与 $y$ 是 $E$ 的同一连通分支中两个不同的点.定义映射 $f:\mathbb{R}^n\to[0,\infty)$ 为 $f(z)=|z-x|$.由于 $f$ 不增大距离，由\cref{lem:1.8} 得 $\mathcal{H}^s(f(E))\leq\mathcal{H}^s(E)<\infty$.由于 $s<1$，$f(E)$ 是 $\mathbb{R}$ 的勒贝格测度为零的子集，特别地其补集稠密.取数 $r$ 使 $r\notin f(E)$ 且 $0<r<f(y)$，则
\[
  E=\{z\in E:|z-x|<r\}\cup\{z\in E:|z-x|>r\}.
\]
这是 $E$ 的一个开分解，$x$ 属于其中一部分，$y$ 属于另一部分，与 $x,y$ 属于 $E$ 的同一连通分支的假设矛盾.
\end{proof}

下述密度性质的证明是 \textcite{marstrand1954a} 的证明的一个略微缩短的版本.

\begin{theorem}\label{thm:4.2}
若 $E$ 是满足 $0<s<1$ 的 $s$-集，则密度 $D^s(E,x)$ 在 $E$ 的几乎每一点处不存在.
\end{theorem}

\begin{proof}
若结论不成立，则由\cref{cor:2.5}，$E$ 有一个正测度的可测子集，在其上密度存在且至少为 $2^{-s}>\frac{1}{2}$.选取足够小的 $\rho$，可找到 $s$-集 $F\subset E$ 使得当 $x\in F$ 时 $D^s(E,x)$ 存在且
\begin{equation}\label{eq:4.1}
  \mathcal{H}^{s}\bigl(E\cap B_r(x)\bigr)>\frac{1}{2}(2r)^{s}
\end{equation}
对一切 $r\leq\rho$ 成立.由 $\mathcal{H}^s$ 的正则性还可设 $F$ 是闭集.设 $y$ 为 $F$ 的一个聚点，并取数 $\eta$ 使 $0<\eta<1$.若 $A_{r,\eta}$ 表示环形区域 $B_{r(1+\eta)}(y)\setminus B_{r(1-\eta)}(y)$，则当 $r\to 0$ 时
\begin{equation}\label{eq:4.2}
  \begin{aligned}
    (2r)^{-s}\mathcal{H}^{s}(E\cap A_{r,\eta})
    &=(2r)^{-s}\mathcal{H}^{s}\bigl(E\cap B_{r(1+\eta)}(y)\bigr)
    -(2r)^{-s}\mathcal{H}^{s}\bigl(E\cap B_{r(1-\eta)}(y)\bigr)\\
    &\to D^{s}(E,y)\bigl((1+\eta)^{s}-(1-\eta)^{s}\bigr).
  \end{aligned}
\end{equation}
另一方面，对任意小的 $r$ 值可找到 $x\in F$ 使 $|x-y|=r$.于是
\[
  B_{\frac{1}{2}r\eta}(x)\subset A_{r,\eta},
\]
故由 \eqref{eq:4.1}，对任意小的 $r$ 值（见图~\ref{fig:4.1}）有
\[
  \frac{1}{2}r^{s}\eta^{s}<\mathcal{H}^{s}\bigl(E\cap B_{\frac{1}{2}r\eta}(x)\bigr)\leq\mathcal{H}^{s}(E\cap A_{r,\eta}).
\]
利用 \eqref{eq:4.2} 得出
\[
  2^{-(s+1)}\eta^{s}\leq D^{s}(E,y)\bigl((1+\eta)^{s}-(1-\eta)^{s}\bigr)=D^{s}(E,y)\bigl(2s\eta+O(\eta^{2})\bigr)
\]
（$\eta\to 0$）.当 $s<1$ 时这是不可能的，定理由这一矛盾得证.
\end{proof}

\begin{figure}[htbp]
  \centering
  \includegraphics[width=0.34\textwidth]{fig4-1.png}
  \caption{\cref{thm:4.2} 证明中的环形区域}
  \label{fig:4.1}
\end{figure}

\begin{corollary}\label{cor:4.3}
任何满足 $0<s<1$ 的 $s$-集都是非正则的.
\end{corollary}

\textcite{marstrand1954a} 关于下角密度的下述结果，其证明在某些方面与\cref{thm:4.2} 的证明相仿.

\begin{theorem}\label{thm:4.4}
设 $\theta$ 为单位向量，$\phi<\frac{1}{2}\pi$.则若 $E$ 是满足 $0<s<1$ 的 $s$-集，在几乎所有的 $x\in E$ 处有 $D^s(E,x,\theta,\phi)=0$.
\end{theorem}

\begin{proof}
反设不然.则按通常方式可找到 $\rho_0$、$\alpha>0$ 与闭 $s$-集 $F\subset E$，使得当 $x\in F$ 且 $r\leq\rho_0$ 时
\begin{equation}\label{eq:4.3}
  \mathcal{H}^{s}\bigl(E\cap S_r(x,\theta,\phi)\bigr)>\alpha r^{s}.
\end{equation}
由\cref{cor:2.4,cor:2.5} 当然可以选取 $y\in F$ 使 $D^s(E\setminus F,y)=0$ 且 $\bar{D}^s(E,y)=c2^{-s}$，其中 $0<c<\infty$.于是，给定 $\varepsilon>0$，存在 $\rho_1\leq\rho_0$ 使当 $r\leq\rho_1$ 时
\begin{equation}\label{eq:4.4}
  \mathcal{H}^{s}\bigl((E\setminus F)\cap B_r(y)\bigr)<\varepsilon r^{s}
\end{equation}
且
\begin{equation}\label{eq:4.5}
  \mathcal{H}^{s}\bigl(E\cap B_r(y)\bigr)<(c+\varepsilon)r^{s}.
\end{equation}
另外，选取 $\rho\leq\frac{1}{2}\rho_1$ 使
\begin{equation}\label{eq:4.6}
  (c-\varepsilon)\rho^{s}<\mathcal{H}^{s}\bigl(E\cap B_\rho(y)\bigr).
\end{equation}

\begin{figure}[htbp]
  \centering
  \includegraphics[width=0.36\textwidth]{fig4-2.png}
  \caption{\cref{thm:4.4} 证明中的点 $x$ 的选取}
  \label{fig:4.2}
\end{figure}

设 $x$ 为 $F\cap B_\rho(y)$ 中使数量积 $x\cdot\theta$ 最大的点，见图~\ref{fig:4.2}.则
\[
  \bigl(F\cap B_\rho(y)\bigr)\cup\bigl(F\cap S_r(x,\theta,\phi)\bigr)\subset F\cap B_{\rho+r}(y),
\]
并且这个并除点 $x$ 外是不交的.于是
\[
  \mathcal{H}^{s}\bigl(F\cap B_\rho(y)\bigr)+\mathcal{H}^{s}\bigl(F\cap S_r(x,\theta,\phi)\bigr)\leq\mathcal{H}^{s}\bigl(F\cap B_{\rho+r}(y)\bigr),
\]
故
\[
  \mathcal{H}^{s}\bigl(E\cap B_\rho(y)\bigr)+\mathcal{H}^{s}\bigl(E\cap S_r(x,\theta,\phi)\bigr)
  \leq\mathcal{H}^{s}\bigl(E\cap B_{\rho+r}(y)\bigr)+2\mathcal{H}^{s}\bigl((E\setminus F)\cap B_{\rho+r}(y)\bigr).
\]
若 $r\leq\rho$，则 $\rho+r\leq\rho_1$，于是可用 \eqref{eq:4.3}--\eqref{eq:4.6} 得到
\[
  (c-\varepsilon)\rho^{s}+\alpha r^{s}<(c+\varepsilon)(\rho+r)^{s}+2\varepsilon(\rho+r)^{s}.
\]
令 $\gamma=r/\rho$，上式成为
\[
  c-\varepsilon+\alpha\gamma^{s}<(c+3\varepsilon)(1+\gamma)^{s},
\]
对 $0<\gamma<1$ 成立.这对一切 $\varepsilon>0$ 成立，故
\[
  c+\alpha\gamma^{s}\leq c(1+\gamma)^{s}\leq c+cs\gamma
\]
对 $0<\gamma<1$ 成立.这与 $\alpha$ 的正性矛盾.
\end{proof}

若取 $S_r\bigl(x,\theta,\frac{1}{2}\pi\bigr)$ 为开半圆来计算密度，则该结果对 $\phi=\frac{1}{2}\pi$ 仍然成立.Davies 曾观察到，这可作为 $\mathbb{R}$ 中单侧密度结果的直接 $n$ 维类似物得到；后者由 \textcite{besicovitch1968} 在他的倒数第二篇论文中给出了漂亮的证明.

$0<s<1$ 的 $s$-集的切线问题没有特别的意义，因为这些集合如此稀疏，以至于用线段作逼近的任何想法都变得相当无意义.

在一种意义下，切线在所有点处存在是可能的.设 $\Gamma$ 是 $\mathbb{R}^n$ 中由光滑双射 $\psi:\mathbb{R}\to\mathbb{R}^n$ 定义的光滑曲线，并设 $E$ 是 $\mathbb{R}$ 中的 $s$-集.容易看出 $\psi(E)$ 是 $\mathbb{R}^n$ 中的 $s$-集，而且对几乎所有的 $x\in\psi(E)$，$\Gamma$ 在 $x$ 处的切线就是 $\psi(E)$ 在 $x$ 处按 \eqref{eq:3.4} 意义的切线.

\section{\texorpdfstring{$s>1$ 的 $s$-集}{s>1 的 s-集}}

当 $s>1$ 时，\newterm{s-set-gt1}{$s$-集}{$s$-set (s > 1)}的密度与切线性质要复杂得多，我们把详细的论述限于平面的子集.

先处理切线问题.一种途径是重复 Besicovitch 对\cref{thm:3.29} 的证明，推出平面 $s$-集（$1<s<2$）在任意一对对顶角中上角密度之和几乎处处严格为正.Besicovitch 的证明依赖于所论集合与任何可求长曲线的交测度为零；由\cref{lem:3.2} 之后的注记，当 $s>1$ 时这当然成立.于是由切线的定义 \eqref{eq:3.4} 可得：满足 $1<s<2$ 的平面 $s$-集几乎处处没有切线.然而，更容易证明一个更强的结果，而且在涉及圆形密度的问题时，我们需要这种关于切向行为的更细致的信息.

遵循 \textcite{marstrand1954a}，我们定义弱切线.若 $s$-集 $E$ 满足 $D^s(x,E)>0$，且对每个 $\phi>0$，
\begin{equation}\label{eq:4.7}
  \varliminf_{r\to 0}r^{-s}\mathcal{H}^{s}\bigl(E\cap(B_r(x)\setminus S_r(x,\theta,\phi)\setminus S_r(x,-\theta,\phi))\bigr)=0,
\end{equation}
则称 $E$ 在 $x$ 处有方向 $\theta$ 的\newterm{weak tangent}{弱切线}{weak tangent}.（回忆 $S_r(x,\theta,\phi)$ 是由 $B_r(x)$ 中在 $x$ 处与从 $x$ 出发方向 $\theta$ 的半直线夹角至多为 $\phi$ 的那些点构成的扇形.）

与 \eqref{eq:3.4} 定义的切线不同，一个集合在一点处在原则上可以有许多弱切线.然而这种情形通常不会发生.

\begin{lemma}\label{lem:4.5}
设 $E$ 是 $\mathbb{R}^2$ 中满足 $1<s<2$ 的 $s$-集.则对几乎所有的 $x\in E$，
\[
  \bar{D}^{s}(E,x,\theta,\phi)\leq 4\cdot 10^{s}\phi^{s-1}
\]
对一切 $\theta$ 与一切 $\phi\leq\frac{1}{2}\pi$ 成立.
\end{lemma}

\begin{proof}
固定 $\rho>0$，令
\begin{equation}\label{eq:4.8}
  F=\{x\in E:\mathcal{H}^{s}\bigl(E\cap B_r(x)\bigr)<2^{s+1}r^{s}\ \text{对一切}\ r\leq\rho\}.
\end{equation}
取 $x\in F$ 及任意满足 $0<\phi\leq\frac{1}{2}\pi$ 的 $\theta$ 与 $\phi$.对 $i=1,2,\ldots$，令 $A_i$ 为圆环与扇形之交
\[
  A_i=S_{ir\phi}(x,\theta,\phi)\setminus S_{(i-1)r\phi}(x,\theta,\phi),
\]
则 $S_r(x,\theta,\phi)\subset\bigcup_{1}^{m}A_i\cup\{x\}$，其中整数 $m<2/\phi$.当 $r<\rho/20$ 时，每个 $A_i$ 的直径至多为 $10r\phi<\rho$，于是对每个含 $F$ 的点的 $A_i$ 应用 \eqref{eq:4.8} 并求和，
\[
  \mathcal{H}^{s}\bigl(F\cap S_r(x,\theta,\phi)\bigr)\leq 2\phi^{-1}2^{s+1}(10r\phi)^{s},
\]
即当 $r<\rho/20$ 时
\[
  (2r)^{-s}\mathcal{H}^{s}\bigl(F\cap S_r(x,\theta,\phi)\bigr)\leq 4\cdot 10^{s}\phi^{s-1}.
\]
于是在几乎所有的 $x\in F$ 处 $\bar{D}^s(F,x,\theta,\phi)\leq 4\cdot 10^{s}\phi^{s-1}$；又由\cref{cor:2.6}，在几乎所有的 $x\in F$ 处 $\bar{D}^s(E\setminus F,x)=0$，故引理结论在几乎所有的 $x\in F$ 处成立.由\cref{cor:2.5}，$E$ 的几乎所有的点对某个 $\rho>0$ 属于这样的 $F$，故结论得证.
\end{proof}

\begin{corollary}\label{cor:4.6}
若 $E$ 是 $\mathbb{R}^2$ 中满足 $1<s<2$ 的 $s$-集，则在 $E$ 的几乎所有的点处不存在弱切线.
\end{corollary}

\begin{proof}
由\cref{lem:4.5}，对几乎所有的 $x\in E$ 及一切 $\theta$ 与 $\phi<\frac{1}{2}\pi$，
\begin{align*}
  \varliminf_{r\to 0}(2r)^{-s}\mathcal{H}^{s}&\bigl(E\cap(B_r(x)\setminus S_r(x,\theta,\phi)\setminus S_r(x,-\theta,\phi))\bigr)\\
  &\geq\underline{D}^{s}(E,x)-\bar{D}^{s}(E,x,\theta,\phi)-\bar{D}^{s}(E,x,-\theta,\phi)
  \geq\underline{D}^{s}(E,x)-8\cdot 10^{s}\phi^{s-1}.
\end{align*}
在这样的点处，要么 $D^s(E,x)=0$，要么当 $\phi$ 充分小时 \eqref{eq:4.7} 不成立，故不存在弱切线.
\end{proof}

为证明 $1<s<2$ 时 $s$-集是非正则的，我们先考虑角密度而非圆形密度.

\begin{theorem}\label{thm:4.7}
设 $E$ 是 $\mathbb{R}^2$ 中满足 $1<s<2$ 的 $s$-集.则当 $\phi<\frac{1}{2}\pi$ 时，对几乎所有的 $x\in E$，下角密度 $D^s(E,x,\theta,\phi)$ 对某个 $\theta$ 为零.
\end{theorem}

\begin{proof}
固定正数 $\alpha$ 与 $\rho$，令
\begin{equation}\label{eq:4.9}
  F_0=\bigl\{x:\mathcal{H}^{s}\bigl(E\cap S_r(x,\theta,\phi)\bigr)>\alpha r^{s}
  \ \text{对一切}\ r\leq\rho\ \text{与一切}\ \theta\bigr\}.
\end{equation}
我们证明 $\mathcal{H}^s(F_0)=0$.假设不然，则由\cref{cor:2.5} 可找到 $\rho_1\leq\rho$ 与正测度集 $F\subset F_0$，使得当 $x\in F$ 且 $r\leq\rho_1$ 时
\begin{equation}\label{eq:4.10}
  \mathcal{H}^{s}\bigl(E\cap B_r(x)\bigr)<2^{s+1}r^{s}.
\end{equation}
由 $\mathcal{H}^s$ 的正则性还可设 $F$ 为闭集.

设 $y$ 为 $F$ 中使 $E\setminus F$ 的圆形密度为零的一点；由\cref{cor:2.6}，$F$ 的几乎任一点都可以.于是给定 $\varepsilon>0$，可找到 $\rho_2\leq\rho_1$ 使当 $r\leq\rho_2$ 时
\begin{equation}\label{eq:4.11}
  \mathcal{H}^{s}\bigl((E\setminus F)\cap B_r(y)\bigr)<\varepsilon r^{s}.
\end{equation}
现在我们在圆盘 $B_{\rho_2}(y)$ 内工作.首先证明 $B_{\frac{1}{2}\rho_2}(y)$ 中有相对远离 $F$ 的点.设对某个 $\gamma\leq\frac{1}{2}\rho_2$，$B_{\frac{1}{2}\rho_2}(y)$ 的所有点与 $F$ 的距离都不超过 $\gamma$.则若 $x\in B_{\frac{1}{2}\rho_2}(y)$，有 $F$ 的点 $z$ 位于 $B_\gamma(x)$ 内.由 \eqref{eq:4.9}，对任意 $\theta$，
\begin{equation}\label{eq:4.12}
  \alpha\gamma^{s}
  <\mathcal{H}^{s}\bigl(E\cap S_\gamma(z,\theta,\phi)\bigr)
  \leq\mathcal{H}^{s}\bigl(E\cap B_\gamma(z)\bigr)
  \leq\mathcal{H}^{s}\bigl(E\cap B_{2\gamma}(x)\bigr).
\end{equation}
若 $\gamma<\frac{1}{4}\rho_2$，则 $B_{\rho_2}(y)$ 含有 $\bigl(\rho_2/\gamma\bigr)^2/16$ 个圆心在 $B_{\frac{1}{2}\rho_2}(y)$ 中、半径为 $2\gamma$ 的不交圆盘.于是把这些圆盘上的 \eqref{eq:4.12} 求和，并由 \eqref{eq:4.10} 得
\[
  \bigl(\rho_2/\gamma\bigr)^{2}\alpha\gamma^{s}/16
  <\mathcal{H}^{s}\bigl(E\cap B_{\rho_2}(y)\bigr)
  <2^{s+1}\rho_2^{s},
\]
故 $\gamma>c\rho_2$，其中 $c$ 只依赖于 $\alpha$ 与 $s$.于是当 $\gamma\leq c\rho_2$ 时，存在含于 $B_{\rho_2}(y)$ 内且不含 $F$ 的点的半径为 $\gamma$ 的圆盘.从而可找到圆盘 $B_{\rho_3}(w)\subset B_{\rho_2}(y)$，其内部不含 $F$ 的点，但其边界含有 $F_0$ 的一点 $v$，且
\begin{equation}\label{eq:4.13}
  \rho_2\geq\rho_3\geq c\rho_2.
\end{equation}

\begin{figure}[htbp]
  \centering
  \includegraphics[width=0.36\textwidth]{fig4-3.png}
  \caption{\cref{thm:4.7} 证明中的圆盘与扇形}
  \label{fig:4.3}
\end{figure}

设 $\theta$ 为 $B_{\rho_3}(w)$ 在 $v$ 处的内法向方向，并令 $\rho_4$ 为 $B_{\rho_3}(w)$ 中过 $v$ 且与 $\theta$ 成角 $\phi$ 的弦长的一半（见图~\ref{fig:4.3}）.由初等几何，
\begin{equation}\label{eq:4.14}
  \rho_4=\rho_3\cos\phi.
\end{equation}
由于扇形 $S_{\rho_4}(v,\theta,\phi)$ 位于 $B_{\rho_3}(w)$ 内，它除 $v$ 外不含 $F$ 的点.于是由 \eqref{eq:4.11}、\eqref{eq:4.13} 与 \eqref{eq:4.14}（$c_1$ 只依赖于 $\phi,\alpha$ 与 $s$）得
\begin{equation}\label{eq:4.15}
  \mathcal{H}^{s}\bigl(E\cap S_{\rho_4}(v,\theta,\phi)\bigr)
  =\mathcal{H}^{s}\bigl((E\setminus F)\cap S_{\rho_4}(v,\theta,\phi)\bigr)
  \leq\mathcal{H}^{s}\bigl((E\setminus F)\cap B_{\rho_2}(y)\bigr)
  <\varepsilon\rho_2^{s}
  <\varepsilon c_1\rho_4^{s}.
\end{equation}
于是对任意 $\varepsilon>0$ 可找到 $v\in F_0$、满足 $0<\rho_4<\rho$ 的 $\rho_4$ 与 $\theta$ 使 \eqref{eq:4.15} 成立.这与 $F_0$ 的定义 \eqref{eq:4.9} 矛盾，故 $\mathcal{H}^s(F_0)=0$.注意到 $\alpha,\rho>0$ 任意，证明完成.
\end{proof}

\begin{corollary}\label{cor:4.8}
设 $E$ 是 $\mathbb{R}^2$ 中满足 $1<s<2$ 的 $s$-集.则在 $E$ 的几乎每一点处，下角密度 $D^s\bigl(E,x,\theta,\frac{1}{2}\pi\bigr)$ 对某个 $\theta$ 为零.
\end{corollary}

\begin{proof}
取递增趋于 $\frac{1}{2}\pi$ 的序列 $\{\phi_i\}$.由\cref{thm:4.7}，对几乎所有的 $x\in E$ 可找到方向序列 $\{\theta_i\}$ 使 $D^s(E,x,\theta_i,\phi_i)=0$ 对每个 $i$ 成立.抽取收敛子列，可设 $\theta_i\to\theta$.于是对任意 $\phi<\frac{1}{2}\pi$ 有 $D^s(E,x,\theta,\phi)=0$.因此，只要 $x$ 还不是\cref{lem:4.5} 的例外点，便有 $D^s\bigl(E,x,\theta,\frac{1}{2}\pi\bigr)=0$.
\end{proof}

下一步我们考察 $\mathbb{R}^2$ 中满足 $1<s<2$ 的正则 $s$-集的下角密度，作为证明这类集合不存在的一个步骤.\textcite{marstrand1954a} 的下述精巧论证大致上说：若这样的集合在过其一点的一条直线的一侧稀疏，则它在另一侧也必定稀疏.

\begin{lemma}\label{lem:4.9}
设 $E$ 是 $\mathbb{R}^2$ 中的 $s$-集，$1<s<2$.设 $x$ 是 $E$ 的上凸密度等于 $1$ 的正则点，并设对某个 $\theta$ 有 $D^s\bigl(E,x,-\theta,\frac{1}{2}\pi\bigr)=0$.则 $E$ 在 $x$ 处有垂直于 $\theta$ 的弱切线.
\end{lemma}

\begin{proof}
由于 $D^s(E,x)=1$ 且 $\bar{D}^s_{\mathrm{c}}(E,x)=1$，给定任意 $\eta>0$，可找到任意小的 $\rho$ 值使
\begin{equation}\label{eq:4.16}
  \mathcal{H}^{s}\bigl(E\cap B_r(x)\bigr)>2^{s}r^{s}(1-\eta)\qquad(r\leq\rho),
\end{equation}
\begin{equation}\label{eq:4.17}
  \mathcal{H}^{s}(E\cap U)<(1+\eta)|U|^{s}\qquad(x\in U\ \text{且}\ 0<|U|\leq 2\rho),
\end{equation}
并且
\begin{equation}\label{eq:4.18}
  \mathcal{H}^{s}\Bigl(E\cap S_\rho\bigl(x,-\theta,\tfrac{1}{2}\pi\bigr)\Bigr)<2^{s}\eta\rho^{s}.
\end{equation}
取 $0<\phi<\frac{1}{2}\pi$，设 $L$ 为过 $x$ 且垂直于 $\theta$ 的直线，$M$ 与 $M'$ 为从 $x$ 出发与 $\theta$ 成角 $\phi$ 的半射线.对固定的正整数 $m$，归纳地构造以 $x$ 为圆心、基于 $L$ 的 $m+1$ 个半径为 $r_i$ 的半圆 $S_{r_i}$，其中 $\rho=r_0>r_1>\cdots>r_m$.每个半圆 $S_{r_i}$ 的端点 $y_i$ 与 $y_i'$ 在 $L$ 上，且与 $M$ 和 $M'$ 分别交于 $z_i$ 与 $z_i'$.设 $S_{r_i}$ 已经构造，则如下指定 $S_{r_{i+1}}$：取 $y_{i+1}$ 为 $[x,y_i]$ 上使其到 $y_i'$ 的距离等于它到 $y_i$ 与到 $z_i'$ 的距离之和的那一点.（一个直截了当的连续性论证表明这样的点存在.）以 $y_{i+1}$ 为圆心、过 $z_i'$ 作弧交 $L$ 于 $y_{i+1}'$.由对称性，以 $y_{i+1}'$ 为圆心、过 $z_i$ 的弧交 $L$ 于 $y_{i+1}$（见图~\ref{fig:4.4}）.记 $S_{r_i}$ 被这两条弧截出的（阴影）凸部分为 $U_i$.则
\begin{equation}\label{eq:4.19}
  |U_i|=2r_{i+1}.
\end{equation}

\begin{figure}[htbp]
  \centering
  \includegraphics[width=0.44\textwidth]{fig4-4.png}
  \caption{\cref{lem:4.9} 证明中的半圆与区域 $U_i$}
  \label{fig:4.4}
\end{figure}

我们估计 $E$ 介于由 $M$ 与 $M'$ 界定的两个相邻扇形之间的部分的测度.由
\[
  S_{r_i}(x,\theta,\phi)\setminus S_{r_{i+1}}(x,\theta,\phi)
  \subset U_i\setminus S_{r_{i+1}}\bigl(x,\theta,\tfrac{1}{2}\pi\bigr)
  \subset U_i\cup S_\rho\bigl(x,-\theta,\tfrac{1}{2}\pi\bigr)\setminus B_{r_{i+1}}(x),
\]
得
\begin{align*}
  \mathcal{H}^{s}&\bigl(E\cap S_{r_i}(x,\theta,\phi)\bigr)-\mathcal{H}^{s}\bigl(E\cap S_{r_{i+1}}(x,\theta,\phi)\bigr)\\
  &\leq\mathcal{H}^{s}(E\cap U_i)+\mathcal{H}^{s}\Bigl(E\cap S_\rho\bigl(x,-\theta,\tfrac{1}{2}\pi\bigr)\Bigr)
    -\mathcal{H}^{s}\bigl(E\cap B_{r_{i+1}}(x)\bigr)\\
  &<(1+\eta)|U_i|^{s}+2^{s}\eta\rho^{s}-2^{s}r_{i+1}^{s}(1-\eta)\\
  &\leq 2^{s+1}\eta r_{i+1}^{s}+2^{s}\eta\rho^{s}
  \leq 2^{s+2}\eta\rho^{s},
\end{align*}
这里用了 \eqref{eq:4.16}--\eqref{eq:4.19}.把这一估计对全部 $m$ 个扇形求和，由 \eqref{eq:4.17} 得
\begin{align*}
  \mathcal{H}^{s}\bigl(E\cap S_\rho(x,\theta,\phi)\bigr)
  &<2^{s+2}\eta\rho^{s}m+\mathcal{H}^{s}\bigl(E\cap S_{r_m}(x,\theta,\phi)\bigr)\\
  &<2^{s+2}\eta\rho^{s}m+2^{s}(1+\eta)r_m^{s}.
\end{align*}
由构造，$r_m/\rho$ 只依赖于 $\phi$ 与 $m$，且当 $m\to\infty$ 时该比值趋于零.于是给定 $\varepsilon>0$，可找到与 $\eta$ 无关的 $m$ 使
\[
  \mathcal{H}^{s}\bigl(E\cap S_\rho(x,\theta,\phi)\bigr)<\rho^{s}\bigl(2^{s+2}\eta m+\varepsilon(1+\eta)\bigr).
\]
这对任意小的 $\rho$ 值与任意 $\eta>0$ 成立，故当 $\phi<\frac{1}{2}\pi$ 时 $D^s(E,x,\theta,\phi)=0$.由于 $D^s\bigl(E,x,-\theta,\frac{1}{2}\pi\bigr)=0$，由 \eqref{eq:4.7} 可得 $E$ 在 $x$ 处有垂直于 $\theta$ 的方向的弱切线.
\end{proof}

\begin{corollary}\label{cor:4.10}
设 $E$ 是 $\mathbb{R}^2$ 中满足 $1<s<2$ 的 $s$-集.则 $E$ 是非正则的.
\end{corollary}

\begin{proof}
由\cref{thm:2.3}，在 $E$ 的几乎所有的点处上凸密度等于 $1$；由\cref{cor:4.8}，对某个 $\theta$ 有 $D^s\bigl(E,x,\theta,\frac{1}{2}\pi\bigr)=0$.于是由\cref{lem:4.9}，集合 $E$ 在其几乎所有的正则点处有弱切线.这与\cref{cor:4.6} 不相容，除非正则点之集测度为零；此时 $E$ 是非正则 $s$-集.
\end{proof}

在发表上述结果后不久，\textcite{marstrand1955} 把该方法推广，证明了更强的事实：$\mathbb{R}^2$ 中（$1<s<2$）$s$-集的密度几乎处处不存在.

为完整起见，我们把本章的主要结果汇集在一起.

\begin{theorem}\label{thm:4.11}
平面 $s$-集除非 $s$ 是整数，必是非正则的.
\end{theorem}

\begin{proof}
合并\cref{cor:4.3} 与\cref{cor:4.10}.
\end{proof}

\section{高维集合}

主要的高维结果与\cref{thm:4.11} 类似.

\begin{theorem}\label{thm:4.12}
$\mathbb{R}^n$ 中的 $s$-集除非 $s$ 是整数，必是非正则的.
\end{theorem}

$0<s<1$ 的情形已在\cref{thm:4.2} 中证明，\cref{cor:4.10} 的证明稍作调整即可涵盖 $n-1<s<n$ 的情形.然而对中间的 $s$ 值，需要实质性的新思想，该定理最终由 \textcite{marstrand1964} 证明.事实上，Marstrand 的工作远不止于此：它不仅适用于更一般的测度，还证明了当 $s$ 非整数时密度几乎处处不存在.其证明包含对正则集的弱切向性质深入得多的研究.

\begin{exercises}
\item 调整\cref{thm:4.2} 的证明，证明若 $E$ 是满足 $0<s<1$ 的 $s$-集，则对几乎所有的 $x$ 有 $D^s(E,x)\leq\bigl(1+2^{s/(s-1)}\bigr)^{s-1}$.
\item 设 $E$ 是 $s$-集（$0<s<1$）.由\cref{thm:4.4} 推出：对任意单位向量 $\theta$，只要密度用开半圆或开半球计算，则在几乎所有的 $x\in E$ 处\newterm{density-hemispherical}{下半球密度}{density, hemispherical} $D^s\bigl(E,x,\theta,\frac{1}{2}\pi\bigr)=0$.举一个例子说明对闭半圆这 不成立.
\item 证明\cref{thm:4.12} 的如下特殊情形：设 $E$ 是 $\mathbb{R}^3$ 中位于光滑（比如说无穷可微）曲面上的 $s$-集.则除非 $s$ 是整数，$E$ 是非正则的.
\end{exercises}

% !TeX root = ../main.tex
% 第 5 章 可比较的网测度 (对应原书 pp. 64--74)

\chapter{可比较的网测度}\label{chap:5}

\section{网测度的构造}

\newterm{comparable-net-measure}{可比较网测度}{comparable net measure}是研究 Hausdorff 测度的极为有用的工具.\newterm{net-measure}{网测度}{net measure}的性质比 Hausdorff 测度方便得多，但仍可构造得对许多用途而言与之等价.网测度最先由 \textcite{besicovitch1952} 使用，用以证明 $\mathcal{H}^s$-测度无穷的闭集包含测度为正但有限的子集；后来 \textcite{marstrand1954b} 在研究集合的 Descartes 积的 Hausdorff 测度时使用了它.本章我们特别关注这两个应用.

这里我们把注意力限于 Euclid 空间上的一种基本形式的网测度.更广泛的讨论见 \textcite{rogers1970} 的 2.7~节，那里描述了 Davies 与 Rogers 的工作.

网测度的构造方式与 Hausdorff 测度类似，只是在定义中使用一个受限的覆盖集合类 $\mathcal{N}$ 而非全体集合.类 $\mathcal{N}$ 被选为一个集合的``网''，其本质性质是：若 $U_1,U_2\in\mathcal{N}$，则 $U_1\cap U_2=\emptyset$、$U_1\subset U_2$、$U_2\subset U_1$ 三者必居其一.此外，我们假定 $\mathcal{N}$ 的每个集合只含于有限多个其他集合之中.特别地，给定 $\mathcal{N}$ 中集合的任意族，可以通过去掉那些含于其他集合之中的集合，找到具有同样并集的不交子族.正如我们将看到的，可以构造与 Hausdorff 测度``可比较''的网测度，即两个测度之比上下都有界的网测度.

为构造 $\mathbb{R}^n$ 上的 $s$ 维网测度 $\mathcal{M}^s$，令 $\mathcal{N}$ 为所有 $n$ 维\newterm{binary-cubes}{半开二进方体}{binary cubes, half open}之族，即形如
\[
  [2^{-k}m_1,2^{-k}(m_1+1))\times[2^{-k}m_2,2^{-k}(m_2+1))\times\cdots\times[2^{-k}m_n,2^{-k}(m_n+1))
\]
的集合，其中 $k$ 为非负整数，$m_1,\ldots,m_n$ 为整数.（若 $n=1$ 或 $2$，网 $\mathcal{N}$ 由半开二进区间或正方形组成.）设 $E\subset\mathbb{R}^n$，$\delta>0$，定义
\begin{equation}\label{eq:5.1}
  \mathcal{M}^{s}_{\delta}(E)=\inf\sum_{i=1}^{\infty}|S_i|^{s},
\end{equation}
其中下确界取遍 $E$ 的所有由 $\mathcal{N}$ 中集合 $\{S_i\}$ 构成的可数 $\delta$-覆盖.由网的性质，只需考虑 $E$ 的由 $\mathcal{N}$ 中两两不交的集合构成的覆盖.于是 $\mathcal{M}^s_{\delta}(E)$ 是 $\mathbb{R}^n$ 上在 $\mathbb{R}^n$ 的有界子集上取有限值的外测度.令
\[
  \mathcal{M}^{s}(E)=\lim_{\delta\to 0}\mathcal{M}^{s}_{\delta}(E)=\sup_{\delta>0}\mathcal{M}^{s}_{\delta}(E),
\]
便得到 $\mathbb{R}^n$ 上的度量外测度 $\mathcal{M}^s$.由\cref{thm:1.5}，Borel 集是 $\mathcal{M}^s$-可测的（事实上 Souslin 集亦然）.与\cref{thm:1.6} 的证明相仿可证 $\mathcal{M}^s$ 是正则外测度.

我们来证明 $\mathcal{M}^s$ 与 $\mathcal{H}^s$ 确实是可比较的测度；这本质上是 \textcite{besicovitch1952} 的工作.

\begin{theorem}\label{thm:5.1}
存在只依赖于维数 $n$ 的常数 $b_n$，使得对每个 $E\subset\mathbb{R}^n$，当 $0<\delta<1$ 时
\begin{equation}\label{eq:5.2}
  \mathcal{H}^{s}_{\delta}(E)\leq\mathcal{M}^{s}_{\delta}(E)\leq b_n\mathcal{H}^{s}_{\delta}(E),
\end{equation}
且
\begin{equation}\label{eq:5.3}
  \mathcal{H}^{s}(E)\leq\mathcal{M}^{s}(E)\leq b_n\mathcal{H}^{s}(E).
\end{equation}
\end{theorem}

\begin{proof}
由外测度的定义立得 $\mathcal{H}^s_{\delta}(E)\leq\mathcal{M}^s_{\delta}(E)$.（Hausdorff 测度的下确界取遍更大的覆盖集合类.）

设 $U$ 为满足 $0<|U|\leq\delta$ 的任意集合，令 $k$ 为满足
\begin{equation}\label{eq:5.4}
  2^{-k-1}\leq|U|<2^{-k}
\end{equation}
的整数，并设 $S$ 为与 $U$ 相交的边长 $2^{-k}$ 的二进方体.则 $U$ 含于由 $S$ 及其直接邻域组成的 $3^n$ 个边长 $2^{-k}$、直径 $2^{-k}n^{1/2}$ 的二进方体之族.把每个这样的方体再剖分为 $2^{n^2}$ 个更小的方体，则 $U$ 含于 $b_n=3^n2^{n^2}$ 个直径为
\begin{equation}\label{eq:5.5}
  2^{-k}n^{1/2}2^{-n}\leq 2^{1-n}n^{1/2}|U|\leq|U|\leq\delta
\end{equation}
的二进方体之中，这里用了 \eqref{eq:5.4}.现在设 $\{U_i\}$ 是 $E$ 的由任意集合构成的 $\delta$-覆盖.对每个 $i$ 有 $U_i\subset\bigcup_{j=1}^{b_n}S_{ij}$，其中 $\{S_{ij}\}_{j=1}^{b_n}$ 是由 $b_n$ 个直径至多为 $|U_i|\leq\delta$ 的方体构成的族.

于是 $E\subset\bigcup_{i}\bigcup_{j}S_{ij}$，且
\[
  \sum_{i=1}^{\infty}\sum_{j=1}^{b_n}|S_{ij}|^{s}\leq b_n\sum_{i=1}^{\infty}|U_i|^{s},
\]
故 \eqref{eq:5.2} 右端的不等式由 $\mathcal{M}^s_{\delta}$ 与 $\mathcal{H}^s_{\delta}$ 的定义得出.

令 $\delta\to 0$ 即得不等式 \eqref{eq:5.3}.
\end{proof}

本节最后证明两条关于集合的单调序列的技术性引理，本章后面将会用到.

\begin{lemma}\label{lem:5.2}
设 $\{E_j\}$ 是 $\mathbb{R}^n$ 的紧子集的递减序列.则对任意 $\delta>0$，
\begin{equation}\label{eq:5.6}
  2^{s}\mathcal{H}^{s}_{\delta}\Bigl(\lim_{j\to\infty}E_j\Bigr)\geq\lim_{j\to\infty}\mathcal{H}^{s}_{2\delta}(E_j).
\end{equation}
\end{lemma}

\begin{proof}
设 $\{U_i\}$ 是 $\lim E_j$ 的任意 $\delta$-覆盖.对每个 $i$，令 $V_i$ 为包含 $U_i$ 的开集，$|V_i|\leq 2|U_i|$，并令 $V$ 表示开集 $\bigcup_{i}V_i$.我们断言对某个整数 $j$ 有 $E_j\subset V$.否则 $\{E_j\setminus V\}$ 是非空紧集的递减序列，由紧性的一个初等推论，它有非空的极限集（$\lim E_j\setminus V$）.于是
\[
  2^{s}\sum|U_i|^{s}\geq\sum|V_i|^{s}\geq\mathcal{H}^{s}_{2\delta}(E_j)\geq\lim_{j\to\infty}\mathcal{H}^{s}_{2\delta}(E_j),
\]
考虑 $\lim E_j$ 的所有此类 $\delta$-覆盖即得 \eqref{eq:5.6}.
\end{proof}

若 $\{E_j\}$ 是 $\mathbb{R}^n$ 的子集的任意递增序列，则由\cref{lem:1.3} 有 $\lim_{j\to\infty}\mathcal{M}^s(E_j)=\mathcal{M}^s\bigl(\lim_{j\to\infty}E_j\bigr)$.然而这一结论对外测度 $\mathcal{M}^s_{\delta}$ 也成立.它是\newterm{increasing-sets-lemma}{``递增集引理''}{increasing sets lemma}的一个特殊情形；该引理在 Hausdorff 测度理论中相当重要，\textcite{rogers1970}（2.6~节）与 \textcite{davies1970} 及其所引文献中有详细讨论.我们需要如下特殊情形.

\begin{lemma}\label{lem:5.3}
设 $\{E_j\}$ 是 $\mathbb{R}^n$ 的子集的递增序列，且每个 $E_j$ 都是二进方体的有限并.则
\[
  \mathcal{M}^{s}_{\delta}\Bigl(\lim_{j\to\infty}E_j\Bigr)=\lim_{j\to\infty}\mathcal{M}^{s}_{\delta}(E_j).
\]
\end{lemma}

\begin{proof}
记 $E=\lim_{j\to\infty}E_j=\bigcup_jE_j$.由于对一切 $j$ 有 $\mathcal{M}^s_{\delta}(E_j)\leq\mathcal{M}^s_{\delta}(E)$，只需在右端有限的假设下证明 $\mathcal{M}^s_{\delta}(E)\leq\lim_{j\to\infty}\mathcal{M}^s_{\delta}(E_j)$.由于 $E_j$ 是二进方体的有限并，$\mathcal{M}^s_{\delta}(E_j)$ 定义中的下确界可由某个由二进方体构成的有限不交族取得.设对每个 $j$，$\mathcal{S}_j$ 是这样一个 $E_j$ 的有限 $\delta$-覆盖，使
\begin{equation}\label{eq:5.7}
  \sum_{C\in\mathcal{S}_j}|C|^{s}=\mathcal{M}^{s}_{\delta}(E_j),
\end{equation}
并设 $\mathcal{S}_j$ 是具有该性质的数值上最小的方体族之一.若 $S\in\mathcal{S}_j$，则 $S$ 必含 $E_j$ 的一点.该点也属于 $E_{j+1}$，从而属于某个 $T\in\mathcal{S}_{j+1}$.由网的性质，要么 $S\subset T$，要么 $T\subset S$.若 $T$ 是 $S$ 的真子集，则要么把 $S$ 换成 $\mathcal{S}_{j+1}$ 中含于 $S$ 的那些方体以缩小 $\sum_{C\in\mathcal{S}_j}|C|^s$，要么把 $\mathcal{S}_{j+1}$ 中含于 $S$ 的方体换成单个方体 $S$，以缩小 $\sum_{C\in\mathcal{S}_{j+1}}|C|^s$ 或该和式的项数.这两种情形都与 $\mathcal{S}_j$ 或 $\mathcal{S}_{j+1}$ 的极小性矛盾.于是我们得出：若 $S\in\mathcal{S}_j$，则存在 $T\in\mathcal{S}_{j+1}$ 使 $S\subset T$.

设 $\{C_1,C_2,\ldots\}$ 是从 $\bigcup_{j=1}^{\infty}\mathcal{S}_j$ 中去掉所有含于该族其他方体之中的方体后所得的方体集.则对每个 $j$ 有 $E_j\subset\bigcup_{i=1}^{\infty}C_i$，从而 $E\subset\bigcup_{i=1}^{\infty}C_i$，且
\begin{equation}\label{eq:5.8}
  \mathcal{M}^{s}_{\delta}(E)\leq\sum_{i=1}^{\infty}|C_i|^{s}.
\end{equation}
但由前一段的结论，每个方体 $C_i$ 属于一切充分大的 $j$ 的 $\mathcal{S}_j$.于是给定 $k$，可找到 $j(k)$ 使方体 $C_1,\ldots,C_k$ 都在 $\mathcal{S}_{j(k)}$ 中.利用 \eqref{eq:5.8} 与 \eqref{eq:5.7},
\[
  \mathcal{M}^{s}_{\delta}(E)
  \leq\lim_{k\to\infty}\sum_{i=1}^{k}|C_i|^{s}
  \leq\lim_{k\to\infty}\sum_{C\in\mathcal{S}_{j(k)}}|C|^{s}
  =\lim_{k\to\infty}\mathcal{M}^{s}_{\delta}\bigl(E_{j(k)}\bigr)
  \leq\lim_{j\to\infty}\mathcal{M}^{s}_{\delta}(E_j),
\]
如所需.
\end{proof}

\section{有限测度的子集}

\label{idx:finite-measure-subsets}
\addterm{有限测度}{finite measure, subsets of}{\pageref{idx:finite-measure-subsets}}
本节结果的一个推论是：$s$-集（即具有正的有限 $\mathcal{H}^s$-测度的集合）极为丰富.事实上，给定任何拓扑上像样的、满足 $\mathcal{H}^s(E)>0$ 的集合 $E$，可以找到 $E$ 的紧子集 $F$ 使 $0<\mathcal{H}^s(F)<\infty$.若 $E$ 不是 $\sigma$-有限的 $\mathcal{H}^s$-测度集，这绝非平凡.（$\sigma$-有限集是指可以表示为有限测度集的可数并的集合.）这类结果最先由 \textcite{besicovitch1952} 对闭集得到，并立即被 \textcite{davies1952b} 推广到 Souslin 集.

\begin{theorem}\label{thm:5.4}
设 $E$ 是 $\mathbb{R}^n$ 中满足 $\mathcal{H}^s(E)=\infty$ 的闭子集.

(a) 设 $c$ 为正数.则存在 $E$ 的紧子集 $F$ 使 $\mathcal{H}^s(F)=c$.

(b) 存在 $E$ 的紧子集 $F$ 使 $\mathcal{H}^s(F)>0$，且对某个常数 $b$ 有
\[
  \mathcal{H}^{s}\bigl(B_r(x)\cap F\bigr)\leq br^{s}\qquad(x\in\mathbb{R}^n,\ r\leq 1).
\]
\end{theorem}

\begin{proof}
为使记号相对简单，我们对 $n=1$ 给出证明；高维情形步骤完全相同，只是把二进区间换成二进方体.不妨设 $E$ 有界.由于 $\mathcal{M}^s(E)=\infty$，可找到整数 $m$ 使
\[
  \mathcal{M}^{s}_{2^{-m}}(E)\geq 2^{s}cb_1,
\]
其中 $b_1$ 是\cref{thm:5.1} 中的常数.

归纳地定义 $E$ 的闭子集的递减序列 $\{E_k\}_{m}^{\infty}$.令 $E_m=E$.对 $k\geq m$，通过对每个长度为 $2^{-k}$ 的二进区间指定交来定义 $E_{k+1}$.设 $I$ 是这样的区间.若 $\mathcal{M}^{s}_{2^{-(k+1)}}(E_k\cap I)\leq 2^{-sk}$，就简单地取 $E_{k+1}\cap I=E_k\cap I$.于是
\begin{equation}\label{eq:5.9}
  \mathcal{M}^{s}_{2^{-(k+1)}}\bigl(E_{k+1}\cap I\bigr)=\mathcal{M}^{s}_{2^{-k}}\bigl(E_k\cap I\bigr)\leq 2^{-sk},
\end{equation}
因为在计算 $\mathcal{M}^s_{2^{-k}}$ 时用 $I$ 作为覆盖区间所给出的估计至少不小于用更短区间所得者.另一方面，若 $\mathcal{M}^{s}_{2^{-(k+1)}}(E_k\cap I)>2^{-sk}$，则选取 $E_{k+1}\cap I$ 为 $E_k\cap I$ 的满足 $\mathcal{M}^{s}_{2^{-(k+1)}}(E_{k+1}\cap I)=2^{-sk}$ 的紧子集.这样的子集存在，因为 $\mathcal{M}^{s}_{2^{-(k+1)}}\bigl(E_{k+1}\cap I\cap(-\infty,u]\bigr)$ 关于 $u$ 连续.由于 $\mathcal{M}^{s}_{2^{-k}}(E_k\cap I)=2^{-sk}$，\eqref{eq:5.9} 仍然成立.现在把 \eqref{eq:5.9} 对所有长度为 $2^{-k}$ 的二进区间 $I$ 求和.用于计算 $\mathcal{M}^s_{2^{-k}}$ 与 $\mathcal{M}^s_{2^{-(k+1)}}$ 的覆盖区间必含于某个这样的 $I$ 中，故得
\[
  \mathcal{M}^{s}_{2^{-(k+1)}}(E_{k+1})=\mathcal{M}^{s}_{2^{-k}}(E_k)\qquad(k\geq m).
\]
迭代得
\begin{equation}\label{eq:5.10}
  \mathcal{M}^{s}_{2^{-k}}(E_k)=\mathcal{M}^{s}_{2^{-m}}(E_m)\qquad(k\geq m).
\end{equation}
设 $I$ 是长度为 $2^{-k}$ 的二进区间.若 $m\leq k<r$，则 $E_r\subset E_{k+1}$，故由 \eqref{eq:5.9},
\[
  \mathcal{M}^{s}_{2^{-(k+1)}}(E_r\cap I)\leq\mathcal{M}^{s}_{2^{-(k+1)}}(E_{k+1}\cap I)\leq 2^{-sk}.
\]
于是，在计算 $\mathcal{M}^s_{2^{-k}}(E_r)$ 时，任何长度为 $2^{-k}$ 的区间 $I$ 都可以换成至多 $2^{-(k+1)}$ 长度的区间而不增大下确界的值.故
\[
  \mathcal{M}^{s}_{2^{-(k+1)}}(E_r)=\mathcal{M}^{s}_{2^{-k}}(E_r)\qquad(m\leq k<r),
\]
从而迭代并结合 \eqref{eq:5.10} 得
\begin{equation}\label{eq:5.11}
  \mathcal{M}^{s}_{2^{-m}}(E_r)=\mathcal{M}^{s}_{2^{-r}}(E_r)=\mathcal{M}^{s}_{2^{-m}}(E_m)\qquad(r\geq m).
\end{equation}
令 $F$ 为紧集 $F=\bigcap_{m}^{\infty}E_k$.于是利用\cref{thm:5.1} 与 \eqref{eq:5.10},
\begin{equation}\label{eq:5.12}
  \mathcal{H}^{s}(F)\leq\mathcal{M}^{s}(F)=\lim_{k\to\infty}\mathcal{M}^{s}_{2^{-k}}(F)
  \leq\varliminf_{k\to\infty}\mathcal{M}^{s}_{2^{-k}}(E_k)=\mathcal{M}^{s}_{2^{-m}}(E_m)<\infty.
\end{equation}
另一方面，利用\cref{lem:5.2}、\cref{thm:5.1} 与 \eqref{eq:5.11},
\[
  2^{s}\mathcal{H}^{s}(F)\geq 2^{s}\mathcal{H}^{s}_{2^{-(m+1)}}(F)\geq\lim_{k\to\infty}\mathcal{H}^{s}_{2^{-m}}(E_k)
  \geq b_1^{-1}\lim_{k\to\infty}\mathcal{M}^{s}_{2^{-m}}(E_k)
  =b_1^{-1}\mathcal{M}^{s}_{2^{-m}}(E_m)\geq 2^{s}c.
\]
合并这些不等式得
\[
  c\leq\mathcal{H}^{s}(F)<\infty.
\]
若左端是严格不等式，则利用测度 $\mathcal{H}^s$ 的上、下连续性（\cref{thm:1.1}），形如 $F\cap[u,\infty)$ 的集合将有测度恰好为 $c$.

上述方法的一个微小变体证明 (b).设 $J$ 是长度为 $2^{-r}$ 的闭二进区间，$r\geq m$.若把 \eqref{eq:5.9} 对位于 $J$ 内的长度为 $2^{-k}$ 的二进子区间 $I$ 求和（$k\geq r$）并照前进行，则在 \eqref{eq:5.10} 的位置得到
\[
  \mathcal{M}^{s}_{2^{-k}}(E_k\cap J)=\mathcal{M}^{s}_{2^{-m}}(E_m\cap J)\qquad(k\geq r\geq m),
\]
从而在 \eqref{eq:5.12} 的位置得到
\[
  \mathcal{H}^{s}(F\cap J)\leq\mathcal{M}^{s}_{2^{-r}}(E_r\cap J)\leq\mathcal{M}^{s}_{2^{-r}}(J)\leq|J|^{s}.
\]
任何区间 $J_0$ 都可被至多六个长度至多为 $|J_0|$ 的相邻二进区间覆盖，由此即得结论.(b) 的完整细节留给读者.
\end{proof}

鉴于\cref{thm:5.4}，一个显然要问的问题是：这类结果对更一般的集合类是否成立.Besicovitch 在其原始文章中把该定理推广到了 $F_{\sigma\delta\sigma}$-集.\textcite{davies1952b} 利用下述中间结果把这一工作推进到其自然的终点，即对一切 Souslin 集证明该定理.

\begin{theorem}\label{thm:5.5}
设 $E$ 是 $\mathbb{R}^n$ 中满足 $\mathcal{H}^s(E)=\infty$ 的任意 Souslin 子集.则 $E$ 包含一个 $\mathcal{H}^s$-测度无穷的闭子集.
\end{theorem}

\begin{proof}
见 \textcite{davies1952b} 或 \textcite{rogers1970}.
\end{proof}

与\cref{thm:5.4} 结合可得：

\begin{theorem}\label{thm:5.6}
设 $E$ 是 $\mathbb{R}^n$ 中满足 $\mathcal{H}^s(E)=\infty$ 的 Souslin 子集.

(a) 对任意正常数 $c$，存在 $E$ 的紧子集 $F$ 使 $\mathcal{H}^s(F)=c$.

(b) 存在 $E$ 的紧子集 $F$ 使 $\mathcal{H}^s(F)>0$，且对某个常数 $b$ 有
\[
  \mathcal{H}^{s}\bigl(B_r(x)\cap F\bigr)\leq br^{s}\qquad(x\in\mathbb{R}^n,\ r\leq 1).
\]
\end{theorem}

本节最后提到一个更具病态性质的结果：若 $E$ 是 $\mathcal{H}^s$-测度非 $\sigma$-有限的紧集，则 $E$ 大到可以分解为连续统多个 $\mathcal{H}^s$-测度非 $\sigma$-有限的子集.\textcite{davies1968} 与 \textcite{rogers1970}（2.8 节）提供了这一思想及相关思想的进一步细节.

\section{集合的 Descartes 积}

为记号简便，本节始终把注意力限于平面的子集，尽管这些工作可以毫无困难地推广到更高维.这里 $(x,y)$ 表示 $\mathbb{R}^2$ 中的 Descartes 坐标.若 $E$ 是 $\mathbb{R}^2$ 的任意子集，用 $E_x$ 表示由 $E$ 中第一坐标等于 $x$ 的那些点构成的截口.

若 $E$ 是 $\mathbb{R}^2$ 的勒贝格可测子集，\newterm{fubinis-theorem}{Fubini 定理}{Fubini's theorem}的一个直接推论是：$E_x$ 对几乎所有的 $x$ 是 $\mathcal{L}^1$-可测的，且
\[
  \mathcal{L}^{2}(E)=\int\mathcal{L}^{1}(E_x)\,\mathrm{d}\mathscr{L}^{1}(x),
\]
其中 $\mathscr{L}^1$ 与 $\mathcal{L}^2$ 分别是 $1$ 维与 $2$ 维勒贝格测度.自然会问：把勒贝格测度换成分数维数的 Hausdorff 测度后，是否有任何类似的结果.例如，若 $A$ 与 $B$ 是两条坐标轴的子集，$\mathcal{H}^{s+t}(A\times B)$ 与 $\mathcal{H}^{s}(A)\mathcal{H}^{t}(B)$ 是否相等，至少相差一个常数倍？事实上，正如后面的例子所示，这类结果远非成立.一种初等的证明尝试或许是：用长度为 $\delta$ 的区间覆盖 $A$ 与 $B$，诱导出 $A\times B$ 的由边长 $\delta$ 的正方形构成的覆盖，从而估计各个 Hausdorff 测度.遗憾的是，使 $A$ 的覆盖``有效''的区间长度可能使 $B$ 的覆盖``无效''，使这类估计毫无用处.（当然，乘积测度 $\mathcal{H}^s\times\mathcal{H}^t$ 未必是 Hausdorff 测度.）不过，得到如下形式的不等式是可能的：
\[
  \mathcal{H}^{s+t}(A\times B)\geq b\,\mathcal{H}^{s}(A)\mathcal{H}^{t}(B).
\]
\textcite{besicovitchmoran1945} 在一定条件下证明了这一点；其他基于密度思想的早期证明由 Moran（\citeyear{moran1946,moran1949}）、\textcite{freilich1950}、\textcite{eggleston1950a} 与 \textcite{eggleston1953b} 给出；另见 \textcite{eggleston1950b}.然而这些证明相当技术化，而且对 $A$ 与 $B$ 附加了限制.下面这个使用可比较网测度的非常一般而漂亮的方法归功于 \textcite{marstrand1954b}.

证明依赖于下述组合引理.

\begin{lemma}\label{lem:5.7}
设 $A$ 是 $\mathbb{R}$ 的任意子集，$\{I_i\}$ 是 $A$ 的由二进区间构成的可数 $\delta$-覆盖，$\{a_i\}$ 是正数序列.设 $c$ 是使
\begin{equation}\label{eq:5.13}
  \sum_{\{i:x\in I_i\}}a_i>c
\end{equation}
对一切 $x\in A$ 成立的常数.则
\begin{equation}\label{eq:5.14}
  \sum_{i}a_i|I_i|^{s}\geq c\,\mathcal{M}^{s}_{\delta}(A).
\end{equation}
\end{lemma}

\begin{proof}
先设族 $\{I_i\}$ 与 $\{a_i\}$ 有限.把 $a_i$ 稍微减小，可取每个 $a_i$ 为有理数而不失去 \eqref{eq:5.13}；再乘以一个通分因子，可进一步设 $a_i$ 为整数.于是对每个 $i$ 取 $a_i$ 份区间 $I_i$，便可设 $a_i=1$ 对一切 $i$ 成立.

在这些假设下，\eqref{eq:5.13} 蕴含 $A$ 中每一点 $x$ 至少属于 $\lceil c\rceil$ 个区间，其中 $\lceil c\rceil$ 是大于 $c$ 的最小整数.由于 $A\subset\bigcup_i I_i$ 且 $\{I_i\}$ 是具有网性质的二进区间，可以取所有不被族中其他区间包含的区间，得到包含 $A$ 的互不重叠的区间组，记作 $\{I_i\}_{i\in\mathcal{S}_1}$.剩下的区间仍至少 $\lceil c\rceil-1$ 次覆盖 $A$，故可用同样的方法选出它们的不重叠子族 $\{I_i\}_{i\in\mathcal{S}_2}$，它也覆盖 $A$.如此进行，得到 $\lceil c\rceil$ 个互不重叠的区间族 $\{I_i\}_{i\in\mathcal{S}_j}$（$j=1,2,\ldots,\lceil c\rceil$），每个族都覆盖 $A$，且指标集 $\mathcal{S}_j$ 两两不交.由于对一切 $i$ 有 $|I_i|\leq\delta$，
\[
  \sum_{i\in\mathcal{S}_j}|I_i|^{s}\geq\mathcal{M}^{s}_{\delta}(A),
\]
对所有 $j$ 求和，便得到 $\{I_i\}$ 为有限区间族情形的 \eqref{eq:5.14}.

最后，设 $\{I_i\}_{1}^{\infty}$ 是无穷的二进区间族.对每个 $k$ 令
\[
  A_k=\Bigl\{x\in\mathbb{R}:\sum_{\substack{i:\,x\in I_i\\ i\leq k}}a_i>c\Bigr\}.
\]
由有限情形，
\[
  \sum_{1}^{k}a_i|I_i|^{s}\geq c\,\mathcal{M}^{s}_{\delta}(A_k).
\]
但每个 $A_k$ 都是二进区间的有限并，序列 $\{A_k\}$ 递增，且 $A\subset\bigcup_k A_k=\lim A_k$.由\cref{lem:5.3},
\[
  \sum_{1}^{\infty}a_i|I_i|^{s}
  \geq c\lim_{k\to\infty}\mathcal{M}^{s}_{\delta}(A_k)
  =c\,\mathcal{M}^{s}_{\delta}\Bigl(\lim_{k\to\infty}A_k\Bigr)
  \geq c\,\mathcal{M}^{s}_{\delta}(A),
\]
无穷情形证毕.
\end{proof}

\begin{theorem}\label{thm:5.8}
设 $E$ 是平面集合，$A$ 是 $x$ 轴的任意子集.设若 $x\in A$ 则 $\mathcal{H}^1(E_x)>c$，$c$ 为某常数.则
\[
  \mathcal{H}^{s+t}(E)\geq bc\,\mathcal{H}^{s}(A),
\]
其中 $b$ 只依赖于 $s$ 与 $t$.
\end{theorem}

\begin{proof}
由\cref{thm:5.1}，把结果中的 `H' 全部换成 `M' 来证明即可.

给定 $\delta>0$，设 $\{S_i\}_{1}^{\infty}$ 是构成 $E$ 的 $2^{\frac{1}{2}}\delta$-覆盖的二进正方形族.对每个 $x\in A$，
\[
  E_x\subset\bigcup_{1}^{\infty}(S_i)_x,
\]
故
\[
  \mathcal{M}^{t}_{\delta}(E_x)\leq\sum_{1}^{\infty}|(S_i)_x|^{t}.
\]
若 $A_{\delta}=\{x\in A:\mathcal{M}^{t}_{\delta}(E_x)>c\}$，则对 $x\in A_{\delta}$，
\[
  c<\sum_{1}^{\infty}|(S_i)_x|^{t}
  =2^{-\frac{1}{2}t}\sum_{\{i:\,x\in\operatorname{proj}S_i\}}|S_i|^{t},
\]
其中 $\operatorname{proj}S_i$ 是把 $S_i$ 投影到 $x$ 轴上所得的长度至多为 $\delta$ 的二进区间.于是
\[
  \sum_{1}^{\infty}|S_i|^{s+t}
  =\sum_{1}^{\infty}|S_i|^{t}|S_i|^{s}
  =2^{\frac{1}{2}s}\sum_{1}^{\infty}|S_i|^{t}|\operatorname{proj}S_i|^{s}
  \geq 2^{\frac{1}{2}(s+t)}c\,\mathcal{M}^{s}_{\delta}(A_{\delta}),
\]
其中在\cref{lem:5.7} 中取 $I_i=\operatorname{proj}S_i$，$a_i=|S_i|^t$.这对 $E$ 的任何由二进正方形 $\{S_i\}$ 构成的 $2^{\frac{1}{2}}\delta$-覆盖都成立，故
\[
  bc\,\mathcal{M}^{s}_{\delta}(A_{\delta})\leq\mathcal{M}^{s+t}_{2^{\frac{1}{2}}\delta}(E)\leq\mathcal{M}^{s+t}(E),
\]
其中 $b=2^{\frac{1}{2}(s+t)}$.由于当 $\delta$ 递减到 $0$ 时 $A_{\delta}$ 递增到 $A$，故当 $\delta\leq\rho$ 时
\[
  \mathcal{M}^{s}_{\delta}(A_{\rho})\leq\mathcal{M}^{s}_{\delta}(A_{\delta})\leq b^{-1}c^{-1}\mathcal{M}^{s+t}(E).
\]
于是对 $\rho>0$，
\[
  \mathcal{M}^{s}(A_{\rho})\leq b^{-1}c^{-1}\mathcal{M}^{s+t}(E),
\]
再由测度 $\mathcal{M}^s$ 的连续性，
\[
  \mathcal{M}^{s}(A)\leq b^{-1}c^{-1}\mathcal{M}^{s+t}(E). \qedhere
\]
\end{proof}

现在取 $E$ 为\newterm{cartesian-product}{Descartes 积}{Cartesian product} $A\times B$.

\begin{corollary}\label{cor:5.9}
对 $\mathbb{R}$ 的任意子集 $A$ 与 $B$，
\[
  \mathcal{H}^{s+t}(A\times B)\geq b\,\mathcal{H}^{s}(A)\mathcal{H}^{t}(B).
\]
\end{corollary}

可以用 Hausdorff 维数来解释这一结果：

\begin{corollary}\label{cor:5.10}
对 $\mathbb{R}$ 的任意子集 $A$ 与 $B$，
\[
  \dim(A\times B)\geq\dim A+\dim B.
\]
\end{corollary}

注意，当 $t=0$（此时 $\mathcal{H}^t(E_x)$ 等于 $E_x$ 中点的个数）时这些结果也成立.常数 $b$ 的值由 \textcite{ernstfreilich1976} 讨论.

一般而言，\cref{thm:5.8} 及其推论中的不等式不可反向，尽管 \textcite{besicovitchmoran1945} 与 \textcite{taylor1952} 给出了使\cref{cor:5.10} 中等号成立的关于 $A,B$ 的若干充分条件.下述关于该不等式可以严格的演示，作为 \textcite{besicovitchmoran1945} 例子的一种简单替代.

\begin{theorem}\label{thm:5.11}
存在 Hausdorff 维数为 $0$ 的 Borel 子集 $A,B\subset\mathbb{R}$，使 $\mathcal{H}^1(A\times B)>0$.
\end{theorem}

\begin{proof}
设 $\{s_j\}$ 是递减于 $0$ 的数列，并设 $0=m_0<m_1<m_2<\cdots$ 是增长足够快以保证
\begin{equation}\label{eq:5.15}
  \left.
  \begin{aligned}
    &(m_1-m_0)+(m_3-m_2)+\cdots+(m_{2j-1}-m_{2j-2})\leq s_jm_{2j}\\
    &(m_2-m_1)+(m_4-m_3)+\cdots+(m_{2j}-m_{2j-1})\leq s_jm_{2j+1}
  \end{aligned}
  \right\}
\end{equation}
的整数列.设 $A$ 为 $[0,1]$ 中如下数构成的子集：当 $m_j+1\leq r\leq m_{j+1}$ 且 $j$ 为奇数时，其第 $r$ 位小数为零.类似地，取 $B$ 为当 $m_j+1\leq r\leq m_{j+1}$ 且 $j$ 为偶数时第 $r$ 位小数为零的数之集.取显然的覆盖：用 $10^k$ 个长度为 $10^{-m_{2j}}$ 的区间覆盖 $A$，其中
\[
  k=(m_1-m_0)+(m_3-m_2)+\cdots+(m_{2j-1}-m_{2j-2}),
\]
则由 \eqref{eq:5.15}，当 $s>0$ 时 $\mathcal{H}^s(A)=0$；类似地 $\mathcal{H}^s(B)=0$.

设 $\operatorname{proj}$ 表示从平面到直线 $L$（即 $y=x$）的正交投影.则 $\operatorname{proj}(x,y)$ 是 $L$ 上与原点的（带符号）距离为 $2^{-\frac{1}{2}}(x+y)$ 的点.若 $u\in[0,1]$，则可找到 $x\in A$ 与 $y\in B$ 使 $u=x+y$（$u$ 的某些小数位由 $x$ 提供，其余由 $y$ 提供）.于是 $\operatorname{proj}(A\times B)$ 是 $L$ 的长度为 $2^{-\frac{1}{2}}$ 的子区间.利用正交投影不增大距离、从而由\cref{lem:1.8} 也不增大 Hausdorff 测度这一事实，
\[
  2^{-\frac{1}{2}}=\mathcal{H}^{1}\bigl(\operatorname{proj}(A\times B)\bigr)\leq\mathcal{H}^{1}(A\times B).
\]
如果需要，可以通过添加可数点集把 $A$ 与 $B$ 变成紧集.
\end{proof}

本节的全部结果都可以在更一般的框架中陈述；见 \textcite{larman1967a} 与 Wegmann（\citeyear{wegmann1969a,wegmann1969b,wegmann1971a}），他们考虑了非常一般的空间中集合的乘积.特别地，\cref{cor:5.9} 的高维类似物可用类似方法证明，即：

\begin{theorem}\label{thm:5.12}
设 $A\subset\mathbb{R}^n$，$B\subset\mathbb{R}^m$.则对某个常数 $b$ 有
\[
  \mathcal{H}^{s+t}(A\times B)\geq b\,\mathcal{H}^{s}(A)\mathcal{H}^{t}(B),
\]
其中 $A\times B\subset\mathbb{R}^{n+m}$.
\end{theorem}

\newterm{cylinder-set}{柱集}{cylinder set}是 $\mathbb{R}^3$ 中形如 $E\times I$ 的子集，其中 $E$ 是 $\mathbb{R}^2$ 的 $1$ 维子集，$I\subset\mathbb{R}$ 是单位区间.\textcite{randolph1936} 提出的一个重要问题是确定何时
\[
  \mathcal{H}^{1}(E)=\frac{1}{4}\pi\,\mathcal{H}^{2}(E\times I),
\]
即何时``面积 = 底 $\times$ 高''.有兴趣的读者可参阅 \textcites{besicovitchmoran1945,freilich1965,ward1967,larman1967c} 了解进一步细节.

\begin{exercises}
\item 设 $E\subset\mathbb{R}$ 是 $s$-集，$I$ 是单位区间.证明 $E\times I$ 是 $\mathbb{R}^2$ 中的 $(s+1)$-集.
\item 记号如\cref{thm:5.8} 所设，证明若 $E$ 是 $\mathbb{R}^2$ 中的 Borel 集，则
\[
  \mathcal{H}^{s+t}(E)\geq b\int\mathcal{H}^{t}(E_x)\,\mathrm{d}\mathcal{H}^{s}(x).
\]
\item 设 $E\subset\mathbb{R}$ 是由 Cantor 型过程构成的 $s$-集，$E=\bigcap_{1}^{\infty}E_j$，其中每个 $E_j$ 有等长区间构成的 $\delta$-覆盖 $\{U_i\}$ 使 $\sum|U_i|^s\leq c$，且当 $j\to\infty$ 时 $\delta\to 0$.设 $F=\{x-y:x,y\in E\}$ 为 $E$ 的差集.证明 $\dim F\leq 2\dim E$.（提示：考虑 $E\times E$ 到直线 $x+y=0$ 上的投影.）
\item 由\cref{exe:1.7} 推出：集合
\[
  G=\{(x,y):y-x\ \text{是无理数}\}\subset\mathbb{R}^2
\]
具有全平面勒贝格测度，但不含任何形如 $A\times B$ 的子集，其中 $A$ 与 $B$ 都是正测度的勒贝格可测集.
\end{exercises}

% !TeX root = ../main.tex
% 第 6 章 投影性质 (对应原书 pp. 75--94)

\chapter{投影性质}\label{chap:6}

\section{引言}

本章建立如下\newterm{projection}{投影}{projection}性质的结果：若 $E$ 是平面中的 $s$-集，则 $E$ 到几乎全部方向的直线上的正交投影具有维数 $t$，其中当 $s\leq 1$ 时 $t=s$，当 $s\geq 1$ 时 $t=1$.更一般地，若 $\Pi$ 是 $\mathbb{R}^n$ 的 $k$ 维子空间，$\operatorname{proj}_{\Pi}$ 表示从 $\mathbb{R}^n$ 到 $\Pi$ 上的\newterm{orthogonal-projection}{正交投影}{orthogonal projection}，我们用 $E$ 的 Hausdorff 测度来估计 $\operatorname{proj}_{\Pi}E$ 的 Hausdorff 测度.下一条引理给出一个方向上显然的不等式.

\begin{lemma}\label{lem:6.1}
设 $E$ 是 $\mathbb{R}^n$ 的任意子集，$\Pi$ 是任意子空间.则
\[
  \mathcal{H}^{s}(\operatorname{proj}_{\Pi}E)\leq\mathcal{H}^{s}(E).
\]
\end{lemma}

\begin{proof}
投影映射不增大距离，即 $|\operatorname{proj}_{\Pi}(x)-\operatorname{proj}_{\Pi}(y)|\leq|x-y|$（$x,y\in E$），故本引理是\cref{lem:1.8} 的直接推论.
\end{proof}

本章的主要部分关心相反方向的估计.这些工作自然地分为一般情形与特殊情形.对一般的 $s$-集 $E$，我们不再沿用 \textcite{marstrand1954a} 原先的几何证明，而是采用一种涉及集合容量的途径.在 $E$ 是 $s$ 为整数的 $s$-集这一特殊情形，结果呈现出微妙的平衡：当 $E$ 正则与 $E$ 非正则时会出现截然不同的现象，此时我们遵循 \textcite{besicovitch1939} 的几何方法.

\section{Hausdorff 测度与容量}

本节把分析中两个表面上极不相同的思想联系起来，同时也为研究 $s$-集的正交投影铺路.

集合的容量的思想最初为处理静电问题而发展.这一理论在被称为\newterm{potential-theory}{位势论}{potential theory}的数学分支中被推广到更普遍的吸引定律，早期的大量工作表述于 \textcite{frostman1935} 的著名论文之中.（更近的论述见 \textcites{taylor1961,carleson1967,haymankennedy1976,hille1973}.）事实证明，集合的 Hausdorff 维数与其容量是相关的，在研究维数性质时使用后者有时更为方便.

Borel 测度 $\mu$ 的\newterm{support}{支集}{support}（有时写作 $\operatorname{supp}\mu$）定义为使 $\int f\,\mathrm{d}\mu=0$ 对每个在 $S$ 上为零的连续函数 $f$ 成立的最小闭集 $S$；直观上可以把支集看作 $\mu$ 集中于其上的集合.

为便于以后引用，我们引述把测度与线性泛函等同起来的著名的 Riesz 表示定理的一个适当版本.（完整的论述见 \textcite{kingmantaylor1966}（9.5 节）或 \textcite{rudin1970}（第 2 章）.）设 $S$ 是 $\mathbb{R}^n$ 的紧子集.若 $\mathscr{F}$ 是 $S$ 上连续函数的空间，则称 $\psi:\mathscr{F}\to\mathbb{R}$ 为 $\mathscr{F}$ 上的\newterm{positive-linear-functional}{正线性泛函}{functional, positive linear}，如果
\[
  \psi(\lambda_1 f_1+\lambda_2 f_2)=\lambda_1\psi(f_1)+\lambda_2\psi(f_2)
  \qquad(\lambda_1,\lambda_2\in\mathbb{R},\ f_1,f_2\in\mathscr{F})
\]
且
\[
  \psi(f)\geq 0\qquad\text{若 } f(x)\geq 0\ \text{对一切}\ x\in S.
\]
显然，对 $S$ 上任何有限 Borel 测度 $\mu$，映射 $\psi(f)=\int_S f\,\mathrm{d}\mu$ 是 $\mathscr{F}$ 上的正线性泛函.Riesz 表示定理告诉我们逆命题也成立.

\begin{theorem}[\newterm{riesz-representation-theorem}{Riesz 表示定理}{Reisz representation theorem}]\label{thm:6.2}
设 $S$ 是 $\mathbb{R}^n$ 的紧子集，$\psi$ 是 $S$ 上连续函数空间 $\mathscr{F}$ 上的正线性泛函.则存在由 $S$ 支撑的 Borel 测度 $\mu$，$\mu(S)<\infty$，使
\[
  \psi(f)=\int_S f\,\mathrm{d}\mu
\]
对一切 $f\in\mathscr{F}$ 成立.
\end{theorem}

$\mathbb{R}^n$ 上具有紧支集且 $0<\mu(\mathbb{R}^n)<\infty$ 的 Borel 测度 $\mu$ 称为\newterm{mass distribution}{质量分布}{mass distribution}.由质量分布 $\mu$ 产生的在点 $x$ 处的 \newterm{t-potential}{$t$-位势}{t-potential}定义为
\[
  \phi_t(x)=\int\frac{\mathrm{d}\mu(y)}{|x-y|^{t}}.
\]
$\mu$ 的 \newterm{t-energy}{$t$-能量}{t-energy}由
\begin{equation}\label{eq:6.1}
  I_t(\mu)=\int\phi_t(x)\,\mathrm{d}\mu(x)=\iint\frac{\mathrm{d}\mu(x)\,\mathrm{d}\mu(y)}{|x-y|^{t}}
\end{equation}
给出.若 $E$ 是 $\mathbb{R}^n$ 的紧子集，$E$ 的 $t$-\newterm{t-capacity}{容量}{t-capacity}记作 $C_t(E)$，定义为
\begin{equation}\label{eq:6.2}
  C_t(E)=\sup_{\mu}\Bigl\{\frac{1}{I_t(\mu)}:\ E\ \text{支撑}\ \mu\ \text{且}\ \mu(E)=1\Bigr\}.
\end{equation}
（注意位势与能量可能为无穷；我们采用约定 $1/\infty=0$.）对 $\mathbb{R}^n$ 的任意子集 $E$，定义
\begin{equation}\label{eq:6.3}
  C_t(E)=\sup\{C_t(F):F\ \text{紧},\ F\subset E\}.
\end{equation}

一些作者给出稍有不同的容量定义，不过这类定义在下述重要意义下等价：具有零容量的集合相同.位势论的一个基本结果是 \eqref{eq:6.2} 等价于
\[
  C_t(E)=\sup_{\mu:\,\mu(E)=1}\Bigl(\sup_{x\in\operatorname{supp}\mu}\{\mu\ \text{在}\ x\ \text{处的}\ t\text{-位势}\}\Bigr)^{-1}.
\]
此外，这些不同的界都可以取到.这里我们不需要这类结果；细节包含在例如 \textcite{haymankennedy1976} 之中.

显然若 $E\subset E'$，则对任意 $t$ 有 $C_t(E)\leq C_t(E')$.又若 $\mu$ 是满足 $I_t(\mu)=\infty$ 的任何质量分布，则对任意 $s>t$ 有 $I_s(\mu)=\infty$.于是若 $C_t(E)=0$，则当 $s>t$ 时 $C_s(E)=0$.

由定义我们注意到下述基本事实：$C_t(E)>0$ 当且仅当存在支集含于 $E$ 中且 $I_t(\mu)<\infty$ 的质量分布 $\mu$.

\begin{lemma}\label{lem:6.3}
设 $E$ 是 $\mathbb{R}^n$ 中满足 $\mathcal{H}^s(E)<\infty$ 的紧子集.设 $\mu$ 是由 $E$ 支撑的质量分布，并令
\[
  F_0=\Bigl\{x\in E:\varlimsup_{r\to 0}\mu\bigl(B_r(x)\bigr)/r^{s}=0\Bigr\}.
\]
则 $\mu(F_0)=0$.
\end{lemma}

\begin{proof}
固定 $\alpha,\rho>0$，令
\[
  F=\{x\in E:\mu(B_r(x))/r^{s}<\alpha\ \text{对一切}\ r\leq\rho\}.
\]
若 $\{U_i\}$ 是 $F$ 的任意 $\delta$-覆盖，$\delta\leq\rho$，则设每个 $U_i$ 都含有 $F$ 的某个点，便存在圆心在 $F$ 中的球 $\{B_i\}$ 使 $|B_i|\leq 2|U_i|\leq 2\rho$ 且对每个 $i$ 有 $U_i\subset B_i$.于是
\[
  \mu(F)\leq\sum_{i}\mu(B_i)<2^{-s}\alpha\sum_{i}|B_i|^{s}\leq\alpha\sum_{i}|U_i|^{s}.
\]
这对任意 $\delta$-覆盖 $\{U_i\}$ 成立，故 $\mu(F)\leq\alpha\mathcal{H}^s_\delta(F)$，从而 $\mu(F)\leq\alpha\mathcal{H}^s(F)\leq\alpha\mathcal{H}^s(E)$.由于 $\rho$ 与 $\alpha$ 可以任意小，故 $\mu(F_0)=0$.
\end{proof}

现在我们得到 Hausdorff 测度与容量之间的基本关系.\cref{thm:6.4}(a) 由 \textcite{frostman1935} 用一种可比较网测度的早期版本证明，另见 \textcite{wallin1969}.(b) 归功于 \textcite{erdosgillis1937}，并被 \textcite{ugaheri1942} 与 \textcite{kametani1945} 进一步推广.

\begin{theorem}\label{thm:6.4}
设 $E$ 是 $\mathbb{R}^n$ 的子集.

(a) 若 $E$ 是满足 $C_t(E)=0$ 的 Souslin 集，则对一切 $s>t$ 有 $\mathcal{H}^s(E)=0$.

(b) 若 $\mathcal{H}^s(E)<\infty$，则 $C_s(E)=0$.
\end{theorem}

\begin{proof}
(a) 设 $E$ 是满足 $\mathcal{H}^s(E)>0$ 的 Souslin 集.我们通过构造一个由 $E$ 的紧子集支撑且 $I_t(\mu)<\infty$ 的质量分布 $\mu$，证明当 $t<s$ 时 $C_t(E)>0$.

由\cref{thm:5.6}(b)，存在紧集 $F\subset E$，$0<\mathcal{H}^s(F)<\infty$，使得对某个常数 $b$ 有
\[
  \mathcal{H}^{s}\bigl(B_r(x)\cap F\bigr)\leq br^{s}\qquad(x\in\mathbb{R}^n,\ r<1).
\]
令 $\mu$ 为 $\mathcal{H}^s$ 在 $F$ 上的限制，则 $\mu$ 是由 $F$ 支撑的质量分布.取 $x\in\mathbb{R}^n$，令
\begin{equation}\label{eq:6.4}
  m(r)=\mu\bigl(B_r(x)\bigr)=\mathcal{H}^{s}\bigl(B_r(x)\cap F\bigr)\leq br^{s}\qquad(r\leq 1).
\end{equation}
于是
\[
\begin{aligned}
  \phi_t(x)
  &=\int_{|x-y|\leq 1}\frac{\mathrm{d}\mu(y)}{|x-y|^{t}}+\int_{|x-y|>1}\frac{\mathrm{d}\mu(y)}{|x-y|^{t}}\\
  &\leq\int_{0}^{1}r^{-t}\,\mathrm{d}m(r)+\mu(\mathbb{R}^n)\\
  &=[r^{-t}m(r)]_{0+}^{1}+t\int_{0}^{1}r^{-(t+1)}m(r)\,\mathrm{d}r+\mu(\mathbb{R}^n)\\
  &\leq b+bt\int_{0}^{1}r^{s-t-1}\,\mathrm{d}r+\mu(\mathbb{R}^n)\\
  &=b\bigl(1+t/(s-t)\bigr)+\mathcal{H}^{s}(F),
\end{aligned}
\]
其中作了分部积分并用了 \eqref{eq:6.4}.于是 $\phi_t(x)$ 在 $\mathbb{R}^n$ 上一致有界，故
\[
  I_t(\mu)=\int\phi_t(x)\,\mathrm{d}\mu(x)<\infty,
\]
如所需.

(b) 由 \eqref{eq:6.3}，只需对 $E$ 为紧集的情形证明.设 $\mu$ 是由 $E$ 支撑的任何质量分布；我们证明 $I_s(\mu)=\infty$.令
\[
  E_0=\Bigl\{x\in E:\varlimsup_{r\to 0}\mu(B_r(x))/r^{s}>0\Bigr\}.
\]
若 $x\in E_0$，则可找到递减于 $0$ 的序列 $\{r_i\}$ 使
\[
  \mu\bigl(B_{r_i}(x)\bigr)\geq\varepsilon r_i^{s}
\]
对某个 $\varepsilon>0$ 成立.除非 $\mu(\{x\})>0$（此时显然 $I_s(\mu)=\infty$），由 $\mu$ 的连续性存在 $q_i$ 使 $0<q_i<r_i$ 且
\[
  \mu(A_i)\geq\frac{1}{2}\varepsilon r_i^{s}\qquad(i=1,2,\ldots),
\]
其中 $A_i$ 是环形区域 $B_{r_i}(x)\setminus B_{q_i}(x)$.必要时取子列，可设对一切 $i$ 有 $r_{i+1}<q_i$，使 $\{A_i\}$ 两两不交.于是当 $x\in E_0$ 时
\[
  \phi_s(x)=\int\frac{\mathrm{d}\mu(y)}{|x-y|^{s}}
  \geq\sum_{i}\int_{A_i}\frac{\mathrm{d}\mu(y)}{|x-y|^{s}}
  \geq\sum_{i=1}^{\infty}\frac{1}{2}\varepsilon r_i^{s}r_i^{-s}=\infty.
\]
由\cref{lem:6.3}，Borel 集 $E_0$ 包含 $E$ 的 $\mu$-几乎所有的点，故 $\mu(E_0)>0$，从而
\[
  I_s(\mu)=\int\phi_s(x)\,\mathrm{d}\mu(x)=\infty.
\]
由于这对由 $E$ 支撑的任何质量分布 $\mu$ 都成立，故 $C_s(E)=0$.
\end{proof}

\begin{corollary}\label{cor:6.5}
若 $E$ 是 $\mathbb{R}^n$ 的 Souslin 子集，则
\[
  \dim E=\inf\{t:C_t(E)=0\}=\sup\{t:C_t(E)>0\}.
\]
\end{corollary}

\begin{proof}
由于当 $s<\dim E$ 时 $\mathcal{H}^s(E)=\infty$，当 $s>\dim E$ 时 $\mathcal{H}^s(E)=0$（见 \eqref{eq:1.14}），推论立即得到.
\end{proof}

\cref{cor:6.5} 常以下列形式使用：

\begin{corollary}\label{cor:6.6}
设 $E$ 是 $\mathbb{R}^n$ 的 Souslin 子集.

(a) 若对某个由 $E$ 支撑的质量分布 $\mu$ 有 $I_t(\mu)<\infty$，则 $t\leq\dim E$.

(b) 若 $t<\dim E$，则存在支集含于 $E$ 中的质量分布 $\mu$ 使 $I_t(\mu)<\infty$.
\end{corollary}

注意，Hausdorff 测度与容量并非完全等价的概念.若 $\dim E=t$，$C_t(E)$ 与 $\mathcal{H}^t(E)$ 的相对大小有种种可能；一些例子见 \textcite{garnett1970} 或 \textcite{mattila1984c}.

如今，Fourier 变换在位势论中起着至关重要的作用（变换理论的一般论述见 \textcite{rudin1973}）.这里我们只需要定义，以及 Plancherel 关于平方可积函数的变换结果的一个推论.若 $f$ 与 $\mu$ 分别是 $\mathbb{R}$ 上的可积函数与质量分布，它们的（$1$ 维）\newterm{fourier-transforms}{Fourier 变换}{Fourier transforms}为
\[
  \left.
  \begin{aligned}
    \hat{f}(p)&=(2\pi)^{-\frac{1}{2}}\int f(u)\,\mathrm{e}^{\mathrm{i}up}\,\mathrm{d}u,\\
    \hat{\mu}(p)&=(2\pi)^{-\frac{1}{2}}\int\mathrm{e}^{\mathrm{i}up}\,\mathrm{d}\mu(u).
  \end{aligned}
  \right\}\qquad(p\in\mathbb{R})
\]
一个初等结果（\newterm{riemann-lebesgue-lemma}{Riemann--Lebesgue 引理}{Riemann-Lebesgue lemma}）断言 $\hat{f}$ 与 $\hat{\mu}$ 是有界连续函数.

\begin{lemma}\label{lem:6.7}
设 $\mu$ 是 $\mathbb{R}$ 上满足
\[
  0<\int|\hat{\mu}(p)|^{2}\,\mathrm{d}p<\infty
\]
的质量分布.则 $\mu$ 的支集具有正勒贝格测度.
\end{lemma}

\begin{proof}
\newterm{plancherels-theorem}{Plancherel 定理}{Plancherel's theorem}断言平方可积的 Fourier 变换对应于平方可积的函数.于是 $\mu$ 是这样的 Borel 测度：对任何在某个有界区间之外为零的连续函数 $f$，
\begin{equation}\label{eq:6.5}
  \int f(u)\,\mathrm{d}\mu(u)=\int f(u)g(u)\,\mathrm{d}u,
\end{equation}
其中 $g$ 是满足
\begin{equation}\label{eq:6.6}
  0<\int|\hat{\mu}(p)|^{2}\,\mathrm{d}p=\int|g(u)|^{2}\,\mathrm{d}u
\end{equation}
的 Borel 可测函数.若 $\mu$ 由一个勒贝格测度为零的紧集支撑，\eqref{eq:6.5} 将蕴含 $g(u)=0$ 对几乎所有的 $u$ 成立，由 \eqref{eq:6.6} 这是不可能的.
\end{proof}

\section{任意维数集合的投影性质}

现在证明关于一般维数集合投影的最重要的结果.这类结果最先由 \textcite{marstrand1954a} 直接证明（未用容量）.这里给出的位势论证明本质上属于 \textcite{kaufman1968}.为使基本思想不至被记号淹没，我们在二维情形给出证明，并稍后提及推广.此时 $\operatorname{proj}_{\theta}$ 表示从 $\mathbb{R}^2$ 到 $L_{\theta}$（过原点且与某固定轴成角 $\theta$ 的直线）上的正交投影.$L_{\theta}$ 的子集上的 $1$ 维勒贝格测度按显然的方式记作 $\mathscr{L}^1$.这里``几乎所有的 $\theta$''指按勒贝格测度.

\begin{theorem}\label{thm:6.8}
设 $E$ 是 $\mathbb{R}^2$ 中 $\dim E=s$ 的 Souslin 子集.

(a) 若 $s\leq 1$，则对几乎所有的 $\theta\in[0,\pi)$ 有 $\dim(\operatorname{proj}_{\theta}E)=s$.

(b) 若 $s>1$，则对几乎所有的 $\theta\in[0,\pi)$ 有 $\mathscr{L}^1(\operatorname{proj}_{\theta}E)>0$.
\end{theorem}

\begin{proof}
(a) 由\cref{lem:6.1}，对一切 $\theta$ 有 $\dim(\operatorname{proj}_{\theta}E)\leq s$.若 $t<s<1$，则用\cref{cor:6.6}(b) 可选取支集含于 $E$ 的紧子集中的质量分布 $\mu$ 使 $I_t(\mu)<\infty$.映射 $f\mapsto\int f(x\cdot\theta)\,\mathrm{d}\mu(x)$ 是正线性泛函，故由 Riesz 表示定理（\cref{thm:6.2}），对每个 $\theta$ 可找到 $\mathbb{R}$ 上的质量分布 $\mu_\theta$，使对连续函数 $f$ 有
\begin{equation}\label{eq:6.7}
  \int f(u)\,\mathrm{d}\mu_\theta(u)=\int f(x\cdot\theta)\,\mathrm{d}\mu(x).
\end{equation}
（这里 $\boldsymbol{\theta}$ 表示方向 $\theta$ 的单位向量，``$\cdot$''为通常的数量积，$x$ 视为从原点到 $x$ 的向量.）利用积分的单调收敛定理作通常的逼近，\eqref{eq:6.7} 对任何非负可测函数成立.把 $u\in\mathbb{R}$ 等同于 $L_{\theta}$ 上的点 $u\boldsymbol{\theta}$，便可把 $\mu_\theta$ 视为 $L_{\theta}$ 上的质量分布，其支集含于 $\operatorname{proj}_{\theta}E$.现在
\[
  I_t(\mu_\theta)
  =\iint\frac{\mathrm{d}\mu_\theta(u)\,\mathrm{d}\mu_\theta(v)}{|u-v|^{t}}
  =\iint\frac{\mathrm{d}\mu(x)\,\mathrm{d}\mu(y)}{|x\cdot\boldsymbol{\theta}-y\cdot\boldsymbol{\theta}|^{t}}
  =\iint\frac{\mathrm{d}\mu(x)\,\mathrm{d}\mu(y)}{|(x-y)\cdot\boldsymbol{\theta}|^{t}}.
\]
对 $[0,\pi)\times\mathbb{R}^2\times\mathbb{R}^2$ 应用 Fubini 定理，可见 $I_t(\mu_\theta)$ 是 $\theta$ 的可测函数，且
\[
  \int_{0}^{\pi}I_t(\mu_\theta)\,\mathrm{d}\theta
  =c\iint\frac{\mathrm{d}\mu(x)\,\mathrm{d}\mu(y)}{|x-y|^{t}}
  =c\,I_t(\mu)<\infty,
\]
其中 $c=\int_{0}^{\pi}|\boldsymbol{\theta}\cdot\boldsymbol{\tau}|^{-t}\,\mathrm{d}\theta<\infty$（因为 $t<1$，且该积分与单位向量 $\boldsymbol{\tau}$ 无关）.于是对几乎所有的 $\theta$ 有 $I_t(\mu_\theta)<\infty$，故由\cref{cor:6.6}(a) 得 $\dim(\operatorname{proj}_{\theta}E)\geq t$.这对一切 $t<s$ 成立，故对几乎所有的 $\theta$ 有 $\dim(\operatorname{proj}_{\theta}E)=s$.

(b) 若 $s>1$，则由\cref{cor:6.6}(b)，存在由 $E$ 支撑且 $I_1(\mu)<\infty$ 的质量分布 $\mu$.与 (a) 的证明一样，用 \eqref{eq:6.7} 定义 $L_\theta$ 上的质量分布 $\mu_\theta$，但这次我们考察测度 $\mu_\theta$ 的 Fourier 变换 $\hat{\mu}_\theta$.由 \eqref{eq:6.7},
\begin{align*}
  |\hat{\mu}_\theta(p)|^{2}
  &=\frac{1}{2\pi}\int\mathrm{e}^{\mathrm{i}up}\,\mathrm{d}\mu_\theta(u)\int\mathrm{e}^{-\mathrm{i}vp}\,\mathrm{d}\mu_\theta(v)\\
  &=\frac{1}{2\pi}\iint\mathrm{e}^{\mathrm{i}(u-v)p}\,\mathrm{d}\mu_\theta(u)\,\mathrm{d}\mu_\theta(v)
  =\frac{1}{2\pi}\iint\mathrm{e}^{\mathrm{i}p(x-y)\cdot\boldsymbol{\theta}}\,\mathrm{d}\mu(x)\,\mathrm{d}\mu(y).
\end{align*}
于是
\[
  |\hat{\mu}_\theta(p)|^{2}+|\hat{\mu}_{\theta+\pi}(p)|^{2}
  =\frac{1}{\pi}\iint\cos\bigl(p(x-y)\cdot\boldsymbol{\theta}\bigr)\,\mathrm{d}\mu(x)\,\mathrm{d}\mu(y),
\]
从而
\[
  \int_{0}^{2\pi}|\hat{\mu}_\theta(p)|^{2}\,\mathrm{d}\theta
  =\frac{1}{2\pi}\int_{0}^{2\pi}\iint\cos\bigl(p|x-y|\cos\theta\bigr)\,\mathrm{d}\mu(x)\,\mathrm{d}\mu(y)\,\mathrm{d}\theta,
\]
因为该定积分同样与 $x-y$ 的辐角无关.于是利用 Fubini 定理与 Bessel 函数的定义 $J_0(u)=(2\pi)^{-1}\int_{0}^{2\pi}\cos(u\cos\theta)\,\mathrm{d}\theta$，得
\[
  \int_{0}^{2\pi}|\hat{\mu}_\theta(p)|^{2}\,\mathrm{d}\theta
  =\iint J_0\bigl(p|x-y|\bigr)\,\mathrm{d}\mu(x)\,\mathrm{d}\mu(y).
\]
（Bessel 函数的详细理论见 \textcite{watson1966}.）若 $0<m<\infty$，Fubini 定理给出
\[
\begin{aligned}
  \int_{-m}^{m}\int_{0}^{2\pi}|\hat{\mu}_\theta(p)|^{2}\,\mathrm{d}\theta\,\mathrm{d}p
  &=\iint\int_{-m}^{m}J_0\bigl(p|x-y|\bigr)\,\mathrm{d}p\,\mathrm{d}\mu(x)\,\mathrm{d}\mu(y)\\
  &=\iint\int_{u=-m|x-y|}^{m|x-y|}J_0(u)\,\frac{1}{|x-y|}\,\mathrm{d}u\,\mathrm{d}\mu(x)\,\mathrm{d}\mu(y)\\
  &\leq b\iint_{\mathbb{R}}\frac{\mathrm{d}\mu(x)\,\mathrm{d}\mu(y)}{|x-y|}
  =b\,I_1(\mu),
\end{aligned}
\]
其中 $b$ 与 $m$ 无关.（反常 Riemann 积分 $\int_{-\infty}^{\infty}J_0(u)\,\mathrm{d}u$ 收敛，故定积分 $\int_{-m}^{m}J_0(u)\,\mathrm{d}u$ 当 $m$ 变化时一致有界.）

令 $m\to\infty$ 并再次利用 Fubini 定理（容易看出 $\hat{\mu}_\theta(p)$ 关于 $\theta$ 与 $p$ 都连续），得
\[
  \int_{0}^{2\pi}\int_{-\infty}^{\infty}|\hat{\mu}_\theta(p)|^{2}\,\mathrm{d}p\,\mathrm{d}\theta\leq b\,I_1(\mu)<\infty,
\]
故对 $[0,\pi)$ 中几乎所有的 $\theta$ 有 $\int_{-\infty}^{\infty}|\hat{\mu}_\theta(p)|^{2}\,\mathrm{d}p<\infty$.由\cref{lem:6.7}，对任何这样的 $\theta$，$\mu_\theta$ 的支集具有正勒贝格测度.而 $\mu_\theta$ 的支集正是 $\mu$ 的支集的投影（紧集的投影是紧集），它含于 Souslin 集 $\operatorname{proj}_{\theta}E$ 之中.

因此对几乎所有的 $\theta$，$\operatorname{proj}_{\theta}E$ 可测且具有正勒贝格测度.
\end{proof}

\cref{thm:6.8} 依投影子空间的维数而有许多高维类似物.记 $G_{n,k}$ 为由 $\mathbb{R}^n$ 的所有 $k$ 维子空间构成的 \newterm{grassmann-manifold}{Grassmann 流形}{Grassmann manifold}.$G_{n,k}$ 可以自然地赋予一个规范化的旋转不变测度；当我们说``几乎所有的 $\Pi\in G_{n,k}$''时即指关于该测度.

Kaufman 的方法（使用高维 Fourier 变换）与 Marstrand 的方法（\citeyear{marstrand1954a}）都可毫无困难地推广，见 \textcite{mattila1975b}.

\begin{theorem}\label{thm:6.9}
设 $E$ 是 $\mathbb{R}^n$ 中 $\dim E=s$ 的 Souslin 子集.

(a) 若 $s\leq k$，则对几乎所有的 $\Pi\in G_{n,k}$ 有 $\dim(\operatorname{proj}_{\Pi}E)=s$.

(b) 若 $s>k$，则对几乎所有的 $\Pi\in G_{n,k}$，$\operatorname{proj}_{\Pi}E$ 具有正的 $k$ 维勒贝格测度.
\end{theorem}

得到这类定理的另一种方法是积分投影到子空间上的测度的 Fourier 变换，利用测度 $\mu$ 的能量以其 $n$ 维 Fourier 变换表示的重要表达式
\begin{equation}\label{eq:6.8}
  I_t(\mu)=c\int_{\mathbb{R}^n}\frac{|\hat{\mu}(p)|^{2}}{|p|^{n-t}}\,\mathrm{d}p,
\end{equation}
其中 $c$ 只依赖于 $n$ 与 $t$.形式上这个等式是变换理论中卷积定理的推论，但严格的推导需要一些细心（见 \textcite{carleson1967}，第 23 页）.\textcite{falconer1982} 用这样的方法得到了\cref{thm:6.8,thm:6.9} 中例外集的 Hausdorff 维数的非平凡上界.（\newterm{exceptional-set}{例外集}{exceptional set}是指使定理结论不成立的 $\theta$（相应地 $\Pi$）的集合.）Kaufman 的原始证明与 \textcite{mattila1975b} 的推广给出了例外集大小的其他信息.\textcite{kaufmanmattila1975} 的文章给出了 $\mathbb{R}^n$ 中例外集可以大到可能的最大程度的例子，见\cref{thm:8.17}.

\textcite{mattila1975b,falconer1982} 的近期工作以容量而非 Hausdorff 测度表述，因此比上述结果稍更一般.

Marstrand（\citeyear{marstrand1954a}）对投影定理的原始表述包含一个有时有用的变式，它可以经调整适用于所讨论的任何结果.若 $E$ 是 $\mathbb{R}^2$ 中 $s>1$ 的 $s$-集，则对几乎所有的 $\theta$，不仅 $\operatorname{proj}_{\theta}E$ 具有正勒贝格测度，而且若 $F$ 是 $E$ 的满足 $\mathcal{H}^s(F)>0$ 的任何 Souslin 子集，则 $\operatorname{proj}_{\theta}F$ 具有正勒贝格测度.换言之，$\theta$ 的例外集不依赖于所选的子集 $F$.

\textcite{mattila1975b} 版本的\cref{thm:6.9}(b) 有更强的结论：若 $\dim E>k$，则对几乎所有的 $\Pi\in G_{n,k}$ 有 $\int_{u\in\Pi}\mathcal{H}^0\bigl(E\cap\operatorname{proj}_{\Pi}^{-1}u\bigr)\,\mathrm{d}u=\infty$（$\mathcal{H}^0$ 就是数一数集合中点的个数）.该文的另一个推论是：若 $0<t<\dim E-k$，则对几乎所有的 $\Pi\in G_{n,k}$，使 $\mathcal{H}^t\bigl(E\cap\operatorname{proj}_{\Pi}^{-1}u\bigr)=\infty$ 的 $u\in\Pi$ 之集具有正的 $\mathcal{H}^k$-测度.

Davies（\citeyear{davies1979,davies1982}） 最近利用\newterm{continuum-hypothesis}{连续统假设}{continuum hypothesis}证明，这些结果对仅仅是\newterm{essentially-s-dimensional}{``本质 $s$ 维''}{essentially $s$-dimensional}的集合不成立.他构造了一个不能表示为维数小于 $2$ 的集合的可数并的集合，但其到几乎每条直线上的投影勒贝格测度为零.

\section{整数维数集合的投影性质}

本节在某种意义上承接\cref{chap:3}，讨论 $s$ 为整数时 $s$-集 $E$ 的投影.我们把证明限于 $E$ 是平面的线性可测子集的情形；高维类似物的分析要复杂得多.

整数维数的正则集与非正则集的投影行为有强烈的对照.$\mathbb{R}^2$ 中的正则 $1$-集在几乎所有的方向上投影成正的 $1$ 维测度集合，而非正则 $1$-集在几乎所有的方向上投影成零测度集合.

正则集容易处理.与之前一样，$\operatorname{proj}_{\theta}$ 表示到与某轴成角 $\theta$ 的直线 $L_{\theta}$ 上的正交投影.

\begin{theorem}\label{thm:6.10}
设 $E$ 是 $\mathbb{R}^2$ 中的正则 $1$-集.则除至多一个 $\theta$ 值外，对一切 $\theta$ 有 $\mathscr{L}^1(\operatorname{proj}_{\theta}E)>0$.
\end{theorem}

\begin{proof}
由\cref{thm:3.25}，正则 $1$-集由可求长曲线的可数并的可测子集与一个零测集组成，故只需对 $E$ 是含于可求长曲线 $\Gamma$ 中的 $1$-集的情形证明.

设 $x$ 是 $E$ 中同时为 $E$ 与 $\Gamma$ 的正则点的一点.则给定 $\varepsilon>0$，可找到 $r$ 使
\[
  \mathcal{H}^{1}\bigl(E\cap B_r(x)\bigr)>(1-\varepsilon^{2})2r
  \quad\text{与}\quad
  \mathcal{H}^{1}\bigl(\Gamma\cap B_r(x)\bigr)<(1+\varepsilon)2r
\]
同时成立，从而
\[
  \mathcal{H}^{1}\bigl(E\cap B_r(x)\bigr)>(1-\varepsilon)\mathcal{H}^{1}\bigl(\Gamma\cap B_r(x)\bigr),
\]
进而
\[
  \mathcal{H}^{1}\bigl((\Gamma\setminus E)\cap B_r(x)\bigr)<\varepsilon\,\mathcal{H}^{1}\bigl(\Gamma\cap B_r(x)\bigr).
\]
由于 $\Gamma\cap B_r(x)$ 由至多可数多条不交弧组成，且 $\Gamma$ 可求长，可选取弧 $\Gamma_0\subset\Gamma\cap B_r(x)$ 使
\[
  \mathcal{H}^{1}(\Gamma_0\setminus E)<\varepsilon\,\mathcal{H}^{1}(\Gamma_0)<2\varepsilon|y-z|,
\]
其中 $y$ 与 $z$ 是 $\Gamma_0$ 的端点.设 $L_{\theta}$ 为与 $[y,z]$ 成角 $\phi$ 的直线，其中 $\cos\phi>2\varepsilon$.则投影到 $L_\theta$ 上并用\cref{lem:6.1},
\[
  \mathscr{L}^{1}(\operatorname{proj}_{\theta}E)
  >|y-z|\cos\phi-\mathscr{L}^{1}\bigl(\operatorname{proj}_{\theta}(\Gamma_0\setminus E)\bigr)
  \geq|y-z|\cos\phi-\mathcal{H}^{1}(\Gamma_0\setminus E)
  >|y-z|(\cos\phi-2\varepsilon)>0.
\]
于是除包含在长度为 $2\sin^{-1}2\varepsilon$ 的区间内（对一切 $\varepsilon>0$）的方向集 $\theta$ 之外，$\mathscr{L}^1(\operatorname{proj}_{\theta}E)>0$，即至多除一个方向外成立.
\end{proof}

可以证明比上述结果更具定量性质的结果.例如，$\mathbb{R}^2$ 中正则 $1$-集 $E$ 的 $\mathcal{H}^1$-测度可以通过把 $E\cap L$ 中的点数沿平面内所有直线 $L$ 积分而求得，其中直线之集被赋予适当的不变测度.这类思想是\newterm{integral-geometry}{积分几何}{integral geometry}的基础；有兴趣的读者可参阅 \textcite{santalo1976} 或 \textcite{federer1969} 了解进一步细节.

我们转向平面中不规则 $1$-集的投影.Gillis（\citeyear{gillis1934a,gillis1936a}） 证明：在不规则 $1$-集 $E$ 中几乎所有的 $x$ 处，对几乎所有的 $\theta$ 有 $\lim_{r\to 0}\mathscr{L}^1\bigl(\operatorname{proj}_{\theta}(E\cap B_r(x))\bigr)/\mathcal{H}^1\bigl(E\cap B_r(x)\bigr)=0$；而 \textcite{besicovitch1939} 证明了自然的结果：对几乎所有的 $\theta$ 有 $\mathscr{L}^1(\operatorname{proj}_{\theta}E)=0$.证明的出发点是不规则 $1$-集几乎处处没有切线（\cref{cor:3.30}）.在整个证明中，$I$ 与 $J$ 表示 $[0,\pi)$ 中的方向区间.遵循 Besicovitch，我们从三个定义开始.

若方向 $\theta$ 满足：过 $x$ 沿方向 $\theta$ 的直线 $L_{\theta}(x)$ 在 $x$ 的每个邻域内都与 $E$ 相交无穷多次，则称 $\theta$ 是 $E$ 在点 $x\in E$ 处的\newterm{condensation direction}{第一类凝聚方向}{condensation direction}.

给定 $x\in E$ 与正数 $\rho,\varepsilon,m$，如下定义 $[0,\pi)$ 的子集 $T(x,\rho,\varepsilon,m)$：$\theta\in T(x,\rho,\varepsilon,m)$ 当且仅当存在 $r$（$0<r<\rho$）与某个开区间 $I$（$\theta\in I\subset[0,\pi)$，$|I|<\varepsilon$）使得
\begin{equation}\label{eq:6.9}
  \frac{\mathcal{H}^{1}\bigl(E\cap C_r(x,I)\bigr)}{r}\geq m|I|.
\end{equation}
（$C_r(x,I)$ 表示由 $B_r(x)$ 中位于某条 $L_\theta(x)$（$\theta\in I$）上的点构成的双扇形.）注意，若 $E$ 不规则，则 $\mathcal{H}^1\bigl(E\cap C_r(x,I)\bigr)$ 关于 $r$ 连续，由此容易看出 $T(x,\rho,\varepsilon,m)$ 是 $[0,\pi)$ 的勒贝格可测子集.

若 $\theta\in T$，其中
\begin{equation}\label{eq:6.10}
  T=\bigcap_{\rho>0}\ \bigcap_{\varepsilon>0}\ \bigcap_{0<m<\infty}T(x,\rho,\varepsilon,m),
\end{equation}
则称方向 $\theta$ 是 $E$ 在点 $x\in E$ 处的第二类凝聚方向.

若 $[0,\pi)$ 中几乎所有的 $\theta$（按勒贝格测度）都是 $E$ 在 $x$ 处的凝聚方向，则称 $x$ 为 $E$ 的一个\newterm{point of radiation}{辐射点}{radiation, point of}.

Besicovitch 关于不规则集投影的定理的证明关键包含在下一条引理中，它依赖于勒贝格密度定理的一个精巧应用.

\begin{lemma}\label{lem:6.11}
设 $E$ 是平面中的闭不规则 $1$-集.则 $E$ 的几乎所有的点都是辐射点.
\end{lemma}

\begin{proof}
由\cref{cor:3.30} 知，对 $\mathcal{H}^1$-几乎所有的 $x\in E$ 及任意区间 $I\subset[0,\pi)$，
\begin{equation}\label{eq:6.11}
  \varlimsup_{r\to 0}\frac{\mathcal{H}^{1}\bigl(E\cap C_r(x,I)\bigr)}{2r}>\frac{1}{20}|I|.
\end{equation}
固定这样一个点 $x$；我们证明 $x$ 是 $E$ 的辐射点.

首先注意 $x$ 处第一类凝聚方向的集合 $S$ 是 $[0,\pi)$ 的勒贝格可测子集.事实上，若 $S(r)$ 是使 $L_\theta(x)\cap B_r(x)$ 含有 $E$ 中异于 $x$ 的点的 $\theta$ 之集，则 $S(r)$ 是闭集 $\{\theta:\text{存在}\ y\in L_\theta\ \text{使}\ \delta\leq|x-y|\leq r\}$ 对正有理数 $\delta$ 取并所得，故 $S(r)$ 是 Borel 集.而 $S=\bigcap_{j=1}^{\infty}S(1/j)$，故 $S$ 也是 Borel 集，从而勒贝格可测.

我们的目的是证明：若 $\rho,\varepsilon,m$ 为正，则对几乎所有的 $\theta\notin S$ 有 $\theta\in T(x,\rho,\varepsilon,m)$.设 $\theta\in[0,\pi)$ 是使 $S$ 的勒贝格密度 $\lim_{r\to 0}\mathscr{L}^1\bigl(S\cap[\theta-r,\theta+r]\bigr)/2r$ 为零的方向.（此处 $\mathscr{L}^1$ 是 $[0,\pi)$ 上的勒贝格测度.）由勒贝格定理（\cref{thm:1.13}），这对几乎所有的 $\theta\notin S$ 成立.于是对一切包含 $\theta$ 的充分小区间 $I\subset[0,\pi)$，
\begin{equation}\label{eq:6.12}
  \mathscr{L}^{1}(S\cap I)<\frac{|I|}{20m}.
\end{equation}
设 $I$ 是任一这样的区间，$0<|I|<\varepsilon$.由于当 $r\to 0$ 时 $S(r)\searrow S$，存在 $\rho_1\leq\rho$ 使当 $r<\rho_1$ 时
\begin{equation}\label{eq:6.13}
  \mathscr{L}^{1}\bigl(S(r)\cap I\bigr)<\frac{|I|}{20m}.
\end{equation}
利用 \eqref{eq:6.11} 选取 $r<\rho_1$ 使
\begin{equation}\label{eq:6.14}
  \mathcal{H}^{1}\bigl(E\cap C_r(x,I)\bigr)>\frac{2r|I|}{20}.
\end{equation}
如果碰巧
\begin{equation}\label{eq:6.15}
  \mathcal{H}^{1}\bigl(E\cap C_r(x,I)\bigr)>mr|I|,
\end{equation}
则由定义 $\theta\in T(x,\rho,\varepsilon,m)$.否则，由 \eqref{eq:6.13} 与勒贝格测度的定义，存在不交区间的可数族 $\{I_i\}$ 使
\begin{equation}\label{eq:6.16}
  S(r)\cap I\subset\bigcup_{1}^{\infty}I_i
\end{equation}
且
\begin{equation}\label{eq:6.17}
  \mathscr{L}^{1}\Bigl(\bigcup_{1}^{\infty}I_i\Bigr)<\frac{|I|}{20m}.
\end{equation}
令 $Q$ 为使
\begin{equation}\label{eq:6.18}
  \mathcal{H}^{1}\bigl(E\cap C_r(x,I_i)\bigr)>mr|I_i|
\end{equation}
成立的指标 $i$ 之集.由 \eqref{eq:6.17},
\begin{equation}\label{eq:6.19}
  \sum_{i\notin Q}\mathcal{H}^{1}\bigl(E\cap C_r(x,I_i)\bigr)
  \leq mr\,\mathscr{L}^{1}\Bigl(\bigcup_{i\notin Q}I_i\Bigr)
  <\frac{r|I|}{20}.
\end{equation}
但 $E\cap C_r(x,I)$ 的点都位于方向 $\theta\in S(r)\cap I$ 的直线 $L_\theta(x)$ 上，故由 \eqref{eq:6.16},
\[
  \bigcup_{i=1}^{\infty}\bigl(E\cap C_r(x,I_i)\bigr)\supset E\cap C_r\bigl(x,S(r)\cap I\bigr)\supset E\cap C_r(x,I).
\]
于是由 \eqref{eq:6.19} 与 \eqref{eq:6.14},
\begin{equation}\label{eq:6.20}
  \sum_{i\in Q}\mathcal{H}^{1}\bigl(E\cap C_r(x,I_i)\bigr)
  =\sum_{i=1}^{\infty}\mathcal{H}^{1}\bigl(E\cap C_r(x,I_i)\bigr)
  -\sum_{i\notin Q}\mathcal{H}^{1}\bigl(E\cap C_r(x,I_i)\bigr)
  >\mathcal{H}^{1}\bigl(E\cap C_r(x,I)\bigr)-\frac{r|I|}{20}
  >\frac{r|I|}{20}.
\end{equation}
利用 \eqref{eq:6.18} 可知，通过扩张与合并区间 $\{I_i\}_{i\in Q}$ 的过程，可以得到不交开区间族 $\{J_j\}_{1}^{\infty}$ 使
\[
  J=\bigcup_{j=1}^{\infty}J_j\supset\bigcup_{i\in Q}I_i
\]
且
\begin{equation}\label{eq:6.21}
  \mathcal{H}^{1}\bigl(E\cap C_r(x,J_j)\bigr)=mr|J_j|\qquad(j=1,2,\ldots);
\end{equation}
由于假定 \eqref{eq:6.15} 不成立，可以使 $J\subset I$.由 \eqref{eq:6.20} 与 \eqref{eq:6.21},
\[
  \mathscr{L}^{1}(J)=\sum_{j}|J_j|>\frac{|I|}{20m}.
\]
若 $\theta'\in J$，则 $\theta'\in J_j$ 对某个 $j$ 成立，故由 \eqref{eq:6.21}，$\theta'\in T(x,\rho,\varepsilon,m)$.于是 $J\subset T(x,\rho,\varepsilon,m)\cap I$，故
\[
  \mathscr{L}^{1}\bigl(T(x,\rho,\varepsilon,m)\cap I\bigr)>\frac{|I|}{20m}.
\]
这对一切包含 $\theta$ 的充分小区间成立，于是我们证明了：对几乎所有的 $\theta\notin S$，要么 $\theta\in T(x,\rho,\varepsilon,m)$，要么 $T(x,\rho,\varepsilon,m)$ 在 $\theta$ 处的下勒贝格密度至少为 $1/(20m)$.由勒贝格密度定理（\cref{thm:1.13}），$T(x,\rho,\varepsilon,m)$ 的勒贝格密度在其补集的几乎所有的点处为零，故几乎所有的 $\theta\notin S$ 都属于 $T(x,\rho,\varepsilon,m)$.

这对一切 $\rho,\varepsilon,m$ 的值成立，而 \eqref{eq:6.10} 中的交可以取为可数的（因为 $T(x,\rho,\varepsilon,m)$ 随 $\rho,\varepsilon$ 减小与 $m$ 增大而减小），于是我们得出：几乎所有的 $\theta\notin S$ 属于 $T$.引理证毕.
\end{proof}

对证明的第二段稍作修改，\cref{lem:6.11} 对任何不规则 $1$-集 $E$ 都成立，不过我们不需要用到这一点.

\cref{lem:6.11} 能否加强为``$E$ 的几乎所有的点都是第二类辐射点''（即几乎所有的方向都是第二类凝聚方向），尚属未知.

现在考虑由要求 $(\mathcal{H}^1\times\mathscr{L}^1)(A\times B)=\mathcal{H}^1(A)\mathscr{L}^1(B)$（对 Borel 集 $A$ 与 $B$）所定义的 $E\times[0,\pi)$ 上的乘积测度 $\mathcal{H}^1\times\mathscr{L}^1$.设 $G_1$ 与 $G_2$ 为由
\[
  G_i=\{(x,\theta):\theta\ \text{是}\ E\ \text{在}\ x\ \text{处的第}\ i\ \text{类凝聚方向}\}\qquad(i=1,2)
\]
给出的子集.

\begin{lemma}\label{lem:6.12}
若 $E$ 是闭不规则 $1$-集，则 $G_1$ 与 $G_2$ 是 $E\times[0,\pi)$ 的 Borel 子集，从而是 $(\mathcal{H}^1\times\mathscr{L}^1)$-可测的.
\end{lemma}

\begin{proof}
(a) 对 $0<r<\rho$ 写
\[
  G_{r,\rho}=\{(x,\theta):x\in E\ \text{且}\ r\leq|x-y|\leq\rho\ \text{对某个}\ y\in E\cap L_\theta(x)\}.
\]
则由于 $E$ 闭，$G_{r,\rho}$ 是闭集.由于 $G_{r,\rho}$ 随 $r$ 减小而增大、随 $\rho$ 减小而减小，
\[
  G_1=\bigcap_{\rho>0}\Bigl[\bigcup_{r\leq\rho}G_{r,\rho}\Bigr]
\]
是 Borel 集.

(b) 由于 $E$ 不规则，$\mathcal{H}^1\bigl(E\cap C_r(x,I)\bigr)$ 关于 $x,r$ 与 $I$ 连续.于是若按与 $T(x,\rho,\varepsilon,m)$ 相同的方式定义 $T'(x,\rho,\varepsilon,m)$ 但把 \eqref{eq:6.9} 中的不等号改为严格的，则 $\{(x,\theta):\theta\in T'(x,\rho,\varepsilon,m)\}$ 是 $E\times[0,\pi)$ 中的开集.但
\[
  G_2=\bigcap_{\rho>0}\ \bigcap_{\varepsilon>0}\ \bigcap_{0<m<\infty}\{(x,\theta):\theta\in T'(x,\rho,\varepsilon,m)\}.
\]
（由于 $T(x,\rho,\varepsilon,m)$ 关于 $m$ 单调，这里可以写 $T'$ 而不写 $T$.）这些交可以对 $\rho,\varepsilon,m$ 的有理值取，故 $G_2$ 是 Borel 集.
\end{proof}

现在可以推出投影定理了.

\begin{theorem}\label{thm:6.13}
设 $E$ 是 $\mathbb{R}^2$ 中的不规则 $1$-集.则
\[
  \mathscr{L}^{1}(\operatorname{proj}_{\theta}E)=0\qquad\text{对几乎所有的}\ \theta\in[0,\pi).
\]
\end{theorem}

\begin{proof}
先设 $E$ 是紧集.由\cref{lem:6.12}，集合 $G=G_1\cup G_2$ 是 $E\times[0,\pi)$ 的 $(\mathcal{H}^1\times\mathscr{L}^1)$-可测子集，且由\cref{lem:6.11}，
\[
  \mathscr{L}^{1}\{\theta:(x,\theta)\in G\}=\pi
\]
对几乎所有的 $x\in E$ 成立.Fubini 定理的条件满足，于是我们得出：对 $\mathscr{L}^1$-几乎所有的 $\theta\in[0,\pi)$，$(x,\theta)\in G$ 对 $\mathcal{H}^1$-几乎所有的 $x\in E$ 成立.设 $\theta$ 是这样一个方向，使 $\theta$ 对几乎所有的 $x$ 都是凝聚方向，并设 $L_\phi$ 为过原点且垂直于 $\theta$ 的直线（即 $\phi=\theta+\frac{1}{2}\pi\ (\mathrm{mod}\ \pi)$）.我们断言 $E$ 到这条直线上的投影勒贝格测度为零.

令 $E=E_0\cup E_1\cup E_2$，其中 $\mathcal{H}^1(E_0)=0$，而 $E_i$（$i=1,2$）是使 $\theta$ 为其第 $i$ 类凝聚方向的 $E$ 中的点之集.则：

(a) 由\cref{lem:6.1}，$\mathscr{L}^1(\operatorname{proj}_{\phi}E_0)=0$.

(b) 对 $E_1$ 我们用集合 Descartes 积的\cref{thm:5.8}.若 $x\in E_1$，则 $L_\theta(x)$ 与 $E$ 相交无穷多点，即 $\mathcal{H}^0\bigl(L_\theta(x)\cap E\bigr)=\infty$.把 $\mathbb{R}^2$ 视为乘积 $L_\phi\times L_\theta$ 并在 $s=1$、$t=0$ 时应用\cref{thm:5.8}，对任意大的 $c$ 有
\[
  \infty>\mathcal{H}^{1}(E_1)\geq bc\,\mathcal{H}^{1}(\operatorname{proj}_{\phi}E_1),
\]
从而 $\mathscr{L}^1(\operatorname{proj}_{\phi}E_1)=\mathcal{H}^1(\operatorname{proj}_{\phi}E_1)=0$.

(c) 设 $m$ 为任意正数，并设 $\mathcal{V}$ 为由 $L_\phi$ 的子区间构成的类
\[
  \mathcal{V}=\{J:\mathcal{H}^{1}\{x\in E:\operatorname{proj}_{\phi}x\in J\}>m|J|\}.
\]
于是利用 \eqref{eq:6.9}、$E_2$ 的定义，并注意到若 $\theta\in I$，则 $C_r(x,I)$ 可以被包进一个宽度小于 $|I|r$、一对边沿方向 $\theta$ 的矩形之中，可知 $\mathcal{V}$ 是 $\operatorname{proj}_{\phi}E_2$ 的一个 Vitali 覆盖.由 Vitali 定理（\cref{thm:1.10}(b)），给定 $\varepsilon>0$，存在 $\mathcal{V}$ 中的不交区间 $\{J_i\}$ 使
\[
  \mathscr{L}^{1}(\operatorname{proj}_{\phi}E_2)\leq\sum|J_i|+\varepsilon\leq\frac{1}{m}\mathcal{H}^{1}(E)+\varepsilon.
\]
于是由于 $1/m$ 与 $\varepsilon$ 可以任意小，
\[
  \mathscr{L}^{1}(\operatorname{proj}_{\phi}E_2)=0.
\]
把 (a)--(c) 合并，$\mathscr{L}^1(\operatorname{proj}_{\phi}E)=0$，这对几乎所有的 $\phi$ 成立.最后，若 $E$ 是任意可测 $1$-集，则利用 $\mathcal{H}^1$ 的正则性可写 $E=F\cup G$，其中 $F$ 紧且 $\mathcal{H}^1(G)<\delta$.我们已经证明对几乎所有的 $\theta$ 有 $\mathscr{L}^1(\operatorname{proj}_{\theta}F)=0$，故由\cref{lem:6.1}，
\[
  \mathscr{L}^{1}(\operatorname{proj}_{\theta}E)=\mathscr{L}^{1}(\operatorname{proj}_{\theta}G)\leq\mathcal{H}^{1}(G)<\delta.
\]
但 $\delta$ 可以任意小，故对几乎所有的 $\theta\in[0,\pi)$ 有 $\mathscr{L}^1(\operatorname{proj}_{\theta}E)=0$.\footnote{原书此句印作 $\mathscr{H}^1(\operatorname{proj}_{\theta}E)=0$；由上文（用 \cref{lem:6.1} 以 $\mathcal{H}^1(G)$ 控制投影测度）可知应为 $\mathscr{L}^1$ 之误印.——译者注.}
\end{proof}

下一条推论为判定不规则 $1$-集提供了一个非常有用的判据.

\begin{corollary}\label{cor:6.14}
$\mathbb{R}^2$ 中的 $1$-集 $E$ 是不规则的，当且仅当它在两个不同方向上的投影勒贝格测度为零.
\end{corollary}

\begin{proof}
若 $E$ 不规则，则由\cref{thm:6.13} 当然可以找到两个使投影 $\mathscr{L}^1$-测度为零的方向.另一方面，若 $E$ 有正线性测度的正则部分，则由\cref{thm:6.10}，它投影成正 $\mathscr{L}^1$-测度的集合.
\end{proof}

鉴于\cref{thm:6.13}，人们自然会问：$1$-集能否在所有方向上都投影成零测度集合.\textcite{besicovitch1928a}（第 3 章）的下述构造表明这是可能的；\textcite{morgan1935} 描述了另一种这样的构造.

\begin{theorem}\label{thm:6.15}
$\mathbb{R}^2$ 中存在（具有正 $\mathcal{H}^1$-测度的）到每条直线上的投影 $\mathscr{L}^1$-测度均为零的 $1$-集.
\end{theorem}

\begin{proof}
设 $k\geq 2$ 为整数，并设 $r$ 为使 $1/k>2\pi/r$ 的最小整数.设 $\phi$ 为角 $2\pi/r$.设 $P$ 是边长为 $a$ 与 $a/2k$、内角为 $\frac{1}{2}(\pi\pm\phi)$ 的平行四边形.我们描述一列操作，把 $P$ 换成由许多与 $P$ 相似的小平行四边形组成的子集.

首先在 $P$ 内放入 $2k$ 个与 $P$ 相似的平行四边形：较短边长为 $a/(2k)^2$、等距地排在 $P$ 的较长边上，较长边长为 $a/2k$、与 $P$ 的较长边成角 $\frac{1}{2}(\pi-\phi)$.对每个新的平行四边形做完全同样的操作，得到 $(2k)^2$ 个与 $P$ 相似但相对 $P$ 倾斜角 $\phi$、边长为 $a/(2k)^2$ 与 $a/(2k)^3$ 的平行四边形，见图~\ref{fig:6.1}.记这些平行四边形之并为 $P_1$，并记把 $P$ 换成子集 $P_1$ 的操作为 $T$.我们注意，组成 $P_1$ 的平行四边形的较长边长度之和为 $a$，而且
\[
  \mathscr{L}^{1}(\operatorname{proj}_{\theta}P_1)\leq a/k\qquad\text{若}\ 0\leq\theta\leq\phi,
\]
其中角 $\theta$ 从 $P$ 的基线方向量起.现在对组成 $P_1$ 的每个平行四边形施行操作 $T$，得到由 $(2k)^4$ 个与 $P$ 相似、相对 $P$ 倾斜角 $2\phi$、较长边长为 $a/(2k)^4$ 的平行四边形组成的子集 $P_2$.由构造，当 $\phi\leq\theta\leq 2\phi$ 时 $\mathscr{L}^1(\operatorname{proj}_{\theta}P_2)\leq a/k$；由于 $P_2\subset P_1$，这对 $0\leq\theta\leq 2\phi$ 也成立.

如此继续，对组成 $P_{j-1}$ 的平行四边形施行操作 $T$ 得到 $P_j$（$j=2,\ldots,r$），它是每个都相对 $P$ 倾斜角 $j\phi$ 的平行四边形之并.由于 $\phi=2\pi/r$，组成 $P_r$ 的 $(2k)^{2r}$ 个平行四边形都与 $P$ 位似（相似且同样地放置）.这些平行四边形的底边长度之和为 $a$，而且对 $0\leq\theta\leq 2\pi$ 有 $\mathscr{L}^1(\operatorname{proj}_{\theta}P_r)\leq a/k$.记把 $P$ 换成按此方式构造的子集 $P_r$ 的操作为 $T_k$.

接下来定义一个相伴的操作 $U_k$，它把长度为 $a$ 的线段 $S_0$ 换成大量的平行线段.如前，设 $r$ 是使 $1/k>2\pi/r=\phi$ 的最小整数，并以线段 $S_0$ 为底、以 $a$ 与 $a/2k$ 为边、内角为 $\frac{1}{2}(\pi\pm\phi)$ 竖起平行四边形 $P$.对 $P$ 施行操作 $T_k$ 得到集合 $P_r=Q$，并设 $S$ 为组成 $Q$ 的平行四边形的底线之集.由平行四边形的相应性质可得：$S$ 由总长度为 $a$ 的 $(2k)^{2r}$ 条线段组成，并且 $S$ 在一切方向上投影成至多 $a/k$ 测度的集合.令 $U_k$ 为把线段 $S_0$ 换成线段集 $S$ 的操作.

最后，设 $\{k_j\}$ 是趋于无穷的整数列.取平面上的一条单位线段，对其施行操作 $U_{k_1}$ 得到由线段组成的集合 $S_1$，并以 $Q_1$ 记由相应平行四边形构成的集合.如此继续得到线段集 $S_j$ 与相应的平行四边形集 $Q_j$，其中 $S_j$ 是对 $S_{j-1}$ 中每条线段施行操作 $U_{k_j}$ 所得线段的总体.

令 $E=\bigcap_{j=1}^{\infty}Q_j$.则 $E$ 具有所需的性质：由于对一切 $\theta$ 有 $\mathscr{L}^1(\operatorname{proj}_{\theta}Q_j)\leq 1/k_j$，当然对一切 $\theta$ 有 $\mathscr{L}^1(\operatorname{proj}_{\theta}E)=0$.为看出 Borel 集 $E$ 是 $1$-集，首先注意由于对任意大的 $m$，$E$ 可以用 $m$ 个直径至多为 $2/m$ 的平行四边形覆盖，故 $\mathcal{H}^1(E)\leq 2$.若 $\{V_i\}$ 是 $E$ 的由开凸集构成的覆盖，则由紧性论证，对某个 $j$ 有 $Q_j\subset\bigcup_i V_i$.由构造，$Q_j$ 中与 $V_i$ 相交的平行四边形的底边长度之和至多为 $3|V_i|$，故 $\frac{1}{3}\leq\sum|V_i|$.从而 $\frac{1}{3}\leq\mathcal{H}^1(E)$，$E$ 是 $1$-集.（事实上 $\mathcal{H}^1(E)=2^{-\frac{1}{2}}$.）
\end{proof}

\begin{figure}[htbp]
  \centering
  \includegraphics[width=0.5\textwidth]{fig6-1.png}
  \caption{\cref{thm:6.15} 证明中的平行四边形构造}
  \label{fig:6.1}
\end{figure}

\textcite{gillis1936b} 证明\cref{thm:6.15} 中描述的集合也从平面上的每一点投影成（角）测度为零的集合.

我们已经见过至少在两个方向上投影成单位区间的不规则 $1$-集（\cref{thm:3.32}）.\textcite{marstrand1954a} 证明了不规则 $1$-集可以在一维的 $\theta$ 集上具有 $\mathscr{L}^1(\operatorname{proj}_{\theta}E)>0$，改进了 \textcite{gillis1934a} 更早的结果.这样的集合将在\cref{thm:8.17} 中描述.

下面的\cref{thm:6.16,thm:6.17} 是\cref{thm:6.10,thm:6.13} 到高维的自然推广：

\begin{theorem}\label{thm:6.16}
设 $E$ 是 $\mathbb{R}^n$ 中的正则 $k$-集，$k$ 为整数.则对几乎所有的 $\Pi\in G_{n,k}$ 有 $\mathscr{L}^k(\operatorname{proj}_{\Pi}E)>0$.
\end{theorem}

\begin{theorem}\label{thm:6.17}
设 $E$ 是 $\mathbb{R}^n$ 中的不规则 $k$-集，$k$ 为整数.则对几乎所有的 $\Pi\in G_{n,k}$ 有 $\mathscr{L}^k(\operatorname{proj}_{\Pi}E)=0$.
\end{theorem}

这两个定理由 \textcite{federer1947} 在其浩繁的文章中（对一般 Hausdorff 测度）证明.$n$ 维 $s$-集结构的许多工作都依赖这些投影结果，一个更简单的证明将是文献中非常有价值的补充.

\section{进一步的变式}

我们简述本章所述思想的若干变式.

首先，回顾\newterm{projective-transformation}{射影变换}{projective transformation}的思想.给定平面内任意直线 $L$，可以找到平面向其添加无穷远直线后的扩张上的射影双射 $\psi_L$，把无穷远点映到 $L$ 上.映射 $\psi_L$ 把直线变为直线，并且除被映到无穷远直线上的点外是局部解析的.于是，若 $\operatorname{proj}_x$ 表示从点 $x$ 出发的投影（即把 $y$ 映到从 $x$ 出发过 $y$ 的半射线与以 $x$ 为心的单位圆的交点的映射），则可以在 $\psi_L$ 下变换\cref{thm:6.8,thm:6.10,thm:6.13}，得到关于从直线 $L$ 上的点出发的投影的结果.例如，变换\cref{thm:6.8}(a) 可见：若 $E$ 是 $\dim E=s\leq 1$ 的平面 Souslin 集，则对每条直线 $L$，$\dim(\operatorname{proj}_x E)=s$（作为单位圆的子集）对几乎所有的 $x\in L$ 成立，从而（取一族平行线的完备组）对几乎所有的 $x\in\mathbb{R}^2$ 成立.

\textcite{marstrand1954a}（第 6 节）以展示 $s$-集局部结构的方式考察了集合与过其点的直线的交.他证明了：若 $E$ 是 $s>1$ 的平面 $s$-集，则过 $E$ 的几乎每一点的几乎每条直线与 $E$ 交成维数 $s-1$ 且 $\mathcal{H}^{s-1}$-测度有限的集合.\textcite{mattila1975b} 以自然的方式把它推广到高维，并在后来的文章（\citeyear{mattila1981}）中把其中一些结果以容量表述.

也可以考虑 $s$-集与曲线（而非直线）的\newterm{intersection}{交}{intersection}.例如，若 $f$ 是 Lipschitz 函数使 $f^{-1}(x)$ 是曲线或曲面，则可以得到形如
\[
  \int\mathcal{H}^{s}\bigl(E\cap f^{-1}(x)\bigr)\,\mathrm{d}\mathcal{H}^{t}(x)\leq c\,\mathcal{H}^{s+t}(E)
\]
的不等式；见 \textcite{federer1969}，容量类似物见 \textcites{mattila1984a,mattila1984b}.若 $\Gamma$ 是曲线，$E$ 是 $s>1$ 的平面 $s$-集，则对正测度的刚体运动 $\sigma$ 有 $\dim\bigl(E\cap\sigma(\Gamma)\bigr)=s-1$.\textcite{mattila1984a} 还描述了高维类似物，并证明至少在维数上这类结果对不规则集的全等拷贝的交不成立.另一方面，若 $E$ 本身可求长，我们就得到 \textcite{federer1969} 或 \textcite{santalo1976} 中所述的经典积分几何的结果.\textcite{mattila1982} 还研究了 $E$ 是自相似集（见 8.3~节）时的这些思想.在这一非常特殊的情形，$E$ 与曲线或曲面的交的维数表现出惊人的一致性.

\begin{exercises}
\item 证明容量是次可加的，即对 $E,E'\subset\mathbb{R}^n$ 有 $C_t(E\cup E')\leq C_t(E)+C_t(E')$.举一个例子说明即使 $E$ 与 $E'$ 是不交的 Borel 集等号也未必成立，从而 $C_t$ 不是 Borel 测度.
\item 设 $\Gamma$ 是 $\mathbb{R}^2$ 中的可求长曲线.设 $n(t,\theta)$ 表示 $\Gamma$ 与方向为 $\theta$、距原点距离为 $t$ 的直线的交点个数.得到积分几何的 \newterm{poincares-formula}{Poincaré 公式}{Poincare's formula}：
\[
  \mathscr{L}(\Gamma)=\frac{1}{2}\int_{\theta=0}^{\pi}\int_{t=-\infty}^{\infty}n(t,\theta)\,\mathrm{d}t\,\mathrm{d}\theta.
\]
\item 用显然的方式把单位正方形分成 $k^2$ 个边长为 $1/k$（$k\geq 3$）的小正方形.在阵列的每一列中给一个正方形着阴影，使这些正方形的中心不共线.用 Cantor 型过程构造集合 $E$：反复把每个着阴影的正方形换成整个图形的相似拷贝并取极限集.证明 $E$ 是不规则 $1$-集.
\end{exercises}

% !TeX root = ../main.tex
% 第 7 章 Besicovitch 集与 Kakeya 集 (对应原书 pp. 95--112)

\chapter{Besicovitch 集与 Kakeya 集}\label{chap:7}

\section{引言}

Kakeya 问题有一段有趣的历史.1917 年 Besicovitch 正在研究 Riemann 积分方面的问题，遇到了如下问题：若 $f$ 是定义在平面上的 Riemann 可积函数，是否总能找到一对正交坐标轴，使 $\int f(x,y)\,\mathrm{d}x$ 对一切 $y$ 都作为 Riemann 积分存在，而且所得的 $y$ 的函数也 Riemann 可积？Besicovitch 注意到，如果能构造一个平面勒贝格测度为零、且在每一方向都含有一条线段的紧集 $F$，就会给出一个反例：因为平移 $F$（如有必要）后可设 $F$ 不含与一对固定坐标轴之一平行且有理距离的线段.令 $f$ 为由 $F$ 中至少有一个有理坐标的点构成的集合 $F_0$ 的特征函数.由于 $F$ 在每个方向都含有一条使 $F_0$ 及其补都稠密的线段，每个方向都有一条使 $f$ 不 Riemann 可积的线段.另一方面，$F$ 的不连续点集平面测度为零，故由勒贝格的著名判据，$f$ 在整个平面上 Riemann 可积.

\textcite{besicovitch1919} 成功构造了具有所需性质的集合，即如今所谓的 \newterm{besicovitch-set}{``Besicovitch 集''}{Besicovitch set}.由于当时俄国的动荡局势，他的文章流传有限，该构造后来在 \emph{Mathematische Zeitschrift}（\citeyear{besicovitch1928b}）上重新发表.

大约同一时期，\textcite{kakeya1917} 与 \textcite{fujiwarakakeya1917} 提出了求最小凸集面积的问题：在这个凸集内部可以把一条单位长线段掉头，即把它挪动到原位置但旋转 $180^{\circ}$ 而不离开该集合.他们猜测高为 $1$ 的等边三角形是这样的最小集合，并记录了 Kubota 的一个观察：若去掉凸性条件，则可以有更小的集合，即三尖内摆线.凸情形的猜测由 \textcite{pal1921} 证明，他重申了不带凸性假设的更有趣的问题.这就是后来著名的 \newterm{kakeya-problem}{Kakeya 问题}{Kakeya problem}.

在 Besicovitch 于 1924 年离开俄国后不久，人们便认识到对 Besicovitch 集作简单的修改即可给出测度任意小的 Kakeya 问题的解（\textcite{besicovitch1928b,perron1928}），该问题以出乎意料的方式得到解决.

或许其后最重要的发展再次归功于 \textcite{besicovitch1964a}，他发现了 Besicovitch 集与本书迄今所描述的几何测度论之间的基本联系.利用极对偶技术，他证明了 Besicovitch 集的存在性可以作为不规则线性可测集投影定理的一个简单的对偶结果得到.令人惊讶的是，这一对偶关系此前竟未被 Besicovitch（多年来他同时对这两个专题都感兴趣）或其他研究者注意到.

大量文章描述了这些问题的种种变式.例如，人们构造了包含所有多边形的拷贝或所有半径的圆的零测度平面集，并考虑了高维类似物的可能性.近来，这类思想已被用来回答调和分析中此前困难的问题.

值得一提的是，1958 年美国数学协会选定 Kakeya 问题作为其第一部影片的主题.任何观看这部由 Besicovitch 本人解说的迷人影片的机会都不应错过.

\section{Besicovitch 集与 Kakeya 集的构造}

本节描述 Besicovitch 集的构造，即测度为零且在每个方向都含有一条线段的集合.Besicovitch 原先的构造（\citeyear{besicovitch1919,besicovitch1928b}）已被大大简化（\textcite{perron1928,vanalphen1942,rademacher1962,schoenberg1962b,besicovitch1963c,cunningham1971,fisher1973}）.所有构造背后的基本想法是形成一个 \newterm{perron-tree}{``Perron 树''}{Perron tree}：把一个高为 $1$ 的等边三角形通过分割底边分成许多同高的小三角形，然后使这些基本三角形沿底线 $L$ 滑动不同的距离（见图~\ref{fig:7.2}）.这样的图形当然含有与 $L$ 成角至少 $60^{\circ}$ 的一切方向的单位长线段.困难在于证明用这种方式可以得到测度任意小的集合.

\begin{lemma}\label{lem:7.1}
设 $T_1$ 与 $T_2$ 是底在直线 $L$ 上的一对相邻三角形，底长为 $b$，高为 $h$.取 $\frac{1}{2}<\alpha<1$.若把 $T_2$ 沿 $L$ 滑动距离 $2(1-\alpha)b$ 使之与 $T_1$ 重叠，则所得图形 $S$ 由一个与 $T_1\cup T_2$ 位似（即相似且同样放置）的三角形 $T$（$\mathcal{L}^2(T)=\alpha^2\mathcal{L}^2(T_1\cup T_2)$）以及两个辅助三角形组成.把 $T_1\cup T_2$ 换成 $S$ 所实现的面积缩减为
\[
  \mathscr{L}^{2}(T_1\cup T_2)-\mathscr{L}^{2}(S)=\mathscr{L}^{2}(T_1\cup T_2)(1-\alpha)(3\alpha-1).
\]
（$\mathcal{L}^2$ 为平面面积即勒贝格测度.）
\end{lemma}

\begin{proof}
情形见图~\ref{fig:7.1}.利用初等几何极容易求出两图形的面积差.（注意这一计算是仿射不变的.）
\end{proof}

\begin{figure}[htbp]
  \centering
  \includegraphics[width=0.42\textwidth]{fig7-1.png}
  \caption{\cref{lem:7.1} 中的三角形滑动}
  \label{fig:7.1}
\end{figure}

\begin{theorem}\label{thm:7.2}
设 $T$ 是底在直线 $L$ 上的三角形.把 $T$ 的底分成 $2^k$ 等份，把每个分点与对顶点相连，得到 $2^k$ 个\newterm{elementary-triangle}{基本三角形}{elementary triangle} $T_1,\ldots,T_{2^k}$.只要 $k$ 取得足够大，就可以把基本三角形沿 $L$ 平移到使所得（闭）图形 $S$ 的面积任意小的位置.此外，若 $V$ 是包含 $T$ 的开集，还可以使 $S\subset V$.
\end{theorem}

\begin{proof}
这一构造需要对稍后指定的固定 $\alpha$ 值反复应用上一引理.我们先处理相邻的基本三角形对.对每个 $i$（$1\leq i\leq 2^{k-1}$），把 $T_{2i}$ 沿 $L$ 相对 $T_{2i-1}$ 移动，得到图形 $S_i^1$，它由一个三角形 $T_i^1$（与 $T_{2i-1}\cup T_{2i}$ 位似）及两个辅助三角形组成.由\cref{lem:7.1}，可以使 $\mathcal{L}^2(T_i^1)=\alpha^2\mathcal{L}^2(T_{2i-1}\cup T_{2i})$，并且（比较 $S_i^1$ 与 $T_{2i-1}\cup T_{2i}$）面积缩减了 $(1-\alpha)(3\alpha-1)\mathcal{L}^2(T_{2i-1}\cup T_{2i})$.构造的第二阶段对相邻的 $S_i^1$ 进行.对 $1\leq i\leq 2^{k-2}$，把 $S_{2i}^1$ 相对 $S_{2i-1}^1$ 平移得到 $S_i^2$.由于 $T_{2i-1}^1$ 的一边与 $T_{2i}^1$ 的对边平行且相等，由\cref{lem:7.1}，可以使 $S_i^2$ 包含一个满足 $\mathcal{L}^2(T_i^2)=\alpha^2[\mathcal{L}^2(T_{2i-1}^1)+\mathcal{L}^2(T_{2i}^1)]$ 的三角形 $T_i^2$，且 $T_{2i-1}^1$ 与 $T_{2i}^1$ 的重叠导致面积至少缩减
\[
  (1-\alpha)(3\alpha-1)[\mathcal{L}^{2}(T_{2i-1}^{1})+\mathcal{L}^{2}(T_{2i}^{1})]
  =(1-\alpha)(3\alpha-1)\alpha^{2}\mathcal{L}^{2}(T_{4i-3}\cup T_{4i-2}\cup T_{4i-1}\cup T_{4i}).
\]
如此继续；在第 $(r+1)$ 阶段，通过把 $S_{2i}^r$ 相对 $S_{2i-1}^r$ 移动（$1\leq i\leq 2^{k-r}$），使三角形 $T_{2i-1}^r$ 与 $T_{2i}^r$ 的重叠导致面积至少缩减 $(1-\alpha)(3\alpha-1)\alpha^{2r}$ 乘以移来构成 $S_{2i-1}^r$ 与 $S_{2i}^r$ 的基本三角形的总面积，从而得到 $S_i^{r+1}$.最终得到单个图形 $S_1^k$，我们就取它为 $S$.$k=3$ 时的构造阶段示于图~\ref{fig:7.2}.

比较 $\bigcup_{i}S_i^r$ 在相继的 $r$ 值下的面积，可见
\[
  \mathcal{L}^{2}(S)
  \leq\mathcal{L}^{2}(T)-(1-\alpha)(3\alpha-1)(1+\alpha^{2}+\cdots+\alpha^{2(k-1)})\mathcal{L}^{2}(T)
  =\Bigl(1-\frac{(3\alpha-1)(1-\alpha^{2k})}{(1+\alpha)}\Bigr)\mathcal{L}^{2}(T).
\]
把 $\alpha$ 取得足够接近 $1$（使 $(3\alpha-1)/(1+\alpha)$ 接近 $1$），再取足够大的 $k$，就能使 $\mathscr{L}^2(S)$ 如所需地小.

最后，若 $T$ 的底长为 $b$，则在构造中固定 $T_1$ 并相对 $T_1$ 安放其他基本三角形，任何三角形被移动的距离都不超过 $b$.于是，把原三角形分成底长至多为 $\varepsilon$ 的子三角形并对每个子三角形施行上述操作，便可以在任何基本三角形移动距离不超过 $\varepsilon$ 的前提下得到面积任意小的图形.若 $\varepsilon$ 取得足够小，该图形将含于开集 $V$ 之中.
\end{proof}

\begin{figure}[htbp]
  \centering
  \includegraphics[width=0.62\textwidth]{fig7-2.png}
  \caption{$k=3$ 时 Perron 树的构造阶段}
  \label{fig:7.2}
\end{figure}

一个有趣的问题是：当 $T$ 被分成 $m$ 个基本三角形时，确定沿底线滑动三角形所能得到的最小面积 $A_m$.Schoenberg（\citeyear{schoenberg1962a,schoenberg1962b}） 证明了 $A_m=O(1/\log m)$，但甚至不知道 $A_m$ 是否关于 $m$ 单调递减.不过，要使上述证明奏效，我们只需要 $\lim_{m\to\infty}A_m=0$.我们用\cref{thm:7.2} 来构造 Besicovitch 集.

\begin{theorem}\label{thm:7.3}
存在勒贝格测度为零且在每个方向都含有一条单位长线段的平面集合.
\end{theorem}

\begin{proof}
我们构造一个测度为零、在一个 $60^{\circ}$ 扇形内的一切方向都含有单位长线段的集合 $F$.把 $F$ 与其旋转 $60^{\circ}$ 与 $120^{\circ}$ 的全等拷贝取并，便得到具有所需性质的集合.

构造由反复应用\cref{thm:7.2} 组成.设 $S_1$ 是基于直线 $L$ 的高为 $1$ 的等边三角形.取包含 $S_1$ 的开集 $V_1$ 使 $\mathcal{L}^2(\overline{V}_1)\leq 2\mathcal{L}^2(S_1)$（横线表示闭包）.由\cref{thm:7.2}，可以把 $S_1$ 分成若干个高为 $1$、底在 $L$ 上的基本三角形，并把它们沿 $L$ 平移，形成含于 $V_1$ 内且 $\mathcal{L}^2(S_2)\leq 2^{-2}$ 的闭图形 $S_2$.由于 $S_2$ 是三角形的有限并，可找到开集 $V_2$ 使 $S_2\subset V_2\subset V_1$ 且 $\mathcal{L}^2(\overline{V}_2)\leq 2\mathcal{L}^2(S_2)$.类似地，可以把 $S_2$ 的每个基本三角形进一步分成基本三角形，并沿 $L$ 移动它们使所得图形 $S_3$ 含于 $V_2$ 且 $\mathcal{L}^2(S_3)\leq 2^{-3}$.把 $S_3$ 包进满足 $\mathcal{L}^2(\overline{V}_3)\leq 2\mathcal{L}^2(S_3)$ 的开集 $V_3$，如此继续.我们得到图形序列 $\{S_i\}$——每个都是基于 $L$、高为 $1$ 的基本三角形的有限并——以及开集序列 $\{V_i\}$，使得对每个 $i$，
\[
  S_i\subset V_i\subset V_{i-1}
\]
且
\[
  \mathcal{L}^{2}(\bar{V}_i)\leq 2\mathcal{L}^{2}(S_i)\leq 2^{-i+1}.
\]
令 $F$ 为闭集 $F=\bigcap_{i=1}^{\infty}\overline{V}_i$.我们断言 $F$ 具有所需的性质.当然 $\mathscr{L}^2(F)=0$.由构造，每个 $S_i$（从而每个 $\overline{V}_i$）在与 $L$ 成角 $60^{\circ}$ 或更大的任何方向上含有单位长线段.必须证明这对 $F$ 也成立.设 $\theta$ 是这样一个方向，并对每个 $i$ 取方向 $\theta$ 的含于 $\overline{V}_i$ 的单位长线段 $M_i$.由标准的紧性论证，必要时取子列，可设 $\{M_i\}$ 按显然的意义收敛于方向 $\theta$ 的单位长线段 $M$.由于 $\{\overline{V}_i\}$ 递减，当 $i\geq j$ 时 $M_i\subset\overline{V}_j$，故因 $\overline{V}_j$ 闭，对每个 $j$ 有 $M\subset\overline{V}_j$.于是 $M\subset\bigcap_{i=1}^{\infty}\overline{V}_i=F$，如所需.
\end{proof}

人们还找到了构造 Besicovitch 集的其他方法：\textcite{kahane1969} 注意到把 $x$ 轴上一个 Cantor 集的点与一条平行的 Cantor 型集的点连接起来便能得到这样的集合，另见 \textcite{kinney1970} 与 \textcite{alexander1975}.另一种使用对偶性的方法将在 7.3~节描述.

同样的构造可以用来形成 \newterm{nikodym-set}{Nikodym 集}{Nikodym set}：即单位正方形的子集 $F$，$\mathscr{L}^2(F)=1$，使得对每个 $x\in F$ 都有一条过 $x$ 且与 $F$ 仅交于点 $x$ 的直线（我们称 $F$ 的每一点都是\newterm{linearly-accessible-set}{``线性可达的''}{linearly accessible set}）.\textcite{nikodym1927} 为回答 Banach 的一个问题给出的复杂构造被 \textcite{davies1952a} 简化，后者证明甚至可以找到这样的集合：过其每一点都有不可数多条这样的直线.（这些构造在 de Guzmán（\citeyear{deguzman1975,deguzman1979,deguzman1981}）、\textcite{casas1978} 与 \textcite{casasdeguzman1981} 中也有描述.）\textcite{davies1952a} 用类似想法得到下述令人惊讶的结果，此处引用而不给出证明.

\begin{theorem}\label{thm:7.4}
给定任何具有有限 $\mathscr{L}^2$-测度的平面集 $G$，存在由直线的并表示的集合 $F$ 使 $G\subset F$ 且 $\mathscr{L}^2(G)=\mathscr{L}^2(F)$.
\end{theorem}

\textcite{larman1971a} 也在其\newterm{impossible-set}{``不可能集''}{impossible set}的构造中使用了这类技术.他给出了一个由 $\mathbb{R}^n$（$n\geq 3$）中闭线段的不交并构成的紧集 $F$，使 $\mathscr{L}^n(F)>0$ 而 $\mathscr{L}^n(G)=0$，其中 $G$ 是相应开线段之并.于是 $F$ 的测度集中于线段的端点.（这样的构造在 $\mathbb{R}^2$ 中是不可能的.）

接下来我们用\cref{thm:7.2} 解决 Kakeya 问题.需要下述简单引理.

\begin{lemma}\label{lem:7.5}
设 $L_1$ 与 $L_2$ 是平面中的平行直线.则给定 $\varepsilon>0$，存在包含 $L_1$ 与 $L_2$ 且 $\mathscr{L}^2(E)<\varepsilon$ 的集合 $E$，使一条单位长线段可以连续地从 $L_1$ 移动到 $L_2$ 而不离开 $E$.
\end{lemma}

\begin{proof}
设 $x_1$ 与 $x_2$ 分别是 $L_1$ 与 $L_2$ 上的点.令 $E$ 为由 $L_1$、$L_2$、连接 $x_1$ 与 $x_2$ 的线段 $M$ 以及以 $x_i$ 为心、介于 $L_i$ 与 $M$ 之间的单位扇形（$i=1,2$）构成的集合，见图~\ref{fig:7.3}.把 $x_1$ 与 $x_2$ 取得相距足够远，$E$ 的总面积便可任意小.进一步，一条单位长线段可以通过在第一个扇形中旋转、沿 $M$ 平移、再在第二个扇形中旋转而从 $L_1$ 移动到 $L_2$.
\end{proof}

\begin{figure}[htbp]
  \centering
  \includegraphics[width=0.5\textwidth]{fig7-3.png}
  \caption{\cref{lem:7.5} 中的集合 $E$}
  \label{fig:7.3}
\end{figure}

\begin{theorem}\label{thm:7.6}
给定 $\varepsilon>0$，存在满足 $\mathscr{L}^2(E)<\varepsilon$ 的\newterm{kakeya-set}{Kakeya 集}{Kakeya set} $E$，在其中一条单位长线段可以被连续移动到原位置但旋转过 $180^{\circ}$.
\end{theorem}

\begin{proof}
我们构造满足 $\mathscr{L}^2(E_0)<\frac{1}{3}\varepsilon$ 的集合 $E_0$，在其内部可以把一条单位长线段挪动到与其原方向成 $60^{\circ}$ 角的位置.取这一构造的三份拷贝便可得到所需的集合 $E$.（若需要在三个组成部分之间作平移，可能要用到\cref{lem:7.5}.）

若 $T$ 是基于直线 $L$ 的高为 $1$ 的等边三角形，则可把 $T$ 分成 $m=2^k$ 个基本三角形并沿 $L$ 把它们滑动到位置 $T_1',\ldots,T_m'$，使 $\mathscr{L}^2(\bigcup_1^m T_i')<\varepsilon/6$（见\cref{thm:7.2}）.令 $E_1=\bigcup_1^m T_i'$.注意对每个 $i$，$T_i'$ 的一边与 $T_{i+1}'$ 的对边平行.于是由\cref{lem:7.5}，对每个 $i$ 可以向 $E_1$ 添加一个测度至多为 $\varepsilon/6m$ 的集合，使一条单位长线段能从 $T_i'$ 移动到 $T_{i+1}'$.这就得到测度至多为 $\frac{1}{6}\varepsilon+\frac{1}{6}(m-1)\varepsilon/m<\frac{1}{2}\varepsilon$ 的集合 $E_0$，在其中一条单位长线段可旋转过 $60^{\circ}$，如所需.
\end{proof}

\section{对偶方法}

本节把 Besicovitch 集与\cref{chap:3,chap:6} 所述的线性可测集理论联系起来.基本想法是用点来参数化直线，使得点集 $E$ 在某方向上的投影在几何上相似于由 $E$ 参数化的直线与某条固定直线的交.于是可以通过把 $E$ 沿各个方向投影来考察由这些直线之并构成的集合（``直线集''）.

\textcite{besicovitch1964a} 以这种方式使用了\newterm{polar-reciprocity}{极对偶}{polar reciprocity}技术.设 $L(x)$ 为与原点距离 $1/|x|$ 且垂直于径向量 $x$ 的直线；即 $L(x)$ 是 $x$ 关于单位圆 $C$ 的极线.无穷远点按自然方式表示过原点的直线.极对偶的基本性质是：$x\in L(y)$ 当且仅当 $y\in L(x)$.容易看出，若 $\operatorname{proj}_{\theta}$ 表示到过原点、方向 $\theta$ 的直线 $L_\theta$ 上的投影，则对任何集合 $E$，直线集 $\{L(x):x\in E\}$ 与 $L_\theta$ 的交恰是集合 $\operatorname{proj}_{\theta}E$ 关于 $C$ 的几何反演.因此\cref{chap:6} 的投影定理可用于研究直线族的交的性质.

出于技术上的原因，我们引入\newterm{duality}{对偶性}{duality}的另一种形式，尽管基本思想相同.暂时我们在平面内工作，以 $(x,y)$ 等表示 Descartes 坐标.

若 $(a,b)\in\mathbb{R}^2$，令 $L(a,b)$ 表示直线 $y=a+bx$ 上的点之集.若 $E$ 是 $\mathbb{R}^2$ 的子集，令 $L(E)$ 为直线集 $\bigcup_{(a,b)\in E}L(a,b)$.设 $c$ 为常数，$L_c$ 为直线 $x=c$，则
\[
  L(a,b)\cap L_c=(c,\,a+bc)=(c,\,(a,b)\cdot(1,c)),
\]
其中 ``$\cdot$'' 表示 $\mathbb{R}^2$ 中通常的数量积.于是若 $E\subset\mathbb{R}^2$，
\[
  L(E)\cap L_c=\{(c,\,(a,b)\cdot(1,c)):(a,b)\in E\},
\]
故集合 $L(E)\cap L_c$ 与 $\operatorname{proj}_{\theta}E$ 几何相似，相似比为 $(1+c^2)^{\frac{1}{2}}$，其中 $c=\tan\theta$.（照例，$\operatorname{proj}_{\theta}$ 表示到过原点、与 $x$ 轴成角 $\theta$ 的直线 $L_\theta$ 上的投影.）特别地，
\begin{equation}\label{eq:7.1}
  \mathscr{L}^{1}\bigl(L(E)\cap L_c\bigr)=0
  \quad\Longleftrightarrow\quad
  \mathscr{L}^{1}(\operatorname{proj}_{\theta}E)=0.
\end{equation}
进一步，考虑到 $y$ 轴上的投影可见：若 $(a,b)\in\mathbb{R}^2$，则 $\operatorname{proj}_{\pi/2}(a,b)=b$ 恰是直线 $L(a,b)$ 的斜率.于是若 $E$ 是任意点集且 $b\in\operatorname{proj}_{\pi/2}E$，则 $L(E)$ 含有一条斜率为 $b$ 的直线.

下面的定理给出线性可测集与直线集之间的基本关系.

\begin{theorem}\label{thm:7.7}
设 $E$ 是 $\mathbb{R}^2$ 中的 $1$-集.则 $L(E)$ 是 $\mathbb{R}^2$ 的 $\mathscr{L}^2$-可测子集.若 $E$ 不规则，则 $\mathscr{L}^2(L(E))=0$；若 $E$ 正则，则 $\mathscr{L}^2(L(E))>0$.
\end{theorem}

\begin{proof}
唯一稍显麻烦的是验证 $L(E)$ 的可测性，不过照例这并不构成严重的困难（见 \textcite{croft1965} 在其对 Besicovitch 文章的评论中的说明）.

容易看出，若 $E$ 是开集则 $L(E)$ 是开集，若 $E$ 是闭集则 $L(E)$ 是闭集.于是若 $E$ 是 $G_{\delta}$-集（开集的可数交）或 $F_{\sigma}$-集（闭集的可数并），则 $L(E)$ 也是同一类型的集合.

设 $\mathcal{H}^1(E)=0$.由 $\mathcal{H}^1$ 的正则性（\cref{thm:1.6}），可找到 $G_{\delta}$-集 $E_0$ 使 $E\subset E_0$ 且 $\mathcal{H}^1(E_0)=0$.由\cref{lem:6.1}，对一切 $\theta$ 有 $\mathscr{L}^1(\operatorname{proj}_{\theta}E_0)=0$，故由 \eqref{eq:7.1}，对一切 $c$ 有 $\mathscr{L}^1(L(E_0)\cap L_c)=0$.但 $L(E_0)$ 是 $G_{\delta}$-集，从而是平面可测的.Fubini 定理适用，故 $\mathscr{L}^2(L(E_0))=0$，由此得 $L(E)\subset L(E_0)$ 是 $\mathscr{L}^2$-可测的且 $\mathscr{L}^2(L(E))=0$.

现在设 $E$ 是任意 $1$-集.由正则性，$E=E_0\cup F$，其中 $E_0$ 是 $F_{\sigma}$-集且 $\mathcal{H}^1(F)=0$.于是 $L(E)=L(E_0)\cup L(F)$ 作为两个可测集之并是 $\mathscr{L}^2$-可测的.

若 $E$ 是不规则 $1$-集，则由\cref{thm:6.13}，对几乎所有的 $\theta$ 有 $\mathscr{L}^1(\operatorname{proj}_{\theta}E)=0$.由对偶原理 \eqref{eq:7.1}，对几乎所有的 $c$ 有 $\mathscr{L}^1\bigl(L(E)\cap L_c\bigr)=0$.由于 $L(E)$ 是 $\mathscr{L}^2$-可测的，Fubini 定理蕴含 $\mathscr{L}^2(L(E))=0$.

若 $E$ 是正则 $1$-集，则由\cref{thm:6.10}，对几乎所有的 $\theta$ 有 $\mathscr{L}^1(\operatorname{proj}_{\theta}E)>0$，故由 \eqref{eq:7.1}，对几乎所有的 $c$ 有 $\mathscr{L}^1\bigl(L(E)\cap L_c\bigr)>0$.此时由 $L(E)$ 的可测性与 Fubini 定理得 $\mathscr{L}^2(L(E))>0$.
\end{proof}

\textcite{besicovitch1964a} 进一步证明：若 $E$ 是正则 $1$-集，则只要把直线集覆盖的重数考虑在内，$L(E)$ 就具有无穷平面测度.然而 \textcite{davies1965} 构造了一个满足 $\mathscr{L}^2(L(E))<\infty$ 的正则 $1$-集 $E$，表明重数假设是关键的.（虽然这些文章用极对偶表示直线，不同的参数化在一个局部实解析双射下等价，因此集合的正则性与所用的参数化无关.）

现在我们用对偶性给出 Besicovitch 集的一个简单构造.这加强了\cref{thm:7.3}，因为我们得到的是在每个方向都含有一条完整直线（而不仅仅是一条线段）的零测度集合.

\begin{theorem}\label{thm:7.8}
存在平面中勒贝格测度为零且在每个方向都含有一条直线的子集.
\end{theorem}

\begin{proof}
设 $E$ 是 $\mathbb{R}^2$ 中的不规则 $1$-集，且 $E$ 到 $y$ 轴上的投影 $\operatorname{proj}_{\pi/2}E$ 包含线段 $-1\leq y\leq 1$.例如，$E$ 可以取为\cref{thm:3.32} 的集合连同其关于 $x$ 轴的反射.考虑直线集 $L(E)$；由\cref{thm:7.7}，$\mathscr{L}^2(L(E))=0$.另一方面，\cref{thm:7.7} 之前的注记表明 $L(E)$ 含有斜率介于 $-1$ 与 $1$ 之间的一切直线.把 $L(E)$ 与其旋转 $\frac{1}{2}\pi$ 的全等拷贝取并，便得到具有所需性质的 $\mathbb{R}^2$ 的子集.
\end{proof}

我们对偶性的下一个应用归功于 \textcite{davies1971}，它表明 Besicovitch 集必然相当大.

\begin{theorem}\label{thm:7.9}
设 $F$ 是平面中在每个方向都含有一条直线的子集.则 $F$ 的 Hausdorff 维数为 $2$.
\end{theorem}

\begin{proof}
任何集合都含于一个同维数的 $G_{\delta}$-集，故不失一般性可设 $F$ 是 $G_{\delta}$-集.于是集合
\[
  E=\{(a,b):L(a,b)\subset F\}
  =\bigcap_{r=1}^{\infty}\{(a,b):L(a,b)\cap B_r(0)\subset F\cap B_r(0)\}
\]
也是 $G_{\delta}$-集，从而是 Borel 可测的.由于 $F$ 在每个方向都含直线，$\operatorname{proj}_{\pi/2}E$ 是整条 $y$ 轴，故由\cref{lem:6.1} 有 $\mathcal{H}^1(E)=\infty$.由投影定理（\cref{thm:6.8}(a)），对几乎所有的 $\theta$ 有 $\dim(\operatorname{proj}_{\theta}E)=1$，故由对偶原理，对几乎所有的 $c$ 有 $\dim\bigl(L(E)\cap L_c\bigr)=1$.由\cref{thm:5.8}，$\dim L(E)=2$，于是由 $L(E)\subset F$ 得 $\dim F=2$.
\end{proof}

把上述证明稍作修改可以证明：若 $F$ 在每个方向含有一条线段，则 $\dim F=2$（见 \textcite{davies1971}）.事实上还有更强的结果：若 $\dim F<2$，则 $F$ 在几乎所有的方向上与一切直线交成 $1$ 维测度为零的集合.这可以用相同的基本思想并采用 6.3~节末尾提到的投影定理的 Marstrand 变式来证明.（这一技术在\citeyear{falconer1980b} 中被用于高维.）

\textcite{besicovitchrado1968}（在 Besicovitch 的最后一篇文章中）与 \textcite{kinney1968} 考虑了另一个曲线填装问题，给出了包含每个半径的圆的零测度平面集的直接构造.而正是 \textcite{davies1972} 认识到这一问题适合对偶方法.

\begin{theorem}\label{thm:7.10}
存在平面中勒贝格测度为零且包含每个半径的圆的子集.
\end{theorem}

\begin{proof}
设 $E$ 是一个有界不规则 $1$-集（例如\cref{thm:3.32} 所描述者），其到 $x$ 轴上的投影 $\operatorname{proj}_{\theta}E$ 包含区间 $0\leq x\leq 1$.如下定义映射 $\psi:\mathbb{R}^2\to\mathbb{R}^2$：
\[
  \psi(x,y)=\bigl(x(1+y^{2})^{\frac{1}{2}},\,y\bigr).
\]
对每个满足 $0\leq x\leq 1$ 的 $x$，存在 $y$ 使 $(x,y)\in E$.于是对任何这样的 $x$ 值，可以在 $\psi(E)$ 中找到形如 $\bigl(x(1+y^2)^{\frac{1}{2}},y\bigr)$ 的点，即找到 $(a,b)\in\psi(E)$ 使 $a^2-x^2b^2=x^2$ 或 $x=|a|(1+b^2)^{-\frac{1}{2}}$.而这恰是直线 $L(a,b)$ 到原点的垂直距离，故直线集 $L(\psi(E))$ 含有与原点距离介于 $0$ 与 $1$ 之间的一切直线.进一步，由于 $\psi$ 是实解析映射，$\psi(E)$ 是不规则 $1$-集（$\psi$ 把可求长曲线映为可求长曲线），故由\cref{thm:7.7}，$\mathscr{L}^2(L(\psi(E)))=0$.现在令 $F$ 为把 $L(\psi(E))$ 关于原点反演所得的集合（换言之，$F$ 是 $L(\psi(E))$ 在由极坐标 $(r,\theta)\to(1/r,\theta)$ 给出的变换下的像）.反演把零测度集变为零测度集，并把距原点 $x$ 的直线变为过原点、半径为 $1/(2x)$ 的圆.于是 $F$ 测度为零，且含有半径大于 $\frac{1}{2}$ 的一切圆.取与 $F$ 相似、相似比趋于零的集合的可数并，便得到具有所需性质的集合.
\end{proof}

鉴于这一构造，自然会问：测度为零的平面集是否可能包含圆心位于每一点的圆.这个有趣而困难的问题至今未解.

Davies 的定理（\cref{thm:7.4}）有一个令人惊讶的对偶.实质上，我们可以找到一个集合，使它在每个方向上具有任何想要的投影或``影子''.这推广了 \textcite{talagrand1980} 的结果.

\begin{theorem}\label{thm:7.11}
对每个 $\theta\in[0,\pi)$ 设 $G_\theta$ 是 $\mathbb{R}$ 的子集，并设集合 $\{(\theta,y):y\in G_\theta\}$ 是平面勒贝格可测的.则存在集合 $E\subset\mathbb{R}^2$，使得对几乎所有的 $\theta$ 有 $\{u\boldsymbol{\theta}:u\in G_\theta\}\subset\operatorname{proj}_{\theta}E$ 且 $\mathscr{L}^1(G_\theta)=\mathscr{L}^1(\operatorname{proj}_{\theta}E)$，其中 $\boldsymbol{\theta}$ 是方向 $\theta$ 的单位向量.
\end{theorem}

\begin{proof}
令
\begin{equation}\label{eq:7.2}
  G=\bigl\{(c,y):(1+c^{2})^{-\frac{1}{2}}y\in G_{\theta},\ \text{其中}\ c=\tan\theta\bigr\}.
\end{equation}
由\cref{thm:7.4}，可找到直线集（记作 $L(E)$）包含 $G$ 且 $\mathscr{L}^2(G)=\mathscr{L}^2(L(E))$.由于 $G$ 平面可测，由 Fubini 定理，对几乎所有的 $c$ 有 $\mathscr{L}^1(G\cap L_c)=\mathscr{L}^1\bigl(L(E)\cap L_c\bigr)$，且 $G\cap L_c\subset L(E)\cap L_c$.由对偶原理，集合 $L(E)\cap L_c$ 与 $\operatorname{proj}_{\theta}E$ 相似（$c=\tan\theta$），相似比为 $(1+c^2)^{\frac{1}{2}}$.于是由 \eqref{eq:7.2}，对几乎所有的 $\theta$，$\operatorname{proj}_{\theta}E$ 与 $G_\theta$ 至多相差一个零测集.
\end{proof}

最后，我们用对偶性考察 Besicovitch 集的高维类似物的可能性.

\begin{theorem}\label{thm:7.12}
$\mathbb{R}^3$ 中 $\mathscr{L}^3$-测度为零的子集不可能包含每个平面的一个平移.
\end{theorem}

\begin{proof}
对 $\mathbb{R}^3$ 中的平面作适当的参数化后，证明与\cref{thm:7.9} 的证明非常相似.对 $(a,b,c)\in\mathbb{R}^3$，令 $\Pi(a,b,c)$ 表示平面 $z=a+bx+cy$ 上的点之集；若 $E\subset\mathbb{R}^3$，令 $\Pi(E)=\bigcup_{(a,b,c)\in E}\Pi(a,b,c)$.若 $d,e$ 为常数，$L_{d,e}$ 为直线 $x=d,\ y=e$，则
\[
  \Pi(a,b,c)\cap L_{d,e}=(d,e,\,a+bd+ce)=(d,e,\,(a,b,c)\cdot(1,d,e)),
\]
其中 ``$\cdot$'' 是 $\mathbb{R}^3$ 中的数量积.于是若 $E\subset\mathbb{R}^3$，
\[
  \Pi(E)\cap L_{d,e}=\{(d,e,\,(a,b,c)\cdot(1,d,e)):(a,b,c)\in E\},
\]
故 $\Pi(E)\cap L_{d,e}$ 与 $\operatorname{proj}_{\theta}E$ 几何相似，后者是 $E$ 到方向为 $(1,d,e)$ 的直线上的投影.于是
\begin{equation}\label{eq:7.3}
  \mathscr{L}^{1}\bigl(\Pi(E)\cap L_{d,e}\bigr)=0
  \quad\Longleftrightarrow\quad
  \mathscr{L}^{1}(\operatorname{proj}_{\theta}E)=0.
\end{equation}
现在设集合 $F\subset\mathbb{R}^3$ 包含每个平面的一个平移.不失一般性可设 $F$ 是 $G_{\delta}$-集，此时容易看出集合 $E=\{(a,b,c):\Pi(a,b,c)\subset F\}$ 是 Borel 集.若 $F$ 包含与所有单位向量垂直的平面，且 $\Pi$ 为平面 $x=0$，则 $\operatorname{proj}_{\Pi}E$ 是整个 $\Pi$，故由\cref{lem:6.1}，$\mathcal{H}^2(E)=\infty$.由\cref{thm:6.9}(b)，对几乎所有的 $\theta$ 有 $\mathscr{L}^1(\operatorname{proj}_{\theta}E)>0$，从而由 \eqref{eq:7.3}，对几乎所有的 $(d,e)\in\mathbb{R}^2$ 有 $\mathscr{L}^1\bigl(\Pi(E)\cap L_{d,e}\bigr)>0$.与\cref{thm:7.9} 的证明一样，$E$ 的可测性蕴含 $\Pi(E)$ 是 $\mathscr{L}^3$-可测的，故由 Fubini 定理，$\mathscr{L}^3(\Pi(E))>0$.由于 $\Pi(E)\subset F$，集合 $F$ 必有正测度.
\end{proof}

把上述证明稍作修改并不难证明：$\mathbb{R}^3$ 中任何零测度子集与几乎所有方向（按球面测度）垂直的一切平面都交成平面测度为零的集合.

Besicovitch 集的高维类似物是否存在是颇有意义的问题.称 $\mathbb{R}^n$ 的子集为 $(n,k)$-Besicovitch 集，如果它 $n$ 维勒贝格测度为零，却包含 $\mathbb{R}^n$ 的每个 $k$ 维子空间的一个平移.我们在\cref{thm:7.8} 中证明了 $(2,1)$-Besicovitch 集存在，而且对任意 $n$，这样的集合与 $\mathbb{R}^{n-2}$ 的 Descartes 积是一个 $(n,1)$-Besicovitch 集.\textcite{marstrand1979b} 用直接的覆盖方法证明了 $(3,2)$-集不可能存在；同时 \textcite{falconer1980a} 用 Fourier 变换证明：当 $k>\frac{1}{2}n$ 时 $(n,k)$-Besicovitch 集不存在.\textcite{oberlinstein1982} 也联系泛函分析提到了 $(n,n-1)$ 情形.\cref{thm:7.12} 足以推广到 $(n,n-1)$-Besicovitch 集不存在的结论，但在一般的 $(n,k)$ 情形使用对偶比较麻烦.\textcite{falconer1980b} 通过使用一种对偶形式并结合投影定理（投影经由一个公共的中间空间分解），证明了当 $2\leq k\leq n-1$ 时 $(n,k)$-Besicovitch 集不存在.这里我们陈述这一结果的最一般形式：

\begin{theorem}\label{thm:7.13}
设 $2\leq k\leq n-1$，$F$ 是 $\mathbb{R}^n$ 中 $n$ 维勒贝格测度为零的子集.则对几乎所有的 $k$ 维子空间 $\Pi\subset\mathbb{R}^n$（按 Grassmann 流形 $G_{n,k}$ 上通常的不变测度），$\Pi$ 的每个平移与 $F$ 交成 $k$ 维测度为零的集合.
\end{theorem}

\section{推广}

本节综述由 Besicovitch 与 Kakeya 构造引发的众多推广与变式.（其中一些材料在 Croft 未刊的油印讲义以及 \textcite{davies1971,cunningham1974,taylor1975} 中也有讨论.）

一类一般形式的问题是：给定一族几何图形，是否可能找到一个零测度平面集，包含该族中每个图形的一个平移，或哪怕只是一个全等拷贝？Besicovitch 集对全体直线的平移给出了肯定答案.

\textcite{kinney1968} 观察到，由于 Cantor 集 $E$ 勒贝格测度为零且包含 $(0,1)$ 中的一切距离，``Cantor 格子布''$(E\times[0,1])\cap([0,1]\times E)$ 平面测度为零，且包含边长小于 $1$ 的每个矩形的全等拷贝.取相似集合的可数并便得到包含所有矩形拷贝的零测度集.Ward（\citeyear{ward1970a,ward1970b}） 构造了包含所有多边形的全等拷贝的零测度集；事实上他造出了一个 Hausdorff 维数为 $2-1/n$ 的集合，包含每个 $n$ 边多边形的一个全等拷贝.\textcite{davies1971} 随后修改了 Besicovitch 的原始构造，给出包含每条多边形折线的平移的零测度集——由\cref{thm:7.9} 它必有维数 $2$.\textcite{marstrand1972} 在一篇适用于任何拓扑测度空间的非常一般的文章中证明了这类最强的结果.他证明了，除其他事项外，存在包含每条直线的可数并的平移的零测度集，还存在 Hausdorff 维数小至 $1$、包含每个这样的并的全等拷贝的集合.

另一个自然会考虑的问题是\newterm{curve-packing}{曲线填装}{curve packing}：把曲线弧的拷贝填装进零测度集.我们已经见过（\cref{thm:7.10} 及相关注记）存在包含所有半径的圆的零测度集，也提到了所有圆心的问题尚未解决.其他显然的问题是：零测度集能否包含所有椭圆、所有二次曲线、或每个度数 $k>1$ 的平面曲线的拷贝.虽然人们普遍相信这样的集合不存在，这些问题相当困难，部分原因是其中涉及非线性，迄今一无所证.\textcite{falconer1982} 给出了关于此类集合的维数以及在 $\mathbb{R}^n$（$n\geq 3$）中填装高维曲面的一些相对较弱的结果.其他这样的填装结果（通常在三维或更高维）可以从调和分析的极大定理推出，见 7.5~节.

毫无疑问，这一领域最重要的发展归功于 \textcite{marstrand1979a}，他解决了\newterm{worm-problem}{``虫''问题}{worm problem}（``一条单位长的虫最小的舒适住所是什么？''），证明了对 $\mathbb{R}^n$（$n\geq 2$）中任何 $n$ 维测度为零的子集 $F$，存在某条长度为 $1$ 的光滑（$C^\infty$）曲线，其全等拷贝都不含于 $F$.事实上，存在这样的曲线，它不能被任何在包含该曲线的区域内定义的可逆实解析映射变换成 $F$ 的子集.复杂的证明使用了全有界度量空间的熵的思想，即对小的 $\varepsilon$ 估计覆盖该空间所需的直径为 $\varepsilon$ 的集合的个数.

还应简要提及相反类型的填装问题：是否存在不包含某些指定集合的全等或相似拷贝的正测度集合？这类结果中最著名的归功于 \textcite{steinhaus1920}：直线的任何正测度子集包含 $(0,c)$ 中的一切距离，其中 $c$ 为某个正数.（见\cref{exe:1.7}；各种推广见 \textcites{besicovitchmiller1948,besicovitch1948,eggleston1949,besicovitchtaylor1952}.）测度为零的 Cantor 集也具有这一性质.Steinhaus 结果由勒贝格密度定理（\cref{thm:1.13}）推出，后者还蕴含直线的任何正长度子集包含任何有限点集的相似拷贝.Erdős 提出了诱人的问题：这对可数（收敛）点列是否仍成立；\textcite{falconer1984} 给出了部分解答.\textcite{daviesmarstrandtaylor1960} 与 \textcite{darstgoffman1970} 考察了这一主题的其他若干变式.

求包含某些类型图形的所有拷贝的最小凸集的问题实际上属于凸性而非几何测度论.我们仅指出：在激发本章所描述的许多研究的 \textcite{pal1921} 的文章中，证明了等边三角形是包含每个方向的单位长线段的最小平面凸集；高维的类似问题仍未解决.\textcite{eggleston1957,besicovitch1965b,wetzel1973} 等人考虑了相关的``凸性''问题.

我们回到在一个尽可能小的集合内把单位长线段转动 $180^{\circ}$ 的 Kakeya 问题.最终由 \textcite{cunningham1971} 解决，他构造了含于单位圆盘内的测度任意小的单连通 Kakeya 集.同一篇文章考虑了极小星形 Kakeya 集，并证明其面积的下界介于 $\pi/108$ 与 $(5-2\sqrt{2})\pi/24$ 之间，取代了 \textcite{walker1952,blank1963,cunninghamschoenberg1965} 更早的估计.（称集合是星形的，如果它有一点可用含于该集合内的线段连接到集合的所有其他点.）极小凸 Kakeya 集仍是高为 $1$ 的等边三角形（\citeyear{pal1921}），三维类似物未知.

\textcite{cunningham1974} 考虑了值得一提的三个问题.首先他求出球面上可以在其中把一段大圆弧掉头的最小子集的面积.极小面积（对小于半圆的弧，此前已由 \textcite{wilker1971} 证明为零）取决于弧所张的角.文中讨论的第二个问题涉及求出小的（平面）集合，使一条定角而半径可变的圆弧可以在其中收缩成一点.Cunningham 的最后一个变式把``鸟''定义为一条中央线段（``身体'')，其两端各非刚性地连接一段（``翅膀''）.文中证明可以把鸟连续地移过任何有界集 $E$，使身体经过 $E$ 的全部，而翅膀保持在一个测度任意小的集合内.这一问题在调和分析中有应用，并与 \textcite{davies1952a} 关于线性可达集的工作密切相关.

\section{与调和分析的关系}

若不提及与调和分析的关系，本章便不完整.一方面，有时可以从某些积分的更一般的范数界得到填装结果.另一方面，Kakeya 构造被用来为调和分析中的一些重大猜想提供反例.我们宽泛地使用\newterm{harmonic-analysis}{``调和分析''}{harmonic analysis}一词，用以描述包括 Fourier 分析以及微分理论、极大算子与乘子理论在内的广阔数学领域.我们只能给出少量例子来说明这种关系.

在第一个例子中，我们从泛函分析不等式推出\cref{thm:7.12}（高维 Besicovitch 集的不存在性）.设 $f(x)$ 是定义在 $\mathbb{R}^3$ 上的可测函数.照例记 $\|f\|_p=\bigl[\int_{\mathbb{R}^3}|f|^p\,\mathrm{d}x\bigr]^{1/p}$，$1\leq p<\infty$.令 $F(t,\theta)$ 表示 $f$ 在垂直于单位向量 $\theta$、距原点距离为 $t$ 的平面上的积分；由 Fubini 定理的简单应用，对一切 $\theta$，$F(t,\theta)$ 对几乎所有的 $t$ 存在.把 $f$ 的 Fourier 变换关于 $t$ 沿所有单位向量 $\theta$ 积分，并利用可积函数的变换有界这一事实，不难证明
\begin{equation}\label{eq:7.4}
  \int_{\theta}\operatorname*{ess\,sup}_{t}|F(t,\theta)|\,\mathrm{d}\theta\leq c\bigl(\|f\|_{1}+\|f\|_{2}\bigr)
\end{equation}
对某个与 $f$ 无关的常数 $c$ 成立.（见 \textcite{falconer1980a} 或 \textcite{oberlinstein1982}.）

现在设 $E$ 是零测度集，并记 $F_0(\theta)$ 为 $E$ 与垂直于 $\theta$ 的平面的交的平面外测度的上确界.常规方法可证 $F_0$ 关于球面测度可测.由勒贝格测度的正则性，可找到包含 $E$ 且测度小于 $\varepsilon$ 的开集 $V$；于是若 $f$ 是 $V$ 的特征函数，则 $\|f\|_1=\|f\|_2<\varepsilon$.由于 $V$ 开，
\[
  F_0(\theta)\leq\operatorname*{ess\,sup}_{t}F(t,\theta),
  \quad\text{故由 \eqref{eq:7.4}，}\quad
  \int F_0(\theta)\,\mathrm{d}\theta\leq 2\varepsilon c.
\]
由于 $\varepsilon$ 任意且 $F_0(\theta)$ 非负，我们得出 $F_0(\theta)=0$ 对几乎所有的 $\theta$ 成立，换言之，几乎所有方向上的所有平面与 $E$ 交成平面测度为零的集合.

这类论证也可以反过来用.我们知道（\cref{thm:7.3}）存在在每个方向含有一条直线的零测度平面集，因此若形如 \eqref{eq:7.4} 的估计对 $\mathbb{R}^2$ 上的函数成立（此时 $\theta$ 表示平面内的单位向量），与上述类似的论证将导致矛盾.于是 Besicovitch 集为``\eqref{eq:7.4} 对平面函数成立''这一看似合理的猜想提供了反例.

如\cref{thm:7.13} 及其前的注记所述，$(n,k)$-Besicovitch 集在 $k>1$ 时不存在.有趣的是，上述方法经调整可以在 $k>\frac{1}{2}n$ 时证明这一点，但此外则不然（见 \textcite{falconer1980a}）.这或许暗示：当 $2\leq k\leq\frac{1}{2}n$ 时，$(n,k)$-Besicovitch 集的不存在性是一种内在的几何性质，而非更一般的泛函分析结果的推论.

我们转向球面平均与球的填装问题.若 $f$ 是 $\mathbb{R}^n$ 上的可测函数，令 $F_r(x)$ 表示 $f$ 在以 $x$ 为心、$r$ 为半径的球面上的平均，即 $f$ 在该球面上的积分除以球的 $(n-1)$ 维曲面面积.令
\begin{equation}\label{eq:7.5}
  F(x)=\sup_{r>0}F_r(x).
\end{equation}
\textcite{stein1976} 证明了
\begin{equation}\label{eq:7.6}
  \|F\|_{p}\leq c\|f\|_{p},
\end{equation}
只要 $n/(n-1)<p$ 且 $n\geq 3$.（如果我们改为把 $F_r(x)$ 定义为 $f$ 在以 $x$ 为心、$r$ 为半径的球体上的平均，则 $F(x)$ 就是 Hardy--Littlewood \newterm{maximal-operator}{极大算子}{maximal operator}，而 \eqref{eq:7.6} 就是经典的极大不等式，它对一切 $n$ 与 $1<p<\infty$ 成立，见 \textcite{deguzman1975}.）

我们使用与上面完全相同的论证，但从 \eqref{eq:7.6} 而非 \eqref{eq:7.4} 出发，推出：若 $n\geq 3$ 且 $E$ 是 $\mathbb{R}^n$ 中 $n$ 维勒贝格测度为零的任意子集，则对几乎所有的 $x\in\mathbb{R}^n$，以 $x$ 为心的一切球与 $E$ 交成 $(n-1)$ 维曲面测度为零的集合.

若 \eqref{eq:7.6} 对 $n=2$ 与某个 $p$ 值得证，则将蕴含零测度平面集不可能包含以每点为心的圆，从而解决前文提到的问题.另一方面，若构造出这样的零测度圆填装，则将蕴含 \eqref{eq:7.6} 对一切 $p$ 都不成立，而不只是对 $p\leq 2$ 不成立——后者是目前在 $\mathbb{R}^2$ 中仅知的.

\eqref{eq:7.6} 的证明关键依赖于球面曲率的不为零.对其他``曲率良好''曲面上的平均也有类似结果（见 \textcite{steinwainger1978}），并随之得到填装结果；但对立方体的表面这类结果则不成立，例如.

勒贝格密度定理的许多推广涉及球以外的集合上的平均.一个看似合理的猜想（对连续函数当然成立）或许是：若 $f$ 是平面上的``合理''函数，则
\begin{equation}\label{eq:7.7}
  \inf\Bigl\{\frac{1}{\mathscr{L}^{2}(R)}\int_{R}f\,\mathrm{d}\mathscr{L}^{2}:\ R\ \text{是中心在}\ x\ \text{处且}\ 0<|R|\leq r\ \text{的矩形}\Bigr\}
\end{equation}
当 $r\to 0$ 时对几乎所有的 $x$ 收敛于 $f(x)$，对上确界也有类似结果.然而，Nikodym 集的存在性（在\cref{thm:7.3} 之后讨论）的一个推论是：这类结果对集合的特征函数甚至都不成立.（我们说全体矩形之类不构成一个微分基.）事实上，设 $E$ 是 Nikodym 集，即单位正方形中测度为 $1$、且每一点都位于某条与 $E$ 仅交一点的直线上的子集.由正则性，可找到 $E$ 的闭子集 $F$ 具有同样的可达性质且 $\mathscr{L}^2(F)>0$.令 $f$ 为 $F$ 的特征函数.由于 $F$ 的补是开集，容易看出：取长轴沿``例外直线''的细矩形，\eqref{eq:7.7} 的下确界当 $r\to 0$ 时对一切 $x\in F$ 趋于零.

Besicovitch 集在 \newterm{fourier-multiplier}{Fourier 乘子理论}{Fourier multiplier}中提供了一个重要而惊人的反例.设 $T$ 是 $\mathbb{R}^2$ 上 $p$ 次方可积函数空间上的线性变换，用 Fourier 变换记号定义为
\[
  \widehat{(Tf)}(\xi)=\chi_{\mathrm{B}}(\xi)\hat{f}(\xi),
\]
其中 $\chi_{\mathrm{B}}$ 是单位圆盘的特征函数.已知
\begin{equation}\label{eq:7.8}
  \|Tf\|_{p}\leq c\|f\|_{p}
\end{equation}
对 $p=2$ 成立，且人们曾普遍预期这一不等式对一切 $\frac{4}{3}<p<4$ 也成立.然而 \textcite{fefferman1971} 精妙地利用 Besicovitch 集证明：\eqref{eq:7.8} 当 $p\neq 2$ 时不成立.有兴趣的读者可参阅 Fefferman 行文清晰的文章了解进一步细节.

\textcite{mitjaginnikisin1973} 遵循类似的过程得到：对 $p<2$，存在单位正方形上的 $p$ 次方可积函数 $f$，其\newterm{fourier-series}{Fourier 级数的部分和}{Fourier series}几乎处处无界，即对几乎所有的 $(x,y)\in\mathbb{R}^2$，
\[
  \varlimsup_{r\to\infty}\Bigl|\sum_{k^{2}+j^{2}<r}a_{kj}\exp\bigl(2\pi\mathrm{i}(xk+yj)\bigr)\Bigr|=\infty,
\]
其中 $\{a_{kj}\}$ 是 $f$ 的二维 Fourier 系数.这与一维情形形成鲜明对比：Carleson--Hunt 定理断言，若 $f$ 对某个 $1<p<\infty$ 是 $p$ 次方可积的，则 $f$ 的 Fourier 级数几乎处处逐点收敛.

最后我们提到在\newterm{functional-analysis}{泛函分析}{functional analysis}中的一个应用：$T$ 是 Hilbert 空间上的 \newterm{hypernormal-operator}{超正规算子}{hypernormal operator}（即 $T^{*}T-TT^{*}\geq 0$ 且 $T$ 没有非平凡的约化空间使其上 $T$ 正规），当且仅当 $T$ 的谱在其每一点的任何邻域中都具有正测度.\textcite{putnam1974} 因此用 Nikodym 集证明了一个正面结果，而不仅仅是提供反例.

要更深入地探索调和分析的相关领域，读者可参阅 \textcite{steinwainger1978} 的综述文章、由 \textcite{waingerweiss1979} 编辑的美国数学会调和分析会议文集，以及 de Guzmán（\citeyear{deguzman1975,deguzman1981}） 的书.几何测度论与调和分析看来必将继续相互产生巨大的影响.

\begin{exercises}
\item 证明存在 $\mathbb{R}^n$（$n\geq 2$）的子集 $E$，在其中有 $\mathscr{L}^n(E)$ 任意小的测度，且在其内部一条单位长线段可以被挪动到任意方向.
\item 设 $S$ 是 $[0,\pi)$ 的 Borel 子集.设 $F$ 是 $\mathbb{R}^2$ 中在 $S$ 的每个方向都含有一条直线的子集.用对偶性证明 $\dim F\geq 1+\dim S$.
\item 证明存在 $\mathbb{R}^2$ 的 $\mathscr{L}^2$-测度为零的子集，包含过 $x$ 轴每一点的一条互不相同的直线.
\item 证明若 $E$ 是 $\mathbb{R}^2$ 中 $\dim E>1$ 的 Borel 子集，则 $E$ 中各点的极线之集（关于任何固定的圆）覆盖一个平面勒贝格测度为无穷的集合.
\item 若 $f$ 是平面可测函数，令 $F_r(x)$ 表示 $f$ 在以 $x$ 为心、$r$ 为半径的圆周上的平均.可以证明（\citeyear{wainger1979}）若 $E$ 是 Cantor 集，则对某个常数 $c$ 有
\[
  \bigl\|\sup_{r\in E}F_{1+r}(x)\bigr\|_{2}\leq c\|f(x)\|_{2}.
\]
进而证明：勒贝格测度为零的平面集不可能在 $\mathbb{R}^2$ 的每一点处包含一个半径为 $1+r$（对某个 $r\in E$）的圆.
\end{exercises}

% !TeX root = ../main.tex
% 第 8 章 分形集杂例 (对应原书 pp. 113--149)

\chapter{分形集杂例}\label{chap:8}

\section{引言}

本章综述由特定构造产生、或出现于数学或物理其他分支中的分数维数集合的例子，并把它们与本书前面的内容联系起来.这些专题的选择完全出于作者的一时兴致，而非因为它们代表了分形集合最重要的出现场合.每一节中选取若干有意义的结果加以证明，其余则仅加引用.希望这一方式能鼓励读者在别处更深入地跟进其中某些专题.大部分例子来自拥有浩繁文献的数学领域，因此本章的参考文献只给出主要出处以及包含进一步综述与文献的近期文章和书籍.

\section{分数维数的曲线}

本节在 $(x,y)$ 坐标平面内工作，考察 $\Gamma$ 的 Hausdorff 维数（\newterm{curve-dimension}{曲线的维数}{curve, dimension of}），即定义在（比如说）单位区间上的函数 $f$ 的图像所构成的点集 $\{(x,f(x))\}$ 的 Hausdorff 维数.

若 $f$ 是有界变差函数，即对一切分割 $0=x_0<x_1<\cdots<x_m=1$，$\sum_{i=1}^{m}|f(x_i)-f(x_{i-1})|$ 有界，则我们实际上回到了 3.2~节的情形：$\Gamma$ 是可求长曲线，从而是正则 $1$-集.然而，若 $f$ 充分不规则（尽管连续），$\Gamma$ 的维数可以大于 $1$.此时仅由 $f$ 的知识来计算 $\Gamma$ 的 Hausdorff 维数与测度可能很难.不过，若 $f$ 满足 Lipschitz 条件，则容易得到一个上界.

\begin{theorem}\label{thm:8.1}
设
\begin{equation}\label{eq:8.1}
  |f(x+h)-f(x)|\leq ch^{2-s}
\end{equation}
对一切 $x$ 及一切满足 $0<h\leq h_0$ 的 $h$ 成立，其中 $c$ 与 $h_0$ 为正常数.则 $\mathcal{H}^s(\Gamma)<\infty$.
\end{theorem}

\begin{proof}
设 $I$ 是 $x$ 轴上长度 $h<h_0$ 的任意区间.由 \eqref{eq:8.1}，在 $I$ 上方取一列正方形可见：$\Gamma$ 的部分 $\{(x,f(x)):x\in I\}$ 可以用至多 $h^{-1}ch^{2-s}+1$ 个边长为 $h$ 的正方形覆盖.于是把 $[0,1]$ 分成 $m$ 等份，每份长 $h=1/m<h_0$，可见
\[
  \mathcal{H}^{s}_{2^{\frac{1}{2}}h}(\Gamma)\leq m\bigl(h^{-1}ch^{2-s}+1\bigr)\bigl(2^{\frac{1}{2}}h\bigr)^{s}\leq c\,2^{1+s/2}
\]
对一列趋于 $0$ 的 $h$ 值成立.从而 $\mathcal{H}^s(\Gamma)\leq c\,2^{1+s/2}<\infty$.
\end{proof}

得到图像具有精细结构的函数的最容易办法是把一系列振荡越来越快的函数相加.于是若 $\sum_{1}^{\infty}|a_i|<\infty$ 且 $\lambda_i\to\infty$，由三角级数
\begin{equation}\label{eq:8.2}
  f(x)=\sum_{i=1}^{\infty}a_i\sin(\lambda_i x)
\end{equation}
定义的函数，在 $a_i$ 与 $\lambda_i$ 选取得当的情况下，可以预期其图像的维数大于 $1$.PLACEHOLDER-NEVER-5
\begin{equation}\label{eq:8.3}
  f(x)=\sum_{i=1}^{\infty}\lambda^{(s-2)i}\sin(\lambda^{i}x),
\end{equation}
其中 $1<s<2$，$\lambda>1$；它是 \newterm{weierstrass-function}{Weierstrass 函数}{Weierstrass function}，连续却处处不可微.Weierstrass 函数的图像很可能具有维数 $s$，但这似乎尚未被严格证明.（维数不能超过 $s$，见\cref{exe:8.3}.）\textcite{mandelbrot1977} 引入了 Weierstrass 函数的一个变式
\[
  f(x)=\sum_{i=-\infty}^{\infty}\lambda^{(s-2)i}\bigl(1-\cos\lambda^{i}x\bigr),
\]
该函数具有尺度性质：对一切 $x$ 有 $f(\lambda x)=\lambda^{2-s}f(x)$；同样其图像的维数应当是 $s$.\textcite{berrylewis1980} 给出了这类函数的计算机实现，并引用了各种物理应用的文献.

为便于计算，把 \eqref{eq:8.2} 中的正弦函数换成形式稍异的周期函数会更方便.设 $g$ 是 $\mathbb{R}$ 上周期为 $4$ 的``锯齿''函数：
\[
  g(4k+x)=
  \begin{cases}
    x & (0\leq x<1)\\
    2-x & (1\leq x<3)\\
    x-4 & (3\leq x<4),
  \end{cases}
\]
其中 $k$ 为整数且 $0\leq x<4$.（特别注意 $g$ 在所有非整数 $x$ 处导数的模为 $1$.）于是我们研究函数
\begin{equation}\label{eq:8.4}
  f(x)=\sum_{i=1}^{\infty}a_i g(\lambda_i x)
\end{equation}
以代替 \eqref{eq:8.2}.\textcite{besicovitchursell1937} 对这类函数中某些函数的图像建立了 Hausdorff 维数：

\begin{theorem}\label{thm:8.2}
设 $\Gamma$ 是函数
\begin{equation}\label{eq:8.5}
  f(x)=\sum_{i=1}^{\infty}\lambda_i^{s-2}g(\lambda_i x)
\end{equation}
（$x\in[0,1]$，$1<s<2$）的图像.设 $\{\lambda_i\}$ 是正数序列，$\lambda_{i+1}/\lambda_i$ 递增于无穷且 $\log\lambda_{i+1}/\log\lambda_i\to 1$.则 $\dim\Gamma=s$.
\end{theorem}

\begin{proof}
若
\begin{equation}\label{eq:8.6}
  \lambda_{k+1}^{-1}\leq h<\lambda_k^{-1},
\end{equation}
则
\begin{align*}
  |f(x+h)-f(x)|
  &\leq\sum_{1}^{k}\lambda_i^{s-2}|g\bigl(\lambda_i(x+h)\bigr)-g(\lambda_i x)|+2\sum_{k+1}^{\infty}\lambda_i^{s-2}\\
  &\leq h\sum_{1}^{k}\lambda_i^{s-1}+2\sum_{k+1}^{\infty}\lambda_i^{s-2}
  \leq 2h\lambda_k^{s-1}+4\lambda_{k+1}^{s-2},
\end{align*}
只要 $k$ 充分大.于是由 \eqref{eq:8.6}，当 $h$ 充分小时
\[
  |f(x+h)-f(x)|\leq 2h^{2-s}+4h^{2-s}=6h^{2-s}.
\]
由\cref{thm:8.1} 得 $\mathcal{H}^s(\Gamma)<\infty$.

$\Gamma$ 维数的下估计则相当棘手.设 $S$ 是边长为 $h$、边平行于坐标轴的正方形.设 $I$ 是 $S$ 到 $x$ 轴上投影的区间.我们证明集合 $E=\{x:(x,f(x))\in S\}$ 的勒贝格测度不能太大.

定义部分和
\[
  f_k(x)=\sum_{1}^{k}\lambda_i^{s-2}g(\lambda_i x).
\]
以下始终设 $k$ 大到足以保证 $\lambda_{k+1}\geq 2\lambda_k\geq 2$，使
\begin{equation}\label{eq:8.7}
  |f(x)-f_k(x)|=\Bigl|\sum_{k+1}^{\infty}\lambda_i^{s-2}g(\lambda_i x)\Bigr|\leq\sum_{k+1}^{\infty}\lambda_i^{s-2}\leq 2\lambda_{k+1}^{s-2},
\end{equation}
并且（对非例外的 $x$）
\begin{equation}\label{eq:8.8}
  |f_k'(x)|=\Bigl|\sum_{1}^{k}\lambda_i^{s-1}g'(\lambda_i x)\Bigr|\geq\lambda_k^{s-1}-\sum_{1}^{k-1}\lambda_i^{s-1}\geq\frac{1}{2}\lambda_k^{s-1}.
\end{equation}

先设正方形 $S$ 的边 $h=\lambda_k^{-1}$ 对某个这样的 $k$ 成立.设 $m$ 是满足
\begin{equation}\label{eq:8.9}
  \lambda_{k+m}^{s-2}\leq h=\lambda_k^{-1}<\lambda_{k+m-1}^{s-2}
\end{equation}
的整数.当然 $m\geq 1$.另一方面，由于 $\lambda_{i+1}/\lambda_i$ 递增，
\[
  \Bigl(\frac{\lambda_{k+1}}{\lambda_k}\Bigr)^{(m-1)(2-s)}\lambda_k^{2-s}
  <\Bigl(\frac{\lambda_{k+m-1}}{\lambda_{k+m-2}}\cdots\frac{\lambda_{k+1}}{\lambda_k}\lambda_k\Bigr)^{2-s}
  =\lambda_{k+m-1}^{2-s}<\lambda_k,
\]
故
\[
  \Bigl(\frac{\lambda_{k+1}}{\lambda_k}\Bigr)^{(m-1)(2-s)}
  <\lambda_k^{s-1}
  <\Bigl(\frac{\lambda_k}{\lambda_{k-1}}\cdots\frac{\lambda_2}{\lambda_1}\lambda_1\Bigr)^{s-1}
  <\Bigl(\frac{\lambda_{k+1}}{\lambda_k}\Bigr)^{(k-1)(s-1)}\lambda_1^{s-1}.
\]
于是取对数得
\begin{equation}\label{eq:8.10}
  m\leq ak,
\end{equation}
其中 $a$ 与 $k$ 无关.

若 $m=1$，则由 \eqref{eq:8.7}，$(x,f(x))$ 能落在 $S$ 中仅当 $(x,f_k(x))$ 落在把 $S$ 向上、下各扩大距离 $2\lambda_{k+1}^{s-2}\leq 2h$ 所得矩形 $S_1$ 之中.导数 $f_k'(x)$ 在区间 $I$ 内至多变号一次.在 $f_k'(x)$ 恒号的每一段上 $|f_k'(x)|\geq\frac{1}{2}\lambda_k^{s-1}$，故 $(x,f_k(x))$ 能落在 $S_1$ 中的 $x$ 只能属于长度至多为 $2\lambda_k^{1-s}$ 乘以 $S_1$ 的高度的区间.于是
\begin{equation}\label{eq:8.11}
  \mathscr{L}^{1}(E)\leq 2\cdot 2\lambda_k^{1-s}\cdot 5h=20h^{s}.
\end{equation}
若 $m>1$，则可把 $I$ 分成至多两部分，使 $f_k'(x)$ 在每部分恒号.由 \eqref{eq:8.9}，$S_1$ 的高为 $h+4\lambda_{k+1}^{s-2}\leq 5\lambda_{k+1}^{s-2}$，故在每一部分中寻找 $E$ 的点时只需考虑属于长度为 $2\lambda_k^{1-s}\cdot 5\lambda_{k+1}^{s-2}$ 的子区间的 $x$.再把每个这样的子区间进一步分成 $f_{k+1}'(x)$ 也恒号的部分.这样，从每个旧区间至多得到 $2\lambda_k^{1-s}\cdot 5\lambda_{k+1}^{s-2}\cdot\frac{1}{2}\lambda_{k+1}+1\leq 6\bigl(\lambda_{k+1}/\lambda_k\bigr)^{s-1}$ 个新区间.

若 $m>2$，重复该过程，从上一组区间中的每个至多得到 $6\bigl(\lambda_{k+2}/\lambda_{k+1}\bigr)^{s-1}$ 个使 $f_{k+2}'(x)$ 也恒号的区间.如此进行，最终看到 $E$ 可以被至多
\[
  2\cdot 6^{m-1}\Bigl(\frac{\lambda_{k+1}}{\lambda_k}\cdot\frac{\lambda_{k+2}}{\lambda_{k+1}}\cdots\frac{\lambda_{k+m-1}}{\lambda_{k+m-2}}\Bigr)^{s-1}
  =2\cdot 6^{m-1}\Bigl(\frac{\lambda_{k+m-1}}{\lambda_k}\Bigr)^{s-1}
\]
个使 $f_{k+m-1}'(x)$ 恒号的区间覆盖.由 \eqref{eq:8.7}，在每个这样的区间上，$(x,f(x))\in S$ 仅当 $(x,f_{k+m-1}(x))\in S_2$，其中 $S_2$ 是把 $S$ 向上、下各扩大距离 $2\lambda_{k+m}^{s-2}$ 所得的矩形.$S_2$ 的高为 $h+4\lambda_{k+m}^{s-2}\leq 5h$，于是考虑 $f_{k+m-1}$ 在每个这样的区间上的梯度并利用 \eqref{eq:8.8} 与 \eqref{eq:8.9}，得
\[
  \mathscr{L}^{1}(E)
  \leq 2\cdot 6^{m-1}\Bigl(\frac{\lambda_{k+m-1}}{\lambda_k}\Bigr)^{s-1}\cdot 5h\cdot 2\lambda_{k+m-1}^{1-s}
  \leq 20\cdot 6^{m-1}h^{s}\leq 20\cdot 6^{ak}h^{s}.
\]
于是存在常数 $b$ 与 $c$，使当 $h=\lambda_k^{-1}$ 时 $\mathscr{L}^1(E)\leq cb^kh^s$.

现在设 $S$ 是边长 $h$ 满足 $\lambda_{k+1}^{-1}<h\leq\lambda_k^{-1}$ 的正方形.由上面的结果，当 $t<s$ 时
\[
  \mathscr{L}^{1}(E)\leq cb^{k}\lambda_k^{-s}
  =c\lambda_{k+1}^{-t}\frac{\lambda_{k+1}^{t}}{\lambda_k^{(s+t)/2}}\frac{b^{k}}{\lambda_k^{(s-t)/2}}.
\]
从而
\begin{equation}\label{eq:8.12}
  \mathscr{L}^{1}(E)\leq c_1h^{t},
\end{equation}
因为由 $\lambda_k$ 的增长条件，当 $k$ 充分大时 $\lambda_k^{(s+t)/2}$ 比 $\lambda_{k+1}^{t}$ 增长得快，且 $\lambda_k^{(s-t)/2}$ 比 $b^k$ 增长得快.

若 $\{U_i\}$ 是 $\Gamma$ 的任意覆盖，把每个 $U_i$ 放进边长为 $|U_i|$ 的正方形 $S_i$ 中.记 $E_i=\{x:(x,f(x))\in S_i\}$，则必有 $[0,1]\subset\bigcup_i E_i$.于是由 \eqref{eq:8.12}，
\[
  \sum|U_i|^{t}=\sum 2^{-\frac{1}{2}t}|S_i|^{t}\geq c_1^{-1}\sum\mathscr{L}^{1}(E_i)\geq c_1^{-1}.
\]
故当 $t<s$ 时 $\mathcal{H}^t(\Gamma)>c_1^{-1}>0$，于是 $\dim\Gamma=s$.
\end{proof}

注意，同样的论证表明：在关于 $\lambda_i$ 的相同条件下，函数
\[
  f(x)=\sum_{i=1}^{\infty}\lambda_i^{s-2}g(\lambda_i x+\theta_i)
\]
的图像的维数为 $s$，其中诸 $\theta_i$ 是固定的``相位''.

\textcite{besicovitchursell1937} 用非常类似的方法证明：若在 \eqref{eq:8.5} 中对所有 $i$ 有 $\lambda_{i+1}=\lambda_i^{\mu}$，其中 $\mu=(s-1)(2-t)/(2-s)(t-1)$，$1<t<s$，则 $\dim\Gamma=t$.\textcites{loveyoung1937,kline1945} 给出了一些进一步的计算.

人们还引入了曲线``维数''的各种其他定义，并与 Hausdorff 维数比较.例如，可以考虑当分割的最大区间长度趋于零时 $\sum_{i=1}^{m}|f(x_i)-f(x_{i-1})|^s$ 的下极限或上极限（见 \textcite{ville1936}）.另外，若 $\Gamma(x,x+h)$ 表示 $\Gamma$ 对应于区间 $[x,x+h]$ 的部分，则当 $h\to 0$ 时 $\mathcal{H}^s\bigl(\Gamma(x,x+h)\bigr)/h^s$ 的极限有一些有趣的性质（Besicovitch（\citeyear{besicovitch1964b,besicovitch1965a,besicovitch1967}））.

\section{自相似集}

许多经典的分形集都是``自相似的''，即由与整个集合几何相似但尺度更小的部分构成.最简单的例子是 Cantor 集（见 1.4~节）：$E\cap[0,\frac{1}{3}]$ 与 $E\cap[\frac{2}{3},1]$ 都与 $E$ 相似但按因子 $\frac{1}{3}$ 缩放；$E\cap[0,\frac{1}{9}]$、$E\cap[\frac{2}{9},\frac{1}{3}]$、$E\cap[\frac{2}{3},\frac{7}{9}]$ 与 $E\cap[\frac{8}{9},1]$ 都以比例因子 $\frac{1}{9}$ 与 $E$ 相似，如此等等.平面中的两个例子描绘于图~\ref{fig:8.1} 与图~\ref{fig:8.2}；图~\ref{fig:8.1} 显示的是熟悉的 \newterm{koch-curve}{Koch 曲线}{Koch curve}即雪花曲线.Mandelbrot（\citeyear{mandelbrot1977,mandelbrot1982}） 的书与 \textcite{dekking1982} 的文章给出了许多其他美丽自相似集的计算机实现，\textcite{stepney1984} 给出了绘制这类集合的计算机程序.

整个世纪以来人们都知道自相似集的例子，但只是近来才有人尝试把它们的理论置于系统的基础之上.\textcite{dekking1982} 描述了利用自由群中词的自同态做到这一点的一般方法.这里我们给出 \textcite{hutchinson1981} 优雅处理的一个版本（另见 \textcite{moran1946} 与 \textcite{marion1979}）.

\begin{figure}[htbp]
  \centering
  \includegraphics[width=0.66\textwidth]{fig8-1.png}
  \caption{Koch 曲线及其生成元}
  \label{fig:8.1}
\end{figure}

\begin{figure}[htbp]
  \centering
  \includegraphics[width=0.62\textwidth]{fig8-2.png}
  \caption{两个平面自相似集及其生成元}
  \label{fig:8.2}
\end{figure}

称映射 $\psi:\mathbb{R}^n\to\mathbb{R}^n$ 为\newterm{contraction}{压缩}{contraction}，若 $|\psi(x)-\psi(y)|\leq c|x-y|$ 对一切 $x,y\in\mathbb{R}^n$ 成立，其中 $c<1$.显然任何压缩都是连续映射.称使该不等式对一切 $x,y$ 成立的 $c$ 的下确界为该压缩的\newterm{ratio}{压缩比}{ratio of a contraction}.把 $\mathbb{R}^n$ 的每个子集都变为几何相似集合的压缩称为\newterm{similitude}{相似压缩}{similitude}.于是相似压缩是伸缩、旋转、平移以及或许一个反射的复合；其压缩比恰是相似压缩的尺度因子.

就本节而言，称集合 $E\subset\mathbb{R}^n$ 对一组压缩 $\psi_1,\ldots,\psi_m$ 是不变的（\newterm{invariant}{不变集}{invariant set}），若 $E=\bigcup_{1}^{m}\psi_j(E)$.若此外这些压缩都是相似压缩，且对某个 $s$ 有 $\mathcal{H}^s(E)>0$ 而 $\mathcal{H}^s\bigl(\psi_i(E)\cap\psi_j(E)\bigr)=0$（$i\neq j$），则称 $E$ 是\newterm{self-similar}{自相似的}{self-similar set}.（这一测度条件保证 $E$ 的自相似特征不因重叠而丧失.）我们将证明：对任意有限压缩组，存在唯一非空紧不变集 $E$；而且对满足``开集条件''的相似压缩组，集合 $E$ 是一个 $s$-集，$s$ 可由压缩比算出.为明确思路：Cantor 集是实直线上相似压缩
\[
  \psi_1(x)=x/3,\qquad \psi_2(x)=(2+x)/3
\]
下不变的唯一紧集.

设 $\{\psi_j\}_{1}^{m}$ 是一组压缩，令 $\psi$ 表示由
\[
  \psi(F)=\bigcup_{1}^{m}\psi_j(F)
\]
定义的 $\mathbb{R}^n$ 子集的变换.记 $\psi$ 的迭代为 $\psi^0(F)=F$ 与 $\psi^{k+1}(F)=\psi(\psi^k(F))$（$k\geq 0$）.

注意，本节的工作对定义在 $\mathbb{R}^n$ 的紧子集上的压缩同样有效.

\begin{theorem}\label{thm:8.3}
给定 $\mathbb{R}^n$ 上压缩比为 $r_j<1$ 的压缩组 $\{\psi_j\}_{1}^{m}$，存在唯一非空紧集 $E$ 使
\begin{equation}\label{eq:8.13}
  E=\psi(E)=\bigcup_{1}^{m}\psi_j(E).
\end{equation}
进一步，若 $F$ 是 $\mathbb{R}^n$ 的任何非空紧子集，则当 $k\to\infty$ 时迭代 $\psi^k(F)$ 在 Hausdorff 度量下收敛于 $E$.
\end{theorem}

\begin{proof}
设 $\mathcal{C}$ 是 $\mathbb{R}^n$ 的所有非空紧子集构成的类.由 Blaschke 选择定理（\cref{thm:3.16}）证明中的完备性部分，$\mathcal{C}$ 赋予 Hausdorff 度量 $\delta$ 后成为完备度量空间.若 $F_1,F_2\in\mathcal{C}$，则由 $\delta$ 的定义，
\[
  \delta\bigl(\psi(F_1),\psi(F_2)\bigr)
  =\delta\Bigl(\bigcup_{1}^{m}\psi_j(F_1),\bigcup_{1}^{m}\psi_j(F_2)\Bigr)
  \leq\max_{j}\delta\bigl(\psi_j(F_1),\psi_j(F_2)\bigr)
  \leq\bigl(\max_{j}r_j\bigr)\delta(F_1,F_2).
\]
（由于 $\psi_j$ 是压缩，若 $F_1$ 含于 $F_2$ 的 $\delta$-平行体之中，则 $\psi_j(F_1)$ 含于 $\psi_j(F_2)$ 的 $r_j\delta$-平行体之中.）由于 $\max_j r_j<1$，$\psi$ 是 $\mathcal{C}$ 上的压缩映射.由完备度量空间的压缩映射定理，存在唯一的 $E\in\mathcal{C}$ 使 $\psi(E)=E$，而且对任意 $F\in\mathcal{C}$ 有 $\delta(\psi^k(F),E)\to 0$（$k\to\infty$）.

可以证明，不变集 $E$ 是映射 $\psi_{j_1}\circ\psi_{j_2}\circ\cdots\circ\psi_{j_k}$（对满足 $1\leq j_i\leq m$ 的一切有限序列 $\{j_1,j_2,\ldots,j_k\}$ 取）的不动点之集的闭包.（我们采用约定 $(\psi_1\circ\psi_2)(x)=\psi_1(\psi_2(x))$ 等.）
\end{proof}

现在设 $s$ 是满足
\[
  \sum_{1}^{m}r_j^{s}=1
\]
的数.把它 $k$ 次乘方，得到有用的恒等式
\begin{equation}\label{eq:8.14}
  \sum_{j_1\ldots j_k}(r_{j_1}r_{j_2}\cdots r_{j_k})^{s}=1,
\end{equation}
其中求和取遍满足 $1\leq j_i\leq m$ 的一切 $k$ 元组 $\{j_1\ldots j_k\}$.对任何这样的序列与任何集合 $F$，写
\[
  F_{j_1\ldots j_k}=\psi_{j_1}\circ\cdots\circ\psi_{j_k}(F).
\]

下面的引理可视为\cref{thm:8.3} 在测度上的类似物；事实上 \textcite{hutchinson1981} 用压缩映射定理给出了一个证明.

\begin{lemma}\label{lem:8.4}
存在 Borel 测度 $\mu$，其支集含于 $E$，$\mu(\mathbb{R}^n)=1$，且对任何可测集 $F$ 有
\begin{equation}\label{eq:8.15}
  \mu(F)=\sum_{j=1}^{m}r_j^{s}\mu\bigl(\psi_j^{-1}(F)\bigr).
\end{equation}
\end{lemma}

\begin{proof}
取 $x\in E$，并记 $x_{j_1\ldots j_k}=\psi_{j_1}\circ\cdots\circ\psi_{j_k}(x)$.对 $k=1,2,\ldots$，在 $E$ 上连续函数空间 $\mathscr{F}$ 上如下定义正线性泛函：
\begin{equation}\label{eq:8.16}
  \phi_k(f)=\sum_{j_1\ldots j_k}(r_{j_1}\cdots r_{j_k})^{s}f\bigl(x_{j_1\ldots j_k}\bigr).
\end{equation}
若 $f\in\mathscr{F}$，则 $f$ 在 $E$ 上一致连续，故给定 $\varepsilon>0$，可找到 $p$ 使得当 $k\geq p$ 时 $f$ 在每个 $E_{j_1\ldots j_k}$ 上的变动小于 $\varepsilon$（回忆 $|E_{j_1\ldots j_k}|\leq(\max_j r_j)^k|E|$）.由于当 $k\geq p$ 时 $x_{j_1\ldots j_k}\in E_{j_1\ldots j_p}$，由 \eqref{eq:8.16} 与 \eqref{eq:8.14} 得：当 $k,k'\geq p$ 时 $|\phi_k(f)-\phi_{k'}(f)|\leq\varepsilon$.由收敛的一般原理，$\{\phi_k(f)\}_k$ 对每个 $f$ 收敛，其极限在 $\mathscr{F}$ 上定义了一个正线性泛函.由 Riesz 表示定理（\cref{thm:6.2}），存在 Borel 测度 $\mu$ 使
\begin{equation}\label{eq:8.17}
  \int f\,\mathrm{d}\mu=\lim_{k\to\infty}\phi_k(f)
\end{equation}
（$f\in\mathscr{F}$）.取 $f\equiv 1$，由 \eqref{eq:8.16} 显然有 $\mu(\mathbb{R}^n)=1$.对任意 $f\in\mathscr{F}$，
\begin{align*}
  \phi_k(f)
  &=\sum_{j_1}r_{j_1}^{s}\sum_{j_2\ldots j_k}(r_{j_2}\cdots r_{j_k})^{s}f\bigl(x_{j_1\ldots j_k}\bigr)\\
  &=\sum_{j_1}r_{j_1}^{s}\sum_{j_2\ldots j_k}(r_{j_2}\cdots r_{j_k})^{s}f\bigl(\psi_{j_1}(x_{j_2\ldots j_k})\bigr)
  =\sum_{j}r_j^{s}\phi_{k-1}(f\circ\psi_j).
\end{align*}
令 $k\to\infty$，利用 \eqref{eq:8.17} 得：若 $f\in\mathscr{F}$，则
\[
  \int f\,\mathrm{d}\mu=\sum_{j}r_j^{s}\int f\circ\psi_j\,\mathrm{d}\mu.
\]
利用单调收敛定理作通常的逼近，这对一切非负函数 $f$ 也成立，故 \eqref{eq:8.15} 得证.

最后，若 $f$ 是在 $E$ 上为零的任何连续函数，则由 \eqref{eq:8.16} 对一切 $k$ 有 $\phi_k(f)=0$，故由 \eqref{eq:8.17} 有 $\int f\,\mathrm{d}\mu=0$.于是 $\mu$ 由 $E$ 支撑.
\end{proof}

由此，记
\[
  \mu_{j_1\ldots j_k}(F)=\mu\bigl((\psi_{j_1}\circ\cdots\circ\psi_{j_k})^{-1}(F)\bigr)
  =\mu\bigl(\psi_{j_k}^{-1}\circ\cdots\circ\psi_{j_1}^{-1}(F)\bigr),
\]
则测度 $\mu_{j_1\ldots j_k}$ 的支集含于 $E_{j_1\ldots j_k}$.又由 \eqref{eq:8.15},
\begin{equation}\label{eq:8.18}
  \mu_{j_1\ldots j_k}=\sum_{j}r_j^{s}\mu_{j_1\ldots j_k j}.
\end{equation}

称开集条件（\newterm{open set condition}{开集条件}{open set condition}）对压缩 $\{\psi_j\}_{1}^{m}$ 成立，若存在有界开集 $V$ 使
\begin{equation}\label{eq:8.19}
  \psi(V)=\bigcup_{j=1}^{m}\psi_j(V)\subset V
\end{equation}
且这个并不交.经 $\psi_{j_1\ldots j_k}$ 变换得
\[
  \bigcup_{j=1}^{m}V_{j_1\ldots j_k j}\subset V_{j_1\ldots j_k},
\]
同样是不交并.于是集合 $\{V_{j_1\ldots j_k}\}$（$k$ 任意）构成一个网：该族中任意两个集合要么不交，要么一个含于另一个.

若 \eqref{eq:8.19} 成立，则 $E\subset\overline{V}$（横线表示闭包）；事实上
\begin{equation}\label{eq:8.20}
  E=\bigcap_{k=0}^{\infty}\psi^{k}(\overline{V}),
\end{equation}
因为 $\{\psi^k(\overline{V})\}_k$ 是由\cref{thm:8.3} 在 Hausdorff 度量下收敛于 $E$ 的紧集递减序列，若 $E$ 在 $\overline{V}$ 之外有点则这是不可能的.在 $\psi_{j_1}\circ\cdots\circ\psi_{j_k}$ 下取像得 $E_{j_1\ldots j_k}\subset\overline{V}_{j_1\ldots j_k}$.（当然，由于 $\psi_j$ 连续，有 $\psi_j(\overline{V})\subset\overline{\psi_j(V)}$ 等.）

现在设 $\{\psi_j\}$ 是相似压缩.我们将证明：若开集条件成立，则不变集 $E$ 自相似，且 $E$ 的 Hausdorff 维数与相似维数相等.\newterm{similarity dimension}{相似维数}{similarity dimension}易于计算，它是使
\[
  \sum_{1}^{m}r_j^{s}=1
\]
的唯一正数 $s$.开集条件保证集合 $\psi_j(E)$ 不能重叠太多.

\begin{lemma}\label{lem:8.5}
设 $\{V_i\}$ 是 $\mathbb{R}^n$ 的不交开子集之族，每个 $V_i$ 含有一个半径为 $c_1\rho$ 的球且含于一个半径为 $c_2\rho$ 的球.则任何半径为 $\rho$ 的球 $B$ 至多与 $(1+2c_2)^n c_1^{-n}$ 个 $\overline{V}_i$ 相交.
\end{lemma}

\begin{proof}
若 $\overline{V}_i$ 与 $B$ 相交，则 $\overline{V}_i$ 含于一个与 $B$ 同心、半径为 $(1+2c_2)\rho$ 的球.若 $\{\overline{V}_i\}$ 中有 $q$ 个与 $B$ 相交，则把相应内含球的体积求和得 $q(c_1\rho)^n\leq(1+2c_2)^n\rho^n$，给出所述的 $q$ 的界.
\end{proof}

\begin{theorem}\label{thm:8.6}
设开集条件对压缩比为 $r_j$（$1\leq j\leq m$）的相似压缩 $\psi_j$ 成立.则相伴的紧不变集 $E$ 是一个 $s$-集，其中 $s$ 由
\begin{equation}\label{eq:8.21}
  \sum_{1}^{m}r_j^{s}=1
\end{equation}
确定；特别地 $0<\mathcal{H}^s(E)<\infty$.
\end{theorem}

\begin{proof}
迭代 \eqref{eq:8.13} 得 $E=\bigcup_{j_1\ldots j_k}E_{j_1\ldots j_k}$.由 \eqref{eq:8.14},
\[
  \sum_{j_1\ldots j_k}|E_{j_1\ldots j_k}|^{s}
  =\sum_{j_1\ldots j_k}|E|^{s}(r_{j_1}\cdots r_{j_k})^{s}
  =|E|^{s}.
\]
由于当 $k\to\infty$ 时 $|E_{j_1\ldots j_k}|\leq(\max_j r_j)^k|E|\to 0$，我们得出 $\mathcal{H}^s(E)\leq|E|^s<\infty$.

下界用两条引理得到.设使 \eqref{eq:8.19} 成立的开集 $V$ 含有半径为 $c_1$ 的球并含于半径为 $c_2$ 的球.取任意 $\rho>0$.对每个无穷序列 $\{j_1,j_2,\ldots\}$（$1\leq j_i\leq m$），在该序列于使
\begin{equation}\label{eq:8.22}
  \bigl(\min_{j}r_j\bigr)\rho\leq r_{j_1}\cdots r_{j_k}\leq\rho
\end{equation}
的最小 $k\geq 1$ 处截断，并记 $\mathcal{S}$ 为由此得到的有限序列之集.由开集的网性质，$\{V_{j_1\ldots j_k}:j_1\ldots j_k\in\mathcal{S}\}$ 是不交族.每个这样的 $V_{j_1\ldots j_k}$ 含有半径为 $c_1r_{j_1}\cdots r_{j_k}$ 从而由 \eqref{eq:8.22} 半径为 $c_1(\min_j r_j)\rho$ 的球，类似地含于半径为 $c_2r_{j_1}\cdots r_{j_k}$ 从而半径为 $c_2\rho$ 的球.由\cref{lem:8.5}，任何半径为 $\rho$ 的球 $B$ 至多与族 $\{\overline{V}_{j_1\ldots j_k}:j_1\ldots j_k\in\mathcal{S}\}$ 中的 $q=(1+2c_2)^n c_1^{-n}(\min_j r_j)^{-n}$ 个集合相交.又对任何 $\{j_1,\ldots,j_k\}$ 有 $\mu_{j_1\ldots j_k}(\mathbb{R}^n)=1$ 且 $\operatorname{support}(\mu_{j_1\ldots j_k})\subset E_{j_1\ldots j_k}\subset\overline{V}_{j_1\ldots j_k}$.按需要迭代 \eqref{eq:8.18} 可见
\[
  \mu=\sum_{j_1\ldots j_k\in\mathcal{S}}(r_{j_1}\cdots r_{j_k})^{s}\mu_{j_1\ldots j_k},
\]
故 $\mu(B)\leq\sum(r_{j_1}\cdots r_{j_k})^s\mu_{j_1\ldots j_k}(\mathbb{R}^n)$，其中求和取遍 $\mathcal{S}$ 中使 $\overline{V}_{j_1\ldots j_k}$ 与 $B$ 相交的序列 $\{j_1,\ldots,j_k\}$.于是利用 \eqref{eq:8.22}，对任何满足 $|B|<|V|$ 的球有 $\mu(B)\leq q\rho^s=q2^{-s}|B|^s$.但给定 $E$ 的任何覆盖 $\{U_i\}$，可用满足 $|B_i|\leq 2|U_i|$ 的球 $\{B_i\}$ 覆盖 $E$，故
\[
  1=\mu(E)\leq\sum\mu(B_i)\leq q2^{-s}\sum|B_i|^{s}\leq q\sum|U_i|^{s}.
\]
可以选取 $\{U_i\}$ 使 $\sum|U_i|^s$ 任意接近 $\mathcal{H}^s(E)$，故 $\mathcal{H}^s(E)\geq q^{-1}>0$，如所需.
\end{proof}

\begin{corollary}\label{cor:8.7}
若开集条件成立，则 $\mathcal{H}^s\bigl(\psi_i(E)\cap\psi_j(E)\bigr)=0$（$i\neq j$），特别地 $E$ 是自相似的.
\end{corollary}

\begin{proof}
利用 $\psi_j$ 是相似压缩这一事实，
\[
  \sum_{j=1}^{m}\mathcal{H}^{s}\bigl(\psi_j(E)\bigr)=\sum_{j=1}^{m}r_j^{s}\mathcal{H}^{s}(E)=\mathcal{H}^{s}(E).
\]
由于 $0<\mathcal{H}^s(E)<\infty$，由 \eqref{eq:8.13} 与 $\mathcal{H}^s$ 的可加性，这只有在 $\mathcal{H}^s\bigl(\psi_i(E)\cap\psi_j(E)\bigr)=0$（$i\neq j$）时才可能.
\end{proof}

若 $\psi_j$ 是压缩而非相似压缩，这一技术仍可用于估计 Hausdorff 维数.这里给出一个有用的情形（另见\cref{exe:8.5}）.

\begin{theorem}\label{thm:8.8}
设 $\{\psi_j\}_{1}^{m}$ 是 $\mathbb{R}$ 上的压缩，开集条件 \eqref{eq:8.19} 对开区间 $V$ 成立.设对每个 $j$ 与一切 $x,y\in\overline{V}$ 有
\begin{equation}\label{eq:8.23}
  q_{j}|x-y|\leq|\psi_{j}(x)-\psi_{j}(y)|\leq r_{j}|x-y|.
\end{equation}
则 $s\leq\dim E\leq t$，其中 $s$ 与 $t$ 由 $\sum_{1}^{m}q_j^{s}=1=\sum_{1}^{m}r_j^{t}$ 定义.
\end{theorem}

\begin{proof}
\cref{thm:8.6} 的方法给出上界.为得下界，照那里的方法进行，但不使用\cref{lem:8.5}，而使用如下事实：若 $\{V_i\}$ 是长度至少为 $2c_1\rho$ 的不交区间之族，则任何长度为 $2\rho$ 的区间 $B$ 至多与 $c_1^{-1}+2$ 个 $\overline{V}_i$ 相交.
\end{proof}

当 $\psi_j$ 仅仅是压缩时，通过对 $k>1$ 的映射 $\{\psi_{j_1\ldots j_k}\}$ 进行工作并用更紧的估计替换 \eqref{eq:8.23}，通常可以得到 $\dim E$ 更好的估计.

1.5~节的方法也可用于估计这类集合的维数.

\textcite{mandelbrot1982} 引入了自相似集的\newterm{generator}{生成元}{generator}的概念，它方便地指定一组相似压缩.一个生成元由若干直线段与两个特别标定的点组成.我们把映射这两个特殊点到线段端点的相似压缩与每条线段相联系.利用生成元，通过迭代``把每条线段换成生成元的相似拷贝''的过程，可以构造逼近自相似集的集合序列.这样得到的集合序列正是 $\psi^k(G)$，其中 $G$ 是生成元.注意在 $\mathbb{R}^2$ 中相似压缩只能在相差一个反射的意义下确定，但其定向可以通过展示第一次迭代来指定.由生成元利用\cref{thm:8.6} 容易算出集合的维数.这一思想最好用例子说明.图~\ref{fig:8.1}(b) 与图~\ref{fig:8.2}(b) 显示了图~\ref{fig:8.1}(a) 与图~\ref{fig:8.2}(a) 中曲线的生成元及逼近的前几个阶段.

作为一个变式，可以在这类构造中引入随机因素.若在迭代过程的每一步按某种概率分布随机选取 $\psi_j$，也可能发展出一套理论.此时所得集合本身不是自相似的，但在统计意义上是自相似的：集合与其分支子集具有相同的概率分布.在合理条件下，这类集合以概率 $1$ 达到一个特定的维数，它可以用 $\psi_j$ 的压缩比的期望表示.

\section{密切填装与 Apollonian 填装}

设 $V$ 是 $\mathbb{R}^n$ 的有界开子集.称闭球序列 $\{B_i\}_{1}^{\infty}$ 为 $V$ 的一个\newterm{osculatory packing}{密切填装}{osculatory packing}，如果对每个 $j$，$B_j$ 是含于 $V\setminus\bigcup_{1}^{j-1}B_i$ 中的最大球（或最大球之一）.最有趣的情形称为 \newterm{Apollonian packing}{Apollonian 填装}{Apollonian packing}，当 $V$ 是 $\mathbb{R}^2$ 中一个曲边三角形的内部时得到（这一填装的前几个阶段见图~\ref{fig:8.3}）.

特别有意义的是闭集 $E=V\setminus\bigcup_{1}^{\infty}B_i$ 的 Hausdorff 维数，它称为填装的\newterm{residual set}{剩余集}{residual set}.对 Apollonian 填装显然有 $1\leq\dim E\leq 2$，但证明这两个不等式都是严格的就出奇地困难，更不用说给出 $\dim E$ 的好的估计了.\textcite{hirst1967} 最先得到非平凡的估计，这里我们概述他的论证，并加入源自 \textcite{boyd1973b} 的一些改进.某些繁琐却直截了当的代数计算的细节从略.

\begin{figure}[htbp]
  \centering
  \includegraphics[width=0.5\textwidth]{fig8-3.png}
  \caption{Apollonian 填装}
  \label{fig:8.3}
\end{figure}

利用初等三角学可以证明下述关于曲边三角形内切圆半径与外接圆半径的公式.内切圆半径的表达式即著名的 \newterm{soddy-formula}{Soddy 公式}{Soddy's formula}，其证明见 \textcite{coxeter1961}（第 12 页），并被物理学家 \textcite{soddy1936} 以韵文形式表述.

\begin{lemma}\label{lem:8.9}
设 $T$ 是由半径为 $a,b,c$ 的相互外切的圆构成的曲边三角形.则 $T$ 的内切圆半径 $r$ 与外接圆半径 $\rho$ 由
\[
  r^{-1}=a^{-1}+b^{-1}+c^{-1}+2\bigl(b^{-1}c^{-1}+a^{-1}c^{-1}+a^{-1}b^{-1}\bigr)^{\frac{1}{2}}
\]
与
\[
  \rho^{2}=(abc)/(a+b+c)=\bigl(b^{-1}c^{-1}+a^{-1}c^{-1}+a^{-1}b^{-1}\bigr)^{-1}
\]
给出.
\end{lemma}

设 $T$ 是任意曲边三角形，边为半径 $a,b,c$ 的圆，其内切圆盘为 $B$.设 $T_1,T_2,T_3$ 是从 $T$ 中去掉 $B$ 后剩下的三个曲边三角形.设 $T,T_j$ 的外接圆半径分别为 $\rho,\rho_j$.

\begin{lemma}\label{lem:8.10}
若 $s\geq\log 3/\log\bigl(1+2\cdot 3^{-\frac{1}{2}}\bigr)=1.4311\ldots$，则 $\rho^{s}\geq\rho_1^{s}+\rho_2^{s}+\rho_3^{s}$.
\end{lemma}

\begin{proof}
对 $T$ 与 $T_1$ 应用\cref{lem:8.9}，经一些计算得
\[
  \rho/\rho_1=1+\rho\bigl(b^{-1}+c^{-1}\bigr)
\]
（$T$ 与 $T_1$ 公共边的半径为 $b$ 与 $c$），对 $T_2,T_3$ 有类似公式.于是
\[
  \rho_1^{s}+\rho_2^{s}+\rho_3^{s}
  =\rho^{s}\bigl[\bigl(1+(b^{-1}+c^{-1})\rho\bigr)^{-s}+\bigl(1+(a^{-1}+c^{-1})\rho\bigr)^{-s}
  +\bigl(1+(a^{-1}+b^{-1})\rho\bigr)^{-s}\bigr],
\]
其中 $\rho=(b^{-1}c^{-1}+a^{-1}c^{-1}+a^{-1}b^{-1})^{-\frac{1}{2}}$.方括号中的表达式关于 $a,b,c$ 齐次，故为证它对任意 $a,b,c$ 都以 $1$ 为上界，可设 $\rho=1$.这是一个约束下的优化问题；只要 $s\geq\log 3/\log\bigl(1+2\cdot 3^{-\frac{1}{2}}\bigr)$，不难得到所需的界，见 \textcite{hirst1967}.
\end{proof}

现在设 $T$ 是边为单位半径圆的等边曲边三角形.从 $T$ 中去掉闭内切圆盘 $B$ 后剩下三个曲边三角形 $T_1,T_2,T_3$.去掉这三个三角形的内切圆盘 $B_1,B_2,B_3$ 后留下九个曲边三角形 $T_{j_1j_2}$（$1\leq j_1,j_2\leq 3$），其内切圆盘为 $B_{j_1j_2}$.如此进行，每一阶段从 $T_{j_1\ldots j_k}$ 中去掉内切圆盘 $B_{j_1\ldots j_k}$，得到曲边三角形 $T_{j_1\ldots j_k1}$、$T_{j_1\ldots j_k2}$、$T_{j_1\ldots j_k3}$.（我们一致地分配指标 $1,2,3$ 以指示三角形相对 $B_{j_1\ldots j_k}$ 的位置.）注意，如此得到的曲边三角形的内部构成一个集合之网.对满足 $1\leq j_i\leq 3$ 的一切有限序列 $\mathbf{j}=\{j_1,\ldots,j_k\}$ 取圆盘族 $\{B_{j_1\ldots j_k}\}$，便得到 $T$ 的 Apollonian 填装.

\begin{theorem}\label{thm:8.11}
若 $E$ 是 Apollonian 填装的剩余集，则 $\dim E\leq\log 3/\log\bigl(1+2\cdot 3^{-\frac{1}{2}}\bigr)=1.4311\ldots$.
\end{theorem}

\begin{proof}
给定 $\delta>0$，可选取足够大的 $k$ 使对所有长度为 $k$ 的序列 $\mathbf{j}$ 有 $|T_{\mathbf{j}}|\leq\delta$.由于 $E\subset\bigcup_{\mathbf{j}}T_{\mathbf{j}}$（并对一切这样的 $k$ 元组取并），
\[
  \mathcal{H}^{s}_{\delta}(E)\leq\sum_{j_1\ldots j_k}|T_{j_1\ldots j_k}|^{s}\leq 2^{s}\sum_{j_1\ldots j_k}\rho_{j_1\ldots j_k}^{s},
\]
其中 $\rho_j$ 是 $T_j$ 的外接圆半径.若 $s>\log 3/\log\bigl(1+2\cdot 3^{-\frac{1}{2}}\bigr)$，则由\cref{lem:8.10} 的反复应用，
\[
  \sum_{j_1\ldots j_k}\rho_{j_1\ldots j_k}^{s}
  \leq\sum_{j_1\ldots j_{k-1}}\rho_{j_1\ldots j_{k-1}}^{s}
  \leq\cdots\leq\sum_{j_1}\rho_{j_1}^{s}\leq\rho^{s},
\]
其中 $\rho=3^{-\frac{1}{2}}$ 是 $T$ 的外接圆半径.于是对一切 $\delta>0$ 有 $\mathcal{H}^s_{\delta}(E)\leq(2\cdot 3^{-\frac{1}{2}})^{s}$，从而 $\mathcal{H}^s(E)\leq(2\cdot 3^{-\frac{1}{2}})^{s}$.
\end{proof}

为得到剩余集维数的下界，我们估计曲边三角形 $T_j$ 的内切圆半径 $r_j$.记 $T_j$ 各边（圆）的曲率（即半径的倒数）为 $\alpha_j,\beta_j,\gamma_j$（同样用一致的定向约定），并令
\begin{equation}\label{eq:8.24}
  \sigma_j=\bigl(\beta_j\gamma_j+\alpha_j\gamma_j+\alpha_j\beta_j\bigr)^{\frac{1}{2}}.
\end{equation}
至少在原则上，可以通过反复应用 Soddy 公式算出这些半径与曲率.设 $M_1,M_2,M_3$ 为矩阵
\[
  M_1=\begin{pmatrix}1&0&0&0\\1&1&0&1\\1&0&1&1\\2&0&0&1\end{pmatrix},\qquad
  M_2=\begin{pmatrix}1&1&0&1\\0&1&0&0\\0&1&1&1\\0&2&0&1\end{pmatrix},\qquad
  M_3=\begin{pmatrix}1&0&1&1\\0&1&1&1\\0&0&1&0\\0&0&2&1\end{pmatrix}.
\]
由\cref{lem:8.9}，行向量 $(\alpha_{\mathbf{j}},\beta_{\mathbf{j}},\gamma_{\mathbf{j}},\sigma_{\mathbf{j}})$ 满足关系
\[
  (\alpha_{\mathbf{j}p},\beta_{\mathbf{j}p},\gamma_{\mathbf{j}p},\sigma_{\mathbf{j}p})=(\alpha_{\mathbf{j}},\beta_{\mathbf{j}},\gamma_{\mathbf{j}},\sigma_{\mathbf{j}})M_p\qquad(p=1,2,3),
\]
其中 $\mathbf{j}p$ 表示序列 $\{j_1,\ldots,j_k,p\}$.（容易验证此时 \eqref{eq:8.24} 对把 $\mathbf{j}$ 换成 $\mathbf{j}p$ 也成立.）于是迭代得
\begin{equation}\label{eq:8.25}
  (\alpha_{\mathbf{j}},\beta_{\mathbf{j}},\gamma_{\mathbf{j}},\sigma_{\mathbf{j}})=(1,1,1,\sqrt{3})\,M_{j_1}M_{j_2}\cdots M_{j_k},
\end{equation}
由此可把 $B_{\mathbf{j}}$ 的半径表示为
\begin{equation}\label{eq:8.26}
  r_{\mathbf{j}}^{-1}=\alpha_{\mathbf{j}}+\beta_{\mathbf{j}}+\gamma_{\mathbf{j}}+2\sigma_{\mathbf{j}}
  =(1,1,1,\sqrt{3})\,M_{j_1}M_{j_2}\cdots M_{j_k}\,(1,1,1,2)' ,
\end{equation}
其中撇号表示向量的转置.

于是填装圆半径趋于零的速率直接取决于矩阵乘积 $M_{j_1}M_{j_2}\cdots M_{j_k}$ 当 $k\to\infty$ 时的行为.可以把乘积中的矩阵看作以相等概率从 $M_1,M_2,M_3$ 中随机选取，而所需的信息便包含在 $k$ 充分大时矩阵乘积诸元素的值的概率分布之中.随机矩阵乘积的理论正在发展之中，见 Kingman（\citeyear{kingman1973,kingman1976}）.我们概述一种用这些思想估计 $\dim E$ 下界的方法，而不是 \textcite{hirst1967,boyd1973a} 更复杂但计算上更有效的方法.利用 \textcite{kingman1976} 定理 2.2 证明的一个修改，可以证明：存在数 $\lambda$，使得矩阵 $M_{j_1}M_{j_2}\cdots M_{j_k}$ 的每个元素的 $k$ 次根当 $k\to\infty$ 时对几乎每个序列 $\{j_1,j_2,\ldots\}$（$1\leq j_i\leq 3$）（关于无穷序列集合上的自然乘积测度）都收敛于 $\lambda$.关键是：对所给的矩阵 $M_1,M_2,M_3$，有 $1<\lambda<3$.于是对任意 $\varepsilon>0$，可找到足够大的 $m$，使得对比例为 $1-\varepsilon$ 的所有长度为 $m$ 的序列 $\{j_1,\ldots,j_m\}$，$M_{j_1}\cdots M_{jm}$ 的一切元素都小于 $(\lambda+\varepsilon)^m$.当然这一事实可以直接计算验证.对给定的 $\varepsilon$，令 $\mathcal{U}$ 表示具有该性质的 $m$ 项序列之集.

\begin{lemma}\label{lem:8.12}
设 $t\leq\bigl(\log 3+m^{-1}\log(1-\varepsilon)\bigr)/\log(\lambda+\varepsilon)$.设 $\mathcal{S}$ 是有限序列的有限集，使三角形 $\{T_{\mathbf{j}}:\mathbf{j}\in\mathcal{S}\}$ 覆盖剩余集 $E$.则
\begin{equation}\label{eq:8.27}
  \sum_{\mathbf{j}\in\mathcal{S}}r_{\mathbf{j}}^{t}\geq c,
\end{equation}
其中 $c$ 是正常数.
\end{lemma}

\begin{proof}
只需在 $\mathcal{S}$ 中每个序列的项数都是 $m$ 的倍数的假设下证明引理.因为把每个 $T_{\mathbf{j}}$ 换成 $T_{\mathbf{j}'}$（其中 $\mathbf{j}'$ 由 $\mathbf{j}$ 的前 $mq$ 项组成，$q$ 是 $\mathbf{j}$ 的项数除以 $m$ 的商），由 \eqref{eq:8.26} 使 \eqref{eq:8.27} 中和式的每一项至多增大 $5^{mt}$ 倍.由三角形的网性质，还可设 $\{T_{\mathbf{j}}:\mathbf{j}\in\mathcal{S}\}$ 除公共顶点外是不交族.设 $\mathbf{j}_1,\mathbf{j}_2,\ldots,\mathbf{j}_q$ 是 $\mathcal{S}$ 中最长的序列之一，其中每个 $\mathbf{j}_i$ 都是 $m$ 元组.由于这些三角形覆盖 $E$，序列 $\mathbf{j}_1,\mathbf{j}_2,\ldots,\mathbf{j}_{q-1},\mathbf{j}$ 必属于 $\mathcal{S}$ 对每个 $m$ 元组 $\mathbf{j}$ 成立.但若 $\mathbf{j}\in\mathcal{U}$，则由 \eqref{eq:8.26},
\begin{align*}
  r_{\mathbf{j}_1\ldots\mathbf{j}_{q-1}\mathbf{j}}^{-1}
  &=(1,1,1,\sqrt{3})\,M_{\mathbf{j}_1}\cdots M_{\mathbf{j}_{q-1}}M_{\mathbf{j}}(1,1,1,2)'\\
  &\leq(1,1,1,\sqrt{3})\,M_{\mathbf{j}_1}\cdots M_{\mathbf{j}_{q-1}}(1,1,1,2)'\,(\lambda+\varepsilon)^{m}\\
  &\leq(\lambda+\varepsilon)^{m}r_{\mathbf{j}_1\ldots\mathbf{j}_{q-1}}^{-1}.
\end{align*}
于是对所述的 $t$，
\[
  \sum_{\mathbf{j}}r_{\mathbf{j}_1\ldots\mathbf{j}_{q-1}\mathbf{j}}^{t}
  \geq\sum_{\mathbf{j}\in\mathcal{U}}r_{\mathbf{j}_1\ldots\mathbf{j}_{q-1}\mathbf{j}}^{t}
  \geq(1-\varepsilon)3^{m}(\lambda+\varepsilon)^{-mt}r_{\mathbf{j}_1\ldots\mathbf{j}_{q-1}}^{t}
  \geq r_{\mathbf{j}_1\ldots\mathbf{j}_{q-1}}^{t},
\]
故把三角形 $T_{\mathbf{j}_1\ldots\mathbf{j}_{q-1},\mathbf{j}}$ 换成单个三角形 $T_{\mathbf{j}_1\ldots\mathbf{j}_{q-1}}$ 不会增大 \eqref{eq:8.27} 中的和.可以重复这一过程直至到达单个三角形 $T$，此时 \eqref{eq:8.27} 显然.故 \eqref{eq:8.27} 对曲边三角形子族构成的 $E$ 的任何覆盖成立.
\end{proof}

\begin{lemma}\label{lem:8.13}
设 $\{V_i\}$ 是满足 $E\subset\bigcup_i V_i$ 的开圆盘有限集.则存在有限序列的有限集 $\mathcal{S}$ 使 $E\subset\bigcup_{\mathbf{j}\in\mathcal{S}}T_{\mathbf{j}}$ 且
\begin{equation}\label{eq:8.28}
  \sum_{i}|V_i|^{t}\geq 4^{-t}\sum_{\mathbf{j}\in\mathcal{S}}r_{\mathbf{j}}^{t}.
\end{equation}
\end{lemma}

\begin{proof}
显然可设每个 $V_i$ 都与 $T$ 重叠，且没有圆盘包含任何其他圆盘.依次考察每个圆盘 $V_i$.设 $k$ 是使 $V_i$ 恰与一个项数为 $k$ 的序列 $\mathbf{j}$ 对应的三角形 $T_{\mathbf{j}}$ 相交的最大整数.分两种情形.

(a) 若 $V_i$ 与 $T_{\mathbf{j}1}$、$T_{\mathbf{j}2}$、$T_{\mathbf{j}3}$ 都相交，则 $V_i$ 必与构成这三个三角形边界的四个圆中至少三个相交，故 $V_i$ 的半径至少是这三个圆构成的曲边三角形的内切圆半径.于是注意到 $T_{\mathbf{j}}$ 各边（圆）的半径大于 $r_{\mathbf{j}}$ 并用\cref{lem:8.9}，得
\begin{equation}\label{eq:8.29}
  \frac{1}{2}|V_i|\geq\min\{r_{\mathbf{j}1},r_{\mathbf{j}2},r_{\mathbf{j}3},r_{\mathbf{j}}\}\geq\frac{1}{8}r_{\mathbf{j}},
\end{equation}
并把 $\mathbf{j}$ 放入 $\mathcal{S}$.

(b) 设 $V_i$ 只与 $T_{\mathbf{j}1},T_{\mathbf{j}2},T_{\mathbf{j}3}$ 中的两个相交，不失一般性设为 $T_{\mathbf{j}1}$ 与 $T_{\mathbf{j}2}$.设 $w$ 为 $T_{\mathbf{j}1}$ 与 $T_{\mathbf{j}2}$ 的公共点.若 $w\notin V_i$，则正如 (a) 中那样，
\begin{equation}\label{eq:8.30}
  \frac{1}{2}|V_i|\geq r_{\mathbf{j}3}\geq\frac{1}{8}r_{\mathbf{j}},
\end{equation}
并把 $\mathbf{j}$ 放入 $\mathcal{S}$.否则 $w\in V_i$.记 $B^0=B_{\mathbf{j}1}$，$B^1=B_{\mathbf{j}12}$，$B^2=B_{\mathbf{j}122}$，$B^3=B_{\mathbf{j}1222},\ldots$.设 $m$ 是使 $B^m$ 与 $V_i$ 相交的最小整数.则 $B^{m+1}\subset V_i$，故照前使用\cref{lem:8.9} 得
\begin{equation}\label{eq:8.31}
  |V_i|\geq|B^{m+1}|\geq\frac{1}{8}|B^{m}|.
\end{equation}
设 $\mathbf{j}'\in\mathcal{S}$，其中 $\mathbf{j}'=\mathbf{j}1222\ldots 2$ 是对应于 $B^m$ 的序列.对形如 $B_{\mathbf{j}211\ldots 1}$ 的球作对称的论证，找到对应于含于 $T_{\mathbf{j}2}$ 内的三角形的序列 $\mathbf{j}''$，也把它放入 $\mathcal{S}$.则 $E\cap V_i\subset T_{\mathbf{j}'}\cup T_{\mathbf{j}''}$.

我们注意到，按此构造 $E\subset\bigcup_{\mathbf{j}\in\mathcal{S}}T_{\mathbf{j}}$，且每个 $E\cap V_i$ 被 $\{T_{\mathbf{j}}:\mathbf{j}\in\mathcal{S}\}$ 中一个或两个三角形覆盖.合并 \eqref{eq:8.29}--\eqref{eq:8.31} 即得 \eqref{eq:8.28}.
\end{proof}

\begin{theorem}\label{thm:8.14}
对曲边等边三角形的 Apollonian 填装，$\dim E>1$，其中 $E$ 是填装的剩余集.
\end{theorem}

\begin{proof}
按上文选取 $\varepsilon$ 与 $m$ 使 $t=\bigl(\log 3+m^{-1}\log(1-\varepsilon)\bigr)/\log(\lambda+\varepsilon)>1$.于是合并\cref{lem:8.12,lem:8.13} 可见：对 $E$ 的任何由开圆盘 $\{V_i\}$ 构成的有限覆盖有 $\sum_{i}|V_i|^{t}\geq 4^{-t}c$.若 $\{U_i\}$ 是 $E$ 的由任意集合构成的覆盖，则把每个 $U_i$ 放进一个开圆盘 $V_i$ 并利用 $E$ 的紧性，推得 $\sum_{i}|U_i|^{t}\geq 2^{-t}4^{-t}c$.故 $\mathcal{H}^t(E)>0$，如所需.
\end{proof}

当然，由此得 $\dim E\geq\log 3/\log\lambda$，其中 $\lambda$ 是与矩阵乘积相联系的极限数.

\textcite{hirst1967} 通过比上文更仔细地估计仅四个矩阵的乘积 $M_{j_1}\cdots M_{j_4}$ 的效应，证明了 $\dim E>1.001$.\textcite{boyd1973a} 发展了把 $\dim E$ 求到任意精度的数值方法；现已知道 $1.300<\dim E<1.315$.\textcite{larman1967b} 用一种截然不同的方法证明了任何用圆盘填装一个区域所得剩余集的维数至少为 $1.03$.

注意，任何（凹边的）曲边三角形的 Apollonian 填装在复记法下形如 $z\to(az+b)/(cz+d)$ 的 \newterm{mobius-transformation}{Möbius 变换}{Mobius transformation}下彼此等价.这是因为这类映射（可以选取它们把任何曲边三角形映到任何别的曲边三角形）具有把圆变为圆、例外情形变为直线的性质.特别地，由\cref{lem:1.8} 可得所有曲边三角形的 Apollonian 填装的剩余集具有相同的维数.进一步，用前面的记号，可找到把 $T$ 分别映到 $T_1,T_2,T_3$ 的 Möbius 变换 $\psi_1,\psi_2,\psi_3$.于是剩余集 $E$ 是满足 $E=\bigcup_{j=1}^{3}\psi_j(E)$ 的紧集.因此尽管 $\psi_j$ 只是非扩张而非压缩，8.3~节的思想仍可经调整用于研究 Apollonian 填装.这一问题的一个推广是求形如 $\psi_{j_1}\circ\cdots\circ\psi_{j_k}$ 的所有映射的极限点之集的 Hausdorff 维数，其中 $\psi_{j_i}$ 取自给定的一组 Möbius 变换；见 \textcites{patterson1976,sullivan1979,mandelbrot1983}.

与任何密切填装 $\{B_i\}$ 相联系的一个重要数字是填装的\newterm{exponent}{指数}{exponent of packing} $e$，定义为使 $\sum_{i=1}^{\infty}|B_i|^t<\infty$ 的 $t$ 的下确界.事实表明，对 Apollonian 填装，指数 $e$ 等于剩余集 $E$ 的维数.证明 $\dim E\leq e$ 相对容易（见\cref{exe:8.6}；\textcite{larman1966} 给出了更一般的结果.）\textcite{boyd1973b} 证明的相反不等式要难得多.需要仔细考察矩阵乘积以证明：若 $t<e$，则存在常数 $c$，使得若曲边三角形 $\{T_{\mathbf{j}}:\mathbf{j}\in\mathcal{S}\}$ 覆盖 $E$，则 $\sum_{\mathbf{j}\in\mathcal{S}}r_{\mathbf{j}}^t\geq c$.结果随后由\cref{lem:8.13} 得出.

\textcite{eggleston1953a} 研究了用取向相反的较小等边三角形填装一个等边三角形所得的剩余集.图~\ref{fig:8.4} 所示的构形留下维数为 $\log 3/\log 2$ 的剩余集，这是任何此类填装所能得到的最小值（见\cref{exe:8.7}）.

\begin{figure}[htbp]
  \centering
  \includegraphics[width=0.42\textwidth]{fig8-4.png}
  \caption{用取向相反的等边三角形作填装}
  \label{fig:8.4}
\end{figure}

\section{一个来自数论的例子}

Hausdorff 测度在\newterm{number-theory}{数论}{number theory}中——特别是在\newterm{diophantine-approximation}{丢番图逼近}{Diophantine approximation}中——找到了大量应用.典型地说：``能被有理数逼近到指定精度的数的集合，其 Hausdorff 维数不能太大''.这里我们考虑这类结果中的具体结果，它们就\cref{chap:6} 的投影定理提供了一个有启发性的例子.

按数论的惯例，记 $\|x\|$ 为实数 $x$ 到最近整数的距离.我们需要下述定理的一维与二维两个版本.

\begin{theorem}\label{thm:8.15}
取 $\alpha>0$ 及正整数序列 $n_1,n_2,\ldots$ 使
\begin{equation}\label{eq:8.32}
  n_{j+1}\geq n_j^{j}.
\end{equation}

(a) 设 $F$ 是满足
\begin{equation}\label{eq:8.33}
  \|n_j x\|\leq n_j^{-\alpha}\qquad(j=1,2,\ldots)
\end{equation}
的实数 $x$ 之集.则 $\dim F=1/(1+\alpha)$.

(b) 设 $E$ 是满足
\[
  \|n_j x\|,\ \|n_j y\|\leq n_j^{-\alpha}\qquad(j=1,2,\ldots)
\]
的 $(x,y)\in\mathbb{R}^2$ 之集.则 $\dim E=2/(1+\alpha)$.
\end{theorem}

\begin{proof}
(a) 记
\begin{equation}\label{eq:8.34}
  F_j=\bigcup_{i=-\infty}^{\infty}\bigl[i n_j^{-1}-n_j^{-1-\alpha},\ i n_j^{-1}+n_j^{-1-\alpha}\bigr],
\end{equation}
则
\begin{equation}\label{eq:8.35}
  F=\bigcap_{j=1}^{\infty}F_j.
\end{equation}
首先注意，$F$ 与任何区间 $I$ 的交对每个 $j$ 都可以用至多 $|I|n_j+2$ 个长度为 $2n_j^{-1-\alpha}$ 的区间覆盖.由于当 $j\to\infty$ 时这些长度趋于零，$F\cap I$ 具有有限的 $\mathcal{H}^{1/(1+\alpha)}$-测度，故 $\dim F\leq 1/(1+\alpha)$.

为求维数的下界，取 $s<1/(1+\alpha)$，并把 $F$ 稍加修改以配合\cref{thm:1.15}，证明 $\mathcal{H}^s(F)>0$.不妨设对一切 $j$ 有 $n_{j+1}n_j^{-1-\alpha}\geq 2$、$n_j^{\alpha}\geq 4$ 且 $n_1\geq 4$；若非如此，利用 \eqref{eq:8.32} 忽略小的 $j$ 值并重新编号即可.令
\begin{equation}\label{eq:8.36}
  m_j=\bigl[2n_{j+1}n_j^{-1-\alpha}-2\bigr]\geq n_{j+1}n_j^{-1-\alpha}\geq 2,
\end{equation}
其中 $[\ ]$ 表示``不超过的最大整数''.注意可找到 $c>0$ 使
\begin{equation}\label{eq:8.37}
  m_1\cdots m_{j-1}n_j^{-s(1+\alpha)}\geq c\qquad(j\geq 2).
\end{equation}
事实上，把 $j$ 换成 $j+1$ 时该不等式左端增大了一个因子
\[
  m_j n_{j+1}^{-s(1+\alpha)}n_j^{s(1+\alpha)}
  \geq n_{j+1}^{1-s(1+\alpha)}n_j^{(s-1)(1+\alpha)}
  \geq n_j^{j[1-s(1+\alpha)]+(s-1)(1+\alpha)},
\]
这里用了 \eqref{eq:8.36} 与 \eqref{eq:8.32}；而当 $j$ 充分大时它大于 $1$，因为 $1-s(1+\alpha)>0$.

$F_j$ 的每个区间依两端重叠的情况含有 $m_j$、$m_j+1$ 或 $m_j+2$ 个 $F_{j+1}$ 的完整区间.为顾及这些端点效应，我们如下把 $F$ 缩减为 $F'$.令 $F_0'$ 为单位区间.若 $j\geq 0$ 且 $F_j'$ 已经指定，则从 $F_j'$ 的每个区间中恰好选取 $m_j$ 个 $F_{j+1}$ 的相邻完整区间来定义 $F_{j+1}'$.令 $F'=\bigcap_{j=0}^{\infty}F_j'$，则 $F'\subset F\cap[0,1]$.

现在按\cref{thm:1.15} 之前描述的 Cantor 型过程构造 $[0,1]$ 的子集 $E$.令 $E_0=[0,1]$，并按那里描述的方式构造 $E_1,E_2,\ldots$，使 $E_j$ 的每个区间含有 $m_j$ 个 $E_{j+1}$ 的区间，于是 \eqref{eq:1.23} 成为
\begin{equation}\label{eq:8.38}
  m_j|J|^{s}=|I|^{s},
\end{equation}
其中 $J$ 是 $E_{j+1}$ 的典型区间，$I$ 是 $E_j$ 的区间.由\cref{thm:1.15}，集合 $E=\bigcap_{j=0}^{\infty}E_j$ 的维数为 $s$.

存在从 $F'$ 到 $E$ 的自然双射 $\psi$：若 $x\in F'$ 位于 $F_j'$ 的第 $k_j$ 个区间（从左数起），则 $\psi(x)$ 是对所有 $j$ 都位于 $E_j$ 的第 $k_j$ 个区间中的唯一点.设 $x$ 与 $y$ 是 $F'$ 的不同点.设 $j$ 是使 $x$ 与 $y$ 位于 $F_{j+1}'$ 的不同区间中的最小整数；设 $y$ 位于含 $x$ 的区间右侧第 $k$ 个区间.则
\begin{equation}\label{eq:8.39}
  |y-x|\geq kn_{j+1}^{-1}-2n_{j+1}^{-1-\alpha}\geq\frac{1}{2}kn_{j+1}^{-1}.
\end{equation}
另一方面，$E$ 中的点 $\psi(x)$ 与 $\psi(y)$ 位于 $E_j$ 的同一区间 $I$ 内，但 $\psi(y)$ 位于 $E_{j+1}$ 的在 $\psi(x)$ 所在区间右侧第 $k$ 个的区间中.于是
\[
  |\psi(y)-\psi(x)|
  \leq 2k|I|m_j^{-1}
  \leq 2k(m_1\cdots m_{j-1})^{-1/s}m_j^{-1}
  \leq 2kc_1n_j^{-1-\alpha}m_j^{-1}
  \leq 2kc_1n_{j+1}^{-1},
\]
其中 $c_1=c^{-1/s}$，这里迭代了 \eqref{eq:8.38} 并用了 \eqref{eq:8.37} 与 \eqref{eq:8.36}.与 \eqref{eq:8.39} 结合得
\[
  |\psi(y)-\psi(x)|\leq 4c_1|y-x|.
\]
由\cref{lem:1.8} 得 $s=\dim E\leq\dim F'\leq\dim F$.这对一切 $s<1/(1+\alpha)$ 成立.

(b) 证明 (b) 的一种途径是把\cref{thm:1.15} 与 (a) 的证明修改到二维.另一种途径是：利用 Descartes 积集的\cref{cor:5.10} 连同 (a)，集合 $E=F\times F$ 的维数至少为 $2/(1+\alpha)$.另一方面，若 $S$ 是任何单位边长正方形，则对每个 $j$ 可用 $(n_j+2)^2$ 个直径为 $2^{\frac{1}{2}}n_j^{-1-\alpha}$ 的正方形覆盖 $E\cap S$（从 $F_j\times F_j$ 中取正方形）.于是 $\dim(E\cap S)$，从而 $\dim E$，至多为 $2/(1+\alpha)$.
\end{proof}

注意，满足 \eqref{eq:8.33} 的 $x$ 正是``密切可逼近''的数：即对一切 $j$，存在有理数 $p_j/n_j$ 使 $|x-p_j/n_j|\leq n_j^{-1-\alpha}$.

若把不等式 \eqref{eq:8.33} 要求``对无穷多个 $j$''成立而非``对一切 $j$''成立，相应的结果就更难证明.此时交 \eqref{eq:8.35} 换成一个上极限.我们不加证明地陈述 \newterm{jarniks-theorem}{Jarník 定理}{Jarnik's theorem}的一维版本.其证明都相当技术化，见 \textcites{jarnik1931,besicovitch1934b,eggleston1952} 与 Kaufman（\citeyear{kaufman1970,kaufman1981}）；与连分数的一个有趣联系见 \textcite{good1941}.(b) 部分给出了解增大速度的某种指示.

\begin{theorem}\label{thm:8.16}
取 $\beta>1$.

(a) 使不等式
\[
  \|nx\|\leq n^{-\beta}
\]
对无穷多个 $n$ 成立的实数 $x$ 之集的 Hausdorff 维数为 $2/(1+\beta)$.

(b) 设 $m_1,m_2,\ldots$ 是快速增长（比如说 $m_{j+1}\geq m_j^j$）的整数列.则使
\[
  \|nx\|\leq nm_j^{-1-\beta}\quad\text{且}\quad n\leq m_j
\]
对某个 $n$ 对无穷多个 $j$ 成立的实数 $x$ 之集的 Hausdorff 维数为 $2/(1+\beta)$.
\end{theorem}

现在考虑\cref{thm:8.15}(b) 的平面集 $E$ 到直线上的投影.结果表明，投影的例外方向的集合大到可能出现的最大程度.（比较\cref{thm:6.8}.）

\begin{theorem}\label{thm:8.17}
设 $n_1,n_2,\ldots$ 是快速增长（比如说 $n_{j+1}\geq n_j^j$）的正整数列.对 $s\leq 1$，设 $E$ 是满足 $-1\leq x,y\leq 1$ 且
\[
  \|n_j x\|,\ \|n_j y\|\leq n_j^{1-2/s}\qquad(j=1,2,\ldots)
\]
的 $(x,y)\in\mathbb{R}^2$ 之集.则 $\dim E=s$，且方向集
\begin{equation}\label{eq:8.40}
  \{\theta:\dim(\operatorname{proj}_{\theta}E)<s\}
\end{equation}
的 Hausdorff 维数至少为 $s$.
\end{theorem}

\begin{proof}
由\cref{thm:8.15}(b) 立即得 $\dim E=s$.

对 $\delta>0$ 令
\[
  A_{\delta}=\bigl\{a:\|na\|\leq nn_j^{-2/s}\ \text{对某个满足}\ 1\leq n\leq n_j^{1-\delta}\ \text{的}\ n\ \text{对无穷多个}\ j\bigr\}.
\]
在\cref{thm:8.16}(b) 中取 $m_j=n_j^{1-\delta}$ 与 $\beta=2/s(1-\delta)-1$ 给出
\begin{equation}\label{eq:8.41}
  \dim A_{\delta}=s(1-\delta).
\end{equation}
对 $a\in\mathbb{R}$，设 $\psi_a:\mathbb{R}^2\to\mathbb{R}$ 为线性映射
\begin{equation}\label{eq:8.42}
  \psi_a(x,y)=ax+y.
\end{equation}
若 $a\in A_{\delta}$，则对任意大的 $j$，存在整数 $n$ 与 $b$ 使
\begin{equation}\label{eq:8.43}
  1\leq n\leq n_j^{1-\delta}
\end{equation}
且
\[
  |na-b|\leq nn_j^{-2/s},
\]
故
\begin{equation}\label{eq:8.44}
  |a-b/n|\leq n_j^{-2/s}.
\end{equation}
若 $u,v$ 是满足 $|u|,|v|\leq n_j$ 的整数，则在坐标记号下写
\[
  (u,v)=(qn,-qb)+(r,\,v+qb),
\]
其中 $u=qn+r$，$q,r$ 的选取使 $0\leq r<n$.若 $a'=b/n$，则
\[
  \psi_{a'}(u,v)=\psi_{a'}(qn,-qb)+\psi_{a'}(r,v+qb)=0+\psi_{a'}(r,v+qb).
\]
由于当 $(u,v)$ 变化时 $(r,v+qb)$ 至多取 $cnn_j$ 个不同的值（$c=3(1+|b|)$），$\psi_{a'}(u,v)$ 从而 $\psi_{a'}(un_j^{-1},vn_j^{-1})$ 亦然.但 $E$ 的每一点的这两个坐标都在某对整数 $u,v$ 下距 $\bigl(un_j^{-1},vn_j^{-1}\bigr)$ 不超过 $n_j^{-2/s}$，故 $\psi_{a'}(E)$ 可以用 $cnn_j$ 个长度为 $2(|a'|+1)n_j^{-2/s}$ 的区间覆盖.而由 \eqref{eq:8.44}，当 $|x|,|y|\leq 1$ 时 $|\psi_a(x,y)-\psi_{a'}(x,y)|\leq n_j^{-2/s}$，故 $\psi_a(E)$ 可以用 $cnn_j$ 个长度为 $2(|a'|+2)n_j^{-2/s}$ 的区间覆盖.这可以对任意大的 $j$ 做到，故只要当 $j\to\infty$ 时 $nn_jn_j^{-2t/s}\to 0$——由 \eqref{eq:8.43}，当 $t>(1-\frac{1}{2}\delta)s$ 时正是如此——就有 $\mathcal{H}^t\bigl(\psi_a(E)\bigr)=0$.

于是我们证明了：若 $a\in A_{\delta}$（对任意 $\delta>0$），即 $a\in A=\bigcup_{\delta>0}A_{\delta}$，则 $\dim\bigl(\psi_a(E)\bigr)<s$.由 \eqref{eq:8.41}，$\dim A=s$.由 $\psi_a$ 的定义，当 $a=\cot\theta$ 时集合 $\psi_a(E)$ 与 $\operatorname{proj}_{\theta}(E)$ 几何相似，从而有相同的 Hausdorff 维数.我们得出：使 $\dim\bigl(\operatorname{proj}_{\theta}(E)\bigr)<s$ 的 $\theta$ 之集的维数至少为 $s$.
\end{proof}

上述例子及其高维推广归功于 \textcite{kaufmanmattila1975}，取代了 \textcite{marstrand1954a} 与 \textcite{kaufman1969} 的例子.由 \textcite{mattila1975b} 关于方向例外集的最大维数的结果，集合 \eqref{eq:8.40} 的维数恰为 $s$.从投影的观点看，这类数论例子一般表现出最坏的可能行为.

Hausdorff 测度在数论中的进一步应用见 \textcites{eggleston1952,bakerschmidt1970,rogers1970,baker1978}.丢番图逼近的更广泛的论述见 \textcite{baker1975} 或 \textcite{schmidt1980}.

\section{在凸性中的一些应用}

本节描述 Hausdorff 测度在\newterm{convexity}{凸性}{convexity}（凸集几何）中出现的若干方式.前两个结果表明凸曲面的不规则性不能太大.第三个结果是投影定理的一个应用.

设 $K$ 是 $\mathbb{R}^3$ 中的（紧）凸体.若 $K$ 在 $x$ 处支撑多于一个切平面，则称 $x$ 是 $K$ 的曲面 $\partial K$ 的\newterm{singular-point}{奇点}{singular point}.下面的定理可推广到高维，归功于 \textcite{andersonklee1952}；\textcite{besicovitch1963a} 给出了另一个证明.

\begin{theorem}\label{thm:8.18}
设 $K$ 是 $\mathbb{R}^3$ 中的凸体.则 $\partial K$ 的奇点之集含于可求长曲线的可数并之中，从而具有 $\sigma$-有限的 $\mathcal{H}^1$-测度.
\end{theorem}

\begin{proof}
设 $\Pi$ 是与 $K$ 不相交的平面，$f:\Pi\to\partial K$ 是``最近点''映射.即若 $x\in\Pi$，则 $f(x)$ 是 $K$ 中使 $|f(x)-x|$ 最小的点；由于 $K$ 凸，该点唯一.若 $x,y\in\Pi$，则（可能空间扭歪的）四边形 $x,f(x),f(y),y$ 的角 $\angle xf(x)f(y)$ 与 $\angle yf(y)f(x)$ 必至少为 $\frac{1}{2}\pi$，因为线段 $[f(x),f(y)]$ 含于 $K$ 中且含有作为 $x$ 与 $y$ 的``最近点''候选的点.于是 $f$ 不增大距离，即
\begin{equation}\label{eq:8.45}
  |f(x)-f(y)|\leq|x-y|\qquad(x,y\in\Pi).
\end{equation}
记 $\partial K_{\Pi}$ 为 $\Pi$ 在 $f$ 下的像，粗略地说即 $\partial K$ 上与 $\Pi$ 同侧的那些点.若 $z$ 是 $\partial K_{\Pi}$ 的奇点，则容易看出 $f^{-1}(z)$ 含有一条直线段.事实上，若 $z$ 是锥型奇点（$z$ 是包含 $K$ 的圆锥的顶点），则 $f^{-1}(z)$ 含有一个圆盘.设 $\{L_i\}_{1}^{\infty}$ 为 $\Pi$ 中相对于 $\Pi$ 内某对固定坐标轴具有有理坐标端点的线段构成的可数族.若 $z$ 是 $\partial K_{\Pi}$ 的奇点，则 $f^{-1}(z)$ 含有一条必与某个 $L_i$ 相交的线段.于是 $\partial K_{\Pi}$ 的奇点含于 $\bigcup_i f(L_i)$ 之中.而利用 \eqref{eq:8.45} 可见，$\partial K$ 的子集 $f(L_i)$ 要么是一条可求长曲线，要么是一个单点（见 3.2~节）.于是 $\partial K_{\Pi}$ 的奇点含于可求长曲线的可数族之中，从而具有 $\sigma$-有限的 $\mathcal{H}^1$-测度.

最后，若 $\Pi_1,\ldots,\Pi_4$ 是包住 $K$ 的四面体各面所在的平面，则 $\partial K$ 的每个奇点都是某个 $\partial K_{\Pi_j}$（$1\leq j\leq 4$）的奇点.于是 $\partial K$ 的奇点构成一个 $\sigma$-有限的正则 $1$-集.
\end{proof}

顺便指出，由证明容易看出凸体的锥型奇点之集至多可数.（若 $z$ 是 $\partial K$ 的锥型奇点，则 $f^{-1}(z)$ 含有一个圆盘，而平面 $\Pi$ 只能含可数多个不交圆盘.）

\textcite{klee1957} 问：位于凸曲面内的直线段的方向之集可以有多大.\textcite{mcminn1960} 在三维情形解决了这一问题，这里我们给出 \textcite{besicovitch1963b} 证明的简化版本.

\begin{theorem}\label{thm:8.19}
设 $K$ 是 $\mathbb{R}^3$ 中的凸体.则含于 $\partial K$ 内的线段的方向之集（把方向视为单位球面的子集）具有 $\sigma$-有限的 $\mathcal{H}^1$-测度.
\end{theorem}

\begin{proof}
设 $\Pi_1$ 与 $\Pi_2$ 是与 $\partial K$ 相交成凸曲线 $\Gamma_1$ 与 $\Gamma_2$ 的平行平面.我们先证明：$\partial K$ 内同时截过 $\Gamma_1$ 与 $\Gamma_2$ 的线段的方向之（紧）集的 $\mathcal{H}^1$-测度是 $\sigma$-有限的.如下定义函数 $f_1:\Gamma_1\to\Gamma_2$：设 $x\in\Gamma_1$.对每个使 $\Gamma_1$ 与 $\Gamma_2$ 位于同侧、在 $x$ 处切于 $\Gamma_1$ 且也切于 $\Gamma_2$ 的平面 $\Pi$，$\Pi\cap\Gamma_2$ 要么是一个单点，要么是一条直线段.这些点或线段在所有这样的 $\Pi$ 上的总体是 $\Gamma_2$ 的一个单点或一条闭子弧.令 $f_1(x)$ 为这个点或弧的中点.显然 $f_1$ 是单调函数，从而是有界变差的.于是由向量差 $g_1(x)=f_1(x)-x$ 给出的函数 $g_1:\Gamma_1\to\mathbb{R}^3$ 也是有界变差的.正交投影到 $\Pi_1$ 上，$\operatorname{proj}_{\Pi_1}\circ g_1:\Gamma_1\to\Pi_1$ 有界变差，故由\cref{exe:3.1}，$\operatorname{proj}_{\Pi_1}\bigl(g_1(\Gamma_1)\bigr)$ 是 $\mathcal{H}^1$-测度有限的集合.但 $\operatorname{proj}_{\Pi_1}\bigl(g(\Gamma_1)\bigr)$ 是线段 $\{[x,f_1(x)]:x\in\Gamma_1\}$ 的方向余弦之集，故这类线段的方向之集具有有限的 $\mathcal{H}^1$-测度.

我们以完全相同的方式定义 $f_2:\Gamma_2\to\Gamma_1$，并推得线段 $\{[f_2(y),y]:y\in\Gamma_2\}$ 的方向之集也具有有限的 $\mathcal{H}^1$-测度.

若 $[x,y]$ 是 $\partial K$ 中满足 $x\in\Gamma_1$、$y\in\Gamma_2$ 的线段，且其形式既非 $[x,f_1(x)]$ 又非 $[f_2(y),y]$，则两种可能必居其一，我们依次考虑.首先 $x$ 可以位于 $\Gamma_1$ 的直线段 $L_1$ 上，$y$ 位于 $\Gamma_2$ 中平行的线段 $L_2$ 上.线段 $\{[x,y]:x\in L_1,\ y\in L_2\}$ 的方向之集显然有有限的 $\mathcal{H}^1$-测度.由于凸曲线只能含可数多条直线段，这样的线段对至多可数多对.第二种可能是 $x$ 与 $y$ 分别为 $\Gamma_1$ 与 $\Gamma_2$ 的奇异点（即支撑多于一条切线的点）.这对至多可数多条 $[x,y]$ 可能发生.于是两类``例外''线段 $[x,y]$ 的方向之集都具有 $\sigma$-有限的 $\mathcal{H}^1$-测度.把所有情形合在一起即得：截过 $\Gamma_1$ 与 $\Gamma_2$ 的边界线段的方向之集具有 $\sigma$-有限的 $\mathcal{H}^1$-测度.

最后，取平行于三个坐标平面且与之距离为有理数的全部平面，我们得到平行平面对的可数族，使 $\partial K$ 中的每条线段都与该族中某对平面相交.于是边界线段的方向之集含于 $\sigma$-有限集的可数族之中，从而是 $\sigma$-有限的.
\end{proof}

\textcite{ewaldlarmanrogers1970} 完成了该定理的高维类似物的研究：他们证明，若 $K$ 是 $\mathbb{R}^n$ 中的凸体，则位于 $K$ 边界内的 $k$ 维球体的取向（$G_{n,k}$ 意义下）之集具有 $\sigma$-有限的 $\mathcal{H}^{k(n-k-1)}$-测度.

一个相关问题——凸体表面上位于某条边界线段（或更一般地，某 $k$ 维球体）内的点的测度——由 \textcite{burton1979} 讨论.\textcite{larman1971b} 考察了凸体的面的相对边界之并的测度.

\textcite{klee1959} 的一个著名结果断言：若 $K$ 是 $\mathbb{R}^3$ 中每个平面截面都是多边形的凸体，则 $K$ 必是多胞体，即有限点集的凸包.曾有猜测：若 $K$ 的截面在几乎所有方向上几乎处处是多边形，则 $K$ 至多有可数多个\newterm{extreme-point}{极点}{extreme point}（不在 $K$ 内任何线段的内部的边界点）.这一猜测被 \textcite{dallalarman1980} 推翻，他们证明极点之集的 Hausdorff 维数可以大到 $1$.我们不在此重复他们精巧的构造，而只是从投影定理推出这已是可能出现的最坏情形.

\begin{theorem}\label{thm:8.20}
设 $K$ 是 $\mathbb{R}^3$ 中截面几乎处处是多边形的凸体.则 $K$ 的极点之集 $E$ 的 Hausdorff 维数至多为 $1$.
\end{theorem}

\begin{proof}
这是把 Mattila 对投影定理的改进应用于 Borel 集 $E$ 的简单推论（见 6.3~节的倒数第二段）.若 $\dim E>1$ 且 $0<t<\dim E-1$，则对几乎所有的单位向量 $\boldsymbol{\theta}$，
\[
  \mathcal{H}^{t}\bigl(E\cap\operatorname{proj}_{\theta}^{-1}u\bigr)=\infty
\]
对 $L_\theta$ 上一个正 $\mathscr{L}^1$-测度的 $u$ 集合成立，其中 $\operatorname{proj}_{\theta}$ 表示到直线 $L_\theta$（过原点沿方向 $\boldsymbol{\theta}$）上的投影.对这样的 $\boldsymbol{\theta}$ 与 $u$，平面 $\operatorname{proj}_{\theta}^{-1}u$ 与 $E$ 交成无穷集.$K$ 的任何极点当然是平面凸集 $K\cap\operatorname{proj}_{\theta}^{-1}u$ 的极点，故这些平面截面不可能是多边形.
\end{proof}

\section{动力系统中的吸引子}

撰写本书之时，数学与数学物理的一个增长领域是对\newterm{dynamical-system}{动力系统}{dynamical system}中``混沌''与``奇怪吸引子''的研究.我们尝试用几页篇幅概述这类吸引子如何产生的一些基本思想.

设 $f$ 是度量空间（通常是 $\mathbb{R}^n$，$n\geq 1$ 的子集）到自身的映射.我们关心序列——\newterm{orbit}{轨道}{orbit} $\{f^m(x)\}_{m=1}^{\infty}$——对不同初始点 $x$ 的行为，特别是 $m$ 很大时的行为.（照例 $f^m$ 表示 $f$ 的 $m$ 次迭代.）例如，这些序列可能是周期的，或可能收敛于一条周期轨道.或者 $\{f^m(x)\}$ 可能表现得近乎随机地在度量空间中游荡.人们很快被引导去研究不变集，即使 $f(E)\subset E$ 的集合 $E$.于是若 $x\in E$，迭代 $f^m(x)$ 对一切 $m$ 都被困在 $E$ 内.若 $f$ 连续，不变集的闭包也是不变的，而一个主要问题就是求给定函数 $f$ 的闭不变集.在特别有意义的情形，不变集可能呈现精细结构.也许最简单的例子是 $f:\mathbb{R}\to\mathbb{R}$，$f(x)=\frac{3}{2}(1-|2x-1|)$，Cantor 集对它是不变的，见\cref{exe:8.8}.一些被称为 Julia 集的非凡集合由要求 $f(E)=E$ 定义，其中 $f$ 是复平面的变换 $z\to z^2-\mu$；Mandelbrot（\citeyear{mandelbrot1982}，第 19 节）对参数 $\mu$ 的各种值绘出了这些集合.

然而，实践中最重要的不变集都具有某种稳定性，否则它们不太可能在物理情境或计算机实现中被观察到.这些集合称为\newterm{attractor}{吸引子}{attractor}.吸引子的定义多种多样；本质上，映射 $f$ 的吸引子是满足 $f(E)\subset E$ 的闭集 $E$，而且若 $x\in V$，则 $f^m(x)$ 到 $E$ 的距离当 $m\to\infty$ 时趋于 $0$，其中 $V$ 是包含 $E$ 的``大''集合，例如 $E$ 的开邻域.通常还要求 $E$ 在某种意义上是极小的，比如要求轨道 $\{f^m(x)\}$ 对某个 $x$ 在 $E$ 中稠密.若 $E$ 具有精细结构，或对初始条件有敏感依赖（使两个相近的点在 $f$ 的迭代下未必保持接近），则 $E$ 可称为奇怪吸引子.判定哪些函数产生奇怪吸引子、为何出现这类吸引子以及它们有何结构，仍有许多工作要做.这一学科已变得极其复杂，\newterm{bifurcation-theory}{分歧理论}{bifurcation theory}、\newterm{ergodic-theory}{遍历理论}{ergodic theory}、微分拓扑与泛函分析都在其中起着重要作用.关于一维与二维映射的出色综述见 \textcite{whitley1983}.

对实区间上映射的研究近来进展可观，主要归功于 Feigenbaum（\citeyear{feigenbaum1978,feigenbaum1979}） 与 Jonker 与 Rand（\citeyear{jonkendrand1981a,jonkendrand1981b}）.这里我们详细考察 \textcite{grassberger1981} 研究过的一个函数，它虽然貌似矫揉造作，却易于分析，并且在更一般的背景中具有重要性.

考虑泛函方程，即著名的重正化方程，
\begin{equation}\label{eq:8.46}
  g\bigl(g(x/\alpha)\bigr)=-g(x)/\alpha,
\end{equation}
连同规范化条件 $g(0)=1$.可以证明（见 \textcite{campaninoepstein1981} 与 \textcite{lanford1982}）该方程被某个在 $[-1,1]$ 上的偶实解析函数 $g$ 对 $\alpha=2.50290\ldots$ 所满足.（虽然这个解局部唯一，但尚不知道 \eqref{eq:8.46} 是否对其他 $\alpha$ 值有解.）函数 $g$ 在 $[-1,1]$ 上精确到四位小数为
\begin{equation}\label{eq:8.47}
  \begin{aligned}
    g(x)&=1-1.5276x^{2}+0.1048x^{4}+0.0267x^{6}\\
    &\quad -0.0035x^{8}+0.0008x^{10}+0.0003x^{12},
  \end{aligned}
\end{equation}
其图像示于图~\ref{fig:8.5}.我们用 8.3~节的思想估计 $g$ 的吸引子的维数.

\begin{figure}[htbp]
  \centering
  \includegraphics[width=0.5\textwidth]{fig8-5.png}
  \caption{函数 $y=g(x)$}
  \label{fig:8.5}
\end{figure}

\begin{theorem}\label{thm:8.21}
函数 $g$ 对每个 $k\geq 0$ 有唯一的周期为 $2^k$ 的（不稳定）周期轨道，以及满足 $0.5345<\dim E<0.5544$ 的闭不变吸引子 $E$.若 $x\in[-1,1]$，则要么 $g^m(x)$ 对充分大的 $m$ 是周期点（这只对可数多个 $x$ 发生），要么 $g^m(x)$ 到 $E$ 的距离当 $m\to\infty$ 时趋于 $0$.若 $x\in E$，则轨道 $\{g^m(x)\}_{1}^{\infty}$ 在 $E$ 中稠密.
\end{theorem}

\begin{proof}
迭代 \eqref{eq:8.46} 并利用 $g$ 的偶性，
\begin{equation}\label{eq:8.48}
  g^{2^{k}}\Bigl(\Bigl(-\frac{1}{\alpha}\Bigr)^{k}x\Bigr)=\Bigl(-\frac{1}{\alpha}\Bigr)^{k}g(x)\qquad(k\geq 0).
\end{equation}
记 $x_m=g^m(0)$，可见 $x_{2^k}=(-1/\alpha)^k$（$k\geq 0$），并且经计算得
\[
  x_1=1,\qquad x_2=-0.3995,\qquad x_3=0.7589,\qquad x_4=0.1596.
\]
设 $V$ 为开区间 $(x_2,x_1)$，并在 $\overline{V}$ 上定义映射
\begin{equation}\label{eq:8.49}
  \psi_1(x)=-\frac{1}{\alpha}x,\qquad \psi_2(x)=g^{-1}\Bigl(-\frac{1}{\alpha}x\Bigr),
\end{equation}
其中 $g^{-1}$ 总取正值.注意
\[
  \psi_1(V)=(x_2,x_4),\qquad \psi_2(V)=(x_3,x_1),
\]
故开集条件对压缩 $\psi_1,\psi_2$ 成立（见 8.3~节）.由\cref{thm:8.3,thm:8.8}，存在唯一紧集 $E$ 使 $E=\psi_1(E)\cup\psi_2(E)$，且 $s\leq\dim E\leq t$，其中
\[
  \bigl(\inf|\psi_1'(x)|\bigr)^{s}+\bigl(\inf|\psi_2'(x)|\bigr)^{s}
  =1=\bigl(\sup|\psi_1'(x)|\bigr)^{t}+\bigl(\sup|\psi_2'(x)|\bigr)^{t},
\]
下确界与上确界取遍 $x\in\overline{V}$.由于 $g'$（从而 $(g^{-1})'$）的极端值在区间 $\overline{V}$ 的端点取得，得
\[
  \alpha^{-s}+\bigl(\alpha|g'(x_1)|\bigr)^{-s}=1=\alpha^{-t}+\bigl(\alpha|g'(x_3)|\bigr)^{-t},
\]
计算导数得 $0.5345<\dim E<0.5544$.

剩下考察点的迭代的确界行为.由于 $g[-1,1]\subset\overline{V}$，可以完全在 $\overline{V}$ 内工作.改写 \eqref{eq:8.46} 可见 $g(-(1/\alpha)x)=g\bigl(-(1/\alpha)g(x)\bigr)$（注意 $g(-(1/\alpha)x)>0$），得到在 $[-1,1]$ 上成立的函数恒等式
\begin{equation}\label{eq:8.50}
  g\circ\psi_1=\psi_2\circ g,\qquad g\circ\psi_2=\psi_1.
\end{equation}
于是记 $\psi_{j_1j_2\ldots j_k}=\psi_{j_1}\circ\psi_{j_2}\circ\cdots\circ\psi_{j_k}$，我们有
\begin{equation}\label{eq:8.51}
  \begin{aligned}
    g\circ\psi_{11\ldots1}&=\psi_{22\ldots2}\circ g,\\
    g\circ\psi_{11\ldots12j_1\ldots j_k}&=\psi_{22\ldots21j_i\ldots j_k}\qquad(2\leq i\leq k+1).
  \end{aligned}
\end{equation}
由于 $g(\overline{V})\subset\overline{V}$，得
\begin{equation}\label{eq:8.52}
  g\bigl(\psi^{k}(\overline{V})\bigr)\subset\psi^{k}(\overline{V}),
\end{equation}
其中照 8.3~节，$\psi^k(\overline{V})=\bigcup_{j_1\ldots j_k}\psi_{j_1\ldots j_k}(\overline{V})$.于是 $\psi^k(\overline{V})$ 中的任何点在 $g$ 的迭代下被困于 $\psi^k(\overline{V})$ 之中.进一步，由 \eqref{eq:8.51} 可见，若 $x\in\psi^k(\overline{V})$，则当 $m$ 增大时 $g^m(x)$ 依次访问 $2^k$ 个（不交的）集合 $\psi_{j_1\ldots j_k}(\overline{V})$ 中的每一个.

设 $x_f=0.5493\ldots$ 为 $g$ 在 $[-1,1]$ 中的唯一不动点.由 \eqref{eq:8.48}，点 $(-1/\alpha)^k x_f$ 从而 $g^m\bigl((-1/\alpha)^k x_f\bigr)$ 是周期至多为 $2^k$ 的周期点.然而由上面的注记，$g^m\bigl((-1/\alpha)^kx_f\bigr)$ 在每 $2^k$ 个相继的 $m$ 值中位于不同的 $\psi_{j_1\ldots j_k}(\overline{V})$ 内（注意 $(-1/\alpha)^kx_f\in\psi_{11\ldots1}(\overline{V})$），给出一条周期为 $2^k$ 的轨道.

现在设 $x\in\overline{V}$ 是使 $g^m(x)$ 永不击中这些周期轨道的点.我们证明 $g^m(x)$ 必收敛于 $E$.我们已经知道一旦进入 $\psi^k(\overline{V})$，迭代 $g^m(x)$ 就被困在那里 \eqref{eq:8.52}.我们证明在迭代下 $\psi^k(\overline{V})$ 中的点最终进入 $\psi^{k+1}(\overline{V})$（对每个 $k$）.设 $g^m(x)\in\psi^k(\overline{V})$ 对某个 $k\geq 0$ 成立，设 $g^m(x)=\psi_{j_1\ldots j_k}(y)$，其中 $y\in\overline{V}$.点 $y$ 不能是不动点 $x_f$，否则由 \eqref{eq:8.51}，$g^m(x)$ 将有周期 $2^k$.于是要么 $y\in[x_2,x_4]\cup[x_3,x_1]$，要么 $y\in[x_4,x_f)\cup(x_f,x_3]$.在后一种情形，由于 $g'(x)<-1$ 于 $g^{-1}[x_4,x_3]=[-x_2,x_3]$ 上，$y$ 的迭代离开不稳定不动点 $x_f$ 移动，直至对某个 $r$ 有 $g^r(y)\in[x_2,x_4]\cup[x_3,x_1]$.于是无论哪种情形都可找到 $r\geq 0$ 使 $g^r(y)\in\psi_j(\overline{V})$，$j=1$ 或 $2$.迭代 \eqref{eq:8.51} 得 $g^{r2^k}\circ\psi_{j_1\ldots j_k}=\psi_{j_1\ldots j_k}\circ g^r$，故
\[
  g^{r2^{k}+m}(x)
  =g^{r2^{k}}\bigl(\psi_{j_1\ldots j_k}(y)\bigr)
  =\psi_{j_1\ldots j_k}\bigl(g^{r}(y)\bigr)
  \subset\psi_{j_1\ldots j_k}(\overline{V})
  \subset\psi^{k+1}(\overline{V}).
\]
于是 $\overline{V}$ 中任何不终于周期轨道的点 $x$ 最终进入一切 $\psi^k(\overline{V})$.由 \eqref{eq:8.20}，$E=\bigcap_{k=0}^{\infty}\psi^k(\overline{V})$，故 $g^m(x)$ 到 $E$ 的距离必趋于零.

若 $x\in E$，则对任意 $k$，迭代 $g^m(x)$ 依次访问 $2^k$ 个集合 $\psi_{j_1\ldots j_k}(\overline{V})$ 中的每一个，由此得轨道 $\{g^m(x)\}$ 在 $E$ 中稠密.
\end{proof}

注意，通过估计充分大的 $k$ 的 $\psi_{j_1\ldots j_k}$ 的导数，可把 $\dim E$ 求到任意精度.事实上 $\dim E=0.538\ldots$.\textcite{ruelle1983} 近来证明这类映射的吸引子实际上是 $s$-集，其中实数 $s$ 可以用映射的``压强''表述.

上面描述的映射 $g$ 在某种意义上是一大类呈现混沌行为的 $[-1,1]$ 上映射的代表性映射.设 $f$ 是 $[-1,1]$ 到自身的任何光滑（比如二阶可微）映射，使 $f$ 在 $0$ 处有唯一极大值，$f(0)=1$，$f''(0)<0$，且当 $x<0$ 时 $f'(x)>0$、当 $x>0$ 时 $f'(x)<0$.设 $f_\lambda$ 为 $[-1,1]$ 上的映射 $f_\lambda(x)=\lambda f(x)$，$0<\lambda<1$.当 $\lambda$ 小时映射 $f_\lambda$ 有唯一的稳定不动点（稳定指附近的点在 $f_\lambda$ 迭代下收敛于它）.可以证明，当 $\lambda$ 增大到某个值 $\lambda_1$ 时该不动点变得不稳定，让位于一对周期为 $2$ 的稳定点.继续增大 $\lambda$ 到 $\lambda_2$，这个环分岔成周期为 $4$ 的稳定环，在 $\lambda_3$ 处出现周期为 $8$ 的环，如此等等.（关于这一过程更详细的描述，特别是关于\newterm{logistic-mapping}{``logistic 映射''}{logistic mapping}$f(x)=1-2x^2$——经坐标变换等价于 $[0,1]$ 上的 $f(x)=4x(1-x)$——见 \textcite{may1976,hofstadter1981,feigenbaum1981}.）Feigenbaum（\citeyear{feigenbaum1978,feigenbaum1979,feigenbaum1981}） 发现了序列 $\{\lambda_j\}$ 的一些非凡性质：他证明 $\lambda_j$ 递增于临界值 $\lambda_\infty$，且 $(\lambda_{j+1}-\lambda_j)/(\lambda_{j+2}-\lambda_{j+1})\to\delta$，其中 $\delta=4.669\ldots$ 是一个普适常数，即不依赖于函数 $f$ 的具体形式.已知 $f_{\lambda_\infty}$ 的迭代行为在定性上类似于上文研究的函数 $g$：$f_{\lambda_\infty}$ 对每个 $k$ 总有一条周期为 $2^k$ 的不稳定轨道，以及一个 Cantor 型特征的不变集，它吸引区间的几乎所有点.让 $\lambda$ 越过 $\lambda_\infty$ 后稳定的周期轨道重新占上风，并且当 $\lambda$ 接近对应于周期 $3\cdot 2^j$ 的环的累积的第二个临界值时，它们以类似方式分岔，如此等等.

Feigenbaum 通过考察 $f$ 的周期点当 $j\to\infty$ 的行为证明：对所论类型的任何函数 $f$，
\[
  \lim_{k\to\infty}(-\alpha)^{k}f_{\lambda_\infty}^{2^{k}}\bigl(x/\alpha^{k}\bigr)=g(x)
\]
满足泛函方程 \eqref{eq:8.46}，其中 $\alpha=2.5029\ldots$ 是已知方程有解的那个数.（坐标的线性缩放允许假设 $g(0)=1$.）于是，就 $f_{\lambda_\infty}^m$ 当 $m\to\infty$ 的极限行为而言，$f_{\lambda_\infty}$ 的迭代可以预期以与 $g$ 的迭代类似的方式行动，不仅在定性上，而且在某些方面在定量上也是如此.特别地，人们或许希望 $f_{\lambda_\infty}$ 的吸引子的 Hausdorff 维数总与 $g$ 的相同.\textcite{grassberger1981} 给出了一些计算证据表明情况确实如此：在分岔点的首次累积 $\lambda_\infty$ 处，已知存在的奇怪吸引子具有维数 $0.538\ldots$.

当然，类似的问题也出现在高维，而分析变得相当困难.\textcite{smale1967} 提出的 $\mathbb{R}^2$ 上的变换——\newterm{horseshoe-diffeomorphism}{``马蹄微分同胚''}{horseshoe diffeomorphism}——说明了奇怪吸引子如何作为``折叠与收缩''的结果出现.单位正方形 $S$ 上的映射 $f$ 在定性上如图~\ref{fig:8.6}(a) 所示，其中 $f(a)=a'$ 等，正方形的上下两半被映到马蹄的两臂.$f^m(S)$ 的前几次迭代示于图~\ref{fig:8.6}(b)，显然 $S$ 的所有点的迭代都收敛于 Cantor 型集合 $\bigcap_{m=0}^{\infty}f^m(S)$.通过适当地定义 $f$，可以得到维数介于 $1$ 与 $2$ 之间的任何吸引子.由于 $f$ 的宽度收缩效应，该吸引子在 $f$ 的扰动下稳定.

\begin{figure}[htbp]
  \centering
  \includegraphics[width=0.5\textwidth]{fig8-6.png}
  \caption{Smale 马蹄微分同胚}
  \label{fig:8.6}
\end{figure}

一个更难分析的 $\mathbb{R}^2$ 映射是 \newterm{henon-pomeau-mapping}{Hénon--Pomeau 映射}{Henon-Pomeau mapping}，其坐标形式为
\[
  f(x,y)=(y+1-ax^{2},\,bx),
\]
其中 $a,b$ 为常数.这一映射在一切 $(x,y)$ 处的 Jacobian 为 $-b$，故以一致的速率收缩面积.它受到大量关注，因为在线性坐标变换的意义下它是具有这一性质的最一般的二次映射.它与一维（二次）logistic 映射相关，其吸引子对某些 $(a,b)$ 可以预期具有精细结构.然而，对已知产生分数维数不变集的 $a,b$ 值，结果该集合并不是吸引子.\textcite{henonpomeau1976} 对 $a=1.4$、$b=0.3$ 考察了一些轨道；更近的分析见 \textcite{simo1979} 及 \textcite{whitley1983} 的文章.

至此我们只考虑了离散动力系统，但对连续系统或流（即微分方程组的解曲线）类似的想法也成立.记 $f_t(x)$ 为点 $x$ 在时刻 $t$ 所到达的位置.则
\[
  f_{t_2}\bigl(f_{t_1}(x)\bigr)=f_{t_1+t_2}(x),
\]
这里假定流是自治的（即由不依赖于时间的微分方程导出）.同样，人们寻求闭不变集（使 $f_t(E)\subset E$ 对一切 $t\geq 0$ 成立），它们在如下意义上是吸引子：对大量的 $x$，$f_1(x)$ 对任意大的 $t$ 访问 $E$ 的任何点的每个邻域.

对 Smale 马蹄映射作简单修改提供了一个流产生奇怪吸引子的例子.设 $f_t$ 是方形截面圆环 $V$ 内的一个流，它以恒定速率沿环内侧绕行，绕环一周用单位时间.设 $\Pi$ 是横截于环的平面，并设 $f_t$ 定义得使当 $x\in S=V\cap\Pi$ 时 $f_1(x)=f(x)$，其中 $f$ 是 $S$ 上的马蹄映射.（在 $V$ 上定义具有该性质的流是直截了当的.）于是若 $E$ 是该流的吸引子，则 $E\cap\Pi=\bigcap_{m=0}^{\infty}f^m(S)$ 是马蹄映射的吸引子.为求吸引子的其余部分，只需让这个集合沿环绕行，得 $E=\bigcup_{t=0}^{1}f_t\bigl(\bigcap_{m=0}^{\infty}f^m(S)\bigr)$.于是该吸引子局部上是马蹄吸引子与一条线段的乘积，故由 5.3~节工作的微小变式，其维数等于 $1+\dim\bigcap_{0}^{\infty}f^m(S)$.或者，该吸引子可局部地视为一维 Cantor 集与平面区域的乘积.

实践中，微分方程流的动力学可以极其复杂，这类简单化的处理走不了多远.在寻找奇怪吸引子时被考察最多的微分方程组或许是 \textcite{lorenz1963} 为模拟天气的某些方面而提出的\newterm{lorenz-equations}{方程组}{Lorenz equations}
\begin{equation}\label{eq:8.53}
  \left.
  \begin{aligned}
    \dot{x}&=\sigma(y-x),\\
    \dot{y}&=rx-y-xz,\\
    \dot{z}&=xy-bz,
  \end{aligned}
  \right\}
\end{equation}
其中 $\sigma,r,b$ 为常数.\textcite{sparrow1982} 的书对该方程组的吸引子与分岔作了透彻的论述.混沌行为在某些参数值处出现，通常选取研究的情形是 $\sigma=10$，$r=28$，$b=\frac{8}{3}$.

\textcite{temam1983} 考察了流体力学的 \newterm{navier-stokes-equation}{Navier--Stokes 方程}{Navier-Stokes equation}.特别地，他估计了吸引子的 Hausdorff 维数，以及弱解的奇点作为时空子集的维数.

近来，人们尝试用运动的常数——特别是轨道的 \newterm{lyapunov-exponents}{Lyapunov 特征指数}{Lyapunov exponents}——来表示动力系统吸引子的维数.（Lyapunov 指数是局部化膨胀矩阵的特征值的长期轨道平均.）这些思想的综述及一些例子见 \textcite{frederickson1983}.

关于吸引子及其与湍流等问题的关系的进一步讨论，见会议文集 \textcites{temam1976,bernardratiu1977,haken1982,guckenheimerholmes1983,barenblattioossjoseph1983}.关于吸引子仍有大量工作要做；事实上，许多已被宣称的结果仍有待于建立在严格的基础之上.

\section{布朗运动}

1827 年植物学家 R.~Brown 观察到悬浮在液体中的微小颗粒处于不停的运动之中，并描绘出高度不规则的路径.这后来被解释为分子对颗粒轰击的结果.在空气的烟尘颗粒中也注意到了类似现象.1905 年 Einstein 发表了对这一运动的数学研究，最终导致 Perrin 获得诺贝尔奖的 Avogadro 常数计算.

\textcite{wiener1923} 提出了 \newterm{brownian-motion}{布朗运动}{Brownian motion}的严格概率模型.他构造了 \newterm{wiener-process}{``Wiener 过程''}{Wiener process}，其随机行为与布朗运动非常相似.这里我们概述这一模型的细节，并给出 \textcite{taylor1953} 关于布朗路径以概率 $1$ 具有 Hausdorff 维数 $2$ 的证明.

若 $n\geq 1$，令 $\Omega$ 表示所有满足 $\omega(0)$ 为原点的连续路径 $\omega:[0,\infty)\to\mathbb{R}^n$ 构成的类.把 $\omega(t)$ 想象为描绘路径 $\omega$ 的颗粒在时刻 $t$ 的位置.可以证明：存在概率测度 $p$（$p(\Omega)=1$），定义在 $\Omega$ 的一个大子集系统上，使得：

(a) 路径具有\newterm{independent-increments}{独立增量}{independent increments}，即若 $t_1\leq t_2\leq t_3\leq t_4$，则 $\omega(t_2)-\omega(t_1)$ 与 $\omega(t_4)-\omega(t_3)$ 独立；

(b) $\omega(t+h)-\omega(t)$ 对一切 $t$ 具有零均值、方差为 $h$ 的 Gaussian 分布.特别地，路径增量的分布是\newterm{stationary-process}{平稳的}{stationary process}（$\omega(t+h)-\omega(t)$ 不依赖于 $t$）、\newterm{isotropic-process}{各向同性的}{isotropic process}（与方向无关），且
\begin{equation}\label{eq:8.54}
  p\{\omega:|\omega(t+h)-\omega(t)|\leq\rho\}
  =ch^{-n/2}\int_{0}^{\rho}r^{n-1}\exp(-r^{2}/2h)\,\mathrm{d}r,
\end{equation}
其中 $c$ 是为保证 $p(\Omega)=1$ 而选取的规范化常数.

概率测度 $p$ 可以构造得使一大类路径集合可测，包括 $\Omega$ 的所有形如 $\{\omega:\omega(t_i)\in B_i\ (i=1,2,\ldots,k)\}$ 的子集，其中诸 $B_i$ 是 $\mathbb{R}^n$ 的 Borel 子集.建立这样一个过程的存在性与唯一性相当复杂，有兴趣的读者可参阅例如 \textcite{hida1980}.事实上，测度 $p$ 可以定义在所有（不必连续的）函数 $\omega:[0,\infty)\to\mathbb{R}^n$ 构成的类上，而且由其他条件可得路径以概率 $1$ 连续.

容易看出，把 $\mathbb{R}^n$ 中的布朗过程投影到子空间上给出一个低维布朗过程.

布朗概率分布的重要性在于：它是使路径具有平稳、独立、方差有限的各向同性增量的本质上唯一的分布.这些条件很可能适用于任何物理情境.

可以把布朗运动看作步长趋于零时随机游动的极限情形.该过程\newterm{statistically-self-similar}{在统计意义上是自相似的}{statistically self-similar}：对任意 $\gamma>0$，路径 $\omega(t)$ 与 $\omega(\gamma^2 t)/\gamma$ 具有相同的概率分布.

\begin{lemma}\label{lem:8.22}
若 $0<\lambda<\frac{1}{2}$，则对几乎所有的 $\omega\in\Omega$ 存在 $h_0>0$，使
\[
  |\omega(t+h)-\omega(t)|\leq|h|^{\lambda}
\]
对 $0\leq t\leq 1$ 与 $|h|<h_0$ 成立.
\end{lemma}

\begin{proof}
设 $j$ 与 $m$ 为正整数.则 \eqref{eq:8.54} 蕴含
\[
\begin{aligned}
  p\bigl\{\omega:|\omega(2^{-j}m)-\omega(2^{-j}(m-1))|>2^{-j\lambda}\bigr\}
  &=c\,2^{jn/2}\int_{2^{-j\lambda}}^{\infty}r^{n-1}\exp(-r^{2}2^{j-1})\,\mathrm{d}r\\
  &=c\int_{2^{j(\frac{1}{2}-\lambda)}}^{\infty}r^{n-1}\exp(-r^{2}/2)\,\mathrm{d}r\\
  &\leq c_1\exp\bigl(-2^{j(\frac{1}{2}-\lambda)}\bigr)
  \leq c_2\,2^{-2j},
\end{aligned}
\]
其中作了代换并用了在 $\lambda<\frac{1}{2}$ 时成立的一些粗放的估计，$c_1$ 与 $c_2$ 与 $j,m$ 无关.于是
\begin{align*}
  &p\{\omega:|\omega(2^{-j}m)-\omega(2^{-j}(m-1))|>2^{-j\lambda}\\
  &\qquad\text{对某个满足}\ j>k,\ 1\leq m\leq 2^j\ \text{的}\ j,m\}
  \leq c_2\sum_{j=k}^{\infty}2^{j}2^{-2j}=c_2\,2^{-k}.
\end{align*}
于是对几乎所有的 $\omega\in\Omega$，
\begin{equation}\label{eq:8.55}
  |\omega(2^{-j}m)-\omega(2^{-j}(m-1))|\leq 2^{-j\lambda}\qquad(m=1,2,\ldots,2^{j})
\end{equation}
对一切充分大的 $j$ 成立.任何区间 $(t,t+h)\subset(0,1)$ 都可以表示为形如 $\bigl(2^{-j}(m-1),2^{-j}m\bigr)$（$2^{-j}\leq h$，且每种长度至多两个区间）的相邻二进区间的可数并.（取 $(t,t+h)$ 的所有不被其他这类区间包含的二进子区间.）于是若 $h$ 足够小，$k$ 是满足 $2^{-k}\leq h$ 的最小整数，则利用 \eqref{eq:8.55} 与三角不等式以及 $\omega(t)$ 的连续性，
\[
  |\omega(t+h)-\omega(t)|\leq 2\sum_{j=k}^{\infty}2^{-j\lambda}=2^{-k\lambda}2/(1-2^{-\lambda})\leq h^{\lambda}2/(1-2^{-\lambda}).
\]
只要 $h$ 充分小，这对 $0\leq t<t+h\leq 1$ 与几乎所有的 $\omega\in\Omega$（对任意 $\lambda<\frac{1}{2}$）成立，故引理结论得证.
\end{proof}

\begin{theorem}\label{thm:8.23}
$\mathbb{R}^n$（$n\geq 2$）中的布朗路径以概率 $1$ 具有 Hausdorff 维数 $2$.
\end{theorem}

\begin{proof}
(a) 我们用位势论方法证明对几乎所有的路径 $\omega$ 有 $\dim\omega\geq 2$.固定 $1<s<2$.利用 \eqref{eq:8.54} 得
\[
  \int_{\omega\in\Omega}\frac{\mathrm{d}p(\omega)}{|\omega(t+h)-\omega(t)|^{s}}
  =ch^{-n/2}\int_{0}^{\infty}r^{-s+n-1}\exp(-r^{2}/2h)\,\mathrm{d}r
  =c_1h^{-s/2},
\]
其中在被积函数中作了 $r^2/h$ 的代换，$c_1$ 与 $h,t$ 无关.于是
\[
  \int_{u=0}^{1}\int_{t=0}^{1}\int_{\omega\in\Omega}\frac{\mathrm{d}p(\omega)\,\mathrm{d}t\,\mathrm{d}u}{|\omega(t)-\omega(u)|^{s}}
  =\int_{0}^{1}\int_{0}^{1}\frac{c_1\,\mathrm{d}t\,\mathrm{d}u}{|t-u|^{s/2}}<\infty.
\]
布朗运动的可测性保证 $|\omega(t)-\omega(u)|^{-s}$ 关于 $[0,1]\times[0,1]\times\Omega$ 上的乘积测度可测，故利用 Fubini 定理我们得出：
\[
  \int_{0}^{1}\int_{0}^{1}\frac{\mathrm{d}t\,\mathrm{d}u}{|\omega(t)-\omega(u)|^{s}}<\infty
\]
对几乎所有的 $\omega\in\Omega$ 成立.

路径 $\omega$ 上有自然的质量分布 $\mu_\omega$：$\mu_\omega(E)=\mathscr{L}^{1}\{t:0\leq t\leq 1\ \text{且}\ \omega(t)\in E\}$.于是对几乎所有的 $\omega\in\Omega$，$s$-能量 $I_s(\mu_\omega)$ 有限，故由\cref{cor:6.6}(a)，$s\leq\dim\bigl(\omega[0,1]\bigr)\leq\dim\omega$.这对一切 $s<2$ 成立，故对几乎所有的 $\omega\in\Omega$ 有 $\dim\omega\geq 2$.

(b) 我们用\cref{thm:8.1} 方法的一个变式得到路径维数的上界.取 $s>2$.\cref{lem:8.22} 蕴含对几乎所有的 $\omega\in\Omega$ 可找到 $h_0>0$，使
\[
  |\omega(t+h)-\omega(t)|\leq|h|^{1/s}
  \qquad\text{若}\ |h|\leq h_0\ \text{且}\ 0\leq t<t+h\leq 1.
\]
对这样的路径 $\omega$，当 $h\leq h_0$ 时子弧 $\omega[t,t+h]$ 可以放进半径为 $h^{1/s}$ 的球中.于是若 $m$ 是满足 $1/m\leq h_0$ 的整数，则可以把 $\omega[0,1]=\bigcup_{j=1}^{m}\omega[(j-1)/m,j/m]$ 放进 $m$ 个半径各为 $m^{-1/s}$ 的球中.故当 $\delta\geq 2m^{-1/s}$ 时 $\mathcal{H}^s_{\delta}\bigl(\omega[0,1]\bigr)\leq m2^s(m^{-1/s})^s=2^s$.令 $m\to\infty$，我们得出：若 $s>2$ 则 $\mathcal{H}^s\bigl(\omega[0,1]\bigr)\leq 2^s$；于是事实上由 \eqref{eq:1.14}，$\mathcal{H}^s\bigl(\omega[0,1]\bigr)=0$，从而 $\mathcal{H}^s(\omega)=0$.
\end{proof}

\textcite{taylor1953} 还证明了对 $\mathbb{R}^n$ 中几乎所有的路径有 $\mathcal{H}^s(\omega)=0$.\textcites{ray1963,taylor1964} 进行的更精细的计算表明：以概率 $1$，布朗路径关于测度函数
\[
  h(t)=t^{2}\log\frac{1}{t}\log\log\log\frac{1}{t}
\]
具有正而有限的 Hausdorff 测度.（在 Hausdorff 测度的定义中把 $|U_i|^s$ 换成 $h(U_i)$，这允许更精细的维数定义，见 1.2~节.）

关于布朗路径的二重点、三重点与多重点的几乎必然 Hausdorff 维数也有人提出了类似问题.该理论还被推广到更一般的 \newterm{levy-process}{Lévy 稳定过程}{Levy process}.其中一些专题的进一步讨论与大量文献见 \textcite{rogers1970}（第 133 页）及 \textcites{pruitttaylor1969,taylor1973,pruitt1975,adler1981}，另见 \textcite{mandelbrot1982}（第 39 节）.

\section{结语}

本章我们叙述了分形集合出现的若干情形，以及本书前面讨论的思想可以被应用于其中的场合.Mandelbrot（\citeyear{mandelbrot1975,mandelbrot1977,mandelbrot1982}） 的著述提供了其他此类集合的广阔全景，其中许多尚未经受严格的数学分析.这些著述包含关于更多专题的非常完整的文献，\textcite{rogers1970} 一书关于 Hausdorff 测度应用的最后一章也是如此.

虽然本书所描述的许多工作并不特别新近，但这一学科仍包含许多未解决的问题，有旧有新.分形集几何的理论与应用看来将在未来的许多年里保持为一个活跃的研究领域.借用 Besicovitch 关于 Kakeya 问题的影片中的话：``愿研究这些事情的人们的努力终获成功''.

\begin{exercises}
\item 证明若允许 \eqref{eq:8.1} 中的 $c$ 依赖于 $x$，\cref{thm:8.1} 仍然成立.
\item 证明 $\mathbb{R}$ 上连续可微函数的图像是一个 $1$-集.
\item 证明 Weierstrass 函数 \eqref{eq:8.3} 的图像的维数至多为 $s$.
\item 设 $f:[0,1]\to\mathbb{R}$ 的图像为 $\Gamma$.证明若对一切分割 $0\leq x_0<\cdots<x_m\leq 1$ 有 $\sum_{i=1}^{m}|f(x_i)-f(x_{i-1})|^s\leq c$，则 $\mathcal{H}^s(\Gamma)<\infty$.
\item 设开集条件 \eqref{eq:8.19} 对 $\mathbb{R}^n$ 上的压缩 $\{\psi_j\}_{1}^{m}$ 成立，且对每个 $j$ 与一切 $x,y$ 有
\[
  q_j|x-y|\leq|\psi_j(x)-\psi_j(y)|\leq r_j|x-y|.
\]
若 $E$ 是与这些压缩相伴的不变集，调整\cref{thm:8.6} 的证明以证明 $s'\leq\dim E\leq t$，其中 $s'=s-\log\bigl(\max(r_j/q_j)\bigr)/\log\bigl(\min(1/q_j)\bigr)$，且 $\sum_{1}^{m}q_j^{s}=1=\sum_{1}^{m}r_j^{t}$.
\item 设 $\{B_i\}_{1}^{\infty}$ 是开集 $V\subset\mathbb{R}^n$ 的一个密切填装.证明 $V\subset\bigcup_i B_i'$，其中 $B_i'$ 是与 $B_i$ 同心但半径加倍的球.进而证明剩余集的 Hausdorff 维数至多等于填装的指数.
\item 证明图~\ref{fig:8.4} 所示等边三角形填装的剩余集是 $s=\log 3/\log 2$ 的 $s$-集.（注意剩余集是自相似的.）
\item 证明映射 $f:\mathbb{R}\to\mathbb{R}$，$f(x)=\frac{3}{2}(1-|2x-1|)$ 以 $0$ 为不动点、以 Cantor 集 $E$ 为不变集.证明若 $x\in E\setminus\{0\}$，则迭代 $f^m(x)$ 在 $E$ 中稠密；若 $x\notin E$，则当 $m\to\infty$ 时 $f^m(x)\to-\infty$.（考虑三进展开或许有帮助.）
\end{exercises}


% 参考文献: 模板 v3.2 起 \makebibliography 自动进入后置部分并无编号排版
\renewcommand{\BibliographyNote}{%
  原书各条目末尾方括号中所注为引用该文献的章节，此处以注记形式保留.}
\makebibliography

% 术语索引: 模板 v3.2 起同样自动无编号排版
\printterminology

\end{document}
